\documentclass{article} % Tạo một bản báo cáo
\usepackage{amsmath,amssymb,amsfonts}
\usepackage{cases} 
\usepackage[utf8]{inputenc}
\usepackage[T5]{fontenc} % Để sử dụng Tiếng Việt
\usepackage[fontsize=13pt]{scrextend} % Set fontsize=13pt
\usepackage{amssymb}
%\usepackage[paperheight=29.7cm,paperwidth=21cm,right=2cm,left=3cm,top=2cm,bottom=2.5cm]{geometry}% Chuẩn A4, căn lề phải, trái, trên, dưới.
\usepackage{subfig} %jjj%
\usepackage[paperheight=29.7cm,paperwidth=21cm,right=2cm,left=3cm,top=2cm,bottom=2.5cm,twoside]{geometry}% Chuẩn A4, căn lề phải, trái, trên, dưới.
\usepackage{mathptmx} % Time New Roman
\usepackage{graphicx} % Thư viện chèn ảnh
\usepackage{pdfpages}
\usepackage{booktabs,siunitx}
\usepackage{footnote}
\usepackage{tablefootnote}
\usepackage[flushleft]{threeparttable}
\usepackage[algo2e]{algorithm2e} 
\usepackage{cite}
\usepackage{listings}
\usepackage{algorithmic}
\usepackage{graphicx}
\usepackage{textcomp}
\usepackage{xcolor}
\usepackage{longtable}
\usepackage{flushend}
\usepackage{dblfloatfix}
\usepackage{lipsum}
\usepackage{algorithm}
\usepackage{float}
\usepackage{multirow}
\usepackage{tabularx}
\usepackage{array}
\usepackage{caption}
\usepackage{subcaption}
\usepackage{url}
\usepackage{longtable}
\usepackage{tikz}
\usetikzlibrary{shapes.geometric, arrows.meta, positioning, calc}
\usepackage{booktabs}

\usepackage{float} % Set vị trí chèn ảnh
\usepackage{tikz} % Thư viện tạo khung bìa
\usetikzlibrary{calc} % Thư viện tikz
% --- Bổ sung phần RTOS ---
% Các thư viện TikZ phục vụ vẽ sơ đồ vòng đời Task, kiến trúc đa nhiệm, cơ chế Mutex
\usetikzlibrary{positioning,arrows.meta,shapes.geometric,shapes.misc,fit,backgrounds,chains,calc}
% pgfplots phục vụ vẽ biểu đồ tải CPU và năng lượng (Chương 5)
\usepackage{pgfplots}
\pgfplotsset{compat=1.16}
\usepgfplotslibrary{groupplots}
% --- Hết bổ sung phần RTOS ---
\usepackage{indentfirst} % Thư viện thụt đầu dòng
\renewcommand{\baselinestretch}{1.2} % Giãn dòng 1.2
\setlength{\parskip}{6pt} % Spacing after
\setlength{\parindent}{1cm} % Set khoảng cách thụt đầu dòng mỗi đoạn
\usepackage{titlesec} % Thư viện để set up các kiểu chữ
\setcounter{secnumdepth}{4} % 4 Heading
\titlespacing*{\section}{0pt}{0pt}{30pt} % Heading 1
\titleformat*{\section}{\fontsize{16pt}{0pt}\selectfont \bfseries \centering}

\titlespacing*{\subsection}{0pt}{10pt}{0pt} % Heading 2
\titleformat*{\subsection}{\fontsize{14pt}{0pt}\selectfont \bfseries}

\titlespacing*{\subsubsection}{0pt}{10pt}{0pt} % Heading 3
\titleformat*{\subsubsection}{\fontsize{13pt}{0pt}\selectfont \bfseries \itshape}

\titlespacing*{\paragraph}{0pt}{10pt}{0pt} % Heading 4
\titleformat*{\paragraph}{\fontsize{13pt}{0pt}\selectfont \itshape}


\renewcommand{\figurename}{\fontsize{12pt}{0pt}\selectfont \bfseries Hình}
\renewcommand{\thefigure}{\thesection.\arabic{figure}}
\usepackage[font=bf]{caption}
\captionsetup[figure]{labelsep=space}

\renewcommand{\tablename}{\fontsize{12pt}{0pt}\selectfont \bfseries Bảng}
\renewcommand{\thetable}{\thesection.\arabic{table}}
\captionsetup[table]{labelsep=space}

\usepackage{tabularx}
\newcolumntype{s}{>{\hsize=.3\hsize}X}
\newcolumntype{y}{>{\hsize=.4\hsize}X}
\newcolumntype{d}{>{\hsize=.1\hsize}X}
\newcolumntype{a}{>{\hsize=1.1\hsize}X}
\newcolumntype{g}{>{\hsize=5\hsize}X}
\renewcommand{\tabularxcolumn}[1]{>{\small}m{#1}}

\renewcommand{\theequation}{\thesection.\arabic{equation}} % Thay đổi đánh số phương trình mặc định
\newtheorem{theorem}{Định lý}[section]
\newtheorem{defn}[theorem]{Định nghĩa}
\newtheorem{corollary}[theorem]{Hệ quả}
\newtheorem{lemma}[theorem]{Bổ đề}

\usepackage{lipsum} % Thư viện tạo chữ linh tinh.
\renewcommand{\contentsname}{MỤC LỤC}
\renewcommand{\listfigurename}{DANH MỤC HÌNH VẼ}
\renewcommand{\listtablename}{DANH MỤC BẢNG BIỂU}
\renewcommand{\refname}{TÀI LIỆU THAM KHẢO}

\usepackage[unicode]{hyperref}

% ==== Sửa lỗi hiển thị bookmark/mục lục PDF (chỉ chỉnh phần setup) ====
% Nguyên nhân: mỗi trang CHƯƠNG dùng \section* rồi \addcontentsline nhưng THIẾU
% \phantomsection nên không tạo được "mỏ neo" (anchor) riêng. Hyperref vì thế tính
% sai cấp bậc và lồng nhầm các chương vào mục đứng trước (ví dụ ABSTRACT).
% Cách sửa, không đụng tới nội dung các chương:
%   1) Nạp gói bookmark (sau hyperref) để dựng lại cây bookmark theo đúng cấp.
%   2) Tự chèn \phantomsection cho MỌI mục cấp "section" được ghi ra mục lục, giúp
%      mỗi chương có mỏ neo riêng -> hiển thị tách bạch, đúng liên kết, và các mục
%      con (\subsection, \subsubsection) tự động lồng vào đúng chương của nó.
\usepackage{bookmark}
\makeatletter
\let\hyfix@addcontentsline\addcontentsline
\renewcommand{\addcontentsline}[3]{%
  \def\hyfix@type{#2}\def\hyfix@sec{section}%
  \ifx\hyfix@type\hyfix@sec \phantomsection \fi
  \hyfix@addcontentsline{#1}{#2}{#3}%
}
\makeatother
% ==== Hết phần sửa ====

\usepackage{colortbl}
\definecolor{LightCyan}{rgb}{0.88,1,1}
\usepackage{forloop}
\newcounter{loopcntr}
\newcommand{\rpt}[2][1]{\forloop{loopcntr}{0}{\value{loopcntr}<#1}{#2}}

\begin{document}
\thispagestyle{empty} % Remove page number on this page

\begin{center}
    \vspace{12pt} % line spacing
        \textbf{\fontsize{15pt}{0pt} \selectfont{} ĐẠI HỌC BÁCH KHOA HÀ NỘI \\ \vspace{3pt} TRƯỜNG ĐIỆN - ĐIỆN TỬ}
    \vspace{0.5cm}
    
% INSERT LOGO
\begin{figure}[H]
    \centering
    
\includegraphics[width=2cm,height=3cm]{Image/logodhbk.png}
\end{figure}

% THESIS TYPE
\vspace{48pt}
        \fontsize{25pt}{0pt}\selectfont{} \textbf{ĐỒ ÁN TỐT NGHIỆP}
\vspace{24pt}

 % THESIS TITTLE
        \textbf{\fontsize{17pt}{0pt}\selectfont{} 
        NGHIÊN CỨU THIẾT KẾ VÀ CHẾ TẠO HỆ THỐNG IOT GIÁM SÁT VI KHÍ HẬU CHUỒNG NUÔI GIA CẦM ỨNG DỤNG LORA VÀ ĐIỆN TOÁN BIÊN}
    \vspace{18pt}

% NAME
        \fontsize{14pt}{0pt}\selectfont{} \textbf{NGUYỄN TIẾN ĐẠT }
    \vspace{3pt}

% EMAIL    
        \fontsize{14pt}{0pt}\selectfont{} dat.nt223689@sis.hust.edu.vn
    \vspace{12pt} % line spacing

% MAJOR
        \fontsize{14pt}{0pt}\selectfont{} \textbf{Ngành Kỹ thuật Điện tử Viễn thông}  
\end{center}
    \vspace{48pt}

% INFORMATION: SUPERVISOR NAME, STUDENT NAME,...
\begin{table}[H]
    \centering
        \begin{tabular}{l l c}
            \textbf{Giảng viên hướng dẫn:}    &  PGS.TS Trần Quang Vinh \vspace{6pt} &  \\ 
            &\vspace{3pt} & \fontsize{10pt}{0pt}\selectfont{} Chữ ký của GVHD \\
            \textbf{Khoa:} & Kỹ thuật truyền thông \vspace{3pt}\\ 
            \textbf{Trường:} & Điện - Điện tử
        \end{tabular}
\end{table}


\begin{center}
    \vspace{78pt}
    \fontsize{14pt}{0pt}\selectfont{} Hà Nội, 07/2026
\end{center}

\newpage % Insert a new page

\thispagestyle{empty} % Remove header/footer from this page
\mbox{} 
% \thispagestyle{empty}
% \begin{figure}[H]
%     \begin{center}
%         \includegraphics[width=1.05\linewidth]{Thuc_Hanh/Danhgia/DanhGia-part-1.pdf}
%      \end{center}
%  \end{figure}
%  \newpage
% \thispagestyle{empty}
% \begin{figure}[H]
%     \begin{center}
%         \includegraphics[width=1.05\linewidth]{Thuc_Hanh/Danhgia/DanhGia-part-3.pdf}
%      \end{center}
%  \end{figure}
% \newpage
% \thispagestyle{empty}
% \begin{figure}[H]
%     \begin{center}
%         \includegraphics[width=1.05\linewidth]{Thuc_Hanh/Danhgia/DanhGia-part-5.pdf}
%      \end{center}
%  \end{figure}
% \cleardoublepage

% \includepdf[pages=-]{Thuc_Hanh/Danhgia/DanhGia.pdf}

% ===================================================
%  MACRO: in một phiếu đánh giá CBHD (dùng lại 2 lần)
% ===================================================
\newcommand{\phieuCBHD}{%
  \begingroup
  \newgeometry{top=1.4cm,bottom=1.4cm,left=1.8cm,right=1.5cm}
  \thispagestyle{empty}
  \setlength{\parindent}{0pt}
  \setlength{\parskip}{2pt}
  \renewcommand{\baselinestretch}{1.1}\selectfont

  % ---------- Tiêu đề ----------
  \begin{center}
    {\fontsize{12pt}{14pt}\selectfont ĐẠI HỌC BÁCH KHOA HÀ NỘI}\\[1pt]
    {\fontsize{12pt}{14pt}\selectfont\bfseries\underline{TRƯỜNG ĐIỆN -- ĐIỆN TỬ}}\\[7pt]
    {\fontsize{13pt}{15pt}\selectfont\bfseries ĐÁNH GIÁ ĐỒ ÁN TỐT NGHIỆP}\\[1pt]
    {\fontsize{12pt}{14pt}\selectfont\bfseries (DÀNH CHO CÁN BỘ HƯỚNG DẪN)}
  \end{center}

% ---------- Thông tin chung ----------
\vspace{-10pt}
\begin{table}[H]
\centering
\renewcommand{\arraystretch}{1.14}
\begin{tabular}{p{3.4cm} p{13.5cm}}
Tên đề tài: & {\bfseries  Nghiên cứu thiết kế và chế tạo hệ thống IoT giám sát vi khí hậu chuồng nuôi gia cầm ứng dụng LoRa và Điện toán biên} \\
Họ tên SV: & {\bfseries Nguyễn Tiến Đạt} \hfill MSSV: {\bfseries 20223689} \\
Cán bộ hướng dẫn: & {\bfseries PGS.TS. Trần Quang Vinh} \\
\end{tabular}
\end{table}
\vspace{-6pt}
%----------------------------------------

  \footnotesize
  \renewcommand{\arraystretch}{1.18}
  \setlength{\tabcolsep}{4pt}

  \noindent
  \begin{tabular}{|>{\centering\arraybackslash\bfseries}m{0.62cm}%
                  |>{\centering\arraybackslash}m{2.05cm}%
                  |m{12.55cm}%
                  |>{\centering\arraybackslash}m{1.3cm}|}
  \hline
  \textbf{STT} & \textbf{Tiêu chí (Điểm tối đa)} & \centering\arraybackslash\textbf{Hướng dẫn đánh giá tiêu chí} & \textbf{Điểm tiêu chí} \\
  \hline

  1 &
  \textbf{Thái độ làm việc} \newline\textbf{(2,5 điểm)} &
  Nghiêm túc, tích cực và chủ động trong quá trình làm ĐATN.

  Hoàn thành đầy đủ và đúng tiến độ các nội dung được GVHD giao.
  & 2,5\\
  \hline

  2 &
  \textbf{Kỹ năng viết quyển ĐATN} \newline\textbf{(2 điểm)} &
  Trình bày đúng mẫu quy định, bố cục các chương logic và hợp lý: Bảng biểu, hình
  ảnh rõ ràng, có tiêu đề, được đánh số thứ tự và được giải thích hay để cập đến
  trong đồ án, có căn lề, dấu cách sau dấu chấm, dấu phẩy, có mở đầu chương và
  kết luận chương, có liệt kê tài liệu tham khảo và có trích dẫn, v.v.

  Kỹ năng diễn đạt, phân tích, giải thích, lập luận: Cấu trúc câu rõ ràng, văn
  phong khoa học, lập luận logic và có cơ sở, thuật ngữ chuyên ngành phù hợp, v.v.
  & 2\\
  \hline

  3 &
  \textbf{Nội dung và kết quả đạt được} \newline\textbf{(5 điểm)} &
  Nêu rõ tính cấp thiết, ý nghĩa khoa học và thực tiễn của đề tài, các vấn đề và
  các giả thuyết, phạm vi ứng dụng của đề tài. Thực hiện đầy đủ quy trình nghiên
  cứu: Đặt vấn đề, mục tiêu đề ra, phương pháp nghiên cứu/ giải quyết vấn đề, kết
  quả đạt được, đánh giá và kết luận.

  Nội dung và kết quả được trình bày một cách logic và hợp lý, được phân tích và
  đánh giá thỏa đáng. Biện luận phân tích kết quả mô phỏng/ phần mềm/ thực
  nghiệm, so sánh kết quả đạt được với kết quả trước đó có liên quan.

  Chỉ rõ phù hợp giữa kết quả đạt được và mục tiêu ban đầu đề ra đồng thời cung
  cấp lập luận để đề xuất hướng giải quyết có thể thực hiện trong tương lai. Hàm
  lượng khoa học/ độ phức tạp cao, có tính mới/tính sáng tạo trong nội dung và
  kết quả của đồ án.
  & 5\\
  \hline

  4 &
  \textbf{Điểm thành tích} \newline\textbf{(1 điểm)} &
  Có bài báo KH được đăng hoặc chấp nhận đăng/ đạt giải SV NCKH giải 3 cấp
  Trường trở lên/ Các giải thưởng khoa học trong nước, quốc tế từ giải 3 trở lên/
  Có đăng ký bằng phát minh sáng chế. \textbf{(1 điểm)}

  Được báo cáo tại hội đồng cấp Trường trong hội nghị SV NCKH nhưng không đạt
  giải từ giải 3 trở lên/ Đạt giải khuyến khích trong cuộc thi khoa học trong nước,
  quốc tế/ Kết quả đồ án là sản phẩm ứng dụng có tính hoàn thiện cao, yêu cầu
  khối lượng thực hiện lớn. \textbf{(0,5 điểm)}
  & 0\\
  \hline

  \multicolumn{3}{|r|}{\textbf{Điểm tổng các tiêu chí:}} & 9,5\\
  \hline
  \multicolumn{3}{|r|}{\textbf{Điểm hướng dẫn:}} & 9,5\\
  \hline
  \end{tabular}

  % ---------- Chữ ký ----------
  %\vspace{6pt}
  \begin{flushright}
    \fontsize{12pt}{20pt}\selectfont\textit{Cán bộ hướng dẫn}\\[22pt]
    \fontsize{12pt}{24pt}\selectfont\textbf{PGS.TS. Trần Quang Vinh}
  \end{flushright}

  \noindent{\fontsize{9.5pt}{11pt}\selectfont\textit{%
    Điểm từng tiêu chí cho lẻ đến 0,5.
    Nếu Điểm tổng các tiêu chí $>$ 10 thì Điểm hướng dẫn làm tròn thành 10.}}

  \restoregeometry
  \renewcommand{\baselinestretch}{1.2}\selectfont
  \setlength{\parskip}{6pt}
  \setlength{\parindent}{1cm}
  \endgroup
}% end \phieuCBHD

% ===================================================
%  MACRO: phiếu đánh giá CBPB (dùng lại 2 lần)
% ===================================================
\newcommand{\phieuCBPB}{%
  \begingroup
  \newgeometry{top=1.4cm,bottom=1.4cm,left=1.8cm,right=1.5cm}
  \thispagestyle{empty}
  \setlength{\parindent}{0pt}
  \setlength{\parskip}{2pt}
  \renewcommand{\baselinestretch}{1.1}\selectfont

  % ---------- Tiêu đề ----------
  \begin{center}
    {\fontsize{12pt}{14pt}\selectfont ĐẠI HỌC BÁCH KHOA HÀ NỘI}\\[1pt]
    {\fontsize{12pt}{14pt}\selectfont\bfseries\underline{TRƯỜNG ĐIỆN -- ĐIỆN TỬ}}\\[7pt]
    {\fontsize{13pt}{15pt}\selectfont\bfseries ĐÁNH GIÁ ĐỒ ÁN TỐT NGHIỆP}\\[1pt]
    {\fontsize{12pt}{14pt}\selectfont\bfseries (DÀNH CHO CÁN BỘ PHẢN BIỆN)}
  \end{center}

 % ---------- Thông tin chung ----------
\vspace{-10pt}
\begin{table}[H]
\centering
\renewcommand{\arraystretch}{1.14}
\begin{tabular}{p{3.4cm} p{13.5cm}}
Tên đề tài: & {\bfseries  Nghiên cứu thiết kế và chế tạo hệ thống IoT giám sát vi khí hậu chuồng nuôi gia cầm ứng dụng LoRa và Điện toán biên} \\
Họ tên SV: & {\bfseries Nguyễn Tiến Đạt} \hfill MSSV: {\bfseries 20223689} \\
Cán bộ phản biện: & ................................................................................................................ \\
\end{tabular}
\end{table}
\vspace{-6pt}
%----------------------------------------

  % ---------- Bảng đánh giá ----------
  \footnotesize
  \renewcommand{\arraystretch}{1.18}
  \setlength{\tabcolsep}{4pt}

  \noindent
  \begin{tabular}{|>{\centering\arraybackslash\bfseries}m{0.62cm}%
                  |>{\centering\arraybackslash}m{2.05cm}%
                  |m{12.55cm}%
                  |>{\centering\arraybackslash}m{1.3cm}|}
  \hline
  \textbf{STT} & \textbf{Tiêu chí (Điểm tối đa)} & \centering\arraybackslash\textbf{Hướng dẫn đánh giá tiêu chí} & \textbf{Điểm tiêu chí} \\
  \hline

  1 &
  \textbf{Trình bày quyển ĐATN} \newline\textbf{(4 điểm)} &
  Trình bày đúng mẫu quy định, bố cục các chương logic và hợp lý: Bảng biểu, hình
  ảnh rõ ràng, có tiêu đề, được đánh số thứ tự và được giải thích hay để cập đến
  trong đồ án, có căn lề, dấu cách sau dấu chấm, dấu phẩy, có mở đầu chương và
  kết luận chương, có liệt kê tài liệu tham khảo và có trích dẫn, v.v.

  Kỹ năng diễn đạt, phân tích, giải thích, lập luận: Cấu trúc câu rõ ràng, văn
  phong khoa học, lập luận logic và có cơ sở, thuật ngữ chuyên ngành phù hợp, v.v.
  & \\
  \hline

  2 &
  \textbf{Nội dung và kết quả đạt được} \newline\textbf{(5,5 điểm)} &
  Nêu rõ tính cấp thiết, ý nghĩa khoa học và thực tiễn của đề tài, các vấn đề và
  các giả thuyết, phạm vi ứng dụng của đề tài. Thực hiện đầy đủ quy trình nghiên
  cứu: Đặt vấn đề, mục tiêu đề ra, phương pháp nghiên cứu/ giải quyết vấn đề, kết
  quả đạt được, đánh giá và kết luận.

  Nội dung và kết quả được trình bày một cách logic và hợp lý, được phân tích và
  đánh giá thỏa đáng. Biện luận phân tích kết quả mô phỏng/ phần mềm/ thực
  nghiệm, so sánh kết quả đạt được với kết quả trước đó có liên quan.

  Chỉ rõ phù hợp giữa kết quả đạt được và mục tiêu ban đầu đề ra đồng thời cung
  cấp lập luận để đề xuất hướng giải quyết có thể thực hiện trong tương lai. Hàm
  lượng khoa học/ độ phức tạp cao, có tính mới/tính sáng tạo trong nội dung và
  kết quả của đồ án.
  & \\
  \hline

  3 &
  \textbf{Điểm thành tích} \newline\textbf{(1 điểm)} &
  Có bài báo KH được đăng hoặc chấp nhận đăng/ đạt giải SV NCKH giải 3 cấp
  Trường trở lên/ Các giải thưởng KH quốc tế/ trong nước từ giải 3 trở lên/
  Có đăng ký bằng phát minh sáng chế. \textbf{(1 điểm)}

  Được báo cáo tại hội đồng cấp Trường trong hội nghị SV NCKH nhưng không đạt
  giải từ giải 3 trở lên/ Đạt giải khuyến khích trong các kỳ thi quốc gia và
  quốc tế khác về chuyên ngành. \textbf{(0,5 điểm)}
  & \\
  \hline

  \multicolumn{3}{|r|}{\textbf{Điểm tổng các tiêu chí:}} & \\
  \hline
  \multicolumn{3}{|r|}{\textbf{Điểm phản biện:}} & \\
  \hline
  \end{tabular}

  % ---------- Chữ ký ----------
  \vspace{8pt}
  \begin{flushright}
    \fontsize{12pt}{14pt}\selectfont\textit{Cán bộ phản biện}\\[2pt]
    \fontsize{11pt}{13pt}\selectfont\textit{(Ký và ghi rõ họ tên)}\\[24pt]
    \phantom{PGS.TS. Trần Quang Vinh}
  \end{flushright}

  \vspace{4pt}
  \noindent{\fontsize{9.5pt}{11pt}\selectfont\textit{%
    Điểm từng tiêu chí cho lẻ đến 0,5.
    Nếu Điểm tổng các tiêu chí $>$ 10 thì Điểm hướng dẫn làm tròn thành 10.}}

  \restoregeometry
  \renewcommand{\baselinestretch}{1.2}\selectfont
  \setlength{\parskip}{6pt}
  \setlength{\parindent}{1cm}
  \endgroup
}% end \phieuCBPB

% ===================================================
%  In hai bản phiếu CBHD + hai bản phiếu CBPB
% ===================================================
\phieuCBHD
\cleardoublepage

\phieuCBPB
\cleardoublepage
\newpage
\section*{LỜI NÓI ĐẦU}
\thispagestyle{empty}
Trong bối cảnh chuyển đổi số đang diễn ra mạnh mẽ ở khắp các lĩnh vực, nông nghiệp thông minh (Smart Agriculture) dần trở thành xu hướng tất yếu nhằm nâng cao năng suất và chất lượng nông sản. Tại Việt Nam, ngành chăn nuôi cũng đang có bước chuyển mình rõ rệt từ mô hình hộ gia đình nhỏ lẻ sang các trang trại tập trung, ứng dụng công nghệ cao.

Tuy nhiên, bài toán khó nhất đối với các trang trại quy mô lớn là khâu kiểm soát vi khí hậu chuồng nuôi. Những biến động thất thường về nhiệt độ, độ ẩm hay sự tích tụ khí độc (NH3, H2S, CO2) nếu không được xử lý kịp thời sẽ làm vật nuôi bị stress, sinh bệnh và gây thiệt hại kinh tế nặng nề. Rõ ràng, cách giám sát thủ công truyền thống vừa tốn sức người, vừa đứt quãng và không còn đáp ứng được yêu cầu thực tế.

Xuất phát từ nhu cầu đó, em quyết định chọn đề tài: "NGHIÊN CỨU THIẾT KẾ VÀ CHẾ TẠO HỆ THỐNG IOT GIÁM SÁT VI KHÍ HẬU CHUỒNG NUÔI GIA CẦM ỨNG DỤNG LORA VÀ ĐIỆN TOÁN BIÊN". Mục tiêu của đồ án là phát triển một hệ thống giám sát tự động, kết hợp thiết bị IoT và mạng truyền thông vô tuyến tầm xa (LoRa) để thu thập dữ liệu môi trường liên tục. Điểm cốt lõi là dữ liệu được xử lý ngay tại thiết bị biên theo hướng Điện toán biên (Edge Computing), giúp người chăn nuôi theo dõi và phản ứng nhanh với sự cố theo thời gian thực, kể cả khi đường truyền Internet bị gián đoạn.

Báo cáo này sẽ trình bày chi tiết toàn bộ quá trình thực hiện đồ án, từ bước khảo sát lý thuyết, lựa chọn linh kiện, thiết kế mạch phần cứng thu thập dữ liệu, cho đến khâu lập trình và kiểm thử hệ thống. Dù đã nỗ lực hoàn thiện, nhưng do giới hạn về thời gian và kinh nghiệm thực tiễn, đồ án chắc chắn không tránh khỏi những thiếu sót. Em rất mong nhận được những nhận xét và góp ý quý báu từ quý thầy cô để đề tài được hoàn thiện hơn.

Em xin chân thành cảm ơn!

\cleardoublepage
\section*{LỜI CAM ĐOAN}
\thispagestyle{empty}
Tôi tên là NGUYỄN TIẾN ĐẠT, mã số sinh viên 20223689, sinh viên lớp Tài năng Điện tử Viễn thông khóa K67. Người hướng dẫn là PGS.TS. TRẦN QUANG VINH. Tôi xin cam đoan toàn bộ nội dung được trình bày trong đồ án \textbf{NGHIÊN CỨU THIẾT KẾ VÀ CHẾ TẠO HỆ THỐNG IOT
GIÁM SÁT VI KHÍ HẬU CHUỒNG NUÔI GIA CẦM ỨNG
DỤNG LORA VÀ ĐIỆN TOÁN BIÊN} là kết quả tìm hiểu nghiên cứu của tôi. Mọi thông tin trích dẫn đều tuân thủ các quy định về sở hữu trí tuệ; các tài liệu tham khảo được liệt kê rõ ràng. Tôi xin chịu hoàn toàn trách nhiệm với những nội dung được viết trong đồ án này.

\vspace{6pt}
\hspace{8.3cm}Hà Nội, ... tháng ..., 2026

\hspace{9cm}\textbf{Người cam đoan}

\vspace{2cm}
\hspace{8.85cm}\textbf{NGUYỄN TIẾN ĐẠT}
\cleardoublepage

\addtocontents{toc}{\protect\thispagestyle{empty}}
\tableofcontents %tạo mục lục tự động
\thispagestyle{empty}
\cleardoublepage
\pagenumbering{roman} %đánh số thứ tự la mã


\pagenumbering{roman} % Đánh số thứ tự la mã
\section*{DANH SÁCH TỪ VIẾT TẮT}
\phantomsection \addcontentsline{toc}{section}{\numberline {} DANH SÁCH TỪ VIẾT TẮT}

\begin{longtable}{@{} p{0.20\textwidth} p{0.75\textwidth} @{}}
\toprule
\textbf{Viết tắt} & \textbf{Giải thích đầy đủ (Tiếng Anh / Tiếng Việt)} \\
\midrule
\endfirsthead

\toprule
\textbf{Viết tắt} & \textbf{Giải thích đầy đủ (Tiếng Anh / Tiếng Việt)} \\
\midrule
\endhead

\midrule
\multicolumn{2}{r}{{(Tiếp tục ở trang sau...)}} \\
\bottomrule
\endfoot

\bottomrule
\endlastfoot

\textbf{ACK}    & Acknowledgement (Gói tin xác nhận) \\
\textbf{ADC}    & Analog-to-Digital Converter (Bộ chuyển đổi tương tự sang số) \\
\textbf{AEAD}   & Authenticated Encryption with Associated Data (Mã hóa có xác thực kèm dữ liệu liên kết) \\
\textbf{AES}    & Advanced Encryption Standard (Chuẩn mã hóa tiên tiến) \\
\textbf{AI}     & Artificial Intelligence (Trí tuệ nhân tạo) \\
\textbf{App}    & Application (Ứng dụng di động) \\
\textbf{ASCON}  & Authenticated Encryption Scheme (Thuật toán mật mã hạng nhẹ chuẩn NIST) \\
\textbf{BLE}    & Bluetooth Low Energy (Bluetooth năng lượng thấp) \\
\textbf{BW}     & Bandwidth (Băng thông) \\
\textbf{CC/CV}  & Constant Current / Constant Voltage (Dòng không đổi / Điện áp không đổi) \\
\textbf{CPU}    & Central Processing Unit (Bộ xử lý trung tâm) \\
\textbf{CR}     & Coding Rate (Tỷ lệ mã hóa sửa lỗi) \\
\textbf{CRC}    & Cyclic Redundancy Check (Kiểm tra dư thừa vòng) \\
\textbf{CSS}    & Chirp Spread Spectrum (Điều chế trải phổ chirp) \\
\textbf{EMA}    & Exponential Moving Average (Trung bình trượt mũ) \\
\textbf{EMI}    & Electromagnetic Interference (Nhiễu điện từ) \\
\textbf{ESP}    & Espressif Systems (Tên hãng sản xuất chip) \\
\textbf{FEC}    & Forward Error Correction (Sửa lỗi tiến) \\
\textbf{FIFO}   & First In First Out (Cơ chế hàng đợi vào trước ra trước) \\
\textbf{FR4}    & Flame Retardant 4 (Vật liệu nền mạch in tiêu chuẩn) \\
\textbf{FSM}    & Finite State Machine (Máy trạng thái hữu hạn) \\
\textbf{GCM}    & Galois/Counter Mode (Chế độ mã hóa Galois/Counter) \\
\textbf{GND}    & Ground (Nối đất/Mass) \\
\textbf{GPIO}   & General Purpose Input/Output (Cổng vào/ra đa năng) \\
\textbf{HIL}    & Hardware-In-the-Loop (Kiểm thử phần cứng trong vòng lặp) \\
\textbf{HMI}    & Human-Machine Interface (Giao diện người - máy) \\
\textbf{HVAC}   & Heating, Ventilation, and Air Conditioning (Hệ thống sưởi, thông gió, điều hòa) \\
\textbf{I2C}    & Inter-Integrated Circuit (Giao tiếp nối tiếp đồng bộ) \\
\textbf{IoT}    & Internet of Things (Internet vạn vật) \\
\textbf{IP65}   & Ingress Protection 65 (Cấp bảo vệ chống bụi và nước) \\
\textbf{ISR}    & Interrupt Service Routine (Trình phục vụ ngắt) \\
\textbf{JSON}   & JavaScript Object Notation (Định dạng trao đổi dữ liệu) \\
\textbf{LDO}    & Low-Dropout Regulator (Bộ ổn áp sụt áp thấp) \\
\textbf{LED}    & Light Emitting Diode (Diode phát quang) \\
\textbf{LNA}    & Low Noise Amplifier (Bộ khuếch đại tạp âm thấp) \\
\textbf{LoRa}   & Long Range (Công nghệ truyền thông không dây tầm xa) \\
\textbf{LOS}    & Line-of-Sight (Tầm nhìn thẳng, không vật cản) \\
\textbf{LPWAN}  & Low-Power Wide-Area Network (Mạng diện rộng công suất thấp) \\
\textbf{Lux}    & Illuminance Unit (Đơn vị đo độ rọi ánh sáng) \\
\textbf{MA}     & Moving Average (Trung bình trượt) \\
\textbf{MAC}    & Message Authentication Code (Mã xác thực bản tin) \\
\textbf{MCU}    & Microcontroller Unit (Vi điều khiển) \\
\textbf{MOSFET} & Metal-Oxide-Semiconductor FET (Transistor hiệu ứng trường) \\
\textbf{MOX}    & Metal-Oxide Semiconductor (Cảm biến khí oxit kim loại bán dẫn) \\
\textbf{MQTT}   & Message Queuing Telemetry Transport (Giao thức truyền tin) \\
\textbf{NLOS}   & Non-Line-of-Sight (Có vật cản che khuất tầm nhìn) \\
\textbf{OTA}    & Over-The-Air (Cập nhật firmware từ xa qua sóng) \\
\textbf{PA}     & Power Amplifier (Bộ khuếch đại công suất) \\
\textbf{PCB}    & Printed Circuit Board (Mạch in) \\
\textbf{PDR}    & Packet Delivery Ratio (Tỷ lệ giao gói thành công) \\
\textbf{PLC}    & Programmable Logic Controller (Bộ điều khiển logic lập trình) \\
\textbf{PLR}    & Packet Loss Rate (Tỷ lệ mất gói tin) \\
\textbf{PPM}    & Parts Per Million (Phần triệu - Đơn vị đo nồng độ khí) \\
\textbf{PSRAM}  & Pseudo-Static RAM (Bộ nhớ RAM giả tĩnh mở rộng) \\
\textbf{PWM}    & Pulse Width Modulation (Điều chế độ rộng xung) \\
\textbf{QoS}    & Quality of Service (Chất lượng dịch vụ) \\
\textbf{REST}   & Representational State Transfer (Kiến trúc giao tiếp web) \\
\textbf{RF}     & Radio Frequency (Tần số vô tuyến) \\
\textbf{RH}     & Relative Humidity (Độ ẩm tương đối) \\
\textbf{RSSI}   & Received Signal Strength Indicator (Chỉ số cường độ tín hiệu thu) \\
\textbf{RTC}    & Real Time Clock (Đồng hồ thời gian thực) \\
\textbf{RTOS}   & Real-Time Operating System (Hệ điều hành thời gian thực) \\
\textbf{RTT}    & Round-Trip Time (Thời gian trễ khứ hồi) \\
\textbf{SF}     & Spreading Factor (Hệ số trải phổ) \\
\textbf{SMD}    & Surface-Mount Device (Linh kiện dán bề mặt) \\
\textbf{SMS}    & Short Message Service (Dịch vụ tin nhắn ngắn) \\
\textbf{SNR}    & Signal-to-Noise Ratio (Tỷ số tín hiệu trên nhiễu) \\
\textbf{SNTP}   & Simple Network Time Protocol (Giao thức đồng bộ thời gian qua mạng) \\
\textbf{SoC}    & System-on-Chip (Hệ thống trên một chip) \\
\textbf{SPI}    & Serial Peripheral Interface (Giao diện ngoại vi nối tiếp) \\
\textbf{SPIFFS} & SPI Flash File System (Hệ thống tập tin trên bộ nhớ Flash SPI) \\
\textbf{TCXO}   & Temperature Compensated Crystal Oscillator (Dao động thạch anh bù nhiệt) \\
\textbf{THI}    & Temperature Humidity Index (Chỉ số nhiệt ẩm) \\
\textbf{TLS}    & Transport Layer Security (Bảo mật tầng giao vận) \\
\textbf{UART}   & Universal Asynchronous Receiver-Transmitter (Giao tiếp không đồng bộ) \\
\textbf{USB}    & Universal Serial Bus (Chuẩn kết nối tuần tự đa năng) \\
\textbf{VCC}    & Voltage Common Collector (Điện áp nguồn dương) \\
\textbf{Wi-Fi}  & Wireless Fidelity (Mạng không dây cục bộ) \\
\textbf{WSN}    & Wireless Sensor Network (Mạng cảm biến không dây) \\

\end{longtable}
\newpage
{\let\oldnumberline\numberline
\renewcommand{\numberline}{\figurename~\oldnumberline}
\listoffigures} %tạo mục hình vẽ tự động
\phantomsection \addcontentsline{toc}{section}{\numberline {} DANH MỤC HÌNH VẼ}
\cleardoublepage

{\let\oldnumberline\numberline
\renewcommand{\numberline}{\tablename~\oldnumberline}
\listoftables}
\phantomsection \addcontentsline{toc}{section}{\numberline {} DANH MỤC BẢNG BIỂU}
\cleardoublepage
\section*{TÓM TẮT ĐỒ ÁN}
\phantomsection\addcontentsline{toc}{section}{\numberline {}TÓM TẮT ĐỒ ÁN}

Trong đề tài \textit{Nghiên cứu thiết kế và chế tạo hệ thống IoT giám sát vi khí hậu chuồng nuôi gia cầm ứng dụng LoRa và Điện toán biên}, một hệ thống IoT phân tán được đề xuất nhằm giải quyết "Nghịch lý đám mây", tối ưu hóa độ trễ điều khiển và năng lượng tiêu thụ. Nội dung đồ án bao gồm các chương: 

\textit{Chương 1: Tổng quan về đề tài và cơ sở khoa học}

Trình bày thực trạng chăn nuôi gia cầm và hạn chế của các giải pháp IoT đám mây khi mất kết nối. Từ đó, xác lập mục tiêu thiết kế thiết bị Node ứng dụng Điện toán biên và LoRa với các chỉ tiêu khắt khe về độ trễ, năng lượng và bảo mật.

\textit{Chương 2: Cơ sở lý thuyết}

Hệ thống hóa các lý thuyết cốt lõi gồm: Mạng cảm biến không dây (WSN), sinh lý gia cầm (chỉ số THI), kiến trúc Điện toán biên, lớp vật lý LoRa (kỹ thuật CSS), FreeRTOS và cơ sở thuật toán mã hóa AES-128.

\textit{Chương 3: Phân tích và thiết kế tổng thể hệ thống}

Biện luận lựa chọn giải pháp phần cứng (ESP32-S3, tổ hợp cảm biến, IC nguồn). Đồng thời, thiết kế kiến trúc phần mềm đa nhiệm trên FreeRTOS, lập thuật toán máy trạng thái (FSM) và cơ sở bảo mật hệ thống.

\textit{Chương 4: Triển khai đề tài}

Trình bày chi tiết quá trình hiện thực hóa thiết bị Node. Phần cứng gồm thiết kế sơ đồ nguyên lý và mạch in (PCB) chống nhiễu. Phần mềm bao gồm lập trình Driver, bộ lọc số, mã hóa AEAD, chế độ Hybrid Sleep, lưu trữ ngoại tuyến và công cụ kiểm thử EdgeProfiler.

\textit{Chương 5: Kiểm thử và đánh giá kết quả}

Đo lường định lượng thực nghiệm trên phần cứng và phần mềm. Kết quả về độ trễ xử lý ($<50$ ms), dòng tiêu thụ, tài nguyên RTOS và chất lượng LoRa đều xác nhận hệ thống đáp ứng hoàn toàn các chỉ tiêu kỹ thuật đã cam kết.

\textit{Chương 6: Kết luận và hướng phát triển}

Tổng kết các đóng góp khoa học, thực tiễn và đề xuất hướng mở rộng tương lai như áp dụng trí tuệ nhân tạo tại biên (TinyML), tối ưu năng lượng bằng Pulsed Heating, và triển khai kiến trúc mạng LoRaWAN.

\newpage

\section*{ABSTRACT}
\phantomsection\addcontentsline{toc}{section}{\numberline {}ABSTRACT}

In the project \textit{Research, Design, and Fabrication of an IoT System for Monitoring Poultry Microclimate Using LoRa and Edge Computing}, a distributed IoT system is proposed to address the "Cloud Paradox," optimizing control latency and power consumption. The project content includes the following chapters:

\textit{Chapter 1: Overview and Scientific Basis}

Presents challenges in intensive poultry farming and the limitations of cloud-based IoT solutions during network failures. It establishes the design objectives for an Edge-LoRa Node, specifying strict quantitative requirements for latency, power, and security.

\textit{Chapter 2: Theoretical Basis}

Systematizes the core theoretical foundations, including Wireless Sensor Networks (WSN), poultry physiology (THI index), Edge Computing architecture, LoRa physical layer (CSS), Real-Time Operating Systems (FreeRTOS), and the AES-128 encryption algorithm.

\textit{Chapter 3: System Analysis and Overall Design}

Justifies the selection of hardware technologies (ESP32-S3 SoC, sensors, power ICs). Furthermore, it designs the multitasking software architecture on FreeRTOS, Finite State Machine (FSM) algorithms, and the system's security mechanisms.

\textit{Chapter 4: Project Implementation}

Details the realization of the Edge Node. Hardware implementation covers the schematic design and anti-interference PCB layout. Software includes programming Drivers, digital filters, AEAD encryption, Hybrid Sleep mode, offline storage, and the EdgeProfiler testing tool.

\textit{Chapter 5: Testing and Result Evaluation}

Conducts quantitative empirical measurements on hardware and software. Test results regarding processing latency ($<50$ ms), power consumption, RTOS resource allocation, and LoRa transmission quality confirm that the system fully meets all committed specifications.

\textit{Chapter 6: Conclusion and Future Development}

Summarizes the project's scientific and practical contributions. It also proposes future research directions, such as applying Edge AI (TinyML), further energy optimization using Pulsed Heating, and deploying a LoRaWAN network architecture.

\cleardoublepage
\pagenumbering{arabic} % đánh số 1,2,3...

\newpage

\renewcommand{\figurename}{\fontsize{12pt}{0pt}\selectfont \bfseries Hình}
\renewcommand{\tablename}{\fontsize{12pt}{0pt}\selectfont \bfseries Bảng}
\section*{CHƯƠNG 1. TỔNG QUAN VỀ ĐỀ TÀI VÀ CƠ SỞ KHOA HỌC}
\addcontentsline{toc}{section}{\numberline{}CHƯƠNG 1. TỔNG QUAN VỀ ĐỀ TÀI VÀ CƠ SỞ KHOA HỌC}
\setcounter{section}{1}
\setcounter{figure}{0}
\setcounter{table}{0}
Chương 1 tập trung phân tích tính cấp thiết của việc giám sát vi khí hậu chuồng nuôi và giải quyết rủi ro từ "Nghịch lý đám mây" thông qua kiến trúc Điện toán biên thực thi trực tiếp trên Node. Đề tài xác lập mục tiêu thiết kế thiết bị Node ứng dụng vi điều khiển ESP32-S3 và công nghệ LoRa nhằm đạt được sự tối ưu về độ trễ xử lý, năng lượng tiêu thụ và bảo mật dữ liệu.
\subsection{Lý do chọn đề tài}
\subsubsection{Thực trạng ngành chăn nuôi và các thách thức từ môi trường vi khí hậu}
Trong xu hướng hiện đại hóa nông nghiệp thông minh (Smart Agriculture), việc chuyển dịch từ quy mô hộ gia đình nhỏ lẻ sang các mô hình trang trại tập trung khép kín đang diễn ra mạnh mẽ tại Việt Nam. Tuy nhiên, mật độ nuôi thả cao trong không gian kín đặt ra thách thức lớn đối với khâu kiểm soát các yếu tố vi khí hậu chuồng nuôi, yếu tố mang tính quyết định trực tiếp đến tỷ lệ sinh trưởng và hiệu quả kinh tế.
\begin{figure}[H]
        \centering
        \includegraphics[width=0.9\linewidth]{Thuc_Hanh/Image/ChanNuoiTapTrungGa.jpg}
        \caption{Mô hình chăn nuôi tập trung}
        \label{fig:placeholder}
    \end{figure}
Về mặt đặc tính sinh học, gia cầm là loài không có tuyến mồ hôi, năng lực điều hòa thân nhiệt kém và cơ chế tản nhiệt phụ thuộc hoàn toàn vào hệ thống hô hấp. Ở các giai đoạn tăng trưởng đầu đời như gà con, mức nhiệt độ lý tưởng trong lồng ấp yêu cầu duy trì ngặt nghèo trong khoảng từ 32°C đến 35°C. Nếu nhiệt độ môi trường vượt ngưỡng 35°C, vật nuôi sẽ rơi vào trạng thái stress nhiệt nguy hiểm. Khi mức nhiệt chạm ngưỡng giới hạn 38°C, hiện tượng suy tim diện rộng sẽ xảy ra, dẫn đến tỷ lệ tử vong hàng loạt trong thời gian ngắn.

Bên cạnh nhiệt độ, độ ẩm và nồng độ khí độc tích tụ cũng là những tác nhân gây thiệt hại nặng nề. 
\begin{figure}[H]
        \centering
        \includegraphics[width=1.0\linewidth]{Thuc_Hanh/Image/Chuong1_impactfactor.png}
        \caption{Một số thông số môi trường}
        \label{fig:placeholder}
    \end{figure}
Trong môi trường chuồng trại khép kín có độ ẩm cao vượt mức 75\%, các khối chất thải hữu cơ chưa qua xử lý phân hủy rất nhanh, giải phóng một lượng lớn khí Amoniac ($NH_3$) vào không khí. Khi nồng độ khí $NH_3$ vượt ngưỡng an toàn 25 ppm, niêm mạc đường hô hấp của vật nuôi sẽ bị tổn thương và bào mòn nghiêm trọng. Hiện tượng này làm suy giảm hệ miễn dịch tự nhiên, tạo điều kiện cho các chủng dịch bệnh nguy hiểm bùng phát, gây tổn thất nghiêm trọng cho người chăn nuôi. Do đó, việc giám sát liên tục, chính xác các tham số môi trường là yêu cầu bắt buộc trong quy trình vận hành trang trại công nghiệp.

\subsubsection{Hạn chế của các giải pháp hiện tại và "Nghịch lý đám mây"}
Đối với bài toán giám sát chuồng trại, các phương pháp thủ công truyền thống ngày càng bộc lộ sự lạc hậu: tiêu tốn nhiều nhân lực, khó bảo đảm tính liên tục và dễ mắc sai số. Trên thị trường hiện nay, dù đã xuất hiện một số hệ sinh thái nông nghiệp thông minh toàn diện như giải pháp ONE Farm \cite{7174976} của VNPT, nhưng các hệ thống thương mại này chủ yếu hướng đến phân khúc doanh nghiệp lớn do chi phí đầu tư ban đầu và vận hành tương đối cao.

Xét về mặt kiến trúc hệ thống, các giải pháp thương mại hiện tại bộc lộ những điểm yếu khi đối mặt với điều kiện thực tế tại các vùng nông thôn Việt Nam:
\begin{itemize}
    \item Chi phí duy trì hạ tầng viễn thông: Việc phụ thuộc hoàn toàn vào kết nối mạng di động 3G/4G hoặc hệ thống Wi-Fi diện rộng khiến các hộ chăn nuôi quy mô vừa và nhỏ phải gánh chịu chi phí vận hành cố định hàng tháng tương đối lớn.
    \item Độ ổn định của mạng lưới vô tuyến: Hạ tầng mạng viễn thông tại các vùng nông nghiệp xa trung tâm thường có độ phủ sóng không ổn định, dễ bị gián đoạn do thời tiết hoặc địa hình.
    \item Rủi ro từ "Nghịch lý đám mây" (Cloud Paradox): Việc cấu hình cho các node cảm biến liên tục truyền tải toàn bộ khối lượng dữ liệu thô lên máy chủ đám mây trung tâm để xử lý tính toán gây ra sự lãng phí nghiêm trọng về băng thông vô tuyến, tiêu tốn năng lượng của thiết bị biên và tiềm ẩn các rủi ro lớn về an toàn thông tin. Nguy hiểm hơn, khi xảy ra sự cố đứt cáp hoặc mất kết nối Internet hoàn toàn, toàn bộ luồng dữ liệu bị chặn đứng, khiến hệ thống mất khả năng phân tích và ra quyết định xử lý. Lúc này, các thiết bị chấp hành như quạt làm mát hay hệ thống phun sương không thể tự động kích hoạt, dẫn đến nguy cơ vật nuôi chết ngạt hàng loạt do stress nhiệt tăng vọt.
\end{itemize}
Từ thực tiễn đó, đề tài hướng tới việc nghiên cứu và chế tạo một hệ thống thiết bị nội bộ tự chủ, kết hợp công nghệ truyền thông không dây tầm xa LoRa để loại bỏ chi phí cước viễn thông , đồng thời ứng dụng năng lực Điện toán biên (Edge Computing) xử lý dữ liệu trực tiếp tại lớp Gateway phân cấp. Mô hình này bảo đảm khả năng vận hành độc lập: ngay cả khi toàn bộ kết nối Internet bên ngoài bị gián đoạn, hệ thống vẫn duy trì năng lực tính toán và trực tiếp kích hoạt các kịch bản cứu hộ vật nuôi kịp thời ngay tại chuồng trại.

% --- Bổ sung phần RTOS ---
\subsubsection{Thách thức về giới hạn tài nguyên của Nút biên và bài toán tối ưu hóa RTOS}
Việc dồn toàn bộ năng lực tính toán xuống thiết bị Node (mô hình \textit{Local Binary Offloading}) tuy giải quyết được "Nghịch lý đám mây" nhưng lại đẩy chính Node vào một bài toán kỹ thuật hóc búa: phải gánh đồng thời nhiều luồng công việc dị thể trên một nền tảng tài nguyên cực kỳ eo hẹp. Vi điều khiển ESP32-S3 chỉ sở hữu khoảng $512\,\text{KB}$ SRAM nội và hai lõi xử lý chạy bằng pin, trong khi phải xử lý song song hàng loạt tác vụ có đặc tính thời gian trái ngược nhau: lấy mẫu cảm biến liên tục qua bus I2C/ADC/UART, lọc nhiễu (EMA, trung vị, Kalman), chạy máy trạng thái hữu hạn (FSM) ra quyết định an toàn với ràng buộc tất định dưới $50\,\text{ms}$, mã hóa có xác thực (AEAD), thu phát LoRa với các pha chờ SPI mang tính "chặn" (blocking), ghi log ngoại tuyến và giám sát watchdog. Các tác vụ này cạnh tranh gay gắt ba loại tài nguyên hữu hạn: thời gian CPU, bus ngoại vi dùng chung (đặc biệt là bus I2C) và bộ nhớ RAM.

Mô hình lập trình vòng lặp tuần tự (super-loop) truyền thống không thể bảo đảm tính thời gian thực khi một tác vụ chậm hoặc bị treo (ví dụ nghẽn bus I2C) chiếm dụng CPU \cite{white2011making}. Do đó, thách thức chính về mặt phần mềm của đề tài là \textbf{điều phối đa nhiệm một cách tất định trong giới hạn ngặt nghèo về bộ nhớ và năng lượng}, và biện pháp là xây dựng kiến trúc firmware xoay quanh một Hệ điều hành thời gian thực (RTOS). Từ đó, bên cạnh các mục tiêu về chức năng, đề tài đặt thêm một nhóm mục tiêu tối ưu hóa RTOS có thể kiểm chứng định lượng:
\begin{itemize}
    \item \textbf{An toàn bộ nhớ ngăn xếp (Stack):} cấp phát Stack cho từng tác vụ sao cho mức dùng đỉnh (High-Water Mark) luôn còn dư an toàn, tuyệt đối không xảy ra tràn ngăn xếp (stack overflow) gây sập hệ thống.
    \item \textbf{Ổn định vùng nhớ động (Heap):} dung lượng Heap nội còn trống phải ổn định theo thời gian (không rò rỉ -- memory leak, không phân mảnh nghiêm trọng), duy trì trên một ngưỡng dự phòng an toàn.
    \item \textbf{Dư địa tải CPU:} phân bổ tác vụ và mức ưu tiên giữa hai lõi sao cho tải CPU mỗi lõi không bão hòa, bảo đảm luôn còn dư địa đáp ứng tức thời sự kiện khẩn cấp.
    \item \textbf{Tiết kiệm năng lượng ở mức RTOS:} khai thác cơ chế \textit{Tickless Idle} của FreeRTOS kết hợp Deep/Light Sleep để hạ tối đa năng lượng tiêu thụ trong thời gian rảnh, kéo dài tuổi thọ pin.
    \item \textbf{Chặn nghịch đảo ưu tiên:} bảo vệ tài nguyên dùng chung (bus I2C) bằng Mutex có kế thừa ưu tiên (priority inheritance) để giới hạn thời gian chặn của tác vụ ưu tiên cao.
\end{itemize}
Các chỉ tiêu RTOS nêu trên sẽ được thiết kế ở Chương 3, hiện thực ở Chương 4 và kiểm chứng bằng số liệu đo thực tế ở Chương 5.
% --- Hết bổ sung phần RTOS ---

\subsection{Mục tiêu nghiên cứu và các chỉ tiêu kỹ thuật lượng hóa}
Mục đích cốt lõi của đề tài là nghiên cứu, thiết kế và chế tạo thực nghiệm một hệ thống IoT giám sát vi khí hậu đạt chuẩn công nghiệp, vận hành ổn định trong môi trường chuồng trại khắc nghiệt, có khả năng tính toán tự chủ và bảo mật cao. Để hệ thống thực sự mang lại giá trị ứng dụng thực tiễn, việc thiết kế không chỉ dừng lại ở mức độ mô hình thử nghiệm (prototype) mà phải đáp ứng các ràng buộc khắt khe về mặt kỹ thuật cũng như đặc tính sinh học.

% Thứ nhất, đối với đặc tính sinh lý nhạy cảm của gia cầm, tốc độ phản hồi của hệ thống khi nhiệt độ hoặc nồng độ khí độc tăng vọt là yếu tố sống còn. Do đó, yêu cầu về độ trễ xử lý sự cố được đặt lên hàng đầu, đòi hỏi hệ thống phải đưa ra quyết định điều khiển thiết bị làm mát gần như tức thời ngay tại lớp mạng cục bộ mà không cần phụ thuộc hay chờ đợi dữ liệu khứ hồi từ máy chủ đám mây.

% Thứ hai, trong điều kiện các mô hình chăn nuôi quy mô lớn, khu vực chuồng trại thường phân bố rải rác trên diện tích rộng và gặp nhiều hạn chế về việc kéo hạ tầng điện lưới đến từng điểm đo. Điều này đặt ra bài toán bắt buộc phải sử dụng nguồn năng lượng độc lập (pin sạc) kết hợp với công nghệ truyền thông vô tuyến tầm xa, có khả năng xuyên tín hiệu tốt qua các rào cản vật lý như tường gạch hay vách tôn. Việc tính toán thiết kế phần cứng để tối ưu hóa dòng điện rò ở chế độ ngủ sâu (Deep Sleep) là điều kiện tiên quyết giúp thiết bị bám trụ lâu dài ngoài thực địa, giảm thiểu tối đa chi phí bảo trì định kỳ cho người nông dân.

% Cuối cùng, trong một mạng lưới truyền thông không dây mở như LoRa, các rủi ro về việc dữ liệu bị can thiệp hoặc kẻ gian sử dụng phương thức phát lại gói tin (Replay Attack) để kích hoạt sai rơ-le có thể gây ra những thiệt hại nghiêm trọng cho toàn bộ lứa nuôi. Vì vậy, việc áp dụng các cơ chế mật mã học để bảo vệ tính toàn vẹn và xác thực của đường truyền là một mục tiêu không thể thiếu trong khâu thiết kế phần mềm.

Để phục vụ công tác kiểm thử, đo lường thực nghiệm và làm cơ sở đối chiếu vững chắc khi nghiệm thu đánh giá đồ án, toàn bộ các mục tiêu thiết kế nêu trên đã được quy đổi và lượng hóa chi tiết qua hệ thống các thông số kỹ thuật trình bày tại bảng \ref{Bảng 1.1}:
\begin{table}[H]
\centering
\caption{Bảng chỉ tiêu kỹ thuật lượng hóa của hệ thống}
\label{Bảng 1.1}
\vspace{2mm}
\small

\begin{tabularx}{\textwidth}{>{\raggedright\arraybackslash\bfseries}p{3.2cm} >{\raggedright\arraybackslash\itshape}p{3.5cm} X}
\toprule
\textbf{Phân hệ chức năng} & \textbf{Tham số thiết kế} & \textbf{Chỉ tiêu kỹ thuật } \\
\midrule

\multirow{3}{=}{Cơ khí \& Phần cứng} 
    & Quy cách mạch in & Custom PCB FR4 2 lớp, phủ sơn chống oxy hóa. \\ \cmidrule{2-3}
    & Nguồn cấp & Pin Li-ion 18650 (3200\textit{mAh}), sạc qua cổng USB Type-C. \\ \cmidrule{2-3}
    & Vỏ hộp bảo vệ & Vỏ PETG in 3D ($110 \times 85 \times 75$ \textit{mm}), chuẩn kháng nước/bụi IP65, tích hợp màng lọc PTFE. \\
\midrule

\multirow{3}{=}{Hiệu năng LoRa RF} 
    & Tham số vô tuyến & Băng tần 433 \textit{MHz}; SF = 7; BW = 125 \textit{kHz}; CR = 4/5. \\ \cmidrule{2-3}
    & Cự ly truyền dẫn & 500\textit{m} - 1000\textit{m} trong điều kiện không vật cản (LOS). \\ \cmidrule{2-3}
    & Tỷ lệ suy hao & PLR $<$ 5\% tại khoảng cách 200\textit{m} có vật cản (NLOS). \\
\midrule

\multirow{3}{=}{Điện toán biên \& Bảo mật} 
    & Độ trễ xử lý biên & $<$ 50 \textit{ms} (từ lúc phát hiện sự cố đến khi đóng Rơ-le). \\ \cmidrule{2-3}
    & Cơ chế mật mã & Mã hóa AES-128 và ASCON, kết hợp xác thực MAC và kiểm tra CRC. \\ \cmidrule{2-3}
    & Chống Replay Attack & Khớp nhãn thời gian (Timestamp) từ RTC với cửa sổ $\Delta t < 5s$. \\
\midrule

\multirow{2}{=}{Năng lượng \& Điện năng} 
    & Dòng Deep Sleep & $<$ 20 $\mu$\textit{A}. \\ \cmidrule{2-3}
    & Tuổi thọ pin & $>$ 3 tháng (với chu kỳ lấy mẫu 5 phút/lần). \\
\midrule

% Nền tảng Cloud & Đồng bộ dữ liệu & Giao tiếp Firebase Realtime Database, độ trễ hiển thị $<$ 5\textit{s}, hỗ trợ điều khiển Downlink. \\
% \midrule

% Sản phẩm học thuật & Tài liệu bàn giao & Báo cáo đồ án và 01 bài báo khoa học (tiếng Anh, định dạng chuẩn IEEE). \\
% \bottomrule

\end{tabularx}
\end{table}
\subsection{Đối tượng và phạm vi nghiên cứu}
\subsubsection{Đối tượng nghiên cứu:}
\begin{itemize}
    \item Vi điều khiển đa nhân ESP32-S3 và năng lực xử lý Edge Computing.
    \item Giao thức truyền thông LoRa vật lý (IC SX1278).
    \item Hệ điều hành FreeRTOS điều phối đa nhiệm trên Node.
    \item Thuật toán Điện toán biên thực thi trên Node (Tiền xử lý, lọc nhiễu, máy trạng thái ra quyết định).
    \item Mật mã học bảo mật: AES-128, ASCON, MAC và chống Replay Attack.
\end{itemize}
\subsubsection{Phạm vi nghiên cứu:}
\begin{itemize}
    \item Phần cứng: Tập trung thiết kế, chế tạo PCB và vỏ hộp IP65 cho duy nhất thiết bị Node cảm biến. Khối Gateway sử dụng thiết bị có sẵn để làm trạm trung chuyển dữ liệu lên Cloud.
    \item Phần mềm: Phát triển Firmware cho Node đảm nhiệm 100\% logic tính toán biên. Xây dựng cơ chế Hybrid Deep Sleep (Ngủ định kỳ kết hợp thức dậy bằng sự kiện ngắt ngoài khi có sự cố khẩn cấp).
\end{itemize}
% \subsection{Yêu cầu của đề tài}
% Để đáp ứng mục tiêu nghiên cứu và giải quyết triệt để bài toán thực tiễn tại các trang trại, thiết bị Node phải thỏa mãn các yêu cầu khắt khe được lượng hóa bằng các tham số cụ thể sau:
% \subsubsection{Yêu cầu chức năng}
% \begin{itemize}
%     \item Thu thập và giám sát vi khí hậu: Node phải quét và lấy mẫu liên tục dữ liệu từ tổ hợp cảm biến với sai số đo lường dưới 5\%. Cụ thể, hệ thống phải theo dõi sát sao biểu đồ nhiệt - ẩm và nhận diện ngay lập tức khi môi trường chạm các ngưỡng nguy hiểm: nhiệt độ vượt 35°C (gây stress nhiệt) , độ ẩm trên 75\% , và nồng độ khí độc Amoniac (NH3) vượt ngưỡng 25 ppm.
%     \item Điện toán biên (Edge Computing): Tích hợp Máy trạng thái hữu hạn (FSM) tại vi điều khiển để tự động đối chiếu dữ liệu thô với ma trận ngưỡng an toàn. Khi phát hiện các tham số vượt ngưỡng (ví dụ: nhiệt độ > 35°C), thuật toán biên phải tự động kích hoạt ngõ ra GPIO để đóng Rơ-le điều khiển hệ thống làm mát cục bộ mà không cần chờ lệnh từ Gateway hay Cloud.
%     \item Truyền thông vô tuyến LoRa: Đóng gói dữ liệu Payload (định mức 32-bytes)  và phát sóng định kỳ. Cấu hình phần cứng phải chốt ở tần số 433 MHz, Spreading Factor (SF) = 7, Bandwidth (BW) = 125 kHz và Coding Rate = 4/5 để cân bằng giữa băng thông và tầm xa.
%     \item Lưu trữ ngoại tuyến (Offline Storage): Xây dựng cơ chế phát hiện rớt mạng cục bộ. Khi không nhận được gói tin phản hồi (ACK) từ Gateway, Node phải tự động đẩy dữ liệu log vào vùng nhớ tĩnh (Flash/PSRAM) và đồng bộ lại ngay khi sóng vô tuyến được khôi phục.
%     \item Bảo mật và tính toàn vẹn dữ liệu: Mã hóa toàn bộ khối dữ liệu tải trọng bằng thuật toán đối xứng AES-128. Đồng thời, tích hợp IC thời gian thực (RTC DS3231) để đính kèm nhãn thời gian (Timestamp) vào từng gói tin , thiết lập cửa sổ thời gian chặn Replay Attack với độ lệch $\Delta{t}$ < 5s.
% \end{itemize}
% \subsubsection{Yêu cầu phi chức năng}
% \begin{itemize}
%     \item Độ trễ xử lý Điện toán biên (Edge Latency): Thời gian đáp ứng của luồng điều khiển (tính từ thời điểm cảm biến xuất tín hiệu báo vượt ngưỡng đến khi chân GPIO đóng mạch Rơ-le) bắt buộc phải nhỏ hơn 50 mili-giây.
%     \item Hiệu suất năng lượng và Tuổi thọ pin (Power Profiling): Vận hành trên nguồn pin Li-ion 18650 (dung lượng 3200mAh). Hệ thống phải quản lý dòng tiêu thụ thông qua chế độ Deep Sleep, duy trì dòng rò ở trạng thái chờ dưới 20µA. Với chu kỳ lấy mẫu 5 phút/lần, Node phải có khả năng hoạt động liên tục hơn 3 tháng mà không cần sạc lại.
%     \item Chất lượng đường truyền (Link Quality): Phủ sóng vô tuyến ổn định trong bán kính 500m đến 1000m ở điều kiện không vật cản (Line-of-Sight). Trong môi trường chuồng trại có tường bê tông che khuất (Non-Line-of-Sight), Tỷ lệ mất gói tin (Packet Loss Rate - PLR) phải được duy trì dưới 5% ở mốc khoảng cách 200m.
%     \item An toàn chức năng (Functional Safety): Hệ thống phải hoạt động ổn định 24/7. Yêu cầu bắt buộc phải cấu hình Hardware Watchdog Timer để tự động reset vi điều khiển nếu xảy ra hiện tượng treo bus I2C hoặc rò rỉ vùng nhớ heap (Memory Leak).
% \end{itemize}
% \subsubsection{Yêu cầu kỹ thuật và thiết kế phần cứng}
% \begin{itemize}
%     \item Mạch nguyên lý và khối xử lý cốt lõi: Sử dụng vi điều khiển có năng lực xử lý đa nhân làm khối trung tâm, kết hợp với module truyền thông vô tuyến tầm xa dải tần thấp (Sub-GHz) thông qua giao thức truyền thông nối tiếp tiêu chuẩn. Khối quản lý năng lượng yêu cầu tích hợp mạch sạc và linh kiện ổn áp tuyến tính (LDO) có dòng tĩnh (Quiescent current) cực thấp. Đặc biệt, nhằm tối đa hóa hiệu năng tiết kiệm điện, phần cứng bắt buộc phải tích hợp cấu trúc công tắc tải (Load Switch) sử dụng linh kiện bán dẫn công suất, cho phép vi điều khiển chủ động ngắt hoàn toàn nguồn điện cấp cho toàn bộ tổ hợp cảm biến ngoại vi khi hệ thống tiến vào trạng thái ngủ sâu (Deep Sleep).
%     \item Thiết kế mạch in (Custom PCB): Thiết kế bo mạch in tiêu chuẩn công nghiệp 2 lớp. Quá trình quy hoạch đường mạch (Layout) phải tuân thủ nghiêm ngặt các quy tắc chống nhiễu điện từ (EMI). Cụ thể, cần thiết lập các vùng cách ly (Keep-out zone) không phủ đồng tại khu vực ăng-ten vô tuyến để tránh suy hao phát xạ, đồng thời phân tách vật lý rõ ràng giữa vùng tín hiệu số (Digital) tốc độ cao và đường tín hiệu tương tự (Analog) nhạy cảm để triệt tiêu nhiễu chéo. Bo mạch sau khi hoàn thiện gia công và hàn linh kiện dán (SMD) phải được phủ lớp bảo vệ chuyên dụng để chống ô xi hóa trong môi trường độ ẩm cao và nhiều khí ăn mòn.
% \end{itemize}
\subsection{Phương pháp nghiên cứu}
Để hiện thực hóa các yêu cầu khắt khe tại thiết bị Node từ khâu thiết kế đến đánh giá nghiệm thu, đề tài áp dụng tổ hợp các phương pháp nghiên cứu chuyên sâu bao gồm:
\subsubsection{Phương pháp nghiên cứu lý thuyết và cơ sở toán học}
\begin{itemize}
    \item Phân tích đặc tính sinh học: Nghiên cứu các tài liệu chăn nuôi gia cầm để trích xuất biểu đồ nhiệt - ẩm theo ngày tuổi. Từ đó, lập bảng ma trận trạng thái (State Matrix) làm cơ sở logic nền tảng cho thuật toán Điện toán biên (Edge Computing).
    \item Phân tích kỹ thuật vô tuyến: Dựa trên cơ sở toán học của kỹ thuật điều chế Chirp Spread Spectrum (CSS). Sử dụng công thức tính Quỹ đường truyền (Link Budget) để tính toán lý thuyết sự đánh đổi giữa Spreading Factor (SF), Băng thông (BW) và Khoảng cách truyền sóng.
    \item Nghiên cứu mật mã học: Đọc hiểu tài liệu kỹ thuật của thư viện mbedtls trên nền tảng ESP-IDF. Nắm vững cấu trúc dịch vòng (ShiftRows) và hoán vị (SubBytes) của thuật toán mã hóa đối xứng AES-128 để đảm bảo tối ưu hóa, không gây tràn bộ nhớ RAM khi chạy đa luồng trên vi điều khiển.
\end{itemize}
\subsubsection{Phương pháp đồng thiết kế và mô hình hóa trên máy tính (HW/SW Co-design)}
Đồ án áp dụng chu trình thiết kế song song để đảm bảo tính tương thích tuyệt đối giữa phần cứng IP65 và Firmware xử lý biên. 
\begin{itemize}
    \item Thiết kế Cơ khí (CAD): Sử dụng phần mềm thiết kế 3D để dựng mô hình 3D vỏ hộp. Tính toán buồng khí động học: xác định vị trí đặt các khe hở thông gió và màng lọc PTFE sao cho không khí lưu thông tối ưu bề mặt cảm biến SHT31 và MQ-135, đồng thời đạt chuẩn kháng nước IP65.
    \item Thiết kế Phần cứng (Hardware Layout): Sử dụng Kicad để mô hình hóa sơ đồ nguyên lý (Schematic) và quy hoạch Custom PCB 2 lớp. Ứng dụng phương pháp tính toán dải đồng để phối hợp trở kháng $50 \Omega$ (Impedance Matching) cho đường dẫn ăng-ten vô tuyến, đồng thời thiết lập các vùng cách ly nhiễu (Keep-out zone).
    \item Thiết kế Phần mềm: Sử dụng phương pháp vẽ lưu đồ Máy trạng thái hữu hạn (FSM - Finite State Machine) để quy hoạch kiến trúc phân luồng chạy song song trên FreeRTOS, đảm bảo tính đồng bộ giữa tác vụ thu thập, mã hóa và ra quyết định biên.
\end{itemize}
\subsubsection{Phương pháp chế tạo, thực nghiệm và đo lường định lượng}
\begin{itemize}
    \item Gia công và Lắp ráp (Assembly): Đặt gia công bản mạch in, tự thực hiện quy trình đóng chip bằng máy khò nhiệt (SMD Soldering). Lắp ráp vật lý bo mạch vào vỏ nhựa PETG in 3D kết hợp bơm keo chống ẩm và ép gioăng cao su silicone.
    \item Đo lường định lượng (Quantitative Measurement): Sử dụng các thiết bị đo lường chuẩn công nghiệp để lấy số liệu thực tế tại Node
    \item Phân tích thống kê và Trực quan hóa dữ liệu: Thiết lập các bài test thực địa khắc nghiệt để thu thập hàng ngàn dòng Log Data (bao gồm: RSSI, SNR, Thời gian gửi/nhận). Sử dụng các thuật toán để tiền xử lý dữ liệu, làm sạch dữ liệu nhiễu, tính toán tỷ lệ mất gói (Packet Loss Rate - PLR), chi phí thời gian mã hóa (Encryption Overhead Time).
\end{itemize}
\subsection{Kết luận chương}
Chương 1 đã làm rõ tính cấp thiết của đề tài xuất phát từ hai áp lực
song song: đặc tính sinh lý mỏng manh của gia cầm đòi hỏi kiểm soát vi
khí hậu liên tục và tức thời, trong khi các giải pháp thương mại hiện
tại lại bộc lộ điểm yếu chết người khi mất kết nối Internet -- toàn bộ
luồng điều khiển tê liệt, vật nuôi không được bảo vệ.

Từ đó, đề tài xác lập hướng tiếp cận rõ ràng: đưa năng lực tính toán
xuống trực tiếp thiết bị Node thông qua kiến trúc Điện toán biên với
chiến lược Offloading nhị phân, kết hợp công nghệ LoRa để loại bỏ sự
phụ thuộc vào hạ tầng viễn thông thương mại. Các chỉ tiêu kỹ thuật được
lượng hóa cụ thể, không chỉ là mục tiêu thiết kế
mà còn là thước đo nghiệm thu định lượng xuyên suốt các chương tiếp
theo.
\newpage


\section*{CHƯƠNG 2. CƠ SỞ LÝ THUYẾT}
\addcontentsline{toc}{section}{\numberline {} CHƯƠNG 2. CƠ SỞ LÝ THUYẾT}
\setcounter{section}{2}
\setcounter{subsection}{0}
\setcounter{figure}{0}
\setcounter{table}{0}

Chương này tập trung hệ thống hóa các cơ sở lý thuyết cốt lõi phục vụ cho việc thiết kế hệ thống, bao gồm việc phân tích đặc điểm sinh học của đối tượng giám sát và làm rõ nguyên lý toán học của các công nghệ trọng yếu như truyền thông LoRa, giao thức IoT và kỹ thuật bảo mật dữ liệu

\subsection{Tổng quan về Mạng cảm biến không dây và Vi khí hậu chuồng nuôi}
\subsubsection{Cấu trúc và đặc điểm của Mạng cảm biến không dây (WSN)}
Mạng cảm biến không dây (Wireless Sensor Network - WSN) là một hệ thống phân tán bao gồm hàng chục đến hàng ngàn thiết bị cảm biến nhỏ gọn (gọi là Node) được nhúng sâu vào môi trường cần giám sát. Trong bối cảnh nông nghiệp chính xác (Precision Agriculture), WSN đóng vai trò như hệ thần kinh ngoại biên, thu thập các thông số vật lý và hóa học để cung cấp dữ liệu đầu vào cho hệ thống ra quyết định.
\begin{figure}[H]
        \centering
        \includegraphics[width=0.9\linewidth]{Thuc_Hanh/Image/WSN.png}
        \caption{Kiến trúc Mạng cảm biến không dây (WSN)}
        \label{Hình 2.1}
    \end{figure}
Kiến trúc hệ thống trong Hình \ref{Hình 2.1} được chia thành 4 lớp rõ rệt, đi từ cấp độ thiết bị hiện trường lên đến máy chủ trung tâm:
\begin{itemize}
    \item Lớp Thu thập dữ liệu: Đây là các thiết bị đầu cuối được phân bố rải rác trực tiếp tại các "Khu vực giám sát". Nhiệm vụ chính của chúng là đo đạc và thu thập thông số môi trường cục bộ. 
    \item Lớp Quản lý cục bộ: Thay vì tất cả các nút đều cố gắng truyền dữ liệu đi xa, các nút cảm biến sẽ gửi dữ liệu tầm ngắn về một "Nút quản lý vùng" ở trung tâm cụm. Trạm quản lý này đóng vai trò gom dữ liệu . .
    \item Lớp Truyền dẫn tầm xa: Sau khi được tổng hợp ở Nút quản lý vùng, dữ liệu sẽ được bắn về các trạm thu phát sóng lớn hơn (Trung tâm điều hành).
    \item Lớp Giám sát trung tâm: Các Trung tâm điều hành khu vực sẽ đóng vai trò làm cầu nối (Gateway) cuối cùng, chuyển tiếp gói tin từ mạng vô tuyến nội bộ lên môi trường Internet (qua Wi-Fi, Ethernet hoặc 4G/LTE) và đưa về máy chủ tổng. Tại đây, quản trị viên có thể theo dõi và điều khiển toàn bộ các vùng giám sát khác nhau trên một giao diện hợp nhất.
\end{itemize}
Cấu trúc điển hình của một Node cảm biến bao gồm bốn khối chức năng chính: Khối cảm biến (Sensor Unit) thực hiện việc chuyển đổi đại lượng vật lý thành tín hiệu điện; Khối xử lý trung tâm (Processing Unit) số hóa và xử lý dữ liệu; Khối truyền thông (Transceiver Unit) thực hiện việc phát xạ tín hiệu vô tuyến; và Khối năng lượng (Power Unit) cung cấp dòng điện duy trì hoạt động. Cụ thể về cấu trúc của Node Sensor được thể hiện như Hình \ref{Hình 2.2}.
 \begin{figure}[H]
        \centering
        \includegraphics[width=0.9\linewidth]{Thuc_Hanh/Image/NodeStructure.pdf}
        \caption{Cấu trúc của một node cảm biến}
        \label{Hình 2.2}
    \end{figure}
 Điểm khác biệt lớn nhất của WSN trong môi trường chăn nuôi so với môi trường viễn thông thông thường là tính chất khắc nghiệt của không gian vật lý: độ ẩm bão hòa cao, khí ăn mòn và nhiễu sóng từ các vách ngăn bê tông, đòi hỏi phần cứng phải có độ bền và tính kháng nhiễu vượt trội.

\subsubsection{Phân tích sinh lý học gia cầm và cơ chế tản nhiệt}
Việc kiểm soát vi khí hậu trong chuồng nuôi không đơn thuần là bài toán tự động hóa mà bắt nguồn sâu xa từ đặc tính sinh lý của đối tượng vật nuôi. Gia cầm là động vật hằng nhiệt, cần duy trì nhiệt độ cơ thể ở mức $41^\circ C$ đến $42^\circ C$. Tuy nhiên, đặc điểm sinh học chí mạng của loài này là hoàn toàn không có tuyến mồ hôi. Lớp lông vũ dày bao phủ cơ thể lại hoạt động như một lớp cách nhiệt hoàn hảo, ngăn cản quá trình tỏa nhiệt qua bề mặt da.

Do đó, gia cầm giải phóng lượng nhiệt dư thừa chủ yếu thông qua đường hô hấp (bức xạ và bay hơi nước qua niêm mạc phổi). Khi nhiệt độ môi trường tăng cao tiệm cận với nhiệt độ cơ thể, sự chênh lệch nhiệt độ (Gradient nhiệt) giảm xuống, làm mất đi khả năng tản nhiệt tự nhiên. Vật nuôi buộc phải thở dốc để tăng tốc độ bay hơi. Quá trình thở dốc này tiêu tốn một lượng năng lượng khổng lồ, làm biến đổi độ pH trong máu (gây kiềm hô hấp) và dẫn đến suy kiệt hệ thống tuần hoàn.

\subsubsection{Chỉ số nhiệt ẩm (THI) và hiện tượng Stress nhiệt}
Trong môi trường chuồng trại khép kín, nhiệt độ không hoạt động độc lập mà có sự tương tác cộng hưởng với độ ẩm tương đối (Relative Humidity - RH). Sự kết hợp này được định lượng hóa thông qua Chỉ số nhiệt ẩm (Temperature Humidity Index - THI). Độ ẩm môi trường càng cao, áp suất hơi nước trong không khí càng lớn, làm ức chế quá trình bay hơi nước từ hệ hô hấp của gia cầm, khiến vật nuôi cảm thấy nóng hơn nhiều so với nhiệt độ thực tế đo được bằng nhiệt kế.

Công thức tính toán chỉ số THI \cite{thom1959discomfort, lin2006thermal} được ứng dụng rộng rãi trong các nghiên cứu nông nghiệp được biểu diễn bằng phương trình:
\begin{equation}
    {THI = 0.8 \times T_{db} + \frac{RH}{100} \times (T_{db} - 14.4) + 46.4}.
\end{equation}
Trong đó:
\begin{itemize}
    \item $T_{db}$: Nhiệt độ bầu khô (Dry-bulb temperature) tính bằng độ C.
    \item $RH$: Độ ẩm tương đối tính bằng phần trăm (\%).
\end{itemize}

Dựa trên ma trận THI, gia cầm sẽ rơi vào trạng thái "Stress nhiệt nhẹ" khi THI vượt mức 72, "Stress nhiệt cấp tính" ở mức 78, và có nguy cơ suy tim diện rộng, chết hàng loạt khi THI vượt ngưỡng 82. Việc chỉ sử dụng cảm biến nhiệt độ đơn lẻ là không đủ cơ sở để hệ thống máy tính đưa ra quyết định đóng cắt quạt thông gió, mà bắt buộc phải thiết lập một hàm nội suy kết hợp cả nhiệt độ và độ ẩm ngay tại lớp xử lý biên.
% \subsubsection{Động học phân hủy chất thải và cơ chế sinh khí Amoniac ($NH_3$)}
% Bên cạnh bài toán về nhiệt - ẩm, môi trường chăn nuôi mật độ cao sinh ra một lượng lớn khí độc, đặc biệt là Amoniac ($NH_3$). Quá trình sinh khí $NH_3$ không xuất phát trực tiếp từ vật nuôi mà là kết quả của các phản ứng hóa sinh phức tạp diễn ra tại lớp đệm lót chuồng.

% Trong phân gia cầm có chứa một lượng lớn axit uric ($C_5H_4N_4O_3$) và urea không tiêu hóa hết. Dưới xúc tác của enzyme uricase và urease từ vi sinh vật có trong môi trường, kết hợp với điều kiện độ ẩm cao (RH $> 75\%$) và nhiệt độ ấm, quá trình thủy phân diễn ra mạnh mẽ\cite{ni2000ammonia}:
% \begin{equation}
%     C_5H_4N_4O_3 + 1.5O_2 + 4H_2O \rightarrow 5CO_2 + 4NH_3
% \end{equation}

% $NH_3$ là chất khí không màu, nhẹ hơn không khí nhưng lại có ái lực cực mạnh với nước. Khi khí này bốc lên từ nền chuồng, chúng hòa tan ngay lập tức vào lớp dịch nhầy trên niêm mạc mắt và đường hô hấp của gia cầm, tạo thành dung dịch amoni hydroxit ($NH_4OH$) có tính kiềm mạnh, gây bỏng rát, viêm loét giác mạc và phá hủy lớp lông mao bảo vệ khí quản.

\subsubsection{Tác động độc lý học của vi khí hậu lên vật nuôi}
Theo các tiêu chuẩn chăn nuôi quốc tế, ngưỡng an toàn tối đa của $NH_3$ trong chuồng nuôi không được vượt quá 25 ppm. Ở nồng độ 50 ppm, tỷ lệ tăng trọng của vật nuôi giảm 15-20\% do giảm lượng thức ăn ăn vào và tổn thương hệ tiêu hóa. Ở mức 100 ppm, nguy cơ tử vong tăng đột biến do viêm phổi cục bộ. Sự thay đổi của các tham số này diễn ra rất nhanh, có thể tính bằng phút khi hệ thống thông gió gặp sự cố. Điều này đặt ra yêu cầu thiết kế hệ thống IoT không chỉ dừng lại ở việc "đọc số liệu", mà phải là "phản ứng tức thời" - nền tảng để đưa Điện toán biên vào mô hình thiết kế.
% \subsection{Khảo sát một số nghiên cứu liên quan} % (muc trong, trung lap; phan khao sat day du da o cuoi chuong 2.6)
\subsection{Điện toán biên (Edge Computing) và Kiến trúc IoT phân tán}
Sự bùng nổ của Internet vạn vật (IoT) khiến một lượng dữ liệu khổng lồ (Big Data) được sinh ra từng giây từ vô số thiết bị ngoại vi. Trong hơn một thập kỷ qua, mô hình Điện toán đám mây (Cloud Computing) luôn được xem là xương sống của mọi hệ thống IoT, nơi đám mây đóng vai trò là trung tâm xử lý tập trung, cung cấp năng lực tính toán và không gian lưu trữ gần như vô hạn. Tuy nhiên, khi các ứng dụng IoT tiến sâu vào các lĩnh vực đòi hỏi tính thời gian thực khắt khe và độ tin cậy tuyệt đối như y tế, xe tự hành, tự động hóa công nghiệp và nông nghiệp chính xác, mô hình tập trung truyền thống bắt đầu bộc lộ những giới hạn vật lý và logic không thể vượt qua. Từ bối cảnh đó, Điện toán biên (Edge Computing) ra đời như một sự dịch chuyển tất yếu của kiến trúc mạng, đưa năng lực xử lý, lưu trữ và ra quyết định rời khỏi trung tâm dữ liệu xa xôi để tiến về sát nhất với nguồn phát sinh dữ liệu.
\subsubsection{Nghịch lý đám mây (Cloud Paradox) trong các hệ thống IoT truyền thống}
Trong mô hình kiến trúc IoT truyền thống dựa trên đám mây, các thiết bị đầu cuối (End-nodes) thường được thiết kế dưới dạng các "nút thụ động" (Dummy nodes). Chức năng duy nhất của chúng là đo lường các đại lượng vật lý (như nhiệt độ, độ ẩm, nồng độ khí gas), chuyển đổi thành tín hiệu số và truyền tải toàn bộ dữ liệu thô (Raw Data) qua mạng viễn thông lên máy chủ đám mây. Đám mây sẽ tiếp nhận, chạy các thuật toán phân tích, so sánh với các ngưỡng an toàn và sau đó gửi ngược lại các lệnh điều khiển (Downlink) về các cơ cấu chấp hành (Actuators) tại hiện trường. Mô hình này, mặc dù tận dụng được sức mạnh tính toán vô song của máy chủ, lại vấp phải một rào cản kỹ thuật nghiêm trọng được giới học thuật định danh là "Nghịch lý đám mây" (Cloud Paradox).

Nghịch lý đám mây xoay quanh bốn thách thức cốt lõi không thể khắc phục triệt để nếu chỉ dựa vào việc nâng cấp hạ tầng mạng: độ trễ truyền dẫn, băng thông nghẽn cổ chai, chi phí duy trì năng lượng và rủi ro điểm lỗi đơn lẻ (Single Point of Failure).

Thứ nhất, về độ trễ truyền dẫn (Propagation Latency). Độ trễ toàn trình (Round-Trip Time - RTT) của một hệ thống IoT tập trung được cấu thành từ nhiều thành phần phức tạp, có thể được biểu diễn qua phương trình toán học cơ bản\cite{kurose2017computer}:
\begin{equation}
    {RTT = 2 \times (T_{trans} + T_{prop} + T_{queue} + T_{proc\_net}) + T_{proc\_cloud}}.
\end{equation}
Trong đó, $T_{trans}$ là thời gian phát dữ liệu lên kênh truyền, $T_{prop}$ là thời gian lan truyền tín hiệu điện từ qua các môi trường vật lý (cáp đồng, cáp quang, không gian tự do), $T_{queue}$ là thời gian gói tin phải nằm chờ tại các bộ đệm của bộ định tuyến (Router) trên đường đi, $T_{proc\_net}$ là thời gian xử lý tiêu đề gói tin (Header processing) tại các trạm chuyển tiếp, và $T_{proc\_cloud}$ là thời gian CPU của máy chủ đám mây thực thi thuật toán. 
Đối với các trang trại nằm ở vùng sâu vùng xa, dữ liệu phải đi qua hàng chục trạm trung chuyển trước khi tới được Data Center đặt tại các thành phố lớn hoặc quốc gia khác. Sự biến động của $T_{queue}$ do nghẽn mạng có thể đẩy RTT lên đến hàng giây, hoặc thậm chí hàng chục giây (Jitter). Đối với một hệ thống giám sát vi khí hậu, sự chậm trễ này là một lỗ hổng chí mạng. Khi nồng độ khí độc NH3 tăng vọt, nếu hệ thống thông gió phải chờ đợi lệnh từ Cloud mất vài giây, lượng khí độc tích tụ đã đủ để làm tổn thương vĩnh viễn niêm mạc hô hấp của hàng ngàn cá thể gia cầm.

Thứ hai, giới hạn về băng thông và năng lượng. Việc liên tục phát sóng vô tuyến để đẩy hàng ngàn bản ghi dữ liệu thô (ví dụ: chu kỳ 1 giây/lần) lên mạng 3G/4G tiêu tốn một lượng năng lượng khổng lồ. Năng lượng tiêu thụ cho quá trình phát sóng vô tuyến (RF Transmission) luôn lớn hơn gấp hàng trăm lần so với năng lượng dành cho việc tính toán nội bộ tại vi điều khiển. Việc lạm dụng đám mây làm cạn kiệt nguồn pin của các thiết bị Node một cách nhanh chóng, đi ngược lại triết lý tiêu thụ năng lượng thấp của Mạng cảm biến không dây (WSN).

Thứ ba, rủi ro về điểm lỗi đơn lẻ. Hệ thống IoT phụ thuộc hoàn toàn vào đám mây là một hệ thống thiếu tính kiên cường (Resilience). Nếu đường kết nối Internet (Backhaul link) từ trang trại lên đám mây bị đứt gãy do thời tiết, thiên tai hoặc sự cố nhà mạng, toàn bộ luồng điều khiển sẽ bị tê liệt hoàn toàn. Mặc dù các cảm biến nội bộ vẫn đang đo được mức nhiệt độ gây tử vong là $40^\circ C$, nhưng do không có kết nối Internet để báo cáo lên Cloud, các Node không thể tự nhận thức được sự nguy hiểm, và các quạt thông gió vẫn nằm im lìm. Rủi ro này biến mô hình đám mây trở thành một canh bạc rủi ro cao đối với các ngành công nghiệp trọng yếu.

\subsubsection{Sự dịch chuyển kiến trúc: Từ Cloud Computing đến Edge Computing}
Để giải quyết triệt để "Nghịch lý đám mây", các nhà nghiên cứu kiến trúc mạng đã đề xuất mô hình Điện toán biên (Edge Computing). Khái niệm này không nhằm mục đích thay thế hoàn toàn Cloud, mà tạo ra một sự phân bổ khối lượng công việc hợp lý dọc theo trục topo mạng từ biên đến trung tâm.

Điện toán biên định nghĩa việc đẩy các năng lực tính toán logic, lưu trữ dữ liệu ngắn hạn, phân tích thời gian thực và phân phối ứng dụng ra khu vực rìa của mạng lưới \cite{shi2016edge}. Rìa ở đây có thể là các Gateway công nghiệp, các trạm gốc viễn thông (Base stations), hoặc sâu hơn nữa là trực tiếp tại các vi điều khiển (MCU) của các Node cảm biến cuối cùng. Triết lý của Điện toán biên là: Dữ liệu được sinh ra ở đâu, hãy xử lý nó ngay tại đó, hoặc càng gần đó càng tốt, trước khi quyết định có cần đẩy nó qua mạng diện rộng hay không.

Bảng \ref{Bảng 2.1} so sánh trên cung cấp một cái nhìn đối chiếu toàn diện giữa hai mô hình kiến trúc, làm nổi bật sự đánh đổi (Trade-off) mang tính tất yếu trong thiết kế hệ thống. 

\begin{table}[H]
\centering
\caption{Bảng phân tích so sánh các đặc trưng cốt lõi giữa Điện toán đám mây và Điện toán biên}
\label{Bảng 2.1}
\renewcommand{\arraystretch}{1.5}
\begin{tabular}{p{3.5cm} p{5.5cm} p{5.5cm}}
\hline
\textbf{Tiêu chí đánh giá} & \textbf{Điện toán Đám mây (Cloud Computing)} & \textbf{Điện toán Biên (Edge Computing)} \\ \hline
\textbf{Vị trí xử lý dữ liệu} & Tập trung tại các trung tâm dữ liệu quy mô lớn (Data Centers) nằm xa nguồn dữ liệu. & Phân tán tại các thiết bị ngoại vi (Nodes, Gateways) ngay sát hiện trường vật lý. \\ 
\textbf{Độ trễ toàn trình (Latency)} & Rất cao và khó dự đoán (thường $> 100$ ms đến vài giây) do phụ thuộc vào định tuyến WAN. & Cực kỳ thấp, mang tính tất định (Deterministic), thường ở dải $< 10$ ms. \\ 
\textbf{Sự phụ thuộc Internet} & Phụ thuộc 100\%. Đứt mạng đồng nghĩa với việc mất hoàn toàn khả năng điều khiển. & Hoạt động độc lập (Offline-capable). Vẫn duy trì các vòng lặp an toàn cục bộ khi đứt mạng. \\ 
\textbf{Tiêu thụ băng thông} & Khổng lồ do phải đẩy toàn bộ dữ liệu thô (Raw Data) liên tục. & Rất thấp do chỉ truyền đi dữ liệu đã qua tổng hợp hoặc các bản tin cảnh báo (Metadata). \\ 
\textbf{Năng lực tính toán \& Lưu trữ} & Gần như vô hạn (Sử dụng các cụm Server, GPU, TPU khổng lồ). & Hạn chế nghiêm ngặt bởi vi điều khiển, dung lượng RAM/Flash và nguồn pin cục bộ. \\ 
\textbf{Tính bảo mật \& Quyền riêng tư} & Dữ liệu nhạy cảm phải truyền qua mạng công cộng, rủi ro bị đánh chặn cao. & Dữ liệu nhạy cảm được xử lý và lưu giữ ngay tại mạng nội bộ, giảm thiểu bề mặt tấn công. \\ \hline
\end{tabular}
\end{table}
Thiết bị biên được lập trình các máy trạng thái (State Machines) an toàn dự phòng, đảm bảo tính mạng vật nuôi không bị đe dọa khi cáp quang đứt. Tuy nhiên, sự ưu việt này phải đánh đổi bằng giới hạn về năng lực tính toán. Trong khi Cloud có thể dễ dàng chạy các mô hình Học sâu (Deep Learning) hàng tỷ tham số, thì Edge Devices (đặc biệt là các MCU như ESP32) chỉ sở hữu bộ nhớ RAM tính bằng Kilobyte và không có hệ thống làm mát chủ động. Do đó, các thuật toán Điện toán biên bắt buộc phải được tinh gọn hóa, sử dụng các cấu trúc dữ liệu nhẹ, tránh cấp phát bộ nhớ động (Dynamic memory allocation) để chống phân mảnh, và ưu tiên xử lý bằng các toán tử thao tác bit (Bitwise operations) để giảm chu kỳ máy. Cuối cùng, về mặt bảo mật, việc không đưa dữ liệu thô rời khỏi mạng nội bộ giúp thu hẹp đáng kể bề mặt tấn công (Attack Surface), biến mô hình Edge trở thành lựa chọn hàng đầu cho việc thiết kế kiến trúc chuẩn công nghiệp và an toàn thông tin.
\subsubsection{Phân lớp kiến trúc hệ thống Edge IoT}
Một hệ thống IoT phân tán tích hợp hoàn chỉnh Điện toán biên không phải là một mặt phẳng đồng nhất, mà được cấu trúc thành một hệ thống phân cấp (Hierarchical Architecture) gồm ba lớp phân định rõ ràng như Hình \ref{Hình 2.3}. Mỗi lớp đảm nhiệm một vai trò chuyên biệt, giải quyết các bài toán khác nhau về thời gian đáp ứng và năng lực xử lý.
\begin{figure}[H]
        \centering        \includegraphics[width=0.5\linewidth]{Thuc_Hanh/Image/3TierIoT.png}
        \caption{Kiến trúc phân tầng trong IoT}
        \label{Hình 2.3}
    \end{figure}
Sơ đồ kiến trúc phân lớp thể hiện một hình chóp đặc trưng của luồng dữ liệu và năng lực tính toán. Ở phần đáy rộng nhất là Lớp biên cực trị (Extreme Edge Layer), bao gồm hàng loạt các thiết bị vi điều khiển (MCU), cảm biến đa thông số và rơ-le chấp hành. Lớp này tiếp xúc trực tiếp với thế giới vật lý, có nhiệm vụ "lấy mẫu" môi trường thực tế (Continuous sampling). Đặc tính của lớp này là có số lượng thiết bị vô cùng lớn, nhưng năng lực tính toán trên mỗi thiết bị lại rất hạn hẹp (thường là vi xử lý lõi đơn hoặc lõi kép tần số thấp, chạy bằng pin). Mọi sự kiện thời gian thực cực đoan nhất (như ngắt khẩn cấp khi quá nhiệt độ) bắt buộc phải được khép kín vòng lặp kiểm soát ngay tại lớp này để đảm bảo độ trễ bằng phần nghìn giây.

Tiến lên phần giữa của tháp là Lớp mạng cục bộ / Lớp sương mù (Fog/Network Layer). Lớp này được đại diện bởi các Gateway công nghiệp, các Router Edge hoặc các máy tính nhúng mạnh mẽ hơn (như Raspberry Pi hoặc máy tính công nghiệp IPC). Gateway đóng vai trò là một "nhà ga trung chuyển" và là người phiên dịch giao thức (Protocol Translator). Nó thu thập sóng vô tuyến tầm ngắn/tầm trung (như LoRa, Zigbee) từ vô số Node ở lớp đáy, tổng hợp dữ liệu (Data Aggregation), mã hóa lại thành các giao thức mạng lõi (như MQTT qua TCP/IP) và phân phối luồng đi. Tại lớp này, các tính toán biên ở mức độ vừa (Middle-edge computing) có thể được thực thi, chẳng hạn như tính toán trung bình cộng dữ liệu của toàn bộ 50 Node trong trang trại để đưa ra quyết định bật quạt tổng, trong khi các Node chỉ điều khiển quạt cục bộ của riêng nó.

Đỉnh của hình tháp là Lớp đám mây (Cloud Layer). Được giải phóng hoàn toàn khỏi gánh nặng phải điều khiển trực tiếp các Rơ-le ở cấp độ mili-giây, Cloud giờ đây có thể tập trung vào thế mạnh cốt lõi của mình: Xử lý dữ liệu lớn (Big Data Processing) theo mẻ (Batch processing). Máy chủ đám mây sẽ lưu trữ dữ liệu lịch sử hàng năm của toàn bộ chuỗi trang trại, thực thi các thuật toán khai phá dữ liệu (Data Mining), huấn luyện các mô hình AI dự báo xu hướng dịch bệnh dựa trên biểu đồ vi khí hậu dài hạn, và cung cấp các giao diện giám sát tổng thể (Dashboards) cho người quản trị ở bất kỳ đâu trên thế giới thông qua kết nối Internet.
\subsubsection{Các chiến lược chia tải tác vụ tại biên (Computation Offloading)}
Một trong những bài toán khoa học hóc búa nhất của lĩnh vực Điện toán biên là bài toán Chia tải tác vụ (Computation Offloading). Khi một ứng dụng phần mềm bao gồm nhiều khối logic khác nhau, nhà thiết kế hệ thống phải quyết định toán học xem khối logic nào sẽ được chạy trên bộ vi xử lý yếu ớt tại Node, khối logic nào sẽ được đóng gói gửi qua sóng vô tuyến để chạy trên Gateway, và khối nào sẽ bị đẩy lên tận Cloud. Quyết định này không mang tính cảm tính mà dựa trên việc tối ưu hóa một hàm mục tiêu (Objective Function) liên quan đến tiêu thụ năng lượng và độ trễ.

Xét một tác vụ $i$ cần được thực thi, nó có thể được định nghĩa bởi một bộ ba tham số $T_i = \{D_i, C_i, T_{max\_i}\}$. Trong đó:
\begin{itemize}
    \item $D_i$: Kích thước dữ liệu đầu vào của tác vụ (tính bằng bit hoặc byte).
    \item $C_i$: Khối lượng tính toán yêu cầu để hoàn thành tác vụ (tính bằng số chu kỳ lệnh vi xử lý - CPU cycles).
    \item $T_{max\_i}$: Ràng buộc thời gian tối đa cho phép để hoàn thành tác vụ (Độ trễ tối đa).
\end{itemize}
Hệ thống có hai lựa chọn chiến lược để thực thi tác vụ này. Lựa chọn thứ nhất là Thực thi cục bộ (Local Execution) ngay tại vi điều khiển của Node \cite{mao2017survey} . Thời gian thực thi cục bộ ($t_{local}$) và năng lượng tiêu thụ cục bộ ($E_{ocal}$) được xác định bởi công thức\cite{mach2017mobile}:
\begin{equation}
    {t_{local} = \frac{C_i}{f_{mcu}}}.
\end{equation}
\begin{equation}
    {E_{local} = P_{active} \times t_{local} = \kappa \cdot f_{mcu}^2 \cdot C_i}.
\end{equation}
Trong đó, $f_{mcu}$ là xung nhịp hoạt động của vi điều khiển, $P_{active}$ là công suất tiêu thụ của CPU ở trạng thái hoạt động toàn tải, và $\kappa$ là hệ số điện dung chuyển mạch hiệu dụng tùy thuộc vào cấu trúc bán dẫn của chip.

Lựa chọn thứ hai là Chia tải sang thiết bị khác (Remote Offloading), ví dụ như gửi toàn bộ dữ liệu $D_i$ qua mạng LoRa lên Gateway để Gateway tính toán hộ. Lúc này, tổng thời gian ($t_{offload}$) và tổng năng lượng ($E_{offload}$) \cite{mach2017mobile}đối với Node sẽ là:
\begin{equation}
    t_{offload} = \frac{D_i}{R_{uplink}} + \frac{C_i}{f_{gateway}} + \frac{D_{out}}{R_{downlink}}.
\end{equation}
\begin{equation}
    E_{offload} = P_{tx} \times \frac{D_i}{R_{uplink}} + P_{rx} \times \frac{D_{out}}{R_{downlink}}.
\end{equation}
Với $R_{uplink}$ và $R_{downlink}$ là tốc độ truyền dữ liệu băng thông lên và xuống của mạng LoRa, $P_{tx}$ và $P_{rx}$ là công suất phát và công suất thu của module vô tuyến, $f_{gateway}$ là xung nhịp của CPU tại Gateway (thường lớn hơn rất nhiều so với Node).

Dựa trên các phương trình trên, bộ điều phối sẽ so sánh. Nếu $E_{local} < E_{offload}$ và $t_{local} < T_{max\_i}$, tác vụ sẽ được giữ lại chạy 100\% tại Node. Trong hệ thống giám sát chuồng trại, thuật toán so sánh ngưỡng nhiệt độ (FSM) có khối lượng tính toán $C_i$ vô cùng nhỏ (chỉ vài phép toán so sánh IF-ELSE đơn giản), trong khi kích thước dữ liệu $D_i$ đẩy đi không hề nhỏ, cộng thêm tốc độ vô tuyến LoRa ($R_{uplink}$) rất chậm (chỉ vài trăm bit/giây). Do đó, hàm $t_{offload}$ luôn lớn hơn rất nhiều so với $t_{local}$, và năng lượng phát sóng LoRa $E_{offload}$ cũng đắt đỏ hơn nhiều so với việc tính toán nội bộ. Hệ quả tất yếu về mặt kiến trúc là hệ thống phải áp dụng mô hình Chia tải cục bộ nhị phân (Local Binary Offloading) cho luồng điều khiển, nghĩa là Node sẽ ôm trọn 100\% nhiệm vụ tính toán nhận diện bất thường và ra quyết định đóng Rơ-le, không san sẻ với Gateway để đảm bảo thời gian cực nhanh và tiết kiệm pin tối đa.
\subsubsection{Phương pháp luận đồng thiết kế Phần cứng/Phần mềm (HW/SW Co-design) cho thiết bị biên}
Trong kỹ thuật nhúng truyền thống, phần cứng thường được thiết kế trước, sau đó các kỹ sư phần mềm mới viết mã để cố gắng nhồi nhét vào cấu trúc đã được cố định đó. Tuy nhiên, đối với Edge IoT, cách tiếp cận này sẽ dẫn đến sự lãng phí tài nguyên vô ích hoặc hiệu năng không tốt. Thay vào đó, phương pháp luận Đồng thiết kế Phần cứng/Phần mềm (Hardware/Software Co-design) được áp dụng.
\begin{figure}[H]
    \centering
    \includegraphics[width=1.0\linewidth]{Thuc_Hanh/Image/HWSWCoDesign_UML.pdf}
    \caption{Quy trình đồng thiết kế phần cứng và phần mềm}
    \label{Hình 2.4}
\end{figure}
Hình \ref{Hình 2.4} trên phác họa chu trình phát triển lặp (Iterative process) của phương pháp đồng thiết kế. Điểm khởi đầu của quá trình không phải là chọn chip hay viết code, mà là định lượng hóa các Ràng buộc hệ thống (System Constraints) mang tính sống còn, ví dụ như độ trễ hệ thống không được vượt quá 50ms và thiết bị phải trụ được 3 tháng trên một viên pin 18650. Từ đây, luồng phát triển rẽ nhánh thành hai quá trình chạy song song nhưng tương tác liên tục: Luồng phần cứng (HW) và Luồng phần mềm (SW).

Ở nhánh Thiết kế phần cứng (Hardware Design Flow), quá trình bắt đầu từ việc phân tích kiến trúc CPU/MCU nhằm lựa chọn nền tảng xử lý phù hợp, tiếp theo là thiết kế cơ chế quản lý năng lượng để đáp ứng yêu cầu tiêu thụ điện năng thấp, cuối cùng là cấu hình bộ nhớ ngoài nhằm bảo đảm đủ không gian lưu trữ chương trình, dữ liệu và hỗ trợ các chức năng mở rộng của hệ thống. Song song với đó, nhánh Thiết kế phần mềm (Software Design Flow) tiến hành phân hoạch hệ thống thành các máy trạng thái hữu hạn (FSM Partitioning), cấu hình hệ điều hành thời gian thực FreeRTOS để tổ chức đa nhiệm, sau đó tối ưu hóa bộ nhớ nhằm giảm dung lượng RAM/Flash sử dụng và nâng cao hiệu quả thực thi.

Điểm hội tụ của hai nhánh là bước Mô phỏng đồng thời (Co-Simulation), nơi mô hình phần cứng và phần mềm được đánh giá như một hệ thống thống nhất thay vì kiểm thử riêng lẻ. Sau quá trình mô phỏng, hệ thống được so sánh với các chỉ tiêu hiệu năng đã đặt ra, bao gồm thời gian đáp ứng, mức tiêu thụ năng lượng, khả năng sử dụng bộ nhớ và độ ổn định vận hành. Nếu bất kỳ tiêu chí nào chưa đạt, quy trình sẽ quay trở lại giai đoạn thiết kế, chủ yếu ở nhánh phần mềm để điều chỉnh phương pháp phân hoạch FSM, cấu hình FreeRTOS hoặc tối ưu thuật toán; khi cần thiết, các quyết định ở nhánh phần cứng như lựa chọn kiến trúc MCU, cấu hình bộ nhớ hay cơ chế quản lý năng lượng cũng sẽ được xem xét lại. Chu trình lặp này tiếp tục cho đến khi toàn bộ mục tiêu thiết kế được thỏa mãn, khi đó dự án chuyển sang giai đoạn chế tạo nguyên mẫu phần cứng và nạp chương trình (PCB & Firmware Prototyping) để kiểm chứng trên hệ thống thực tế trước khi tiến hành các bước đánh giá và hoàn thiện sản phẩm.

\subsection{Công nghệ truyền thông vô tuyến LoRa và Phân tích quỹ đường truyền}
Nền tảng để Điện toán biên phân tán có thể hoạt động hiệu quả là một hạ tầng truyền thông vô tuyến đáng tin cậy, đóng vai trò như hệ tuần hoàn kết nối các Node cô lập với trạm Gateway trung tâm. Trong điều kiện chăn nuôi công nghiệp khắc nghiệt, nơi các tín hiệu bị chặn đứng bởi tường bê tông dày, mái tôn kim loại và chịu sự can nhiễu từ các hệ thống điện cơ công suất lớn, các chuẩn vô tuyến truyền thống như WiFi hay Bluetooth hoàn toàn bất lực. Để giải quyết bài toán giao tiếp vô tuyến tầm xa nhưng lại tiêu thụ năng lượng siêu thấp, công nghệ LoRa (Long Range) đã nổi lên như một chuẩn mực công nghiệp không thể thay thế.
\subsubsection{Các công nghệ mạng diện rộng công suất thấp (LPWAN)}
Thế giới mạng không dây IoT được phân chia thành các phổ công nghệ dựa trên hai trục tham số tỷ lệ nghịch với nhau: Tầm xa phủ sóng (Coverage Range) và Tốc độ truyền tải dữ liệu (Data Rate). Ở phổ tốc độ cao nhưng tầm gần, chúng ta có các công nghệ Mạng cục bộ không dây (WLAN/WPAN) như WiFi, Bluetooth Low Energy (BLE), Zigbee, thích hợp cho nhà thông minh nhưng không thể vươn ra các trang trại rộng hàng hecta. Ở phổ tốc độ cao, tầm xa, chúng ta có mạng thông tin di động Cellular (3G/4G/5G). Tuy nhiên, cái giá phải trả cho băng thông di động là năng lượng tiêu thụ cực kỳ lớn và thiết kế phần cứng phức tạp, khiến tuổi thọ pin của thiết bị chỉ đếm bằng ngày hoặc giờ.

Để lấp đầy khoảng trống chiến lược ở phổ thiết bị cần tầm phủ sóng rộng lớn (vài kilomet) nhưng chỉ cần gửi các gói tin dữ liệu nhỏ (vài chục byte) với năng lượng cục pin hoạt động kéo dài nhiều năm, một lớp mạng mới đã được định nghĩa: Mạng diện rộng công suất thấp (Low-Power Wide-Area Network - LPWAN). Trong nhóm LPWAN, ba công nghệ đang dẫn dắt thị trường toàn cầu bao gồm NB-IoT, Sigfox và LoRa, mỗi công nghệ theo đuổi một triết lý thiết kế lớp vật lý hoàn toàn trái ngược nhau.

NB-IoT (Narrowband IoT) là một tiêu chuẩn do tổ chức 3GPP phát triển, hoạt động trên các băng tần phổ được cấp phép (Licensed bands) của mạng LTE hiện hữu. Kỹ thuật này cung cấp chất lượng dịch vụ (QoS) rất cao và độ trễ thấp do được bảo vệ khỏi nhiễu đồng kênh, nhưng yêu cầu doanh nghiệp phải trả phí thuê bao SIM định kỳ và phụ thuộc tuyệt đối vào độ phủ sóng trạm gốc của nhà mạng viễn thông. Sigfox tiếp cận theo hướng hoàn toàn khác, sử dụng kỹ thuật truyền dải siêu hẹp (Ultra-Narrowband - UNB) với tốc độ rùa bò chỉ 100 bps để đạt được khả năng xuyên thấu xa nhất và tiết kiệm pin nhất. Tuy nhiên, kiến trúc của Sigfox bị hạn chế nghiêm trọng về tính năng liên lạc hai chiều (Bi-directional), số lượng bản tin đẩy xuống Node (Downlink) bị giới hạn ngặt nghèo (thường chỉ 4 tin/ngày), khiến nó vô dụng trong việc thực thi các lệnh điều khiển hồi tiếp.

Vượt qua những nhược điểm trên, LoRa (được cấp bằng sáng chế bởi Semtech) sử dụng kỹ thuật trải phổ trên băng tần ISM không cần cấp phép (Unlicensed bands như 433MHz tại Châu Á, 868MHz tại Châu Âu). Sự xuất sắc của LoRa nằm ở việc nó cho phép các kỹ sư tự do xây dựng các mạng cục bộ nội bộ (Private Networks) mà không phải trả bất kỳ khoản phí viễn thông nào. Hơn nữa, nó hỗ trợ giao tiếp hai chiều hoàn hảo, cân bằng tuyệt vời giữa chi phí năng lượng và khả năng xuyên thấu tường chắn, trở thành ứng cử viên duy nhất đáp ứng hoàn hảo cho cấu trúc mạng của dự án.

\subsubsection{Lớp vật lý của LoRa và Kỹ thuật điều chế Chirp Spread Spectrum (CSS) \cite{augustin2016study}}
Sức mạnh vô tuyến của LoRa không nằm ở việc tăng công suất phát (do thiết bị pin không cho phép), mà nằm ở một kỹ thuật điều chế tinh vi ở Lớp vật lý (Physical Layer) có tên là Chirp Spread Spectrum (CSS). Nguồn gốc của CSS xuất phát từ các hệ thống radar quân sự và viễn thông vũ trụ sâu, nơi tín hiệu phải vượt qua mức độ nhiễu nền cực kỳ khủng khiếp.

Khác với kỹ thuật điều chế FSK (Frequency Shift Keying) thông thường (trong đó mức logic 0 và 1 được biểu diễn bởi hai tần số tĩnh rời rạc), CSS biểu diễn dữ liệu bằng cách sử dụng các xung tần số quét tuyến tính liên tục theo thời gian. Một xung như vậy được gọi là "Chirp". Tần số tức thời của một chirp thay đổi liên tục trên toàn bộ băng thông (Bandwidth - BW) được cấp phát. Phương trình tần số tức thời $f(t)$ của một xung Up-chirp (tần số tăng dần) được mô tả bằng hàm toán học\cite{vangelista2015frequency}:
\begin{equation}
    {f(t) = f_0 + \left( \frac{BW}{T} \right) \times t}.
\end{equation}
Trong đó, $f_0$ là tần số bắt đầu ở mép dưới của băng thông, $BW$ là độ rộng dải băng thông (ví dụ 125 kHz), $T$ là khoảng thời gian kéo dài của một ký hiệu (Symbol duration), và tỉ số $\left( \frac{BW}{T} \right)$ chính là tốc độ quét tần số (Sweep rate).

Bằng cách mã hóa dữ liệu vào thời điểm bắt đầu nhảy tần (Cyclic shift) của các xung chirp, tín hiệu chứa thông tin hữu ích bị làm "nhiễu hóa" và trải rộng năng lượng của nó ra toàn bộ băng thông rộng. Hệ quả thần kỳ của kỹ thuật này tuân theo Định lý Shannon-Hartley nổi tiếng về dung lượng kênh truyền \cite{shannon1948mathematical}:
\begin{equation}
    {C = BW \times \log_2 \left( 1 + \frac{S}{N} \right)}.
\end{equation}
Định lý này khẳng định rằng, để truyền được một tốc độ thông tin $C$ không đổi trong điều kiện Tỷ số tín hiệu trên nhiễu (Signal-to-Noise Ratio - $\frac{S}{N}$) cực kỳ thấp (tín hiệu yếu hơn cả nhiễu nền), cách duy nhất là phải mở rộng băng thông $BW$ lên cực kỳ lớn. Máy thu LoRa, nhờ biết trước cấu trúc toán học của các xung chirp, sẽ sử dụng một bộ lọc phối hợp (Matched Filter) tích phân tín hiệu theo thời gian để cộng hưởng lại năng lượng, từ đó khôi phục lại dữ liệu gốc ngay cả khi sóng vô tuyến đo được yếu hơn mức nhiễu nền của môi trường tự nhiên lên tới $20 \text{ dB}$ (nghĩa là tín hiệu chìm sâu dưới nhiễu gấp 100 lần về mặt công suất).
\begin{figure}[H]
    \centering
    \includegraphics[width=0.7\linewidth]{Thuc_Hanh/Image/spectrumLora.png}
    \caption{Up-Chirp và Down-Chirp trong LoRa}
    \label{Hình 2.5}
\end{figure}
Đồ thị phổ tần - thời gian (Spectrogram) như trong Hình \ref{Hình 2.5} là một minh họa trực quan sinh động nhất về cách sóng LoRa khiêu vũ trong không gian điện từ. Trục hoành của đồ thị đại diện cho thời gian ($t$), trong khi trục tung đại diện cho dải tần số hoạt động ($f$). Nhìn vào đồ thị, ta không thấy các điểm tần số rời rạc như công nghệ Bluetooth hay WiFi, mà thấy các đường thẳng chéo liên tục. Một đường thẳng chéo đi lên từ góc dưới bên trái (tần số thấp) lên góc trên bên phải (tần số cao) là một "Up-chirp" nguyên thủy. Độ dốc của đường thẳng này chính là tốc độ quét tần số $\frac{BW}{T}$.

Trong chuẩn LoRa, mỗi phần tử dữ liệu (Symbol) truyền đi không đơn thuần là một bit 0 hoặc 1, mà nó chứa $SF$ bit thông tin. Thông tin này được nhúng vào tín hiệu bằng cách thay đổi điểm xuất phát tần số của đường chéo. Khi đường chéo này chạm đến mép viền trên của băng thông, nó sẽ bị cắt đứt và "vòng" (wrap-around) lại mép viền dưới để tiếp tục quét lên. Sự gián đoạn đột ngột (nhảy bậc) về tần số này chứa đựng giá trị số học của biểu tượng cần truyền. Tại bộ thu Gateway, các mạch nhân tần số sẽ lấy tín hiệu nhận được nhân với một Down-chirp (đường chéo đi xuống) chuẩn cục bộ. Nếu hai xung này đồng bộ, kết quả của phép nhân sẽ tạo ra một xung Dirac duy nhất (một đỉnh nhọn năng lượng) trên miền tần số (thông qua phép biến đổi Fast Fourier Transform - FFT). 
\subsubsection{Các tham số vô tuyến cốt lõi:}
Hiệu năng phát sóng của thiết bị biên không phải là một hằng số cố định, mà có thể được lập trình điều khiển linh hoạt thông qua ba tham số vật lý nền tảng. Sự tinh chỉnh ba tham số này tạo ra một bài toán tối ưu đa biến cực kỳ phức tạp để cân bằng giữa Tầm xa phủ sóng, Thời lượng pin và Tốc độ truyền tải.
\paragraph{Hệ số trải phổ (Spreading Factor - SF):}
SF là tham số linh hồn của LoRa, quy định số lượng xung chirp quét được sử dụng để đại diện cho một biểu tượng (Symbol) thông tin duy nhất. Mối quan hệ toán học được biểu diễn qua công thức: Một Symbol sẽ chứa đựng $SF$ bit dữ liệu và được hình thành từ $2^{SF}$ xung chip. Giá trị SF được chuẩn hóa giới hạn từ SF7 đến SF12.

Khi ta tăng SF lên một đơn vị (ví dụ từ SF7 lên SF8), số lượng chip tạo ra một biểu tượng nhân đôi, đồng nghĩa với việc Thời gian phát sóng (Time on Air - ToA) bị kéo dài gấp đôi. Việc phát sóng chậm lại cho phép bộ thu tích lũy được nhiều năng lượng tín hiệu hơn, làm cải thiện Độ nhạy máy thu (Receiver Sensitivity) thêm trung bình khoảng $2.5 \text{ dB}$, giúp thiết bị có khả năng phát xuyên tường tốt hơn nhiều lần. Tuy nhiên, cái giá phải trả là Tốc độ bit (Bit Rate) giảm thảm hại và lượng điện năng rút ra từ viên pin tăng theo cấp số nhân trong một chu kỳ gửi. Do đó, SF7 thường được cấu hình cho các Node gần Gateway (môi trường Line-of-Sight), trong khi SF10 hoặc SF12 buộc phải sử dụng cho các Node bị khuất sâu bên trong các góc tường bê tông dày.
\paragraph{Băng thông (Bandwidth - BW):}
Băng thông xác định độ rộng của dải phổ mà tín hiệu chirp quét qua. Trong băng tần phổ biến 433MHz, BW thường được thiết lập ở 3 mức: 125 kHz, 250 kHz và 500 kHz. Thời gian để phát một biểu tượng (Symbol Duration - $T_s$) tỷ lệ thuận với hệ số trải phổ và tỷ lệ nghịch với Băng thông\cite{adelantado2017understanding}:
\begin{equation}
    {T_s = \frac{2^{SF}}{BW}}.
\end{equation}
Từ phương trình trên, nếu ta tăng $BW$ gấp đôi, thời gian truyền $T_s$ giảm đi một nửa, tốc độ mạng tăng nhanh. Nhưng theo định lý Shannon, băng thông càng rộng thì lượng nhiễu nhiệt (Thermal noise) lọt vào kênh truyền càng lớn, làm suy giảm nghiêm trọng tỷ số SNR, kéo theo việc giảm tầm xa truyền dẫn. Để tối đa hóa khả năng đâm xuyên cho các ứng dụng đo lường không yêu cầu tốc độ cao, BW 125 kHz luôn là lựa chọn tiêu chuẩn.
\paragraph{Tỷ lệ mã hóa lỗi (Coding Rate - CR):}
Đây là tham số đại diện cho kỹ thuật Sửa lỗi tiến (Forward Error Correction - FEC). Trong môi trường công nghiệp đầy sóng điện từ, các bit dữ liệu gửi đi rất dễ bị lật trạng thái (bit flip) do nhiễu. Thay vì yêu cầu thiết bị gửi lại toàn bộ gói tin khi bị lỗi (tốn pin vô ích), kỹ thuật FEC tự động chèn thêm các bit kiểm tra dư thừa (Redundant bits) vào dòng dữ liệu gốc dựa trên mã Hamming. Tham số $CR$ được tính bằng $\frac{4}{4+n}$, với $n \in \{1, 2, 3, 4\}$. Ví dụ, $CR = \frac{4}{5}$ nghĩa là cứ 4 bit dữ liệu thực tế, chip xử lý vô tuyến sẽ kẹp thêm 1 bit dự phòng để tính toán phát hiện sai lệch. Nếu $CR = \frac{4}{8}$, tỷ lệ dữ liệu thừa lên tới 50\%, khả năng tự hồi phục gói tin lỗi ở phía máy thu là cao nhất, nhưng độ dài chuỗi bit truyền đi tăng gấp đôi, làm giảm thời lượng pin của Node.
\subsubsection{Mô hình toán học đánh giá Quỹ đường truyền (Link Budget) và Độ nhạy máy thu}
Để đảm bảo hệ thống vô tuyến hoạt động xuyên suốt qua các tầng chuồng trại mà không bị mất tín hiệu (Packet Loss), các kỹ sư phải xây dựng một mô hình đánh giá năng lực đường truyền thông qua bài toán Quỹ đường truyền (Link Budget). Link Budget là tổng thể lượng hao hụt suy giảm tối đa (tính bằng Decibel - dB) mà tín hiệu có thể chịu đựng được trên đường bay từ anten phát đến anten thu, mà tại đó máy thu vẫn đủ sức giải mã được thông tin.

Phương trình cân bằng Quỹ đường truyền tổng quát được xây dựng:
\begin{equation}
    {P_{RX} = P_{TX} + G_{TX} - L_{TX} - L_{FS} - L_{Misc} + G_{RX} - L_{RX}}.
\end{equation}
Cụ thể, các tham số kỹ thuật như sau:
\begin{itemize}
    \item $P_{RX}$: Công suất tín hiệu tối thiểu nhận được tại khối khuếch đại LNA của máy thu (dBm). Đây chính là ngưỡng sinh tử; nếu $P_{RX}$ tụt xuống dưới Độ nhạy máy thu ($S$), gói tin sẽ chìm hẳn vào hư vô.
    \item $P_{TX}$: Công suất phát xạ RF do vi điều khiển cấu hình bơm vào đường mạch (ví dụ $+14 \text{ dBm}$ hoặc mức tối đa $+20 \text{ dBm}$). Phát công suất càng cao, pin càng mau cạn kiệt.
    \item $G_{TX}, G_{RX}$: Hệ số khuếch đại (Độ lợi) của anten ở bên phát và bên thu (dBi). Trong dự án này, anten đẳng hướng lò xo thường có mức tăng ích khoảng $2.0 \text{ dBi}$ đến $3.0 \text{ dBi}$.
    \item $L_{TX}, L_{RX}$: Tổn hao tĩnh trên vật liệu bảng mạch in (PCB) ở dạng đường dẫn sóng cao tần, tổn hao do bất đồng kháng (Impedance mismatch) tại cổng hàn SMA.
    \item $L_{Misc}$: Tổng hợp các hệ số tổn hao tạp khác như tán xạ do mưa, sự suy hao hấp thụ do đi qua vật cản là thân nhiệt hoặc hơi nước dày đặc trong chuồng trại chăn nuôi.
    \item $L_{FS}$: Suy hao không gian tự do (Free Space Path Loss) - yếu tố ăn mòn năng lượng tàn khốc nhất.
\end{itemize}
Tham chiếu quan trọng nhất để đánh giá sức mạnh hệ thống không phải là công suất phát, mà là Độ nhạy máy thu (Receiver Sensitivity - $S$). Một máy thu có độ nhạy $-130 \text{ dBm}$ sẽ tốt hơn hàng ngàn lần so với máy thu chỉ có độ nhạy $-100 \text{ dBm}$. Phương trình kinh điển tính toán giới hạn độ nhạy nhiệt động học của một hệ thống giải điều chế LoRa được minh họa chi tiết\cite{bor2016do}:
\begin{equation}
    {S = P_{Thermal} + NF + SNR_{limit}}.
\end{equation}
\begin{equation}
    {S = -174 + 10 \log_{10}(BW) + NF + SNR_{limit}}.
\end{equation}
Trong đó: 
\begin{itemize}
    \item Mức $-174 \text{ dBm/Hz}$ là nhiễu nền nhiệt (Thermal Noise Floor). Đây là giới hạn vật lý tuyệt đối sinh ra do sự chuyển động nhiệt hỗn loạn của các hạt electron bên trong vật dẫn đồng của anten tại nhiệt độ phòng (khoảng $290K$ hoặc $17^\circ C$). Bất kỳ thiết bị vô tuyến nào trên Trái Đất cũng phải chịu một mức nhiễu nền tối thiểu này.
    \item Mảnh ghép $10 \log_{10}(BW)$ là tích phân nhiễu phân bố trên toàn dải băng thông hoạt động. Băng thông càng rộng, miệng phễu hút nhiễu nền càng lớn.
    \item $NF$ (Noise Figure - Hệ số nhiễu phần cứng): Các linh kiện bán dẫn (Transistor) bên trong mạch khuếch đại tần số thấp (LNA) của bản thân máy thu cũng tạo ra nhiễu nhiệt nội bộ. Đối với chip IC LoRa, hệ số này thường dao động ở mức $6 \text{ dB}$.
    \item $SNR_{limit}$ (Tỷ số tín hiệu trên nhiễu cực hạn): Đây là điểm kỳ diệu của thuật toán CSS. Đối với các chuẩn viễn thông cũ, $SNR$ bắt buộc phải là số dương (tín hiệu phải mạnh hơn nhiễu từ 5 đến 10 lần mới đọc được). Nhưng đối với LoRa, nhờ bộ lọc phối hợp toán học, máy thu có thể khôi phục dữ liệu ở mức SNR âm sâu. Với cấu hình SF12, LoRa đòi hỏi SNR chỉ cần đạt $-20 \text{ dB}$.
\end{itemize}
Thế các tham số vào phương trình cho cấu hình tối đa năng lực đâm xuyên ($BW = 125 \text{ kHz}$, $SF = 12$): $S = -174 + 10 \log_{10}(125000) + 6 - 20 = -137 \text{ dBm}$

Một quỹ đường truyền khổng lồ (từ cực đại phát $+20 \text{ dBm}$ đến cực tiểu nhận $-137 \text{ dBm}$) cung cấp cho kiến trúc hệ thống một ngân sách năng lượng lên tới $157 \text{ dB}$, thừa sức xuyên thủng các lớp bê tông cốt thép nhiều tầng mà không cần thêm bộ khuyếch đại nào.
\subsubsection{Suy hao truyền rỗng (FSPL) và mô hình nhiễu vô tuyến trong công nghiệp \cite{rappaport2002wireless}}
Ngay cả khi truyền trong môi trường chân không tuyệt đối không có hạt bụi nào, năng lượng sóng điện từ vẫn bị suy giảm thảm hại do sự lan tỏa cầu của mặt sóng. Năng lượng phát ra từ anten đẳng hướng (Isotropic) lan truyền dưới dạng các bề mặt hình cầu đồng tâm, và diện tích mặt cầu tăng tỷ lệ thuận với bình phương bán kính. Do đó, mật độ công suất trên một đơn vị diện tích suy giảm theo luật bình phương nghịch đảo.

Mô hình này được mô tả bằng định luật Friis thông qua phương trình Suy hao không gian tự do (Free Space Path Loss - $FSPL$)\cite{rappaport2002wireless}:
\begin{equation}
    {FSPL = \left( \frac{4\pi d}{\lambda} \right)^2}.
\end{equation}
Chuyển đổi sang thang đo Decibel (logarit cơ số 10) để tiện cho việc cộng trừ trong ngân sách đường truyền, và thay thế bước sóng $\lambda = \frac{c}{f}$ (với $c$ là vận tốc ánh sáng $\approx 3 \times 10^8 \text{ m/s}$), ta được công thức thực hành:
\begin{equation}
    {FSPL (dB) = 20 \log_{10}(d) + 20 \log_{10}(f) + 32.44}.
\end{equation}
Trong công thức trên, $d$ là khoảng cách truyền sóng (đo bằng km) và $f$ là tần số hoạt động trung tâm (đo bằng MHz). Dựa vào công thức này, ta có thể chứng minh tính chất vật lý của tần số: ở cùng một khoảng cách $d$, sóng mang tần số thấp (như 433 MHz) sẽ bị suy hao ít hơn rất nhiều so với sóng mang tần số cao (như WiFi 2400 MHz). Đây là nền tảng lý thuyết sống còn giải thích tại sao đề tài lại chọn dải tần Sub-GHz cho hệ thống IoT nông nghiệp.

Tuy nhiên, trong thế giới thực của chuồng trại chăn nuôi, sóng vô tuyến hầu như không bao giờ bay thẳng mượt mà trong không gian tự do (Line-of-Sight - LOS). Sóng sẽ va đập vào bề mặt sàn bê tông, mái tôn kẽm và các máng ăn kim loại, tạo ra hiện tượng phản xạ, khúc xạ và tán xạ đa hướng. Các tia sóng dội này, di chuyển theo nhiều quỹ đạo dài ngắn khác nhau, sẽ đến anten thu ở các khoảng thời gian khác nhau (Multipath effect).

Ngoài ra, để đảm bảo liên kết vô tuyến LOS được coi là "hoàn toàn sạch", không chỉ cần một đường thẳng quang học không vật cản nối giữa hai anten, mà còn yêu cầu một khoảng không gian hình chóp nón xung quanh đường thẳng đó không bị xâm phạm. Khoảng không gian này được gọi là Vùng Fresnel (Fresnel Zone). Bán kính của vùng Fresnel thứ nhất $R_1$ tại điểm giữa đường truyền (nơi phình to nhất) được tính theo nguyên lý sóng quang học~\cite{rappaport2002wireless}:
\begin{equation}
    {R_1 = 17.32 \times \sqrt{\frac{d}{4f}}}.
\end{equation}
Nếu các chướng ngại vật (như ngọn cây hoặc dầm thép) lấn sâu vào hơn 40\% diện tích của Vùng Fresnel thứ nhất, hiện tượng nhiễu xạ (Diffraction loss) sẽ phát sinh tàn phá quỹ đường truyền một cách tiêu cực. Do đó, việc am hiểu mô hình suy hao FSPL và hiệu ứng đa đường (Multipath Fading) là điều kiện kiên quyết để nhóm nghiên cứu có thể thiết kế sơ đồ triển khai thực địa: quyết định độ cao treo anten vật lý của Node cảm biến, góc lắp đặt hộp thiết bị chống nước và các chiến lược khôi phục gói tin bằng thuật toán Điện toán biên trong trường hợp đường truyền bị gián đoạn cục bộ bởi nhiễu không gian.

\subsection{Kiến trúc vi điều khiển và Hệ điều hành thời gian thực (RTOS)}
Trong kiến trúc truyền thống, một MCU 8-bit chỉ đảm nhận nhiệm vụ thu thập tín hiệu tuần tự. Tuy nhiên, khi trở thành một Node Điện toán biên, MCU phải gánh vác đồng thời hàng loạt tác vụ nặng nề: tiền xử lý số liệu môi trường, chạy thuật toán Máy trạng thái hữu hạn (FSM) để ra quyết định, giao tiếp vô tuyến nối tiếp, quản lý bộ nhớ cục bộ và mã hóa dữ liệu. Để giải quyết bài toán đa nhiệm phức tạp này trên một thiết bị bị giới hạn nghiêm ngặt về tài nguyên, việc từ bỏ phương pháp lập trình tuần tự truyền thống và áp dụng Hệ điều hành thời gian thực (Real-Time Operating System - RTOS) là một bước tiến bắt buộc về mặt kiến trúc.
\subsubsection{Giới hạn của mô hình lập trình vòng lặp tuần tự}
Trong kỹ thuật lập trình hệ thống nhúng cơ bản (đặc trưng bởi nền tảng Arduino hay PIC sơ khai), kiến trúc phần mềm thường được xây dựng dựa trên mô hình Vòng lặp siêu cấp (Super-loop Architecture), thường được biết đến qua hàm while(1) hoặc loop(). Cơ chế hoạt động của mô hình này là vi điều khiển sẽ thực thi các khối lệnh một cách tuần tự từ trên xuống dưới: đọc cảm biến nhiệt ẩm, đọc cảm biến khí gas, tính toán so sánh, sau đó đẩy dữ liệu ra cổng nối tiếp (UART) tới module vô tuyến, và kết thúc bằng một lệnh chờ (Delay) trước khi lặp lại chu trình.

Hạn chế nhất của Super-loop Architecture là tính chất "Chặn" (Blocking) của luồng thực thi. Khi vi điều khiển giao tiếp với các cảm biến ngoại vi thông qua bus I2C (như quá trình đọc dữ liệu từ cảm biến nhiệt ẩm SHT31 hoặc cảm biến ánh sáng BH1750), nếu xảy ra hiện tượng nhiễu điện từ trường (EMI) trong môi trường chuồng trại làm mất tín hiệu xung nhịp đồng bộ (SCL), toàn bộ bus I2C có thể rơi vào trạng thái treo (Deadlock). Tại thời điểm đó, vi xử lý sẽ bị kẹt vĩnh viễn ở dòng lệnh chờ phản hồi (Polling) từ I2C \cite{white2011making}.

Hệ quả của việc này đối với Điện toán biên là sự vi phạm hoàn toàn các ràng buộc về thời gian. Giả sử trong lúc vi xử lý đang bị treo ở hàm đọc dữ liệu ánh sáng, nồng độ khí độc NH3 và nhiệt độ môi trường đột ngột tăng vượt mức nguy hiểm. Vi điều khiển không thể thoát khỏi hàm bị treo để tiến xuống dòng lệnh bên dưới nhằm kích hoạt chân GPIO điều khiển đóng mạch rơ-le làm mát. Ngay cả khi không xảy ra lỗi, việc một tác vụ truyền thông vô tuyến (thường mất từ hàng chục đến hàng trăm mili-giây) chiếm dụng CPU cũng làm cho hệ thống bị mù tạm thời trước các biến động của môi trường. Nói cách khác, kiến trúc Super-loop thiếu khả năng ưu tiên hóa tác vụ (Task Prioritization), khiến nó hoàn toàn không đủ độ tin cậy để ứng dụng trong các hệ thống giám sát an toàn sinh học có nguy cơ thiệt hại lớn.

\subsubsection{Nguyên lý hoạt động của Hệ điều hành thời gian thực (FreeRTOS)}
Để phá vỡ giới hạn tuyến tính của Super-loop, Hệ điều hành thời gian thực (RTOS), cụ thể là FreeRTOS, được tích hợp trực tiếp vào nhân của vi điều khiển. Khác với các hệ điều hành mục đích chung (General-Purpose OS như Windows hay Linux) được thiết kế để tối đa hóa thông lượng (Throughput) chung của hệ thống, RTOS được thiết kế với mục tiêu tối thượng là tính tất định (Determinism) – khả năng đảm bảo một tác vụ sẽ được thực thi trong một khoảng thời gian dự đoán được, bất kể khối lượng công việc hiện tại của hệ thống \cite{laplante2004real}.

Nguyên lý cốt lõi của FreeRTOS dựa trên cơ chế Lập lịch ưu tiên độc chiếm (Preemptive Priority-based Scheduling). Thay vì một luồng duy nhất, mã nguồn được chia nhỏ thành các Tác vụ (Tasks) độc lập. Mỗi tác vụ được hệ điều hành cấp phát một vùng nhớ ngăn xếp (Stack Memory) riêng biệt và một Khối điều khiển tác vụ (Task Control Block - TCB). Mức độ ưu tiên (Priority) của mỗi tác vụ được kỹ sư gán cứng dựa trên tầm quan trọng của nó đối với sự an toàn của hệ thống.

Thuật toán lập lịch của FreeRTOS hoạt động dựa trên một bộ định thời phần cứng tạo ra các ngắt chu kỳ (Tick Interrupts), thường được cấu hình ở tần số 1000 Hz (nghĩa là 1 mili-giây đánh dấu một nhịp). Tại mỗi nhịp Tick, Bộ lập lịch (Scheduler) sẽ kiểm tra danh sách các tác vụ đang ở trạng thái Sẵn sàng (Ready State). Toán học của quá trình Chuyển đổi ngữ cảnh (Context Switching) được diễn ra khi Scheduler phát hiện một tác vụ có mức ưu tiên cao hơn vừa được giải phóng khỏi trạng thái Chờ (Blocked State) – ví dụ: Task Điện toán biên vừa nhận được một tín hiệu ngắt báo hiệu nhiệt độ quá tải.

Thời gian đáp ứng của hệ thống (System Response Time - $T_{res}$) trong một hệ thống RTOS được tính bằng phương trình vi mô:
\begin{equation}
    {T_{res} = T_{int} + T_{context} + T_{exec}}.
\end{equation}
Trong đó:
\begin{itemize}
    \item $T_{int}$: Thời gian trễ ngắt phần cứng (Interrupt Latency) – khoảng thời gian từ lúc có tín hiệu vật lý đến khi CPU nhận biết.
    \item $T_{context}$: Thời gian chuyển đổi ngữ cảnh (Context Switch Time). CPU phải lưu toàn bộ các thanh ghi (Registers) của tác vụ đang chạy vào vùng Stack của nó, sau đó tải các thanh ghi của tác vụ ưu tiên cao từ Stack của tác vụ đó vào CPU.
    \item $T_{exec}$: Thời gian thực thi thuật toán logic của tác vụ mới.
\end{itemize}
Nhờ cơ chế Preemptive, hệ thống không bao giờ bị nghẽn mạng tại một cảm biến hỏng. Nếu Task đọc I2C bị treo, Task Điều khiển biên (với độ ưu tiên cao nhất) vẫn luôn được Scheduler trao quyền chiếm đoạt CPU (Preempt) ngay lập tức để ngắt điện toàn bộ hệ thống hoặc bật hệ thống làm mát khẩn cấp, đảm bảo Ràng buộc độ trễ (Latency Constraint) luôn nhỏ hơn 50ms.

% --- Bổ sung phần RTOS ---
\subsubsection{Vòng đời và các trạng thái của một Tác vụ (Task Lifecycle)}
Để bộ lập lịch có thể điều phối tất định, mỗi Tác vụ trong FreeRTOS tại mọi thời điểm luôn nằm ở một trong bốn trạng thái logic, và việc dịch chuyển giữa các trạng thái này được kích hoạt bởi các sự kiện xác định. Hiểu rõ vòng đời này (Hình \ref{fig:rtos_task_lifecycle}) là điều kiện tiên quyết để thiết kế đúng mức ưu tiên và cơ chế đồng bộ ở các chương sau.
\begin{figure}[H]
    \centering
    \includegraphics[width=0.8\linewidth]{Thuc_Hanh/Image/TaskCycleFreeRTOS.pdf}
    \caption{Sơ đồ vòng đời và bốn trạng thái của một Tác vụ trong FreeRTOS}
    \label{fig:rtos_task_lifecycle}
\end{figure}
% \begin{figure}[H]
% \centering
% \begin{tikzpicture}[
%     >=Stealth,
%     st/.style={rectangle, rounded corners=3pt, draw, thick, align=center,
%                minimum width=2.7cm, minimum height=1.05cm},
%     lb/.style={font=\scriptsize\itshape, align=center, fill=white, inner sep=1.5pt}
% ]
% \node[st, fill=blue!8]    (ready) at (-3.2,2)    {\textbf{SẴN SÀNG}\\(Ready)};
% \node[st, fill=green!18]  (run)   at (3.2,2)     {\textbf{ĐANG CHẠY}\\(Running)};
% \node[st, fill=orange!20] (block) at (3.2,-1.7)  {\textbf{BỊ CHẶN}\\(Blocked)};
% \node[st, fill=black!12]  (susp)  at (-3.2,-1.7) {\textbf{TẠM DỪNG}\\(Suspended)};

% \draw[->,thick] (ready) to[bend left=12]  node[lb,above]{Được điều phối\\(ưu tiên cao nhất)} (run);
% \draw[->,thick] (run)   to[bend left=12]  node[lb,below]{Bị chiếm quyền /\\hết time-slice} (ready);
% \draw[->,thick] (run)   to                node[lb,right]{Chờ sự kiện:\\Queue / Semaphore /\\vTaskDelay()} (block);
% \draw[->,thick] (block) to[bend left=9]   node[lb,pos=0.55]{Sự kiện đến /\\hết timeout} (ready);
% \draw[->,thick] (ready) to[bend left=15]  node[lb,left]{vTaskSuspend()} (susp);
% \draw[->,thick] (susp)  to[bend left=15]  node[lb,right]{vTaskResume()} (ready);
% \end{tikzpicture}
% \caption{Sơ đồ vòng đời và bốn trạng thái của một Tác vụ trong FreeRTOS cùng các sự kiện chuyển trạng thái}
% \label{fig:rtos_task_lifecycle}
% \end{figure}

Trạng thái \textit{Đang chạy} (Running) là khi mã của tác vụ thực sự được CPU thực thi; trên một lõi đơn, tại mỗi thời điểm chỉ có duy nhất một tác vụ ở trạng thái này. Trạng thái \textit{Sẵn sàng} (Ready) gồm các tác vụ đã đủ điều kiện chạy và đang xếp hàng chờ CPU, nhưng chưa được chọn vì có tác vụ khác ưu tiên cao hơn (hoặc bằng) đang chiếm lõi. Trạng thái \textit{Bị chặn} (Blocked) là khi tác vụ chủ động nhường CPU để chờ một sự kiện -- chẳng hạn chờ dữ liệu trong Hàng đợi, chờ một Cờ hiệu, hoặc chờ hết một khoảng trễ \texttt{vTaskDelay()}; điểm mấu chốt là tác vụ bị chặn \emph{không} tiêu tốn chu kỳ CPU, khác hẳn vòng lặp chờ bận (busy-wait) của super-loop. Cuối cùng, trạng thái \textit{Tạm dừng} (Suspended) là khi tác vụ bị đưa ra khỏi lịch trình một cách tường minh bằng \texttt{vTaskSuspend()} và chỉ quay lại khi được \texttt{vTaskResume()}.


% --- Hết bổ sung phần RTOS ---

\subsubsection{Quản lý đa nhiệm và cơ chế phân luồng (Core Pinning) trên MCU}
 Trong khi các hệ thống nhúng thông thường chỉ có một nhân (Core) phải liên tục chuyển đổi qua lại giữa các tác vụ, dòng vi điều khiển 32-bit Xtensa LX6 cung cấp hai nhân vật lý độc lập có khả năng thực thi hai luồng lệnh song song cùng một chu kỳ thời gian thực \cite{espressif2023esp32}.

Sử dụng bộ API của FreeRTOS (như xTaskCreatePinnedToCore), hệ thống phần mềm của Node có thể ép cứng (Pinning) các tác vụ cụ thể vào một trong hai nhân vật lý: Core 0 (được gọi là PRO\_CPU) và Core 1 (được gọi là APP\_CPU).
\begin{figure}[H]
    \centering
    \includegraphics[width=1.0\linewidth]{Thuc_Hanh/Image/esp32_dual.pdf}
    \caption{Kiến trúc phân chia tác vụ trên bộ xử lý Dual-core của ESP32}
    \label{Hình 2.6}
\end{figure}
Hình \ref{Hình 2.6} minh họa một kiến trúc phân bổ tài nguyên tinh tế nhằm cô lập các tác vụ có tính chất cản trở lẫn nhau. Core 0 (PRO\_CPU) được giao phó toàn bộ khối lượng công việc liên quan đến giao tiếp ngoại vi giao thức vô tuyến (UART) và toán học mã hóa (AES-128). Bản chất của giao tiếp UART với module LoRa họ E32 là quá trình truyền nhận nối tiếp tốc độ tương đối chậm (thường ở mức 9600 bps). Nếu CPU phải ngồi chờ từng byte truyền đi, nó sẽ lãng phí hàng ngàn chu kỳ máy (CPU cycles). Do đó, Core 0 sẽ ôm toàn bộ các ngắt vô tuyến và xử lý các hàm nhân ma trận của thuật toán AES-128.

Trong khi đó, Core 1 (APP\_CPU) được giữ hoàn toàn "sạch sẽ" và chuyên biệt hóa để phục vụ ứng dụng Điện toán biên. Core 1 đóng vai trò là "bộ não" thực thi Máy trạng thái hữu hạn (FSM). Nó liên tục rà soát dữ liệu từ cảm biến SHT31 và MQ-135, đối chiếu với ma trận ngưỡng nhiệt-ẩm. Bởi vì Core 1 không bao giờ bị ngắt quãng bởi các tín hiệu vô tuyến LoRa rác rưởi từ bên ngoài mạng dội vào (vì ngắt vô tuyến đã bị Core 0 chặn đứng và xử lý), thời gian thực thi của thuật toán FSM để ra lệnh đóng mạch Rơ-le bảo vệ là tuyệt đối mang tính tất định (Deterministic), đảm bảo độ trễ (Edge Latency) luôn ghim chặt dưới mức 50 mili-giây. Để hai nhân này có thể chia sẻ dữ liệu môi trường cho nhau một cách an toàn mà không gây ra hiện tượng Cạnh tranh dữ liệu (Race Condition), hệ thống sử dụng cơ chế Giao tiếp liên tiến trình (IPC) của FreeRTOS thông qua các Hàng đợi (Queues) và Cờ hiệu độc quyền (Mutexes).

% --- Bổ sung phần RTOS ---
\subsubsection{Mô hình đa xử lý đối xứng (SMP) và chiến lược ghim lõi}
Cần phân định rõ về mặt học thuật giữa hai mô hình điều phối trên vi xử lý đa nhân. Trong mô hình \textit{Đa xử lý bất đối xứng} (Asymmetric Multiprocessing -- AMP) cổ điển, mỗi lõi chạy một thực thể hệ điều hành riêng biệt. Ngược lại, nhân FreeRTOS được port cho ESP-IDF là một nhân \textit{Đa xử lý đối xứng} (Symmetric Multiprocessing -- SMP): một bộ lập lịch duy nhất quản lý chung danh sách tác vụ Sẵn sàng và phân phối chúng cho cả hai lõi (PRO\_CPU và APP\_CPU) \cite{espressif_freertos_idf}. Trong mô hình SMP, mỗi tác vụ mang một thuộc tính core affinity: nếu để mặc định không ghim (\texttt{tskNO\_AFFINITY}), tác vụ có thể được lập lịch trên bất kỳ lõi nào còn rảnh, giúp cân bằng tải tự động; nếu ghim cứng vào một lõi cụ thể thông qua \texttt{xTaskCreatePinnedToCore()}, tác vụ chỉ chạy trên lõi đó.

Đề tài lựa chọn chiến lược áp đặt một kỷ luật phân hoạch kiểu AMP lên trên nền nhân SMP. Lý do là vì lập lịch SMP tự do, dù tối ưu thông lượng tổng, lại làm tăng tính bất định: một tác vụ có thể di trú (migrate) giữa hai lõi, gây ra chi phí mất hiệu lực bộ nhớ đệm (cache invalidation) và làm khó dự đoán thời điểm thực thi. Bằng cách ghim cứng tác vụ FSM an toàn vào một lõi chuyên biệt, đề tài đánh đổi một phần thông lượng để thu về tính tất định (determinism) -- yếu tố sống còn đối với ràng buộc trễ điều khiển dưới $50\,\text{ms}$.

\subsubsection{Loại trừ tương hỗ trên bus I2C và cơ chế kế thừa ưu tiên}
Trong số các tài nguyên dùng chung, bus I2C là điểm tranh chấp điển hình và nguy hiểm nhất của hệ thống: bốn thiết bị (cảm biến SHT31, BMP180, BH1750 và đồng hồ thời gian thực DS3231) cùng treo trên một đôi dây SDA/SCL duy nhất. Một giao dịch I2C là một chuỗi nguyên tử gồm điều kiện START, gửi địa chỉ, trao đổi nhiều byte và điều kiện STOP. Nếu một tác vụ đang đọc dở SHT31 thì bị một tác vụ khác chiếm quyền (preempt) để khởi tạo một giao dịch tới DS3231, hai chuỗi xung sẽ chồng lấn trên cùng đường dây, phá vỡ trạng thái máy của bus và dẫn đến dữ liệu rác hoặc tệ hơn là khóa cứng bus (bus lockup). Đây chính là một dạng Cạnh tranh dữ liệu (Race Condition) ở mức phần cứng.

Giải pháp là bảo vệ mỗi giao dịch bằng một Mutex. Trước khi truy cập bus, tác vụ phải giành quyền sở hữu Mutex (\texttt{xSemaphoreTake}); sau khi hoàn tất giao dịch mới trả lại (\texttt{xSemaphoreGive}). Cần nhấn mạnh sự khác biệt giữa Mutex và Cờ hiệu nhị phân (Binary Semaphore): tuy cùng có hai trạng thái, nhưng Mutex có khái niệm \textit{quyền sở hữu} (ownership -- chỉ tác vụ đã lấy mới được trả) và quan trọng hơn là tích hợp cơ chế \textit{kế thừa ưu tiên} (priority inheritance), nên Mutex mới là công cụ đúng để bảo vệ tài nguyên, còn Binary Semaphore chỉ nên dùng để báo hiệu sự kiện.

Cơ chế kế thừa ưu tiên ra đời để chống lại hiện tượng Priority Inversion được minh họa ở Hình \ref{fig:rtos_priority_inheritance}.
\begin{figure}[H]
    \centering
    \includegraphics[width=0.95\linewidth]{Thuc_Hanh/Image/Mutex.pdf}
    \caption{Cơ chế kế thừa ưu tiên của Mutex}
    \label{fig:rtos_priority_inheritance}
\end{figure}
Một tác vụ ưu tiên thấp (L) đang giữ Mutex của bus I2C thì bị tác vụ ưu tiên cao (H) đòi cùng Mutex; H buộc phải chờ. Nếu lúc này một tác vụ ưu tiên trung bình (M) -- không liên quan đến bus -- giành CPU và chạy dài, nó sẽ vô tình chặn L hoàn thành giao dịch, khiến H (ưu tiên cao nhất) bị chặn vô thời hạn bởi M (ưu tiên thấp hơn). FreeRTOS giải quyết bằng cách: khi H chặn trên Mutex mà L đang giữ, hệ điều hành tạm thời \textit{nâng} ưu tiên của L lên bằng ưu tiên của H, nhờ đó L không bị M chặn, hoàn tất nhanh giao dịch và trả Mutex; sau đó ưu tiên của L được khôi phục về mức cũ. Cơ chế này \textit{giới hạn cận trên} (bound) thời gian chặn của H đúng bằng độ dài một giao dịch I2C, bảo toàn tính tất định của hệ thống.

% \begin{figure}[H]
% \centering
% \begin{tikzpicture}[font=\scriptsize, >=Stealth]
% % --- truc thoi gian va lan tac vu ---
% \foreach \y/\name in {0/{L (thấp)}, 1.3/{M (trung bình)}, 2.6/{H (cao)}}{
%     \node[anchor=east] at (-0.15,\y+0.25) {\name};
%     \draw[->,gray] (0,\y) -- (10.6,\y);
% }
% \draw[->] (0,-0.6) -- (10.6,-0.6) node[anchor=west,black]{$t$};
% % --- moc thoi gian (duong dut) ---
% \foreach \x in {1,2,3,4,5,7}{ \draw[dashed,black!30] (\x,-0.6) -- (\x,3.2); }
% % --- khoi thuc thi cua L ---
% \fill[blue!25,draw=black] (1,0) rectangle (2,0.5);          % L chay binh thuong
% \fill[blue!25,draw=black] (2,0) rectangle (3,0.5);          % L trong CS
% \fill[red!35,draw=black]  (4,0) rectangle (5,0.5);          % L chay o uu tien ke thua
% \node[font=\tiny] at (2.5,0.78) {giữ Mutex};
% % --- khoi thuc thi cua H ---
% \fill[green!30,draw=black] (3,2.6) rectangle (4,3.1);       % H chay
% \draw[black!50] (4,2.6) rectangle (5,3.1);                  % H cho Mutex (rong)
% \node[font=\tiny] at (4.5,2.85) {chờ Mutex};
% \fill[green!30,draw=black] (5,2.6) rectangle (7,3.1);       % H chay trong CS sau khi lay Mutex
% % --- khoi thuc thi cua M ---
% \fill[orange!35,draw=black] (7,1.3) rectangle (9,1.8);      % M chay sau cung
% % --- chu thich su kien ---
% \draw[->] (2,1.05) -- (2,0.55); \node[font=\tiny,anchor=south] at (2,1.05) {Lock@t=2};
% \draw[->] (5,1.05) -- (5,0.55); \node[font=\tiny,anchor=south] at (5,1.1) {Unlock@t=5};
% \node[font=\tiny,red] at (4.5,-0.35) {Kế thừa ưu tiên: L chạy ở ưu tiên của H};
% \end{tikzpicture}
% \caption{Cơ chế kế thừa ưu tiên của Mutex giúp giới hạn thời gian chặn khi bảo vệ bus I2C dùng chung}
% \label{fig:rtos_priority_inheritance}
% \end{figure}
% --- Hết bổ sung phần RTOS ---

\subsubsection{Kiến trúc phân vùng bộ nhớ và cơ chế lưu trữ ngoại tuyến}
Điện toán biên không chỉ bao hàm năng lực xử lý (Computing) mà còn đòi hỏi khả năng lưu trữ cục bộ (Local Storage). Đối với các thiết bị vi điều khiển nhúng, kiến trúc bộ nhớ là sự kết hợp giữa bộ nhớ truy cập ngẫu nhiên tĩnh/động (SRAM) để lưu các biến thực thi tốc độ cao, và bộ nhớ Flash phi bay hơi (Non-volatile Flash Memory) để lưu trữ mã nguồn và dữ liệu vĩnh viễn. Trong dự án này, bộ nhớ Flash SPI của ESP32 (thường có dung lượng 4MB) đóng vai trò như một "đám mây mini" nội bộ.

Việc lưu trữ dữ liệu thô trên bộ nhớ Flash không thể được thực hiện một cách hỗn loạn bằng cách ghi trực tiếp vào các địa chỉ ô nhớ vật lý, vì bộ nhớ Flash có cấu trúc đặc thù: dữ liệu chỉ có thể được ghi vào các ô nhớ đã bị xóa (Erased), và thao tác xóa chỉ có thể thực hiện theo từng Khối (Block/Sector) lớn chứ không thể xóa từng Byte. Để giải quyết rào cản phần cứng này, một Hệ thống tập tin chuyên biệt dành riêng cho bộ nhớ Flash SPI (ví dụ: SPIFFS hoặc LittleFS) được nhúng vào firmware \cite{espressif2023esp32}. Hệ thống tập tin này cung cấp các hàm API thao tác bậc cao như mở file (open), ghi file (write), đọc file (read) tương tự như trên hệ điều hành máy tính tiêu chuẩn, đồng thời nó tự động cân bằng hao mòn (Wear Leveling) bằng cách phân bổ đều các lệnh ghi/xóa ra khắp các sector, nhằm ngăn chặn việc một sector bị ghi đè quá số lần giới hạn (thường là 100,000 chu kỳ) dẫn đến hỏng bộ nhớ vật lý.
\begin{figure}[H]
    \centering
    \includegraphics[width=0.9\linewidth]{Thuc_Hanh/Image/MemoryPartion.pdf}
    \caption{Phân vùng bộ nhớ Flash và phân bổ không gian lưu trữ của ESP32}
    \label{Hình 2.7}
\end{figure}
Hình \ref{Hình 2.7} minh chứng cho chiến lược kiến tạo cấu trúc phần mềm bền vững đối mặt với "Nghịch lý đám mây". Phân vùng spiffs chiếm đến 1.5 MB dung lượng, cung cấp không gian đủ lớn để chứa hàng ngàn bản ghi (Log records). Mỗi khi quá trình phát sóng vô tuyến thông qua module LoRa thất bại (do mất điện Gateway hoặc nhiễu sóng nghiêm trọng, hàm đọc chân trạng thái AUX báo lỗi), luồng phần mềm tại Node sẽ lập tức rẽ nhánh. Thay vì ném bỏ các giá trị môi trường cực kỳ quan trọng, một tác vụ Ghi ngoại tuyến sẽ mở file log trong phân vùng spiffs và gắn (append) chuỗi dữ liệu (bao gồm nhiệt độ, độ ẩm, nồng độ NH3 và Timestamp). Ngay khi kết nối vô tuyến được khôi phục, một luồng đồng bộ chạy nền ở độ ưu tiên thấp (Background Sync Task) sẽ từ từ đọc từng dòng trong file log này và gửi gối đầu (backlog) về cơ sở dữ liệu trung tâm để khôi phục tính toàn vẹn của biểu đồ thời gian lịch sử, đảm bảo nhà nông không bao giờ mất dấu dữ liệu do hạ tầng mạng đứt gãy.

% --- Bổ sung phần RTOS ---
\subsubsection{Quản lý vùng nhớ Stack và Heap trong FreeRTOS/ESP-IDF}
Nếu bộ nhớ Flash lưu mã nguồn và dữ liệu bền vững thì bộ nhớ RAM (SRAM nội $\approx 512\,\text{KB}$ của ESP32-S3) mới là nơi diễn ra mọi hoạt động thực thi nóng, và nó được phân chia thành hai miền có bản chất hoàn toàn khác nhau: vùng ngăn xếp (Stack) cấp riêng cho từng tác vụ, và vùng nhớ động (Heap) dùng chung. Quản lý chặt chẽ hai miền này là điều kiện tiên quyết cho độ ổn định của một hệ thống chạy liên tục 24/7.

Mỗi khi một tác vụ được tạo ra, FreeRTOS cấp cho nó một vùng Stack riêng cùng một Khối điều khiển tác vụ (TCB). Stack lưu các biến cục bộ, khung lời gọi hàm và toàn bộ thanh ghi được lưu lại tại mỗi lần Chuyển đổi ngữ cảnh. Việc định kích thước Stack là một bài toán đánh đổi sắc bén: cấp quá nhỏ sẽ gây tràn ngăn xếp (Stack Overflow) -- tác vụ ghi đè ra ngoài vùng cấp phát, làm hỏng bộ nhớ của tác vụ khác và sập hệ thống một cách khó tái hiện; cấp quá lớn lại lãng phí RAM vốn rất quý. Để giải bài toán này một cách định lượng thay vì phỏng đoán, FreeRTOS cung cấp hàm \texttt{uxTaskGetStackHighWaterMark()} trả về \textit{mức dùng đỉnh} (High-Water Mark) -- lượng Stack còn trống ít nhất mà tác vụ từng chạm tới kể từ khi khởi động; kết hợp với tùy chọn \texttt{configCHECK\_FOR\_STACK\_OVERFLOW} (đặt một "lính canh" -- canary value -- ở đáy Stack để phát hiện ghi tràn), kỹ sư có thể tinh chỉnh kích thước Stack vừa vặn với mức dư an toàn.

Vùng Heap phục vụ cấp phát động (\texttt{malloc}/\texttt{free}). FreeRTOS chuẩn cung cấp năm sơ đồ quản lý Heap với các đánh đổi khác nhau: \texttt{heap\_1} chỉ cấp phát mà không bao giờ giải phóng (đơn giản, tất định nhất); \texttt{heap\_2} cho phép giải phóng nhưng không gộp các khối trống liền kề (dễ phân mảnh); \texttt{heap\_3} bọc \texttt{malloc}/\texttt{free} của thư viện C; \texttt{heap\_4} (mặc định phổ biến) có gộp khối trống (coalescing) để chống phân mảnh; và \texttt{heap\_5} mở rộng \texttt{heap\_4} cho nhiều vùng nhớ rời rạc. Trên nền tảng đó, ESP-IDF xây dựng một bộ cấp phát đa heap theo năng lực (capability-based allocator): mỗi lần cấp phát, lập trình viên khai báo loại bộ nhớ cần qua các cờ năng lực như \texttt{MALLOC\_CAP\_INTERNAL} (RAM nội tốc độ cao), \texttt{MALLOC\_CAP\_SPIRAM} (PSRAM ngoài dung lượng lớn), hay \texttt{MALLOC\_CAP\_DMA} (vùng nhớ truy cập được bằng DMA cho ngoại vi) \cite{espressif_mem_alloc}. Hai hàm chẩn đoán then chốt là \texttt{heap\_caps\_get\_free\_size()} (tổng dung lượng còn trống của một loại RAM) và \texttt{heap\_caps\_get\_largest\_free\_block()} (khối liền kề lớn nhất còn trống); chênh lệch giữa hai giá trị này chính là thước đo mức độ phân mảnh (fragmentation) của Heap.

Từ phân tích trên rút ra nguyên tắc thiết kế cho firmware Node: hạn chế tối đa việc cấp phát/giải phóng động trong các luồng xử lý nóng (hot path) nhằm tránh phân mảnh và độ trễ bất định; ưu tiên cấp phát tĩnh hoặc cấp phát một lần lúc khởi tạo cho các cấu trúc dữ liệu có vòng đời dài (hàng đợi, bộ đệm gói tin); và giám sát liên tục dung lượng Heap trống để phát hiện sớm rò rỉ bộ nhớ (memory leak) -- một nguyên nhân hàng đầu khiến các thiết bị nhúng sập sau nhiều ngày chạy.
% --- Hết bổ sung phần RTOS ---

\subsubsection{Mô hình quản lý năng lượng phân tầng và cơ chế ngủ sâu (Deep Sleep)}
Một kiến trúc Điện toán biên hoàn hảo về mặt thuật toán sẽ trở nên vô giá trị nếu nó tiêu thụ hết viên pin của thiết bị chỉ trong vài ngày. Năng lượng tiêu thụ tổng cộng của một Node vô tuyến ($E_{total}$) là tích phân của dòng điện theo thời gian, được chia thành ba dải trạng thái chính \cite{gomez2012overview}:
\begin{equation}
    {E_{total} = \int_{0}^{T} V_{bat} \cdot I(t) \, dt = V_{bat} \left( I_{active} \cdot T_{active} + I_{tx} \cdot T_{tx} + I_{sleep} \cdot T_{sleep} \right)}.
\end{equation}
Trong đó, $I_{tx}$ (dòng phát xạ vô tuyến) và $I_{active}$ (dòng hoạt động đỉnh của MCU) là những thông số khổng lồ (từ 100mA đến vài trăm mA). Vì các tác vụ tính toán nhanh (FSM) và phát LoRa chỉ diễn ra trong các khe thời gian chớp nhoáng ($T_{active}$ và $T_{tx}$ cực nhỏ), nên năng lượng hao hụt thực chất nằm ở khoảng thời gian hệ thống chờ đợi các chu kỳ kế tiếp ($T_{sleep}$ chiếm đến 99\% thời gian sống của thiết bị). Việc hạ gục chỉ số $I_{sleep}$ xuống mức tiệm cận 0 là chén thánh của kỹ thuật thiết kế phần cứng WSN.

Vi điều khiển ESP32 được trang bị bộ điều khiển quản lý năng lượng tiên tiến (Power Management Unit) với nhiều dải Sleep khác nhau. Khi gọi hàm tiến vào chế độ Ngủ sâu (Deep Sleep), bộ điều khiển sẽ cắt đứt toàn bộ nguồn điện và xung nhịp (Clock gating) dẫn vào Core 0, Core 1, bộ nhớ SRAM chính và mọi khối ngoại vi kỹ thuật số tốc độ cao (SPI, I2C, UART). Chỉ có duy nhất một bộ Đồng hồ thời gian thực siêu tiết kiệm điện (RTC Controller) và khối nhớ ULP (Ultra-Low Power coprocessor) được giữ lại nguồn để duy trì chức năng đếm giờ (Timer) hoặc chờ kích hoạt ngắt vật lý từ chân cảm biến. Tại trạng thái này, dòng tiêu thụ bản thân của vi điều khiển giảm từ mức $40 \text{ mA}$ xuống chỉ còn xấp xỉ $10 \mu A$.

Tuy nhiên, sự rò rỉ dòng điện tĩnh (Quiescent Current) từ các thành phần ngoại vi trên bo mạch (Module LoRa, các IC ổn áp, và đặc biệt là hệ thống nung nóng vi sợi của cảm biến khí gas MQ-135) vẫn sẽ hút cạn viên pin nếu không bị can thiệp. Để đạt được hiệu năng tối ưu tuyệt đối, kiến trúc phần cứng của Node đã được thiết kế sử dụng một MOSFET kênh P đóng vai trò làm Công tắc tải (Load Switch) \cite{instruments2016load}. Trước khi vi điều khiển gọi hàm Sleep, nó sẽ xuất mức logic đảo tới cực Cổng (Gate) của MOSFET, trực tiếp cắt đứt nguồn vật lý dẫn đến toàn bộ hệ sinh thái cảm biến. Sự kết hợp giữa Deep Sleep bằng phần mềm và Load Switch bằng phần cứng tạo ra một mô hình Hybrid Sleep hoàn mỹ, đè bẹp $I_{sleep}$ tổng thể của toàn hệ thống xuống dưới ngưỡng $20 \mu A$, kéo dài thời gian sống của khối pin 18650 lên đến nhiều tháng trời trong môi trường hoang dã.

\subsection{Bảo mật trong mạng IoT và Thuật toán mã hóa AES-128}
Sự lan tỏa sóng điện từ vô hình trong không gian mang lại sự tự do trong việc triển khai Mạng cảm biến không dây (WSN), nhưng đồng thời nó cũng biến tầng vật lý thành một môi trường mở hoàn toàn, phơi bày mọi dữ liệu nhạy cảm ra trước hàng loạt các mối đe dọa an ninh không gian mạng. Khác với kết nối cáp quang nơi việc nghe lén đòi hỏi phải can thiệp vật lý vào đường truyền, trong mạng vô tuyến công suất thấp như LoRa, bất kỳ đối tượng nào sở hữu một thiết bị thu phát vô tuyến cấu hình tương đương (cùng tần số, băng thông và hệ số trải phổ) đều có khả năng đánh chặn toàn bộ thông lượng thông tin. Đối với một hệ thống Điện toán biên tích hợp điều khiển các hệ thống cơ điện công suất lớn (quạt hút, rèm che nhiệt), một lỗ hổng bảo mật tại Node không đơn thuần dẫn đến sự rò rỉ dữ liệu, mà có thể biến thành một cuộc tấn công phá hoại hạ tầng vật lý (Cyber-Physical Attack).
\subsubsection{Các bề mặt tấn công vô tuyến trong mạng cảm biến}
Hệ thống WSN tại biên mạng nông nghiệp đối mặt với ba bề mặt tấn công thụ động và chủ động cốt lõi \cite{granjal2015security}:
\begin{itemize}
    \item Nghe lén và Rò rỉ thông tin (Eavesdropping / Sniffing): Kẻ tấn công thiết lập một Node giả mạo (Rogue Node) ngay bên ngoài hàng rào trang trại, thiết lập module LoRa ở chế độ liên tục nghe (Continuous Receive Mode). Nếu dữ liệu chuồng trại (nhiệt độ, độ ẩm, nồng độ khí) được gửi dưới dạng bản rõ (Plaintext), kẻ gian sẽ nắm bắt được toàn bộ chu trình sinh học của vật nuôi, từ đó thực hiện hành vi do thám thông tin bí mật thương mại về quy trình chăm sóc.
    \item Sửa đổi và Làm giả dữ liệu (Data Tampering / Forgery): Nghiêm trọng hơn nghe lén, kẻ tấn công bắt được gói tin vô tuyến, giải mã, thay đổi nội dung tải trọng (ví dụ: sửa thông số "Nhiệt độ $25^\circ C$" thành "Nhiệt độ $45^\circ C$") và bơm ngược gói tin đó vào mạng. Việc làm giả mạo bản tin sẽ đánh lừa thuật toán Điện toán biên hoặc đánh lừa bộ ra quyết định tại Gateway, kích hoạt hệ thống làm mát hoặc sưởi ấm sai thời điểm, trực tiếp đe dọa tính mạng toàn đàn gia cầm. Yêu cầu tính Toàn vẹn dữ liệu (Data Integrity) là bắt buộc.
    \item Tấn công phát lại (Replay Attack): Kẻ gian có thể không cần biết cách giải mã thông tin. Bằng cách sử dụng hệ thống Vô tuyến định nghĩa bằng phần mềm (Software Defined Radio - SDR), kẻ tấn công ghi lại dải tín hiệu số của một lệnh điều khiển hợp lệ từ Node (ví dụ: lệnh "Tắt hệ thống thông gió"). Sau đó, chờ đến khi nhiệt độ chuồng trại lên đỉnh điểm vào giữa trưa, chúng phát lại lệnh này liên tục với công suất cực cao. Hệ thống Gateway nhận được lệnh (hoàn toàn hợp lệ về mặt mã hóa) và mù quáng tắt hệ thống thông gió, gây ngạt thở cho vật nuôi.
\end{itemize}
\subsubsection{Nguyên lý của mật mã học khối đối xứng (Symmetric Block Cipher)}
Để thiết lập tường lửa vững chắc ngăn chặn các bề mặt tấn công trên, mã hóa dữ liệu tại nguồn (tại vi điều khiển của Node) trước khi đưa xuống lớp vật lý LoRa là phương pháp duy nhất. Trong lý thuyết mật mã, có hai trường phái chính: Mã hóa bất đối xứng (Asymmetric Cryptography, ví dụ: RSA, ECC) và Mã hóa đối xứng (Symmetric Cryptography).

Mã hóa bất đối xứng sử dụng cặp khóa Công khai/Bí mật, cung cấp mức độ phân quyền rất cao nhưng lại yêu cầu các phép toán số học mô-đun (Modular Arithmetic) khổng lồ trên số nguyên tố lớn, điều này làm vắt kiệt tài nguyên RAM nhỏ bé của MCU ESP32, tiêu tốn hàng giây đồng hồ cho một lần mã hóa, đi ngược lại Ràng buộc thời gian (Latency Constraint) của Điện toán biên. Ngược lại, Mã hóa đối xứng sử dụng chung một khóa bí mật (Secret Key) cho cả thao tác đóng và mở khóa. Cấu trúc toán học của chúng chủ yếu dựa trên các phép toán nhị phân đơn giản như XOR, dịch bit (Shift) và hoán vị (Permutation), mang lại tốc độ thực thi nhanh như chớp và tiêu tốn rất ít chu kỳ lệnh, biến nó thành lựa chọn hoàn hảo nhất cho WSN \cite{roman2011securing}.
\subsubsection{Cấu trúc toán học của chuẩn mã hóa AES-128}
Chuẩn mã hóa tiên tiến (Advanced Encryption Standard - AES) được Viện Tiêu chuẩn và Công nghệ Quốc gia Hoa Kỳ (NIST) công bố là thuật toán mã hóa khối (Block cipher) tiêu chuẩn toàn cầu, dựa trên cơ sở thuật toán Rijndael \cite{daemen2013design}. AES-128 sử dụng một khóa có độ dài 128 bit (16 bytes) để mã hóa một khối dữ liệu thô (Plaintext) cũng có kích thước 128 bit.

Khối dữ liệu 16 bytes đầu vào được tổ chức thành một ma trận trạng thái (State Matrix) $4 \times 4$. Để biến đổi bản rõ thành một bản mã (Ciphertext) hoàn toàn không có tính quy luật thống kê, AES-128 áp dụng quá trình biến đổi lặp đi lặp lại qua 10 vòng (Rounds). Ngoại trừ vòng cuối cùng bị lược bỏ một bước, mỗi vòng lặp thực thi một trình tự gồm 4 hàm toán học phi tuyến tính và tuyến tính đan xen, hoạt động chủ yếu trên nền tảng của toán học Trường hữu hạn Galois $GF(2^8)$ với đa thức bất khả quy $m(x) = x^8 + x^4 + x^3 + x + 1$.

AES-128 sử dụng bốn phép toán chính như trong Hình \ref{Hình 2.8} để thực hiện việc mã hóa và giải mã dữ liệu: SubBytes, ShiftRows, MixColumns và AddRoundKey. Mỗi phép toán này có vai trò quan trọng trong việc bảo vệ dữ liệu khỏi việc bị bóc lột thông tin. SubBytes giúp tạo ra sự thay đổi ngẫu nhiên trong dữ liệu, trong khi ShiftRows giúp phân tán dữ liệu hơn và MixColumns giúp phân tán các thay đổi ngẫu nhiên hơn. Cuối cùng, AddRoundKey giúp xác nhận rằng dữ liệu đã được mã hóa một cách an toàn.

Sự phối hợp chặt chẽ của bốn phép toán này trong mỗi vòng lặp tạo nên cốt lõi bảo mật của thuật toán AES. Chính kiến trúc đa tầng này đã giúp AES an toàn trước hầu hết các phương pháp thám mã hiện đại. 
\begin{figure}[H]
    \centering
    \includegraphics[width=1.0\linewidth]{Thuc_Hanh/Image/aes128.pdf}
    \caption{Biến đổi toán học của một vòng lặp trong thuật toán AES-128}
    \label{Hình 2.8}
\end{figure}

Hàm thay thế byte (SubBytes): Đây là thành phần phi tuyến tính duy nhất của AES. Mỗi byte trong ma trận trạng thái được thay thế bằng một byte khác dựa trên một Bảng tra cứu S-box (Substitution-box). Phép thay thế này thực chất là việc tìm phần tử nghịch đảo nhân của byte đó trong trường $GF(2^8)$, sau đó nhân nó với một ma trận affine và cộng XOR với một hằng số cố định.

Hàm dịch hàng (ShiftRows): Hàm này thực hiện phép hoán vị đơn giản. Hàng đầu tiên của ma trận trạng thái giữ nguyên. Hàng thứ hai dịch vòng (cyclic shift) sang trái 1 byte, hàng thứ ba dịch trái 2 bytes, và hàng thứ tư dịch trái 3 bytes.

Hàm trộn cột (MixColumns): Đây là trung tâm của phép toán Khuếch tán (Diffusion) gây ra hiệu ứng tuyết lở (Avalanche Effect). Mỗi cột của ma trận trạng thái (được xem như một đa thức bậc 3 với các hệ số trong $GF(2^8)$) được nhân với một đa thức tĩnh $a(x) = \{03\}x^3 + \{01\}x^2 + \{01\}x + \{02\} \bmod{x^4+1}$. Phép toán nhân ma trận này có dạng:
\begin{equation}
\begin{bmatrix}
s'_{0,c} \\
s'_{1,c} \\
s'_{2,c} \\
s'_{3,c}
\end{bmatrix}
=
\begin{bmatrix}
02 & 03 & 01 & 01 \\
01 & 02 & 03 & 01 \\
01 & 01 & 02 & 03 \\
03 & 01 & 01 & 02
\end{bmatrix}
\begin{bmatrix}
s_{0,c} \\
s_{1,c} \\
s_{2,c} \\
s_{3,c}
\end{bmatrix}
\end{equation}
% \begin{bmatrix} s'_{0,c} \\ s'_{1,c} \\ s'_{2,c} \\ s'_{3,c} \end{bmatrix} $=$ \begin{bmatrix} 02 & 03 & 01 & 01 \\ 01 & 02 & 03 & 01 \\ 01 & 01 & 02 & 03 \\ 03 & 01 & 01 & 02 \end{bmatrix} \begin{bmatrix} s_{0,c} \\ s_{1,c} \\ s_{2,c} \\ s_{3,c} \end{bmatrix}

Sự kết hợp giữa ShiftRows và MixColumns đảm bảo rằng nếu một bit bất kỳ ở đầu vào bị đổi trạng thái, thì sau 10 vòng, trung bình một nửa số bit ở đầu ra sẽ thay đổi, khiến mọi thuật toán đoán mật khẩu theo tần suất thống kê trở nên vô vọng.

Hàm cộng khóa vòng (AddRoundKey): Ma trận trạng thái hiện tại sẽ được thực hiện phép toán cộng XOR từng bit với Khóa của vòng đó (Round Key). Khóa vòng được sinh ra từ Khóa bí mật (Symmetric Key) ban đầu thông qua Bộ sinh khóa (Key Expansion Schedule).
 
Thông qua việc nhúng thư viện mật mã tiêu chuẩn phần mềm (như mbedtls hoặc thư viện tăng tốc phần cứng bên trong ESP32), các gói tin môi trường và trạng thái FSM từ Node trước khi được đẩy qua module LoRa E32 bằng cổng nối tiếp UART sẽ bị mã hóa thành các đoạn chuỗi thập lục phân hỗn độn vô nghĩa. Ngay cả khi trạm nghe lén bắt được sóng ở băng tần 433 MHz, việc đọc mã là bất khả thi nếu không nắm giữ khóa 128-bit lưu ẩn sâu trong phân vùng tĩnh của thiết bị.
\subsubsection{Cơ chế đồng bộ thời gian thực (Timestamp) và phòng chống tấn công phát lại}
AES-128 giải quyết xuất sắc tính Bảo mật (Confidentiality), nhưng nó bị khuất phục trước kỹ thuật Tấn công phát lại (Replay Attack). Để vá lỗ hổng nghiêm trọng này, một giao thức bảo mật không gian thời gian (Temporal Security) phải được nhúng trực tiếp vào cấu trúc bản tin tại tầng Ứng dụng trước khi bản tin này bị đẩy qua bộ mã hóa AES.
Phương pháp tối ưu và ít tiêu hao băng thông nhất là ứng dụng cơ chế Nhãn thời gian (Timestamp) thay vì sử dụng cơ chế Hỏi-Đáp (Challenge-Response) tốn kém quá nhiều bước bắt tay vô tuyến.

Ngoài ra, cơ chế xác thực dựa trên Timestamp còn giúp hệ thống nâng cao khả năng phát hiện các hành vi bất thường trong mạng truyền thông không dây. Cụ thể, khi Gateway ghi nhận nhiều gói tin có độ lệch thời gian lớn hoặc không hợp lệ trong một khoảng thời gian ngắn, hệ thống có thể kích hoạt các cơ chế cảnh báo hoặc ghi log để phục vụ phân tích an ninh. Điều này đặc biệt quan trọng trong các hệ thống IoT quy mô lớn, nơi các thiết bị hoạt động phân tán và khó kiểm soát trực tiếp.

Bên cạnh đó, việc sử dụng Timestamp kết hợp với thuật toán mã hóa AES-128 còn đảm bảo rằng các gói tin không chỉ được bảo mật về nội dung mà còn đảm bảo tính toàn vẹn và tính mới (freshness). Ngay cả khi tin tặc có khả năng đánh chặn và lưu trữ dữ liệu, việc không thể thay đổi nội dung bên trong gói tin mà không có khóa bí mật khiến cho mọi nỗ lực giả mạo trở nên vô hiệu.

Một ưu điểm khác của phương pháp này là chi phí triển khai thấp và không yêu cầu thay đổi lớn về hạ tầng phần cứng. Các vi điều khiển thông dụng như ESP32 hoặc STM32 hoàn toàn có thể xử lý việc gắn Timestamp và mã hóa AES với độ trễ rất thấp. Đồng thời, IC RTC cung cấp nguồn thời gian ổn định, giúp đảm bảo độ chính xác trong quá trình xác thực.

Tuy nhiên, để hệ thống hoạt động hiệu quả, cần đảm bảo rằng đồng hồ thời gian thực giữa Node và Gateway được đồng bộ ở mức chấp nhận được. Nếu sai lệch thời gian quá lớn, các gói tin hợp lệ cũng có thể bị từ chối. Do đó, trong các triển khai thực tế, cần có cơ chế hiệu chỉnh thời gian định kỳ hoặc sử dụng các giao thức đồng bộ thời gian để duy trì độ chính xác của hệ thống.


\begin{figure}[H]
    \centering
    \includegraphics[width=1.0\linewidth]{Thuc_Hanh/Image/Time-Based Anti-Replay Mechanism in LoRa Communication Using AES-128.pdf}
    \caption{Cơ chế chống Replay Attack dựa trên xác thực dấu thời gian trong LoRa}
    \label{Hình 2.9}
\end{figure}
Hình \ref{Hình 2.9} mô tả chi tiết cơ chế vòng đời của một bản tin an toàn. Ở phần trên của quy trình (Normal Operation), khi Node cảm biến (ESP32) muốn báo cáo dữ liệu môi trường, nó sẽ tiến hành truy vấn IC đồng hồ thời gian thực (RTC DS3231 hoặc bộ định thời cục bộ) để lấy ra một con số Nhãn thời gian hệ thống (T\_{send}, ví dụ đo bằng kỷ nguyên Unix Epoch time). Gói tin mất vài trăm mili-giây bay qua không trung và đến Gateway sau đó Gateway giải mã, thấy nhãn thời gian T\_{send} gửi đi lúc 14:05:10, và thời gian đồng hồ hiện tại của nó là 14:05:11. Độ lệch $\Delta t = 1 \text{ giây}$. Vì $\Delta t$ nằm gọn trong "Cửa sổ an toàn" (Ví dụ: 5 giây), Gateway tin tưởng tuyệt đối đây là gói tin mới (Freshness) và thi hành lệnh.

Ở nửa dưới của sơ đồ, kịch bản tấn công được giả lập. Tên tin tặc dùng phần cứng bắt gọn luồng sóng vô tuyến. Hắn chờ 10 phút sau mới bơm ngược gói tin đó vào mạng nhằm cố ý đảo lộn dữ liệu làm mát. Tuy nhiên, Gateway khi giải mã gói tin rác này vẫn sẽ nhìn thấy cái nhãn thời gian cũ mốc từ 10 phút trước nằm nguyên vẹn bên trong (vì hacker không biết mã AES nên không thể cập nhật nhãn thời gian). Phép trừ toán học $\Delta t$ sẽ lập tức đưa ra con số sai lệch lớn (591 giây). Hệ thống Gateway sẽ đưa ra ngoại lệ (Exception), lập tức bác bỏ gói tin (Drop Packet) và chặn đứng hoàn toàn âm mưu phá hoại. 
\subsection{Kiến trúc và cơ chế cập nhật phần mềm không dây (OTA) trên nền tảng Edge IoT}

Trong bối cảnh các hệ thống mạng cảm biến không dây (WSN - Wireless Sensor Networks) được triển khai tại những vị trí vật lý khắc nghiệt và khó tiếp cận như môi trường cách ly sinh học của chuồng trại chăn nuôi, việc duy trì vòng đời phần mềm (Software Development Life Cycle) thông qua kết nối cáp vật lý (UART/USB) là hoàn toàn bất khả thi. Để giải quyết rào cản này, công nghệ Over-The-Air (OTA) được áp dụng như một cơ chế cốt lõi. 

Về mặt cấu trúc, bộ nhớ Flash SPI ngoại vi phải được phân chia thành Dual-Bank Architecture. Như được trình bày trong Hình \ref{fig:ota_memory_map}, bộ nhớ được chia thành các vùng độc lập: Vùng nạp khởi động (Bootloader), vùng dữ liệu trạng thái (\texttt{otadata}) và ít nhất hai phân vùng ứng dụng (\texttt{APP0} và \texttt{APP1}). Cơ chế này đảm bảo tính nguyên vẹn của hệ thống: tại bất kỳ thời điểm nào, mã lệnh thực thi luôn chạy trên một phân vùng ổn định, trong khi phân vùng còn lại (đang ở trạng thái nhàn rỗi) sẽ được sử dụng làm bộ đệm để ghi khối dữ liệu phần mềm mới tải về.
\begin{figure}[H]
    \centering
    \includegraphics[width=1.0\linewidth]{Thuc_Hanh/Image/Flash (Partition Table).pdf}
    \caption{Giản đồ phân bổ không gian nhớ Flash cho kiến trúc cập nhật Dual-Bank OTA trên ESP32-S3}
    \label{fig:ota_memory_map}
\end{figure}
% \begin{figure}[H]
% \centering
% \begin{tikzpicture}[font=\small, >=stealth, node distance=0.5cm]
%     % Styles
%     \tikzstyle{block} = [rectangle, draw, thick, minimum width=6cm, minimum height=0.8cm, align=center]
%     \tikzstyle{active} = [fill=green!20]
%     \tikzstyle{idle} = [fill=blue!10]
%     \tikzstyle{system} = [fill=gray!20]
    
%     % Memory Map
%     \node[block, system] (boot) {Bootloader (0x1000 - 0x8000)};
%     \node[block, system, below=of boot] (part) {Partition Table (0x8000 - 0x9000)};
%     \node[block, system, below=of part] (otadata) {OTA Data / NVS (0x9000 - 0x10000)};
%     \node[block, active, below=of otadata, minimum height=1.5cm] (app0) {\textbf{APP0 Partition (Factory / Current)}\\ (Trạng thái: Đang thực thi mã lệnh)};
%     \node[block, idle, below=of app0, minimum height=1.5cm] (app1) {\textbf{APP1 Partition (OTA Update)}\\ (Trạng thái: Ghi tuần tự firmware mới)};
    
%     % Arrows and labels
%     \draw[->, thick, red] (otadata.west) -- ++(-1.5,0) |- node[pos=0.25, left, align=center] {Cờ Boot trỏ\\vào APP0} (app0.west);
    
%     \draw[->, thick, blue, dashed] (app1.east) -- ++(1.5,0) |- node[pos=0.25, right, align=center] {Dữ liệu Firmware mới\\từ mạng Wi-Fi ghi vào} (app1.east);
% \end{tikzpicture}
% \caption{Giản đồ phân bổ không gian nhớ Flash (Partition Table) cho kiến trúc cập nhật Dual-Bank OTA trên ESP32-S3}
% \label{fig:ota_memory_map}
% \end{figure}

Sự thách thức lớn nhất khi áp dụng OTA trong mạng cảm biến nông nghiệp của đề tài là nghịch lý về băng thông và năng lượng. Giao thức truyền thông chính của hệ thống là LoRa (Long Range), sở hữu ưu điểm vượt trội về khả năng xuyên thấu và cự ly truyền sóng lên tới hàng kilomet với công suất phát chỉ 20 dBm. Tuy nhiên, tốc độ dữ liệu (Data Rate) của mạng LoRa cực thấp, thường dao động ở mức vài trăm bps. Dung lượng của một tệp tin Firmware chuẩn chứa cả ngăn xếp TCP/IP và thuật toán điều khiển biên thường xấp xỉ 1.2 MB. Nếu dùng kênh LoRa để tải tệp tin này, quá trình cập nhật sẽ kéo dài nhiều ngày, chiếm dụng hoàn toàn băng thông mạng (Duty Cycle) và vắt kiệt năng lượng pin. 

Để hóa giải nút thắt vật lý này, đồ án đã tiến hành thiết kế và lập trình kiến trúc cập nhật phần mềm Hybrid OTA. Cơ chế này sử dụng song song hai chuẩn truyền thông: LoRa làm kênh điều khiển (Control Channel) và Wi-Fi làm kênh truyền tải dữ liệu (Data Channel). Trạng thái mặc định, module Wi-Fi bị vô hiệu hóa hoàn toàn ở mức phần cứng (RF Disabled) để đảm bảo dòng tĩnh của hệ thống duy trì ở mức micro-Ampe.
% \begin{figure}[H]
% \centering
% \includegraphics[]{}
% \caption{Lưu đồ thuật toán điều khiển quá trình Hybrid OTA tích hợp cơ chế bảo vệ Rollback}
% \label{fig:hybrid_ota_flow}
% \end{figure}
\begin{figure}[H]
    \centering
    \includegraphics[width=1.0\linewidth]{Thuc_Hanh/Image/OTA (1).pdf}
    \caption{Lưu đồ thuật toán điều khiển quá trình Hybrid OTA}
    \label{fig:hybrid_ota_flow}
\end{figure}
Như thể hiện trong Hình \ref{fig:hybrid_ota_flow}, ngay sau khi kết nối vào mạng Wi-Fi nội bộ của trang trại, tiến trình tải được thực thi theo phương thức truyền khối (Chunked Transfer Encoding) qua giao thức HTTPS bảo mật. Dữ liệu được ghi trực tiếp vào phân vùng tĩnh thay vì nạp lên RAM nhằm tránh lỗi tràn bộ nhớ (Heap Overflow). Đặc biệt, tính an toàn của hệ thống được đặt lên hàng đầu thông qua cơ chế tự động phục hồi (Auto-Rollback). Nếu bản phần mềm mới tải về chứa lỗi logic (vòng lặp vô hạn, lỗi cấp phát bộ nhớ) khiến hệ điều hành treo cứng, bộ định thời Hardware Watchdog sẽ tự động kích hoạt ngắt Reset. Tại lần khởi động tiếp theo, Bootloader phát hiện cờ \texttt{ESP\_OTA\_IMG\_PENDING\_VERIFY} chưa được gỡ bỏ, nó sẽ tự động đảo ngược con trỏ hệ thống về bản Firmware ổn định cũ. 

\subsection{Khảo sát các công trình nghiên cứu liên quan}
\subsubsection{Kiến trúc Điện toán biên và Phân bổ tài nguyên}
Trong bài báo mang tính nền tảng "Edge Computing: Vision and Challenges" \cite{shi2016edge}, W. Shi và các cộng sự đã chính thức hóa định nghĩa về Điện toán biên. Đóng góp lớn nhất của nghiên cứu này là việc đưa ra khái niệm "Computing Stream" (Luồng tính toán), trong đó dữ liệu không còn chảy một chiều từ thiết bị lên đám mây, mà bị đánh chặn và xử lý tại các "nút thắt" ở rìa mạng. 
\begin{figure}[H]
    \centering
    \includegraphics[width=0.7\linewidth]{Thuc_Hanh/Image/edge Computing diagram.png}
    \caption{Sơ đồ hoạt động của hệ thống điện toán biên trong bài báo}
    \label{Hình 2.11}
\end{figure}
Nhóm tác giả đã sử dụng các kịch bản thử nghiệm giả lập (Case studies) trong nhà thông minh (Smart Home) để đo lường độ trễ. Kết quả chỉ ra rằng, việc xử lý cục bộ làm giảm băng thông tiêu thụ trên mạng trục lõi (Backbone network) lên tới 80\%, đồng thời cắt giảm độ trễ phản hồi từ vài giây xuống mức mili-giây. Tuy nhiên, giới hạn của công trình này là chỉ dừng ở mức lý thuyết vĩ mô, chưa đề xuất các cấu trúc vi mạch cụ thể để triển khai.

Đi sâu hơn vào bài toán tối ưu hóa, Y. Mao và các cộng sự trong công trình "A Survey on Mobile Edge Computing: The Communication Perspective" \cite{mao2017survey} đã xây dựng một nền tảng toán học đồ sộ cho vấn đề Offloading. Phương pháp nghiên cứu của họ là áp dụng lý thuyết tối ưu hóa lồi (Convex Optimization) và Lý thuyết trò chơi (Game Theory) để giải quyết hàm đa mục tiêu: giảm thiểu năng lượng tiêu thụ tại thiết bị di động trong khi vẫn đáp ứng được ràng buộc về thời gian chết (Deadline) của tác vụ. Họ phân loại kiến trúc ra làm hai nhánh: Partial Offloading (cắt nhỏ thuật toán, chạy một nửa ở Node, gửi một nửa lên Edge Server) và Binary Offloading (chọn hoặc chạy 100\% ở Node, hoặc 100\% ở Server). Quá trình mô phỏng bằng Matlab chỉ ra rằng, đối với các thiết bị IoT có lượng dữ liệu cảm biến nhỏ nhưng yêu cầu độ trễ điều khiển máy móc khắt khe, chiến lược Binary Offloading (Thực thi toàn bộ tại Node) mang lại hiệu năng cao nhất do tránh được chi phí thời gian thiết lập kênh truyền vô tuyến.

Các nghiên cứu chỉ ra rằng kiến trúc Điện toán biên là bắt buộc để khắc phục độ trễ. Chiến lược Offloading nhị phân (chạy hoàn toàn tại MCU của Node) là phương pháp tối ưu. Tuy nhiên, việc áp dụng Deep Learning là quá sức đối với MCU chạy pin, định hướng cho chúng ta phải tìm kiếm các thuật toán biên nhẹ nhàng hơn (như Máy trạng thái hữu hạn - FSM).
\subsubsection{Đồng thiết kế phần cứng và Tối ưu hóa năng lượng cảm biến}
Công trình mang tính đột phá nhất liên quan trực tiếp đến cảm biến khí là nghiên cứu của F. Palacio và cộng sự: "Pulsed-Temperature Metal Oxide Gas Sensors for Microwatt Power Consumption" \cite{palacio2020pulsed}. Cảm biến khí oxit kim loại bán dẫn (MOX - như dòng MQ-135 dùng để đo $NH_3$) hoạt động dựa trên nguyên lý thay đổi điện trở bề mặt khi tiếp xúc với khí độc. Tuy nhiên, để phản ứng hóa học xảy ra, lớp màng nhạy khí ($SnO_2$) bắt buộc phải được nung nóng liên tục ở nhiệt độ từ $300^\circ C$ đến $400^\circ C$ bằng một cuộn dây may-so bên trong. Phương pháp thiết kế truyền thống luôn cấp điện liên tục (Continuous Heating) cho cuộn dây này, tiêu tốn một công suất tĩnh khổng lồ khoảng $150 \text{ mW}$ đến $800 \text{ mW}$, hút cạn bất kỳ viên pin nào chỉ trong 2-3 ngày.

Để giải quyết, nhóm của Palacio đã đề xuất phương pháp Điều chế nhiệt độ dạng xung (Pulsed-Temperature). Thay vì nung liên tục, họ sử dụng một công tắc bán dẫn điều khiển bằng xung PWM từ vi điều khiển. Cảm biến sẽ bị ngắt điện hoàn toàn trong thời gian dài (ví dụ 10 phút). Khi cần lấy mẫu, MCU sẽ "đánh thức" cảm biến, bơm một dòng xung cường độ cao trong một khoảng thời gian cực ngắn (chỉ khoảng $50 \text{ ms}$) để làm nóng lớp màng lên $400^\circ C$, đọc giá trị ADC, và ngay lập tức ngắt điện (Power-off). Đánh giá thực nghiệm bằng máy đo nguồn Keithley 2450 cho thấy, phương pháp này giảm công suất tiêu thụ trung bình từ $150 \text{ mW}$ xuống mức siêu vi watt ($16 \mu W$), đồng thời việc tạo sốc nhiệt còn giúp cảm biến loại bỏ khí tồn dư, tăng độ nhạy (Selectivity).
\begin{figure}[H]
    \centering
    \includegraphics[width=1.0\linewidth]{Thuc_Hanh/Image/Low-Power Sensing Cycle.pdf}
    \caption{Chu trình vận hành tiết kiệm năng lượng}
    \label{Hình 2.10}
\end{figure}
Lưu đồ trong Hình \ref{Hình 2.10} mô hình hóa một chu trình sống (Life-cycle) của Node dựa trên triết lý tối ưu năng lượng từ bài báo của F. Palacio. Mấu chốt của quy trình nằm ở khối "Pulsed Heating". Thay vì nối cứng chân VCC của cảm biến khí MQ-135 vào thẳng nguồn pin, hệ thống sử dụng một chân GPIO của ESP32 để kích mở cực cổng (Gate) của một Transistor (MOSFET). Chỉ khi hệ thống thức dậy sau 5 phút ngủ sâu, MOSFET mới đóng mạch cho dòng điện lớn đi qua cuộn nung của cảm biến. Thời gian nung nóng (Pre-heating time) được tinh chỉnh ở mức tối thiểu (khoảng 200ms - vừa đủ để nhiệt độ màng $SnO_2$ đạt đỉnh phản ứng). Ngay sau khi bộ chuyển đổi Tương tự-Số (ADC) chốt được giá trị điện áp, MOSFET lập tức bị ngắt, cô lập cảm biến khỏi nguồn điện. Cơ chế "đánh nhanh rút gọn" này biến một linh kiện vốn ngốn điện khủng khiếp thành một thiết bị tương thích với chuẩn WSN dùng pin.

Bên cạnh đó, nghiên cứu "Low-Power On-Chip Implementation of Enhanced SVM Algorithm..." \cite{chang2022low} của J. Chang chứng minh rằng ta có thể tích hợp thuật toán Học máy cơ bản (Support Vector Machine - SVM) lên các MCU cấp thấp. Họ đã tối ưu hóa hàm Kernel của SVM để nó chỉ tiêu tốn 10KB RAM. Tuy nhiên, việc chạy tính toán dấu phẩy động ma trận vẫn làm CPU mất vài chục mili-giây, tiêu tốn năng lượng hơn việc dùng cấu trúc rẽ nhánh. Cuối cùng, công trình "Hardware Solutions for Low-Power Smart Edge Computing" \cite{wisniewski2022hardware} đã phát triển nền tảng CYSmart, tập trung vào việc áp dụng các IC ổn áp (LDO) có dòng rò tĩnh dưới $1 \mu A$ và sử dụng Sleep mode của phần cứng để đánh bại rào cản năng lượng.

Những khám phá từ nhóm này là kim chỉ nam cho thiết kế mạch nguyên lý. Việc sử dụng MOSFET để cắt nguồn cảm biến (Load Switch) kết hợp cơ chế nung nóng dạng xung (Pulsed Heating) cho MQ-135, kết hợp với các bộ ổn áp tĩnh thấp là những ý tưởng cực kỳ xuất sắc và hoàn toàn khả thi để ứng dụng ngay vào đồ án.

\subsubsection{Mô hình hóa năng lượng truyền thông vô tuyến LoRa}
Việc đưa dữ liệu biên về trung tâm yêu cầu một giao thức không dây tiêu tốn ít năng lượng nhất có thể.

Trong bài báo "Energy Consumption Model for Sensor Nodes Based on LoRa and LoRaWAN" \cite{bouguera2018energy}, nhóm tác giả T. Bouguera đã xây dựng một mô hình toán học giải phẫu chi tiết toàn bộ hồ sơ năng lượng của một Node LoRa. Họ sử dụng nền tảng vi điều khiển STM32L0 kết hợp với IC vô tuyến SX1272. Bằng cách sử dụng Máy phân tích điện năng (Agilent DC Power Analyzer), họ đã đo đạc chính xác dòng điện tiêu thụ ở 4 trạng thái của chip vô tuyến: Sleep (Ngủ), Standby (Chờ), TX (Phát) và RX (Thu).
Nghiên cứu đưa ra công thức tổng quát tính điện năng tiêu thụ cho một gói tin:
\begin{equation}
    E_{packet} = P_{wakeup} \times T_{wakeup} + P_{TX} \times ToA + P_{RX} \times T_{RX} + P_{sleep} \times T_{sleep}
\end{equation}
Phát hiện mang tính bước ngoặt của họ là: việc cấu hình Hệ số trải phổ (Spreading Factor - SF) tác động theo cấp số nhân đến Thời gian bay của sóng (Time on Air - ToA). Ví dụ, truyền một gói tin 50 byte ở SF7 mất $100 \text{ ms}$ với dòng điện $40 \text{ mA}$. Nhưng nếu đổi sang cấu hình SF12 để xuyên tường tốt hơn, gói tin đó sẽ lơ lửng trên không trung mất gần $2 \text{ giây}$, rút cạn một lượng năng lượng lớn gấp 20 lần từ viên pin. Từ đó, họ đề xuất rằng các mạng cảm biến tĩnh không nên mù quáng cài đặt SF tối đa, mà phải có cơ chế cấu hình thích ứng (Adaptive Data Rate) dựa trên khoảng cách. Nếu khoảng cách dưới 500m (Line-of-sight), bắt buộc phải khóa cứng cấu hình ở SF7 hoặc SF8 để bảo toàn pin.


\subsection{Kết luận chương}

Chương 2 đã thiết lập nền tảng lý thuyết theo một logic nhân quả nhất
quán, trong đó mỗi lớp công nghệ được lựa chọn xuất phát từ ràng buộc
của lớp trước.

Đặc tính sinh lý của gia cầm -- không có tuyến mồ hôi, tản nhiệt hoàn
toàn qua hô hấp -- biến chỉ số THI thành thông số sống còn và áp đặt
yêu cầu phản ứng tức thời lên toàn bộ hệ thống. Yêu cầu đó loại bỏ mô
hình Cloud tập trung khỏi phương án khả thi, dẫn đến kiến trúc Điện
toán biên với chiến lược Offloading nhị phân -- toàn bộ logic phán quyết
thực thi ngay tại Node. Để đảm bảo tính tất định cho vòng điều khiển, FreeRTOS với cơ chế ghim lõi (Core Pinning) trên ESP32 là
giải pháp kiến trúc phần mềm bắt buộc.

Song song với bài toán tính toán là bài toán năng lượng. Mô hình toán
học tiêu thụ pin theo Hệ số trải phổ SF chỉ ra rằng tăng SF từ 7 lên 12
kéo dài thời gian phát sóng gần 20 lần, buộc hệ thống phải khóa cứng ở
SF7/SF8 thay vì tối đa hóa chống nhiễu một cách mù quáng. Kết hợp với
kỹ thuật nung nóng xung cho cảm biến khí gas và chế độ ngủ sâu lai, đây
là nền tảng để kéo dài tuổi thọ pin ở mức chấp nhận được cho triển khai
thực địa.

Lớp bảo mật AES-128 kết hợp xác thực Timestamp khép kín thiết kế bằng
cách bịt kín lỗ hổng Replay Attack của mạng vô tuyến mở, trong
khi cơ chế OTA lai đảm bảo khả năng bảo trì từ xa mà không cần can
thiệp vật lý. Toàn bộ cơ sở lý luận này dẫn dắt trực tiếp các quyết
định thiết kế phần cứng và phần mềm được trình bày từ Chương 3 trở đi.
\newpage

\section*{CHƯƠNG 3. PHÂN TÍCH VÀ THIẾT KẾ HỆ THỐNG}
\addcontentsline{toc}{section}{\numberline {} CHƯƠNG 3. PHÂN TÍCH VÀ THIẾT KẾ TỔNG THỂ HỆ THỐNG}
\setcounter{section}{3}
\setcounter{subsection}{0}
\setcounter{figure}{0}
\setcounter{table}{0}
Chương này sẽ tiến hành phân tách bài toán từ trên xuống dưới. Bắt đầu từ việc phân tích sâu sắc các yêu cầu chức năng và phi chức năng để vạch ra biên giới hoạt động của hệ thống, tiếp đó là biện luận lựa chọn các giải pháp công nghệ phần cứng tối ưu nhất. Cuối cùng, hệ thống sẽ được kiến trúc hóa thông qua các sơ đồ khối, kiến trúc phần mềm đa nhiệm trên FreeRTOS và các lưu đồ thuật toán Điện toán biên. 
\subsection{Phân tích yêu cầu hệ thống}
Để đáp ứng mục tiêu nghiên cứu và giải quyết triệt để bài toán thực tiễn tại các trang trại, thiết bị Node phải thỏa mãn các yêu cầu khắt khe được lượng hóa bằng các tham số cụ thể sau:
\subsubsection{Yêu cầu chức năng}
\begin{itemize}
    \item Thu thập và giám sát vi khí hậu: Node phải quét và lấy mẫu liên tục dữ liệu từ tổ hợp cảm biến với sai số đo lường dưới 5\%. Cụ thể, hệ thống phải theo dõi sát sao biểu đồ nhiệt - ẩm và nhận diện ngay lập tức khi môi trường chạm các ngưỡng nguy hiểm: nhiệt độ vượt 35°C (gây stress nhiệt) , độ ẩm trên 75\% , và nồng độ khí độc Amoniac (NH3) vượt ngưỡng 25 ppm.
    \item Điện toán biên (Edge Computing): Tích hợp Máy trạng thái hữu hạn (FSM) tại vi điều khiển để tự động đối chiếu dữ liệu thô với ma trận ngưỡng an toàn. Khi phát hiện các tham số vượt ngưỡng (ví dụ: nhiệt độ > 35°C), thuật toán biên phải tự động kích hoạt ngõ ra GPIO để đóng Rơ-le điều khiển hệ thống làm mát cục bộ mà không cần chờ lệnh từ Gateway hay Cloud.
    \item Yêu cầu cốt lõi ở đây không chỉ là tốc độ phản hồi
mà còn là tính tự chủ của vòng lặp điều khiển -- luồng ra
quyết định phải được khép kín hoàn toàn bên trong Node, hoạt động độc
lập với trạng thái của kênh truyền vô tuyến bên ngoài. Điều đó có nghĩa
là ngay cả khi Gateway mất điện, đường truyền LoRa bị gián đoạn hay
kết nối Internet hoàn toàn đứt gãy, FSM vẫn tiếp tục chu kỳ giám sát
và sẵn sàng kích hoạt cơ cấu chấp hành để bảo vệ vật nuôi trong thời
gian thực. Tính chất này phân biệt rõ ràng kiến trúc đề xuất so với các
giải pháp thương mại phụ thuộc Cloud. Về mặt triển khai phần mềm, FSM
được hiện thực hóa như một tác vụ độc lập chạy trên lõi riêng biệt
(Core Pinning) trong FreeRTOS, đảm bảo không bị gián đoạn bởi các luồng
xử lý truyền thông hay mã hóa đang diễn ra song song. Ràng buộc thời
gian tất định dưới 50 ms được bảo đảm nhờ cơ chế ưu tiên tác vụ
(Preemptive Scheduling) của hệ điều hành thời gian thực.

    \begin{figure}[H]
    \centering
    \includegraphics[width=1.0\linewidth]{Thuc_Hanh/Image/Edge Node System Context.pdf}
    \caption{Edge Node System Context}
    \label{Hình 3.1}
    \end{figure}
    Hình \ref{Hình 3.1} mô tả sơ đồ hệ thống hóa các luồng tương tác giữa thực thể môi trường vật lý (Microclimate, HVAC) và các khối chức năng nội tại của Node biên. Khối FSM Edge Inference đóng vai trò là "Bộ não trung tâm", tách biệt hoàn toàn luồng điều khiển cơ điện (hướng xuống HVAC) khỏi luồng truyền tải dữ liệu (hướng sang mã hóa và LoRa). Cấu trúc rẽ nhánh từ khối LoRa đến khối Storage minh họa hành vi tự nhận thức lỗi của hệ thống: việc ghi log ngoại tuyến không phải là tác vụ chạy liên tục, mà chỉ được đánh thức khi kênh truyền vô tuyến gặp sự cố gián đoạn.
    
    \item Truyền thông vô tuyến LoRa: Đóng gói dữ liệu Payload (định mức 32-bytes)  và phát sóng định kỳ. Cấu hình phần cứng phải chốt ở tần số 433 MHz, Spreading Factor (SF) = 7, Bandwidth (BW) = 125 kHz và Coding Rate = 4/5 để cân bằng giữa băng thông và tầm xa.
    
    \item Lưu trữ ngoại tuyến (Offline Storage): Xây dựng cơ chế phát hiện rớt mạng cục bộ. Khi không nhận được gói tin phản hồi (ACK) từ Gateway, Node phải tự động đẩy dữ liệu log vào vùng nhớ tĩnh (Flash/PSRAM) và đồng bộ lại ngay khi sóng vô tuyến được khôi phục.
    \item Bảo mật và tính toàn vẹn dữ liệu: Mã hóa toàn bộ khối dữ liệu tải trọng bằng thuật toán đối xứng AES-128. Đồng thời, tích hợp IC thời gian thực (RTC DS3231) để đính kèm nhãn thời gian (Timestamp) vào từng gói tin , thiết lập cửa sổ thời gian chặn Replay Attack với độ lệch $\Delta{t}$ < 5s.
\end{itemize}
\subsubsection{Yêu cầu phi chức năng}
\begin{itemize}
    \item Độ trễ xử lý Điện toán biên (Edge Latency): Thời gian đáp ứng của luồng điều khiển (tính từ thời điểm cảm biến xuất tín hiệu báo vượt ngưỡng đến khi chân GPIO đóng mạch Rơ-le) bắt buộc phải nhỏ hơn 50 mili-giây.
    \item Hiệu suất năng lượng và Tuổi thọ pin (Power Profiling): Vận hành trên nguồn pin Li-ion 18650 (dung lượng 3200mAh). Hệ thống phải quản lý dòng tiêu thụ thông qua chế độ Deep Sleep, duy trì dòng rò ở trạng thái chờ dưới 20µA. Với chu kỳ lấy mẫu 5 phút/lần, Node phải có khả năng hoạt động liên tục hơn 3 tháng mà không cần sạc lại.
    \item Chất lượng đường truyền (Link Quality): Phủ sóng vô tuyến ổn định trong bán kính 500m đến 1000m ở điều kiện không vật cản (Line-of-Sight). Trong môi trường chuồng trại có tường bê tông che khuất (Non-Line-of-Sight), Tỷ lệ mất gói tin (Packet Loss Rate - PLR) phải được duy trì dưới 5\% ở mốc khoảng cách 200m.
    \item An toàn chức năng (Functional Safety): Hệ thống phải hoạt động ổn định 24/7. Yêu cầu bắt buộc phải cấu hình Hardware Watchdog Timer để tự động reset vi điều khiển nếu xảy ra hiện tượng treo bus I2C hoặc rò rỉ vùng nhớ heap (Memory Leak).
\end{itemize}
\begin{figure}[H]
    \centering
    \includegraphics[width=1.0\linewidth]{Thuc_Hanh/Image/SystemRequire.pdf}
    \caption{Edge IoT System Requirements and Actors}
    \label{Hình 3.2}
\end{figure}
Hình \ref{Hình 3.2} mô tả sơ đồ yêu cầu hệ thống và các tác nhân (Edge IoT System Requirements \& Actors) của một hệ thống IoT biên (Edge IoT) dùng để giám sát vi khí hậu và điều khiển thông gió trong môi trường chăn nuôi hoặc nhà kính.

Hệ thống bắt đầu từ các yếu tố môi trường (Microclimate Environment) như nồng độ NH3, nhiệt độ và độ ẩm. Các thông số này được thu thập bởi khối Data Acquisition, sử dụng các giao tiếp như I2C hoặc ADC để đọc dữ liệu từ cảm biến. Để tiết kiệm năng lượng, khối này được hỗ trợ bởi cơ chế Pulsed Heating nhằm đánh thức cảm biến khi cần đo đạc.

Dữ liệu thô sau khi thu thập sẽ được chuyển đến khối Edge Computing FSM (Finite State Machine) để xử lý cục bộ tại thiết bị biên. Khối này thực hiện các thuật toán ra quyết định với yêu cầu độ trễ thấp (nhỏ hơn 50 ms), đồng thời đưa ra các hành động điều khiển đối với hệ thống thông gió (Ventilation System) nhằm duy trì điều kiện môi trường phù hợp.

Trước khi truyền đi, dữ liệu được mã hóa bằng AES-128 Encryption để đảm bảo tính bảo mật. Sau đó dữ liệu mã hóa được gửi thông qua mô-đun LoRa TX/RX bằng giao tiếp SPI hoặc UART. Dữ liệu được truyền qua liên kết vô tuyến Sub-GHz RF đến Gateway hoặc Cloud, nơi thực hiện chức năng lưu trữ, giám sát và phân tích dữ liệu.

Để tăng độ tin cậy của hệ thống, khi kết nối mạng hoặc truyền thông LoRa gặp sự cố, cơ chế Offline SPIFFS Storage sẽ được kích hoạt để lưu trữ dữ liệu tạm thời trên bộ nhớ cục bộ, bảo đảm không bị mất dữ liệu và hỗ trợ vận hành theo chế độ dự phòng (Fail-safe).

\begin{table}[H]
\centering
\caption{Định lượng hóa các Yêu cầu Phi chức năng (Non-Functional Specifications)}
\label{Bảng 3.1}
\renewcommand{\arraystretch}{1.5}
\begin{tabular}{p{4.0cm} p{6.5cm} p{4cm}}
\hline
\textbf{Nhóm Yêu cầu} & \textbf{Tham số Định lượng Thiết kế} & \textbf{Kiểm chứng } \\ \hline
\textbf{ Thời gian } & Thời gian đáp ứng từ Sensor $\rightarrow$ Relay $< 50\text{ ms}$. & Đo bằng xung Trigger. \\ 
\textbf{Năng lượng} & Dòng tĩnh Deep Sleep $< 20\mu\text{A}$. Dòng đỉnh Active $< 150\text{ mA}$. & Đo bằng đồng hồ vạn năng. \\ 
\textbf{An toàn và chịu lỗi} & Tự động phục hồi khi I2C treo $< 2\text{ giây}$. Watchdog Timeout = $2\text{s}$. & Ép chập mạch I2C cứng (Short circuit) để test Reset. \\ 
\textbf{Cơ khí} & Vỏ khí động học, màng lọc PTFE. Chuẩn IP65, mạch FR4 phủ keo. & Test ngâm sương ẩm $85\%$ trong $48\text{ giờ}$. \\ \hline
\end{tabular}
\end{table}
Bảng dữ liệu phi chức năng như Bảng \ref{Bảng 3.1} là tiêu chí cho quá trình lựa chọn linh kiện vật lý. Để đạt được dòng rò $< 20\mu\text{A}$, các kỹ sư phải rà soát từng con điện trở và dòng Quiescent của IC nguồn. Để đạt độ trễ $< 50\text{ ms}$, CPU phải duy trì xung nhịp cao ở thời điểm Active. Cột phương pháp kiểm chứng đảm bảo mọi thiết kế trên giấy đều sẽ được đối chứng bằng phần cứng thực tế.
\subsubsection{Yêu cầu kỹ thuật và thiết kế phần cứng}
\begin{itemize}
    \item Mạch nguyên lý và khối xử lý cốt lõi: Sử dụng vi điều khiển có năng lực xử lý đa nhân làm khối trung tâm, kết hợp với module truyền thông vô tuyến tầm xa dải tần thấp (Sub-GHz) thông qua giao thức truyền thông nối tiếp tiêu chuẩn. Khối quản lý năng lượng yêu cầu tích hợp mạch sạc và linh kiện ổn áp tuyến tính (LDO) có dòng tĩnh (Quiescent current) cực thấp. 
    \item Thiết kế mạch in (Custom PCB): Thiết kế bo mạch in tiêu chuẩn công nghiệp 2 lớp. Quá trình quy hoạch đường mạch (Layout) phải tuân thủ nghiêm ngặt các quy tắc chống nhiễu điện từ (EMI). Cụ thể, cần thiết lập các vùng cách ly (Keep-out zone) không phủ đồng tại khu vực ăng-ten vô tuyến để tránh suy hao phát xạ, đồng thời phân tách vật lý rõ ràng giữa vùng tín hiệu số (Digital) tốc độ cao và đường tín hiệu tương tự (Analog) nhạy cảm để triệt tiêu nhiễu chéo. Bo mạch sau khi hoàn thiện gia công và hàn linh kiện dán (SMD) phải được phủ lớp bảo vệ chuyên dụng để chống ô xi hóa trong môi trường độ ẩm cao và nhiều khí ăn mòn.
\end{itemize}

\subsection{Thiết kế kiến trúc và lựa chọn giải pháp công nghệ}
Từ việc phân tích cặn kẽ các ràng buộc hệ thống, quá trình ra quyết định để lựa chọn linh kiện và công nghệ đòi hỏi sự cân nhắc giữa hiệu năng xử lý, khả năng quản lý năng lượng và tính khả thi trong thiết kế.
\subsubsection{Kiến trúc tổng thể hệ thống}
Sơ đồ kiến trúc tổng thể của toàn bộ hệ thống được thể hiện trong Hình \ref{Hình 3.3}
\begin{figure}[H]
    \centering
    \includegraphics[width=1.0\linewidth]{Thuc_Hanh/Image/KienTrucTongThe.pdf}
    \caption{Kiến trúc tổng thể của hệ thống}
    \label{Hình 3.3}
\end{figure}
Kiến trúc tổng thể của hệ thống giám sát vi khí hậu được thiết kế theo mô hình phân lớp (Layered Architecture) tiêu chuẩn của IoT công nghiệp. Hình ảnh minh họa một hệ sinh thái khép kín từ khâu thu thập dữ liệu vật lý đến khâu tương tác với người dùng cuối, được chia thành 5 phân lớp (Layer) chuyên biệt. Việc phân mảnh cấu trúc này không chỉ giúp cô lập lỗi (Fault Isolation) mà còn tối ưu hóa tính linh hoạt và khả năng mở rộng của toàn hệ thống.
\subsubsection{Thiết kế sơ đồ khối tổng thể của Node}
Sau khi đã phân tích lựa chọn cho từng khối chức năng riêng lẻ, các khối này được tích hợp lại thành một kiến trúc phần cứng hoàn chỉnh của thiết bị Node. Sơ đồ khối tổng thể trong Hình \ref{fig:ch3_block} thể hiện rõ ràng năm phân hệ chính và cách chúng trao đổi tín hiệu với khối xử lý trung tâm ESP32-S3.

% ===================================================================
% PlantUML cho Hình 3.8 - render và lưu thành: Thuc_Hanh/Image/Ch3_BlockDiagram.png
% @startuml
% skinparam componentStyle rectangle
% skinparam backgroundColor white
% rectangle "KHOI NGUON\nUSB-C -> TP4056 (sac)\n-> AO3401A -> AP2112K (LDO 3V3)\nPin Li-ion 18650" as PWR
% rectangle "ESP32-S3-WROOM-1\n(Loi kep LX7 240MHz)\n- FreeRTOS\n- FSM Edge\n- AES-128/ASCON" as MCU
% rectangle "CAM BIEN VI KHI HAU\nSHT31 / BMP180 / BH1750\n(I2C)" as ENV
% rectangle "CAM BIEN KHI\nMQ-135 (A0) / CH4 (ADC)\nCO2 (UART)" as GAS
% rectangle "RTC DS3231\n(I2C + SQW)" as RTC
% rectangle "LoRa SX1278\n(SPI, 433MHz)" as LORA
% rectangle "KHOI CHAP HANH\nAct1/Act2/Act3 -> Ro-le\n(Quat/Phun suong/Suoi)" as ACT
% rectangle "RGB LED + Nut BOOT/RESET/USER" as UI
% PWR --> MCU : 3V3
% ENV --> MCU : I2C
% GAS --> MCU : ADC/UART
% RTC --> MCU : I2C/SQW
% MCU --> LORA : SPI
% MCU --> ACT : GPIO
% MCU --> UI
% LORA --> MCU : DIO0 (ngat RX/TX)
% @enduml
% ===================================================================
\begin{figure}[H]
    \centering
    \includegraphics[width=0.95\linewidth]{Thuc_Hanh/Image/BlockDiagramNode.pdf}
    \caption{Sơ đồ khối tổng thể phần cứng của thiết bị Node}
    \label{fig:ch3_block}
\end{figure}

Trung tâm của sơ đồ là vi điều khiển ESP32-S3, đóng vai trò bộ não điều phối mọi luồng tín hiệu. Khối nguồn cung cấp điện áp $3{,}3\,\text{V}$ ổn định cho toàn hệ. Hai tổ hợp cảm biến (vi khí hậu và khí độc) là các giác quan đầu vào: nhóm vi khí hậu giao tiếp số qua I2C, nhóm khí độc giao tiếp qua ADC (MQ-135, CH4) và UART (CO2). Khối RTC DS3231 cung cấp trục thời gian cho cơ chế bảo mật. Khối LoRa SX1278 là cổng giao tiếp vô tuyến hai chiều với Gateway qua SPI, đồng thời báo sự kiện thu/phát qua chân ngắt DIO0. Cuối cùng, khối chấp hành nhận lệnh điều khiển từ GPIO để đóng cắt rơ-le. Điểm mấu chốt cần nhấn mạnh là toàn bộ logic quyết định nằm hoàn toàn bên trong ESP32-S3; các khối ngoại vi chỉ làm nhiệm vụ cảm nhận hoặc chấp hành, không tham gia tính toán. Cách phân bổ này hiện thực hóa đúng mô hình chia tải nhị phân cục bộ (Local Binary Offloading) đã chứng minh là tối ưu ở chương 2.

\subsubsection{Lựa chọn giải pháp công nghệ lõi}
a) \textbf{Khối Vi xử lý trung tâm: Module ESP32-S3-WROOM-1}

Trong khi các dòng chip ESP8266 đã hoàn toàn lỗi thời do bộ nhớ eo hẹp và ESP32 phiên bản tiêu chuẩn thiếu đi tập lệnh tối ưu cho Edge, đồ án quyết định sử dụng module SoC ESP32-S3-WROOM-1 làm khối xử lý trung tâm. ESP32-S3 sở hữu vi kiến trúc lõi kép Xtensa 32-bit LX7, hoạt động ở xung nhịp cực đại $240\text{ MHz}$. 

Hình \ref{fig:espwroom1} trình bày cấu trúc phần cứng bên trong module ESP32-S3-WROOM-1. Có thể thấy module được thiết kế dưới dạng một hệ thống tích hợp hoàn chỉnh (System-on-Module – SoM), trong đó các thành phần cần thiết cho hoạt động của vi điều khiển đều được tích hợp sẵn trên cùng một khối phần cứng.
\begin{figure}[H]
    \centering
    \includegraphics[width=1.0\linewidth]{Thuc_Hanh/Image/blockdiagram.png}
    \caption{Cấu trúc phần cứng của module ESP32-S3-WROOM-1}
    \label{fig:espwroom1}
\end{figure}
Trung tâm của module là chip ESP32-S3, đóng vai trò thực hiện các tác vụ xử lý dữ liệu, điều khiển hệ thống và quản lý truyền thông. Chip này được cấp nguồn hoạt động ở mức điện áp 3.3 V thông qua chân cấp nguồn của module. Tín hiệu EN (Enable) được sử dụng để kích hoạt hoặc khởi động lại vi điều khiển khi cần thiết.

Để đảm bảo hoạt động ổn định và chính xác, ESP32-S3 sử dụng bộ dao động thạch anh ngoài 40 MHz Crystal làm nguồn xung nhịp tham chiếu. Từ tín hiệu đồng hồ này, các bộ PLL nội bộ sẽ tạo ra các mức tần số cao hơn phục vụ cho lõi xử lý Xtensa LX7 và các khối ngoại vi bên trong hệ thống.

Bộ nhớ chương trình được lưu trữ trong khối QSPI Flash kết nối trực tiếp với vi điều khiển thông qua giao diện Quad SPI tốc độ cao. Toàn bộ firmware của hệ thống, các thư viện phần mềm và hệ thống tập tin SPIFFS đều được lưu trữ trên bộ nhớ Flash này. Bên cạnh đó, module còn hỗ trợ mở rộng PSRAM thông qua giao diện SPI chuyên dụng nhằm tăng dung lượng bộ nhớ cho các ứng dụng Edge Computing hoặc xử lý dữ liệu có kích thước lớn.

Khối GPIOs cung cấp các chân vào/ra số cho phép ESP32-S3 giao tiếp với cảm biến, cơ cấu chấp hành và các thiết bị ngoại vi khác. Thông qua các chân này, vi điều khiển có thể triển khai nhiều giao thức truyền thông như UART, SPI, I2C, PWM hoặc ADC phục vụ cho việc thu thập dữ liệu và điều khiển hệ thống.

Đối với chức năng truyền thông không dây, ESP32-S3 tích hợp bộ thu phát RF bên trong chip. Trước khi tín hiệu được đưa tới anten, tín hiệu RF phải đi qua khối RF Matching. Đây là mạng phối hợp trở kháng có nhiệm vụ tối ưu hóa việc truyền năng lượng giữa bộ thu phát và anten, giảm hiện tượng phản xạ tín hiệu, từ đó cải thiện hiệu suất phát sóng và độ nhạy thu của hệ thống Wi-Fi/Bluetooth.

Cuối cùng, tín hiệu RF sau khi được phối hợp trở kháng sẽ được truyền tới anten tích hợp trên module, cho phép ESP32-S3 thực hiện các kết nối không dây như Wi-Fi 2.4 GHz và Bluetooth Low Energy (BLE). Nhờ tích hợp sẵn toàn bộ chuỗi RF từ bộ thu phát đến anten, ESP32-S3-WROOM-1 giúp giảm đáng kể độ phức tạp trong thiết kế phần cứng và nâng cao độ tin cậy của hệ thống.

Từ cấu trúc trên có thể thấy ESP32-S3-WROOM-1 không chỉ là một vi điều khiển đơn thuần mà là một nền tảng phần cứng hoàn chỉnh, tích hợp đầy đủ các thành phần xử lý, bộ nhớ, khối RF và giao tiếp ngoại vi. Điều này giúp giảm chi phí thiết kế, rút ngắn thời gian phát triển sản phẩm và đặc biệt phù hợp với các ứng dụng IoT biên yêu cầu khả năng xử lý thời gian thực, lưu trữ dữ liệu cục bộ và kết nối không dây ổn định.

b) \textbf{Công nghệ Giao tiếp Vô tuyến: Module Lora SX1278 (SPI Bus)}

Để thực hiện truyền dữ liệu từ nút cảm biến đến Gateway với khoảng cách xa và mức tiêu thụ năng lượng thấp, đồ án sử dụng module Ra-02 tích hợp chip thu phát vô tuyến SX1278 của Semtech. Đây là bộ thu phát RF hỗ trợ công nghệ điều chế LoRa (Long Range), cho phép truyền dữ liệu ở khoảng cách lên tới vài kilomet trong điều kiện thực tế, đồng thời vẫn đảm bảo công suất tiêu thụ thấp phù hợp với các hệ thống IoT.

Hình \ref{Hình 3.4} trình bày cấu trúc phần cứng của module SX1278. Có thể thấy chip được tích hợp đầy đủ các khối xử lý tín hiệu RF, bộ điều chế LoRa và các giao diện điều khiển cần thiết để xây dựng một hệ thống truyền thông vô tuyến hoàn chỉnh.

\begin{figure}[H]
\centering
\includegraphics[width=0.95\linewidth]{Thuc_Hanh/Image/LoraSX1278.png}
\caption{Kiến trúc phần cứng của module LoRa SX1278}
\label{Hình 3.4}
\end{figure}

Tại tầng RF, SX1278 tích hợp các khối PA (Power Amplifier) và LNA (Low Noise Amplifier). Khối PA có nhiệm vụ khuếch đại tín hiệu trước khi phát ra anten nhằm tăng phạm vi truyền dẫn, trong khi LNA khuếch đại các tín hiệu thu có công suất rất nhỏ từ môi trường, giúp cải thiện độ nhạy của bộ thu. Theo thông số kỹ thuật của nhà sản xuất, SX1278 có thể đạt công suất phát tối đa khoảng +20 dBm và độ nhạy thu lên tới -148 dBm khi hoạt động ở chế độ LoRa.

Sau tầng RF, tín hiệu được đưa tới khối Mixer và Single-to-Differential Converter để thực hiện quá trình chuyển đổi và xử lý tín hiệu trung tần. Các khối này có vai trò chuyển đổi giữa miền tương tự và miền số, chuẩn bị dữ liệu cho quá trình điều chế và giải điều chế bên trong chip.

Trung tâm của hệ thống thu phát là khối Frac-N Frequency Synthesizer kết hợp với bộ dao động tinh thể ngoài 32 MHz XO (Crystal Oscillator). Bộ tổng hợp tần số này tạo ra tần số sóng mang RF với độ chính xác cao, cho phép SX1278 hoạt động trong các dải tần khác nhau như 169 MHz, 433 MHz, 868 MHz và 915 MHz. Trong đồ án này, module được cấu hình hoạt động tại băng tần 433 MHz, phù hợp với điều kiện triển khai thực tế và có khả năng xuyên vật cản tốt.

Khối LoRa Demodulator/Modulator thực hiện chức năng điều chế và giải điều chế tín hiệu theo kỹ thuật Chirp Spread Spectrum (CSS). Đây là công nghệ cốt lõi của LoRa, cho phép trải phổ tín hiệu trên một băng thông rộng nhằm tăng khả năng chống nhiễu, nâng cao độ nhạy thu và mở rộng khoảng cách truyền dẫn. Ngoài chế độ LoRa, SX1278 còn hỗ trợ các phương thức điều chế truyền thống như FSK, GFSK và OOK để đáp ứng nhiều yêu cầu ứng dụng khác nhau.

Để nâng cao độ tin cậy của truyền thông, chip còn tích hợp các khối RSSI (Received Signal Strength Indicator) và AFC (Automatic Frequency Correction). RSSI được sử dụng để đánh giá cường độ tín hiệu thu, trong khi AFC tự động hiệu chỉnh sai lệch tần số giữa bộ phát và bộ thu, giúp duy trì liên kết ổn định ngay cả khi môi trường truyền dẫn thay đổi.

Dữ liệu truyền và nhận được lưu trữ tạm thời trong Packet Engine và FIFO Memory 256 Byte. Khối này hỗ trợ đóng gói khung dữ liệu, kiểm tra CRC, phát hiện lỗi và quản lý bộ đệm truyền thông. Nhờ đó, vi điều khiển không cần xử lý trực tiếp từng bit dữ liệu mà chỉ cần trao đổi dữ liệu với bộ đệm FIFO thông qua giao diện SPI.

Việc cấu hình và điều khiển SX1278 được thực hiện thông qua giao diện SPI cùng các chân ngắt DIO0–DIO5. Trong hệ thống đề xuất, ESP32-S3 đóng vai trò bộ điều khiển trung tâm, sử dụng giao thức SPI để cấu hình các thanh ghi của SX1278, ghi dữ liệu vào FIFO trước khi phát và đọc dữ liệu nhận được từ bộ thu.

c) \textbf{Thiết kế khối nguồn}

Nhiệm vụ quan trọng đầu tiên của khối nguồn là quản lý chu trình sạc/xả của viên pin này nhằm chống lại hiện tượng phù tụ, đoản mạch hoặc suy giảm tuổi thọ tế bào hóa học . Linh kiện được tích hợp cho nhiệm vụ này là vi mạch quản lý sạc TP4056. 
\begin{figure}[H]
\centering
\includegraphics[width=0.95\linewidth]{Thuc_Hanh/Image/TP4056.png}
\caption{Mạch quản lý sạc TP4056}
\label{Hình 3.5}
\end{figure}
Thuật toán sạc của TP4056 tuân thủ nghiêm ngặt quy trình sạc hai giai đoạn tuyến tính: Dòng điện không đổi / Điện áp không đổi (Constant Current / Constant Voltage - CC/CV). Về mặt toán học, dòng sạc tối đa $I_{CHG}$ được lập trình thông qua một điện trở $R_{PROG}$ nối từ chân PROG xuống mass, tuân theo phương trình $I_{CHG} = 1200 / R_{PROG}$. Với cấu hình $R_{PROG} = 1.2\text{k}\Omega$, hệ thống cho phép dòng sạc danh định là $1\text{A}$. Khi điện áp pin tiệm cận ngưỡng cắt $4.2\text{V}$, IC tự động chuyển sang pha sạc áp không đổi (CV), dòng sạc sẽ giảm dần theo hàm mũ cho đến khi tụt xuống dưới $1/10$ dòng CC ban đầu ($100\text{mA}$), lúc này chu trình sạc được ngắt hoàn toàn để bảo vệ điện cực \cite{texas2018battery}. Để thiết lập một lớp khiên bảo vệ vật lý nhiều tầng, TP4056 được phối hợp với IC bảo vệ DW01 và cụm MOSFET xả/sạc kép. DW01 đóng vai trò như một lính gác điện áp (Voltage Monitor); nếu nó phát hiện điện áp xả của pin rơi xuống dưới ngưỡng $2.4\text{V}$ (Over-discharge) hoặc dòng xả vượt quá $3\text{A}$ (Short Circuit), nó sẽ lập tức ngắt điện áp kích ở cực Gate của MOSFET, cô lập hoàn toàn viên pin khỏi tải tiêu thụ.

Nhiệm vụ thứ hai, và cũng là nguyên nhân gây thất thoát năng lượng lớn nhất, là quá trình hạ áp (Step-down Regulation) từ dải biến thiên $3.0\text{V}-4.2\text{V}$ của pin xuống mức $3.3\text{V}$ phẳng tuyệt đối. Mức điện áp $3.3\text{V}$ này là bắt buộc để duy trì tần số dao động chuẩn xác của bộ thạch anh trên ESP32-S3 và đảm bảo IC LoRa SX1278 không bị cháy. Các bo mạch phát triển giá rẻ thường lạm dụng IC ổn áp tuyến tính AMS1117-3.3V. Tuy nhiên, kiến trúc của AMS1117 sử dụng BJT (Bipolar Junction Transistor) làm phần tử dẫn, sinh ra dòng tĩnh (Quiescent Current - $I_Q$) khổng lồ lên tới $5\text{mA}$ và mức điện áp rơi (Dropout Voltage - $V_{DO}$) $\approx 1.1\text{V}$. Điều này đồng nghĩa với việc, nếu điện áp pin sụt xuống $4.0\text{V}$, đầu ra của AMS1117 chỉ còn $2.9\text{V}$, làm hệ thống sập nguồn ngay cả khi pin vẫn còn đầy $80\%$.

Để phá vỡ rào cản này, bản thiết kế Schematic đã thay thế hoàn toàn bằng IC LDO Low-Power AP2112K-3.3V. Kiến trúc của AP2112K sử dụng P-Channel MOSFET làm phần tử dẫn, cho phép ép mức điện áp rơi $V_{DO}$ xuống chỉ còn $250\text{mV}$ khi chạy đầy tải $600\text{mA}$. Quan trọng hơn, dòng tĩnh $I_Q$ của IC này chỉ xấp xỉ $55 \mu\text{A}$ \cite{diodes2015ap2112}. Tổn hao công suất (Power Dissipation - $P_D$) tỏa nhiệt trên LDO được tính bằng phương trình:
\begin{equation}
    P_D = (V_{IN} - V_{OUT}) \times I_{LOAD} + V_{IN} \times I_Q
\end{equation}
Tuy nhiên, điểm nhấn học thuật gây tranh cãi và đáng chú ý nhất trong thiết kế Schematic này là sự đánh đổi có chủ đích (Trade-off) về mặt động lực học năng lượng: Hệ thống từ chối hoàn toàn việc sử dụng Công tắc tải (Load Switch MOSFET) để cắt điện các cảm biến khí. Các nghiên cứu IoT thông thường thường sử dụng MOSFET để cắt điện các cảm biến ăn dòng lớn khi MCU vào chế độ ngủ (Deep Sleep), gọi là cơ chế đánh xung nhiệt (Pulsed Heating). Nhưng đối với cảm biến khí oxit kim loại (MOX) như MQ-135 (đo $NH_3$) và Cảm biến CH4 DFRobot, bề mặt phản ứng $SnO_2$ bắt buộc phải được nung nóng liên tục ở dải nhiệt độ $>300^\circ\text{C}$ để duy trì trạng thái cân bằng nhiệt động học (Thermal Equilibrium) với các phân tử oxy trong không khí. Nếu cắt nguồn và cấp lại đột ngột, sự biến thiên nhiệt độ (Thermal Shock) sẽ làm đường điện trở nền ($R_0$) của cảm biến bị trôi (Baseline Drift) dữ dội. Thời gian làm nóng để cảm biến ổn định trở lại (Warm-up time) có thể kéo dài từ 5 đến 15 phút, gây sai số đo lường lên tới hàng chục ppm, phá vỡ hoàn toàn độ tin cậy của một hệ thống cảnh báo khí độc công nghiệp \cite{barsan2001conduction}.

Do đó, thiết kế mạch quyết định đấu nối thẳng nguồn +3.3V\_SYS từ AP2112K vào chân cấp điện của toàn bộ mảng cảm biến. Bo mạch chấp nhận việc duy trì một dòng tải liên tục (Continuous Current Draw) $\approx 200\text{mA}$ cho cuộn dây Heater của MQ-135 và CH4. Quyết định kỹ thuật này hy sinh một phần tuổi thọ pin để đổi lấy độ chính xác hóa học tuyệt đối (Absolute Chemical Precision). Để bù đắp lại, vi điều khiển ESP32-S3 được cấu hình vào trạng thái Deep Sleep cực đoan nhất, ngắt toàn bộ lõi CPU và RAM, chỉ tiêu thụ $7\mu\text{A}$ thời gian rảnh, nhường toàn bộ năng lượng pin cho cảm biến khí "thở".

Hình \ref{Hình 3.5} mô tả luồng đi của năng lượng: từ cổng USB Type-C cấp nguồn $5\text{V}$ vào mạch sạc TP4056, sạc cho pin Lithium. Dòng điện từ pin đi qua IC bảo vệ DW01 trước khi tiến vào chân $V_{IN}$ của LDO AP2112K.
\begin{figure}[H]
    \centering
    \includegraphics[width=1.0\linewidth]{Thuc_Hanh/Image/KhoiNguon.pdf}
    \caption{Luồng hoạt động của khối nguồn}
    \label{Hình 3.5}
\end{figure}

d) \textbf{Tổ hợp cảm biến vi khí hậu (SHT31, BMP180, BH1750)}

Để lượng hóa trạng thái vi khí hậu chuồng nuôi, Node tích hợp một tổ hợp cảm biến số giao tiếp qua bus I2C, trong đó mỗi cảm biến đảm nhiệm một nhóm đại lượng vật lý riêng. 
\begin{figure}[H]
    \centering
    \includegraphics[width=0.76\linewidth]{Thuc_Hanh/Image/SHT31.png}
    \caption{Cảm biến nhiệt độ và độ ẩm SHT31 (giao tiếp I2C)}
    \label{fig:ch3_sht31}
\end{figure}
Cảm biến nhiệt độ và độ ẩm SHT31 của Sensirion được chọn làm cảm biến chủ lực cho bài toán Chỉ số nhiệt ẩm (THI). SHT31 cho sai số nhiệt độ điển hình $\pm 0{,}3\,^\circ C$ và sai số độ ẩm $\pm 2\%$RH, nằm trong giới hạn sai số dưới $5\%$ mà đề tài cam kết ở chương 1 \cite{sensirion2018sht31}. Cảm biến tích hợp sẵn mạch chuyển đổi ADC và bộ hiệu chỉnh tuyến tính hóa bên trong, đồng thời đính kèm 8 bit CRC cho mỗi cặp byte dữ liệu trả về; nhờ đó firmware tại Node có thể loại bỏ ngay các khung dữ liệu lỗi do nhiễu bus mà không làm sai lệch khâu tính toán THI phía sau. 
\begin{figure}[H]
    \centering
    \includegraphics[width=0.9\linewidth]{Thuc_Hanh/Image/BMP180.png}
    \caption{Cảm biến áp suất khí quyển BMP180 (giao tiếp I2C)}
    \label{fig:ch3_bmp180}
\end{figure}
Cảm biến áp suất khí quyển BMP180 của Bosch được bổ sung nhằm cung cấp tham chiếu áp suất phục vụ phát hiện xu hướng thời tiết và bù độ cao; BMP180 lưu sẵn 11 hệ số hiệu chỉnh trong bộ nhớ EEPROM nội, cho phép firmware đọc một lần lúc khởi tạo và áp dụng chuỗi công thức bù nhiệt để quy đổi giá trị thô thành áp suất tuyệt đối chính xác \cite{bosch2013bmp180}. 
\begin{figure}[H]
    \centering
    \includegraphics[width=0.9\linewidth]{Thuc_Hanh/Image/BH1750.png}
    \caption{Cảm biến cường độ ánh sáng BH1750 (giao tiếp I2C)}
    \label{fig:ch3_bh1750}
\end{figure}
Cảm biến cường độ ánh sáng BH1750 \cite{rohm2011bh1750} đo độ rọi (lux) phục vụ giám sát chu kỳ chiếu sáng nhân tạo trong chuồng kín, một yếu tố ảnh hưởng trực tiếp đến nhịp sinh học và năng suất của gia cầm.

Một quyết định kiến trúc quan trọng trong sơ đồ nguyên lý là việc bố trí các cảm biến này trên đường bus I2C dùng chung với các đường kéo lên (pull-up) giá trị $4{,}7\,\text{k}\Omega$ (linh kiện R16, R17 trên Schematic). Giá trị điện trở kéo lên được tính toán dựa trên điện dung ký sinh tổng cộng của bus và tốc độ truyền mong muốn theo ràng buộc của chuẩn I2C \cite{nxp2014i2c}: điện trở quá lớn làm sườn lên của tín hiệu bị bo tròn gây lỗi ở tốc độ cao, trong khi điện trở quá nhỏ làm tăng dòng tiêu thụ khi bus ở mức thấp. Mức $4{,}7\,\text{k}\Omega$ là điểm cân bằng tiêu chuẩn cho bus I2C tốc độ 100--400 kHz với vài cảm biến mắc song song. Bảng \ref{Bảng 3.2} tổng hợp đặc tính của tổ hợp cảm biến vi khí hậu.

\begin{table}[H]
\centering
\caption{Đặc tính kỹ thuật của tổ hợp cảm biến vi khí hậu}
\label{Bảng 3.2}
\renewcommand{\arraystretch}{1.4}
\begin{tabular}{p{2.4cm} p{3.0cm} p{2.2cm} p{2.0cm} p{3.4cm}}
\hline
\textbf{Cảm biến} & \textbf{Đại lượng đo} & \textbf{Giao tiếp} & \textbf{Địa chỉ} & \textbf{Sai số điển hình} \\ \hline
SHT31 & Nhiệt độ, độ ẩm & I2C & 0x44 & $\pm 0{,}3\,^\circ C$ / $\pm 2\%$RH \\ 
BMP180 & Áp suất, nhiệt độ & I2C & 0x77 & $\pm 0{,}12$ hPa \\ 
BH1750 & Cường độ sáng & I2C & 0x23 & $\pm 20\%$ (lux) \\ \hline
\end{tabular}
\end{table}

e) \textbf{Tổ hợp cảm biến khí độc (MQ-135, CH4, CO2)}

Trọng tâm an toàn sinh học của hệ thống nằm ở khả năng phát hiện khí độc tích tụ, đặc biệt là Amoniac ($NH_3$) sinh ra từ quá trình phân hủy chất thải đã phân tích ở chương 2. Để đo $NH_3$, đề tài sử dụng cảm biến khí oxit kim loại bán dẫn (Metal-Oxide Semiconductor -- MOX) dòng MQ-135. Nguyên lý hoạt động của MQ-135 dựa trên sự thay đổi điện trở bề mặt của lớp màng nhạy khí $SnO_2$ khi hấp phụ các phân tử khí khử: khi nồng độ $NH_3$ tăng, điện trở cảm biến giảm, và tỉ số $R_s/R_0$ (điện trở trong khí cần đo trên điện trở trong không khí sạch) được quy đổi sang nồng độ ppm thông qua đường đặc tuyến log-log của nhà sản xuất \cite{barsan2001conduction}. Tín hiệu đầu ra của MQ-135 là điện áp tương tự, được đưa về chân ADC của ESP32-S3 (đường A0 trên Schematic) qua đầu nối J9. Bên cạnh $NH_3$, hệ thống dự phòng thêm hai kênh đo khí mở rộng: cảm biến CH4 (đo khí Mêtan, ngõ ra tương tự, đầu nối J11) phục vụ cảnh báo khí dễ cháy sinh ra từ phân hủy yếm khí, và cảm biến CO2 giao tiếp số qua UART (đầu nối J7/J8) cho phép theo dõi mức độ thông thoáng của chuồng.

Tóm tắt mô tả cho các cảm biến khí độc được thể hiện trong Bảng \ref{Bảng 3.3}
\begin{table}[H]
\centering
\caption{Tổ hợp cảm biến khí và giao tiếp tương ứng}
\label{Bảng 3.3}
\renewcommand{\arraystretch}{1.4}
\begin{tabular}{p{2.6cm} p{3.2cm} p{2.6cm} p{4.2cm}}
\hline
\textbf{Cảm biến} & \textbf{Khí mục tiêu} & \textbf{Giao tiếp} & \textbf{Đặc điểm} \\ \hline
MQ-135 & $NH_3$ (Amoniac) & Analog (ADC, A0) & Màng MOX $SnO_2$, cần nung nóng \\ 
CH4 & Mêtan ($CH_4$) & Analog (ADC) & Cảnh báo khí dễ cháy \\ 
CO2 & Cacbonic ($CO_2$) & UART (số) & Theo dõi độ thông thoáng \\ \hline
\end{tabular}
\end{table}

f) \textbf{Khối thời gian thực (RTC DS3231) và khối chấp hành}

Để hiện thực cơ chế chống tấn công phát lại (Replay Attack) dựa trên nhãn thời gian đã trình bày ở chương 2, Node được trang bị IC đồng hồ thời gian thực DS3231 \cite{maxim2015ds3231} giao tiếp qua I2C. 
\begin{figure}[H]
    \centering
    \includegraphics[width=0.9\linewidth]{Thuc_Hanh/Image/DS3231.png}
    \caption{Module đồng hồ thời gian thực DS3231}
    \label{fig:ch3_ds3231}
\end{figure}
DS3231 tích hợp sẵn bộ bù nhiệt bằng tinh thể thạch anh (TCXO), cho sai số cực thấp khoảng $\pm 2$ ppm (tương đương dưới một phút sai lệch mỗi năm), bảo đảm hai đầu Node và Gateway luôn có một mốc thời gian chung đáng tin cậy để cửa sổ xác thực $\Delta t < 5\,\text{s}$ hoạt động chính xác. Chân ngắt/xung vuông (SQW) của DS3231 còn có thể được cấu hình để phát xung định kỳ đánh thức vi điều khiển khỏi chế độ ngủ sâu, đóng vai trò một nguồn đánh thức theo thời gian độc lập với bộ định thời nội của ESP32-S3. Nguồn nuôi cho mạch dao động của DS3231 được duy trì liên tục để đồng hồ không bị mất giờ qua mỗi chu kỳ ngủ sâu của vi điều khiển.

Ở đầu ra điều khiển, Node bố trí ba kênh chấp hành (Act1, Act2, Act3 -- các đầu nối J20, J21, J22 trên Schematic) đưa tín hiệu mức logic tới các module rơ-le bên ngoài để đóng/cắt các thiết bị cơ điện công suất lớn như quạt thông gió, hệ thống phun sương làm mát hoặc đèn sưởi. Việc tách riêng ba kênh cho phép thuật toán Điện toán biên thực thi các kịch bản điều khiển độc lập (ví dụ: kênh 1 cho quạt hút, kênh 2 cho phun sương, kênh 3 cho đèn sưởi), thay vì gộp chung một ngõ ra. Đây chính là mắt xích cuối cùng khép kín vòng điều khiển cục bộ tại biên: từ cảm biến $\rightarrow$ FSM ra quyết định $\rightarrow$ chân GPIO $\rightarrow$ rơ-le, toàn bộ diễn ra trong chưa đầy 50 ms mà không cần bất kỳ kết nối nào ra bên ngoài.





\subsection{Thiết kế kiến trúc phần mềm đa nhiệm trên FreeRTOS}
Phần cứng dù mạnh mẽ cũng chỉ phát huy hiệu quả khi được điều phối bởi một kiến trúc phần mềm hợp lý. Dựa trên cơ sở lý thuyết về hệ điều hành thời gian thực và cơ chế phân luồng đa nhân đã trình bày ở chương 2, phần này cụ thể hóa kiến trúc firmware của Node thành một tập các tác vụ (Task) FreeRTOS chạy song song, được gán cứng vào hai nhân vật lý của ESP32-S3 nhằm bảo đảm tính tất định của luồng điều khiển an toàn.

\subsubsection{Phân rã hệ thống thành các tác vụ}
Toàn bộ chức năng của Node được phân rã thành sáu tác vụ độc lập, mỗi tác vụ có một mức ưu tiên và một vùng ngăn xếp riêng. Việc phân rã tuân theo nguyên tắc tách biệt mối quan tâm (Separation of Concerns): mỗi tác vụ chỉ chịu trách nhiệm một nhóm chức năng đồng nhất, trao đổi dữ liệu với nhau thông qua các hàng đợi (Queue) và cờ hiệu (Semaphore/Mutex) để tránh tranh chấp dữ liệu (Race Condition).

% ===================================================================
% PlantUML cho Hình 3.9 - render và lưu thành: Thuc_Hanh/Image/Ch3_TaskArch.png
% @startuml
% skinparam backgroundColor white
% package "CORE 0 (PRO_CPU) - Vo tuyen & Bao mat" {
%   [Task_LoRa_TX/RX\n(uu tien cao)] as T1
%   [Task_Crypto\nAES-128 / ASCON] as T2
%   [Task_Sync_Offline\n(uu tien thap)] as T3
% }
% package "CORE 1 (APP_CPU) - Dien toan bien" {
%   [Task_Sensor\n(doc I2C/ADC/UART)] as T4
%   [Task_Edge_FSM\n(uu tien cao nhat)] as T5
%   [Task_Telemetry\n(USB JSON)] as T6
% }
% queue "Queue du lieu cam bien" as Q1
% queue "Queue goi gui" as Q2
% T4 --> Q1
% Q1 --> T5
% T5 --> Q2 : goi can gui
% Q2 --> T2
% T2 --> T1 : goi da ma hoa
% T1 --> T3 : khi mat ACK
% T5 --> [GPIO Ro-le] : lenh dieu khien
% T4 ..> T6
% @enduml
% ===================================================================
\begin{figure}[H]
    \centering
    \includegraphics[width=0.95\linewidth]{Thuc_Hanh/Image/PhanLuongDaNhiemj.pdf}
    \caption{Kiến trúc phân luồng đa nhiệm trên ESP32-S3}
    \label{fig:ch3_task}
\end{figure}

Hình \ref{fig:ch3_task} mô tả nguyên tắc gán nhân: Core 0 (PRO\_CPU) gánh toàn bộ các tác vụ liên quan đến vô tuyến và mật mã -- vốn có tính chất "chặn" (blocking) và tiêu tốn nhiều chu kỳ máy khi chờ truyền/nhận trên SPI -- bao gồm tác vụ phát/thu LoRa, tác vụ mã hóa AES-128/ASCON và tác vụ đồng bộ dữ liệu ngoại tuyến chạy nền ở mức ưu tiên thấp. Core 1 (APP\_CPU) được giữ "sạch" để chuyên trách Điện toán biên: tác vụ đọc cảm biến, tác vụ máy trạng thái FSM (mang mức ưu tiên cao nhất hệ thống) và tác vụ phát telemetry. Nhờ sự cô lập này, ngay cả khi một gói LoRa nhiễu rác dội vào và làm bận Core 0, tác vụ FSM trên Core 1 vẫn được lập lịch tức thời để ra lệnh đóng rơ-le, bảo đảm độ trễ điều khiển luôn dưới ngưỡng 50 ms. Bảng \ref{Bảng 3.4} liệt kê chi tiết tập tác vụ cùng mức ưu tiên và nhân được gán.

\begin{table}[H]
\centering
\caption{Phân rã tác vụ FreeRTOS: gán nhân, mức ưu tiên và kích thước Stack dự kiến}
\label{Bảng 3.4}
\renewcommand{\arraystretch}{1.4}
\small
\begin{tabular}{p{2.8cm} p{1.1cm} p{1.3cm} p{1.7cm} p{6.2cm}}
\hline
\textbf{Tác vụ} & \textbf{Nhân} & \textbf{Ưu tiên} & \textbf{Stack (B)} & \textbf{Chức năng} \\ \hline
Task\_Edge\_FSM & Core 1 & 6 & 3072 & Đối chiếu ngưỡng, ra quyết định, điều khiển rơ-le (tới hạn $<50$ ms) \\ 
Task\_LoRa & Core 0 & 5 & 4096 & Phát/thu gói, chờ ACK, đo RSSI/SNR \\ 
Task\_Crypto & Core 0 & 4 & 4096 & Mã hóa/giải mã AES-128-GCM hoặc ASCON-128 \\ 
Task\_Sensor & Core 1 & 3 & 4096 & Đọc I2C/ADC/UART, lọc nhiễu (EMA/Median/Kalman) \\ 
Task\_Telemetry & Core 1 & 2 & 4096 & Phát JSON trạng thái qua USB phục vụ giám sát \\ 
Task\_Sync\_Off & Core 0 & 1 & 3584 & Ghi và đồng bộ lại dữ liệu lưu ngoại tuyến (SPIFFS) \\ \hline
\end{tabular}
\end{table}

Việc gán mức ưu tiên tuân theo nguyên tắc \textit{tầm quan trọng đối với an toàn} kết hợp tinh thần lập lịch đơn điệu theo tần suất (Rate Monotonic) \cite{liu1973scheduling}: tác vụ FSM an toàn -- chịu ràng buộc trễ ngặt nhất ($<50$ ms) -- được gán ưu tiên cao nhất (6); kế đến là các tác vụ vô tuyến nhạy cảm về thời gian (LoRa, Crypto, ưu tiên 5--4); tác vụ đọc cảm biến định kỳ ở mức trung bình (3); còn các tác vụ "mềm" như phát telemetry và đồng bộ nền được đặt ưu tiên thấp nhất (2--1) để không bao giờ cản trở luồng điều khiển tới hạn. Kích thước Stack ở cột thứ tư được ước lượng ban đầu dựa trên độ sâu chuỗi lời gọi hàm và dung lượng bộ đệm cục bộ của mỗi tác vụ (bộ đệm gói tin, trạng thái thuật toán mã hóa, vùng định dạng chuỗi JSON), cộng thêm biên dự phòng an toàn; các giá trị này là \emph{dự kiến} và sẽ được hiệu chỉnh chính xác bằng phép đo High-Water Mark thực tế ở Chương 5 nhằm vừa tránh tràn Stack vừa không lãng phí RAM.

Để thống nhất với các thuật ngữ nêu trong yêu cầu thiết kế, cần làm rõ ánh xạ tên gọi: \texttt{Task\_Sensor} chính là tác vụ đọc cảm biến trên bus I2C (còn gọi là \textbf{Task\_Sensor\_I2C}); cụm xử lý biên gồm lọc Kalman (thực thi trong \texttt{Task\_Sensor}/\texttt{Task\_Edge\_FSM}) phối hợp với mã hóa trong \texttt{Task\_Crypto} tương ứng với khối \textbf{Task\_Kalman\_AES}; và \texttt{Task\_LoRa} đảm nhiệm vai trò phát gói \textbf{Task\_LoRa\_TX}. Ba khối chức năng này tạo thành một đường ống (pipeline) sản xuất -- tiêu thụ nối với nhau bằng các Hàng đợi, sẽ được trình bày ở mục tiếp theo.

\subsubsection{Cơ chế giao tiếp liên tác vụ và đồng bộ}
Để các tác vụ phối hợp nhịp nhàng mà không giẫm chân nhau, firmware sử dụng ba nguyên tắc đồng bộ của FreeRTOS. Thứ nhất, các Hàng đợi (Queue) đóng vai trò băng chuyền dữ liệu một chiều: tác vụ đọc cảm biến đẩy cấu trúc dữ liệu môi trường vào hàng đợi, tác vụ FSM lấy ra để xử lý; tương tự, FSM đẩy gói cần gửi vào hàng đợi cho tác vụ mã hóa và LoRa. Cơ chế hàng đợi vừa truyền dữ liệu vừa tự động đồng bộ nhịp sản xuất -- tiêu thụ, đồng thời sao chép dữ liệu nên loại bỏ rủi ro hai tác vụ cùng truy cập một vùng nhớ. Thứ hai, Cờ hiệu nhị phân (Binary Semaphore) được dùng để báo hiệu sự kiện ngắt: khi chân DIO0 của SX1278 phát ngắt báo đã nhận xong gói, trình phục vụ ngắt (ISR) chỉ "nhả" một cờ hiệu, đánh thức tác vụ LoRa xử lý phần nặng bên ngoài ngữ cảnh ngắt -- một thực hành chuẩn giúp giữ cho ISR cực ngắn. Thứ ba, Mutex (cờ loại trừ tương hỗ) bảo vệ các tài nguyên dùng chung như bus I2C hay biến cấu hình toàn cục, bảo đảm tại một thời điểm chỉ một tác vụ được truy cập, tránh hỏng dữ liệu do truy cập đồng thời.

% --- Bổ sung phần RTOS ---
\subsubsection{Thiết kế cơ chế bảo vệ bus I2C dùng chung bằng Mutex}
Trên cơ sở lý thuyết về loại trừ tương hỗ và kế thừa ưu tiên đã trình bày ở Chương 2, phần này cụ thể hóa thiết kế đường ống xử lý đa nhiệm và cơ chế bảo vệ bus I2C của Node. Hình \ref{Hình 3.14} biểu diễn luồng dữ liệu dạng sản xuất -- tiêu thụ chạy xuyên qua ba khối tác vụ chính, cùng vị trí đặt Mutex bảo vệ tài nguyên dùng chung.


\begin{figure}[H]
    \centering
    \includegraphics[width=1.0\linewidth]{Thuc_Hanh/Image/cochedongbo.pdf}
    \caption{Xử lý đa nhiệm và cơ chế Mutex cho bus I2C}
    \label{Hình 3.14}
\end{figure}
Như sơ đồ thể hiện, bốn thiết bị SHT31, BMP180, BH1750 và DS3231 chia sẻ chung một bus I2C; ngoài ra biến cấu hình toàn cục cũng là tài nguyên dùng chung. Để loại trừ tranh chấp, đề tài thiết kế một Mutex duy nhất (\texttt{i2c\_mutex}) được khởi tạo một lần lúc hệ thống boot bằng \texttt{xSemaphoreCreateMutex()}. Quy ước truy cập được áp đặt nghiêm ngặt: bất kỳ tác vụ nào muốn thực hiện một giao dịch I2C đều phải giành Mutex bằng \texttt{xSemaphoreTake()} \emph{trước} khi phát điều kiện START, và chỉ trả lại bằng \texttt{xSemaphoreGive()} \emph{sau} khi đã phát điều kiện STOP -- toàn bộ giao dịch trở thành một vùng tới hạn (critical section) nguyên tử.

Hai quyết định thiết kế quan trọng được áp dụng. Thứ nhất, lời gọi \texttt{xSemaphoreTake()} luôn kèm một \textit{thời hạn chờ} (timeout, ví dụ $100$ ms) thay vì chờ vô hạn: nếu không giành được Mutex trong thời hạn (dấu hiệu bus bị treo hoặc tác vụ giữ Mutex gặp sự cố), tác vụ đọc cảm biến sẽ bỏ qua chu kỳ đó và đánh dấu dữ liệu là không hợp lệ thay vì treo cứng -- bảo toàn nguyên tắc "không bao giờ chặn luồng điều khiển an toàn". Thứ hai, nhờ Mutex của FreeRTOS có cơ chế kế thừa ưu tiên, nếu tác vụ FSM ưu tiên cao tình cờ cần truy cập bus trong lúc tác vụ cảm biến ưu tiên thấp đang giữ Mutex, ưu tiên của tác vụ cảm biến sẽ được nâng tạm thời để nó nhanh chóng giải phóng bus, giới hạn cận trên thời gian chặn đúng bằng độ dài một giao dịch I2C như trong Hình \ref{fig:rtos_priority_inheritance}.

Việc phân rã hệ thống thành mô hình đường ống (pipeline) với các tác vụ độc lập mang lại ưu điểm vượt trội về khả năng đáp ứng thời gian thực và tối ưu hóa hiệu suất CPU. Tuy nhiên, kiến trúc đa nhiệm này cũng đặt ra thách thức lớn về tính toàn vẹn dữ liệu khi nhiều tiến trình chạy song song có thể cùng yêu cầu truy xuất phần cứng tại một thời điểm. Đặc biệt đối với chuẩn giao tiếp đồng bộ yêu cầu sự chính xác cao về mặt xung nhịp như I2C, bất kỳ sự can thiệp đan xen nào từ một tác vụ khác trong lúc tín hiệu đang truyền tải đều có thể gây ra hiện tượng xung đột (bus collision) hoặc làm hỏng hoàn toàn khung dữ liệu (frame corruption). Do đó, việc thiết lập một cơ chế đồng bộ chặt chẽ không chỉ là giải pháp phòng ngừa lỗi, mà là điều kiện tiên quyết để đảm bảo chuỗi dữ liệu từ khâu sản xuất (đọc cảm biến) đến khâu tiêu thụ (xử lý thuật toán và truyền dẫn) luôn chính xác và đáng tin cậy.
% --- Hết bổ sung phần RTOS ---

\subsection{Thiết kế thuật toán Điện toán biên tại Node}
Phần này thiết kế chi tiết ba thành phần thuật toán cốt lõi: chu trình lấy mẫu tiết kiệm năng lượng, khâu tiền xử lý -- lọc nhiễu dữ liệu cảm biến, và máy trạng thái hữu hạn (FSM) ra quyết định điều khiển. 

\subsubsection{Lưu đồ chu trình hoạt động tổng thể của Node}
Mỗi chu kỳ hoạt động của Node được thiết kế thành một chuỗi trạng thái tuần tự, bắt đầu từ lúc thức dậy và kết thúc bằng việc quay lại trạng thái ngủ. 
\begin{figure}[H]
    \centering
    {\includegraphics[width=0.55\linewidth]{Thuc_Hanh/Image/NodeSequence.pdf}}
    \caption{Lưu đồ chu trình hoạt động một chu kỳ của Node}
    \label{fig:ch3_cycle}
\end{figure}
Việc mô hình hóa chu trình thành một lưu đồ tường minh như trong Hình \ref{fig:ch3_cycle} giúp bảo đảm rằng mọi nhánh xử lý -- kể cả nhánh lỗi -- đều được tính đến, đồng thời cho phép ước lượng chính xác thời gian thức (active time) phục vụ tính toán năng lượng.

% ===================================================================
% PlantUML cho Hình 3.10 - render và lưu thành: Thuc_Hanh/Image/Ch3_SensingCycle.png
% @startuml
% skinparam backgroundColor white
% start
% :Thuc day (timer RTC / ngat su kien);
% :Khoi tao ngoai vi, doc ly do reset;
% :Doc cam bien (SHT31, BMP180, BH1750,\nMQ-135, CH4, CO2);
% :Loc nhieu per-signal (EMA/MA/Median);
% :Tinh THI, doi chieu nguong -> FSM;
% if (Vuot nguong nguy hiem?) then (co)
%   :Dong ro-le chap hanh (<50ms);
% else (khong)
% endif
% :Gan timestamp (DS3231) + dong goi payload;
% :Ma hoa AES-128 / ASCON;
% :Phat LoRa, cho ACK;
% if (Co ACK?) then (co)
%   :Cap nhat thong ke, dong bo backlog neu co;
% else (khong)
%   :Luu ban ghi vao SPIFFS (offline);
% endif
% :Phat telemetry (USB);
% :Vao Deep Sleep / Light Sleep;
% stop
% @enduml
% ===================================================================


Chu trình mở đầu bằng sự kiện thức dậy, có thể đến từ bộ định thời theo chu kỳ (lấy mẫu định kỳ) hoặc từ một ngắt sự kiện khẩn cấp. Sau khi khởi tạo ngoại vi, Node đọc lần lượt toàn bộ cảm biến, đưa từng tín hiệu qua bộ lọc tương ứng, rồi tính chỉ số THI và đối chiếu với ma trận ngưỡng trong khối FSM. Nếu phát hiện vượt ngưỡng nguy hiểm, lệnh đóng rơ-le được phát ra ngay lập tức -- đây là nhánh tới hạn về thời gian, phải hoàn tất dưới 50 ms. Chỉ sau khi đã bảo đảm an toàn cục bộ, Node mới tiến hành gắn nhãn thời gian, đóng gói, mã hóa và phát dữ liệu qua LoRa. Việc đặt nhánh điều khiển an toàn \emph{trước} nhánh truyền thông là một quyết định thiết kế then chốt: nó bảo đảm tính mạng vật nuôi được ưu tiên tuyệt đối so với việc báo cáo số liệu.

\subsubsection{Tiền xử lý và lọc nhiễu dữ liệu theo từng tín hiệu}
Dữ liệu thô từ cảm biến luôn chứa thành phần nhiễu, và nếu đưa trực tiếp vào FSM sẽ gây ra hiện tượng dao động đóng/cắt rơ-le liên tục (chattering) quanh ngưỡng. Vì vậy, trước khi ra quyết định, mỗi tín hiệu được làm sạch bằng một bộ lọc số được lựa chọn riêng theo đúng đặc tính vật lý của nó -- một cách tiếp cận tinh tế hơn nhiều so với việc áp một bộ lọc chung cho mọi đại lượng.

Nhiệt độ và độ ẩm biến thiên chậm, liên tục, nhiễu chủ yếu là nhiễu điện tử biên độ nhỏ; do đó chúng được lọc bằng bộ lọc trung bình mũ (Exponential Moving Average -- EMA), vừa làm mượt tốt vừa chỉ cần lưu một giá trị trạng thái nên rất tiết kiệm RAM:
\begin{equation}
    {y_t = \alpha \, x_t + (1-\alpha)\, y_{t-1}, \qquad 0 < \alpha < 1}.
\end{equation}
Trong đó hệ số làm mượt $\alpha$ được chọn ở mức $0{,}3$ để cân bằng giữa độ trễ đáp ứng và mức khử nhiễu. Áp suất khí quyển rất ổn định trong ngắn hạn nên được lọc bằng trung bình trượt (Moving Average -- MA) cửa sổ $N$ mẫu, cho hiệu quả khử nhiễu Gauss tốt nhất:
\begin{equation}
    {y_t = \frac{1}{N}\sum_{i=0}^{N-1} x_{t-i}}.
\end{equation}
Riêng cường độ ánh sáng có nhiều đột biến xung (bóng che, vật bay qua) nên được lọc bằng bộ lọc trung vị (Median), vốn miễn nhiễm với giá trị ngoại lai:
\begin{equation}
    {y_t = \mathrm{median}\{x_{t-N+1}, \dots, x_{t-1}, x_t\}}.
\end{equation}
Ngoài ra, firmware còn hiện thực bộ lọc Kalman một chiều \cite{kalman1960new} như một tùy chọn nâng cao cho các tín hiệu cần ước lượng tối ưu theo nghĩa thống kê. Bảng \ref{Bảng 3.5} tóm tắt chiến lược lọc theo từng tín hiệu cùng lý do lựa chọn.

\begin{table}[H]
\centering
\caption{Chiến lược lọc nhiễu theo từng tín hiệu cảm biến}
\label{Bảng 3.5}
\renewcommand{\arraystretch}{1.4}
\begin{tabular}{p{2.9cm} p{2.7cm} p{2.0cm} p{6.9cm}}
\hline
\textbf{Tín hiệu} & \textbf{Bộ lọc} & \textbf{Tham số} & \textbf{Lý do lựa chọn} \\ \hline
Nhiệt độ, độ ẩm & EMA & $\alpha = 0{,}3$ & Biến đổi chậm, nhiễu nhỏ; bám xu hướng, ít RAM \\ 
Áp suất & Average & $N = 8$ & Rất ổn định; khử nhiễu Gauss tối ưu \\ 
Ánh sáng & Median & $N = 8$ & Nhiều đột biến xung; trung vị miễn nhiễm ngoại lai \\ 
(Dự phòng) & Kalman 1D & $q, r$ & Ước lượng tối ưu khi cần độ chính xác cao \\ \hline
\end{tabular}
\end{table}

\subsubsection{Máy trạng thái hữu hạn (FSM) ra quyết định điều khiển}
Khối ra quyết định được thiết kế dưới dạng một Máy trạng thái hữu hạn (Finite State Machine -- FSM) gồm ba trạng thái: \textbf{AN TOÀN (SAFE)}, \textbf{CẢNH BÁO (WARNING)} và \textbf{KHẨN CẤP (EMERGENCY)}. Việc dùng FSM thay cho chuỗi câu lệnh điều kiện rời rạc mang lại hai lợi ích: thứ nhất, trạng thái hiện tại được lưu giữ nên hệ thống có "trí nhớ", tránh dao động đóng/cắt; thứ hai, có thể cài đặt cơ chế trễ (hysteresis) ở biên chuyển trạng thái để rơ-le không bị bật/tắt liên tục khi giá trị đo dao động sát ngưỡng.

% ===================================================================
% PlantUML cho Hình 3.11 - render và lưu thành: Thuc_Hanh/Image/Ch3_FSM.png
% @startuml
% skinparam backgroundColor white
% [*] --> SAFE
% state SAFE : Trong nguong an toan\nRo-le TAT
% state WARNING : THI>72 hoac NH3>20ppm\nBat quat muc thap
% state EMERGENCY : T>35C hoac NH3>25ppm\nhoac THI>78 -> Bat toan bo quat/phun suong
% SAFE --> WARNING : vuot nguong canh bao
% WARNING --> EMERGENCY : vuot nguong nguy hiem
% EMERGENCY --> WARNING : ha xuong (co tre hysteresis)
% WARNING --> SAFE : ve binh thuong (co tre)
% EMERGENCY --> EMERGENCY : duy tri lam mat
% @enduml
% ===================================================================
\begin{figure}[H]
    \centering
    \includegraphics[width=1.0\linewidth]{Thuc_Hanh/Image/FSM.pdf}
    \caption{FSM ra quyết định điều khiển vi khí hậu}
    \label{fig:ch3_fsm}
\end{figure}

Logic chuyển trạng thái dựa trên việc đối chiếu đồng thời nhiều đại lượng đã lọc với ma trận ngưỡng an toàn sinh học đã thiết lập ở chương 1 và 2: nhiệt độ ($35\,^\circ C$ gây stress, $38\,^\circ C$ nguy kịch), độ ẩm ($75\%$), nồng độ $NH_3$ ($25$ ppm) và chỉ số tổng hợp THI (các mốc 72, 78, 82). Khi ở trạng thái KHẨN CẤP, FSM kích hoạt đồng thời các kênh chấp hành Act1--Act3 để bật toàn bộ quạt hút và hệ thống phun sương; khi hạ về ngưỡng an toàn, FSM chỉ chuyển trạng thái sau khi giá trị đã nằm dưới ngưỡng một khoảng trễ nhất định, nhằm tránh hiện tượng rơ-le rung. Bảng \ref{Bảng 3.6} cụ thể hóa ma trận ngưỡng và hành động tương ứng của FSM.

\begin{table}[H]
\centering
\caption{Ma trận ngưỡng và hành động của FSM điều khiển biên}
\label{Bảng 3.6}
\renewcommand{\arraystretch}{1.4}
\begin{tabular}{p{2.6cm} p{6.2cm} p{5.5cm}}
\hline
\textbf{Trạng thái} & \textbf{Điều kiện kích hoạt} & \textbf{Hành động chấp hành} \\ \hline
AN TOÀN & $T \le 32\,^\circ C$, $RH \le 70\%$, $NH_3 \le 20$ ppm & Tắt toàn bộ rơ-le \\ 
CẢNH BÁO & $32 < T \le 35\,^\circ C$ hoặc THI $> 72$ hoặc $NH_3 > 20$ ppm & Bật quạt mức thấp (Act1) \\
KHẨN CẤP & $T > 35\,^\circ C$ hoặc $NH_3 > 25$ ppm hoặc THI $> 78$ & Bật toàn bộ quạt + phun sương (Act1--Act3) \\ \hline
\end{tabular}
\end{table}

\subsection{Áp dụng phương pháp đồng thiết kế Phần cứng/Phần mềm}
Dựa trên các ràng buộc cụ thể về độ trễ và năng lượng, các chức năng hệ thống đã được đánh giá và phân hoạch tối ưu giữa phần cứng và phần mềm, với kết quả cuối cùng được tổng hợp tại Bảng \ref{Bảng 3.7}.

\begin{table}[H]
\centering
\caption{Phân hoạch chức năng Phần cứng/Phần mềm cho Node}
\label{Bảng 3.7}
\renewcommand{\arraystretch}{1.4}
\begin{tabular}{p{3.5cm} p{4.5cm} p{6.5cm}}
\hline
\textbf{Chức năng} & \textbf{Phân hoạch} & \textbf{Lý do} \\ \hline
Mã hóa AES-128 & Phần cứng & Engine AES trong ESP32-S3 chạy nhanh, giảm chu kỳ CPU và năng lượng \\ 
Mã hóa ASCON-128 & Phần mềm & Thuật toán hạng nhẹ, không có engine phần cứng; để so sánh đối chứng \\ 
Lọc nhiễu, FSM & Phần mềm (Core 1) & Cần linh hoạt thay đổi ngưỡng/tham số; tải tính toán nhỏ \\ 
Điều chế LoRa CSS & Phần cứng (SX1278) & Xử lý lớp vật lý chuyên dụng, MCU không kham nổi \\ 
Cắt nguồn ngoại vi & Phần cứng (nguồn) & Quyết định không dùng load switch cho cảm biến khí MOX (cân bằng nhiệt) \\ 
Gắn nhãn thời gian & Phần cứng (DS3231) & RTC cấp mốc thời gian chính xác\\ \hline
\end{tabular}
\end{table}

% \begin{figure}[H]
%     \centering
%     \includegraphics[width=0.85\linewidth]{Thuc_Hanh/Image/HWSWCoDesign.png}
%     \caption{Áp dụng chu trình đồng thiết kế HW/SW cho thiết bị Node}
%     \label{fig:ch3_codesign}
% \end{figure}

Trong quá trình thiết kế, vòng lặp tái phân hoạch (Re-partitioning) đã được kích hoạt ở ít nhất hai điểm quan trọng. Thứ nhất, ban đầu thuật toán mã hóa dự kiến chạy hoàn toàn bằng thư viện phần mềm; tuy nhiên khi đo thử thời gian mã hóa AES bằng phần mềm thuần làm tăng đáng kể thời gian thức của Node, nhóm thiết kế đã quyết định đẩy AES-128 sang bộ tăng tốc mã hóa phần cứng tích hợp sẵn trong ESP32-S3, kéo thời gian mã hóa xuống mức dưới mili-giây -- đúng tinh thần "đẩy tác vụ nặng sang phần cứng" của chu trình co-design. Thứ hai, ở khâu quản lý năng lượng, phương án ban đầu theo lý thuyết là dùng công tắc tải (Load Switch) cắt nguồn toàn bộ cảm biến khi ngủ; nhưng phân tích đặc tính nhiệt động học của cảm biến khí MOX cho thấy việc cắt nguồn sẽ phá vỡ cân bằng nhiệt và gây trôi baseline nghiêm trọng, nên thiết kế đã tái phân hoạch theo hướng cấp nguồn liên tục cho mảng cảm biến khí và chỉ tối ưu phần ngủ sâu ở phía vi điều khiển. Hai ví dụ này cho thấy đồng thiết kế không phải là một quy trình tuyến tính một chiều, mà là một vòng lặp tinh chỉnh liên tục giữa hai miền phần cứng và phần mềm cho đến khi đạt điểm cân bằng tối ưu.
\subsection{Thiết kế phân hệ cập nhật phần mềm không dây (OTA) trên nền tảng Edge Computing}

\subsubsection{Bài toán giới hạn băng thông và cơ sở đề xuất kiến trúc Hybrid OTA}
Trong mạng cảm biến không dây (WSN) ứng dụng cho nông nghiệp thông minh, giao thức LoRa được lựa chọn làm xương sống truyền thông nhờ khả năng đâm xuyên vật cản xuất sắc và tầm phủ sóng rộng với năng lượng tối thiểu. Tuy nhiên, theo định lý Shannon-Hartley về dung lượng kênh truyền, để đạt được độ nhạy thu (Receiver Sensitivity) cực thấp ở mức -137 dBm (với Spreading Factor SF = 12), hệ thống phải đánh đổi bằng tốc độ truyền dữ liệu (Data Rate) vô cùng nhỏ, thường chỉ dao động ở mức vài chục đến vài trăm bit/giây (bps). 

Mặt khác, đối với nền tảng điện toán biên (Edge Computing) sử dụng vi điều khiển ESP32-S3, một tệp nhị phân Firmware hoàn chỉnh bao gồm hệ điều hành thời gian thực FreeRTOS, ngăn xếp giao thức mạng (Network Stack) và các thuật toán học máy bù trôi cảm biến thường có kích thước từ 1.2 MB đến 1.5 MB. Phép tính toán học cơ bản cho thấy, nếu truyền tải khối lượng dữ liệu này qua kênh LoRa với tốc độ thực tế 250 bps, thiết bị sẽ mất hơn 13 giờ truyền liên tục để hoàn tất. Điều này không chỉ vi phạm nghiêm ngặt quy định về chu kỳ hoạt động (Duty Cycle, thường giới hạn ở 1\%) của dải tần ISM, mà còn gây cạn kiệt dung lượng pin Lithium và làm sụp đổ toàn bộ mạng lưới do chiếm dụng kênh truyền.

Từ những mâu thuẫn vật lý trên, đồ án đề xuất và thiết kế một kiến trúc cập nhật lai (Hybrid OTA Architecture). Phương pháp này phân tách rõ ràng luồng tín hiệu điều khiển (Control Plane) và luồng dữ liệu (Data Plane):
\begin{itemize}
    \item \textbf{Kênh Điều khiển (LoRa):} Hoạt động 100\% thời gian trong chu kỳ thức. Kênh này chỉ dùng để truyền các lệnh đánh thức (Wake-up triggers), tham số cấu hình và cờ hiệu báo cáo có bản cập nhật mới (OTA Flag).
    \item \textbf{Kênh Dữ liệu (Wi-Fi 2.4 GHz):} Trạng thái mặc định là tắt hoàn toàn (RF Disabled) để tiết kiệm năng lượng. Kênh này chỉ được đánh thức khi có cờ hiệu OTA, thiết lập một kết nối tốc độ cao (lên tới 150 Mbps theo chuẩn 802.11n) vào mạng nội bộ để tải file nhị phân trong thời gian tính bằng giây, sau đó lập tức tự ngắt.
\end{itemize}

\subsubsection{Quy hoạch không gian bộ nhớ Flash (Memory Partitioning)}
Để hệ thống có khả năng tiếp nhận phần mềm mới mà không làm gián đoạn mã lệnh đang thực thi, bộ nhớ Flash SPI ngoại vi (dung lượng 8 MB) của ESP32-S3 phải được quy hoạch theo kiến trúc Ngân hàng kép (Dual-Bank). Thay vì sử dụng một phân vùng ứng dụng nguyên khối, bộ nhớ được cắt lớp theo thiết kế chuẩn công nghiệp, minh họa tại Hình \ref{fig:memory_map}.

\begin{figure}[H]
\centering
\resizebox{1.0\textwidth}{!}{
\begin{tikzpicture}[
    >=stealth, font=\small,
    % Định nghĩa style các khối
    block/.style={rectangle, draw, thick, minimum width=8cm, minimum height=0.9cm, align=center},
    sys/.style={fill=gray!15},
    app_active/.style={fill=blue!10, thick},
    app_idle/.style={fill=white, dashed, thick},
    data/.style={fill=yellow!10}
]

% Tọa độ tuyệt đối (Grid) theo trục dọc
\node[block, sys] (boot) at (0, 0) {\textbf{Bootloader} \\ Offset: \texttt{0x0000} | Size: 32 KB};
\node[block, sys] (part) at (0, -1.2) {\textbf{Partition Table} \\ Offset: \texttt{0x8000} | Size: 4 KB};
\node[block, sys] (otadata) at (0, -2.4) {\textbf{OTA Data (Chứa cờ Boot)} \\ Offset: \texttt{0x9000} | Size: 8 KB};
\node[block, data] (nvs) at (0, -3.6) {\textbf{NVS (Non-Volatile Storage)} \\ Offset: \texttt{0xB000} | Size: 20 KB};

\node[block, app_active, minimum height=1.5cm] (app0) at (0, -5.3) {\textbf{Phân vùng Factory / APP0} \\ Trạng thái: Đang thực thi mã lệnh (Active) \\ Offset: \texttt{0x10000} | Size: 3 MB};

\node[block, app_idle, minimum height=1.5cm] (app1) at (0, -7.3) {\textbf{Phân vùng OTA / APP1} \\ Trạng thái: Nhàn rỗi, chờ ghi dữ liệu (Idle) \\ Offset: \texttt{0x310000} | Size: 3 MB};

% Trục địa chỉ bên trái
\draw[->, thick] (-4.5, 0.5) -- (-4.5, -8.5) node[below] {Địa chỉ bộ nhớ};

% Đường nối mô tả logic
\draw[->, thick, red] (otadata.east) -- ++(1.5,0) |- node[pos=0.25, right] {Trỏ vào phân vùng Active} (app0.east);
\end{tikzpicture}
}
\caption{Quy hoạch phân bổ bộ nhớ Flash phục vụ cơ chế cập nhật OTA (Dual-Bank)}
\label{fig:memory_map}
\end{figure}

% Khối \texttt{OTA Data} đóng vai trò như một thanh ghi định tuyến. Khi thiết bị khởi động, Bootloader không nạp ngay hệ điều hành mà đọc giá trị tại \texttt{OTA Data} để xác định đâu là phân vùng chứa Firmware hợp lệ và mới nhất, sau đó mới ánh xạ (Map) vùng nhớ đó vào bộ đệm lệnh (Instruction Cache) của CPU.
Như minh họa tại Hình \ref{fig:memory_map}, bộ nhớ Flash của ESP32-S3 được quy hoạch theo kiến trúc Ngân hàng kép (Dual-Bank) nhằm đảm bảo an toàn tuyệt đối cho tiến trình cập nhật OTA.

Ở các dải địa chỉ thấp, Bootloader và Bảng phân vùng (Partition Table) làm nhiệm vụ khởi tạo phần cứng và định tuyến bộ nhớ. Phân vùng NVS (Non-Volatile Storage) được dùng để lưu trữ các thông số bất biến như cấu hình vô tuyến và trạng thái máy trạng thái hữu hạn (FSM).

Thành phần cốt lõi điều phối quá trình OTA là phân vùng OTA Data (\texttt{0x9000}). Khởi động hệ thống, Bootloader bắt buộc phải đọc \texttt{OTA Data} để tìm vị trí firmware hợp lệ. Ở trạng thái mặc định (mũi tên đỏ), cờ Boot trỏ vào APP0 (Phân vùng Active), nơi vi điều khiển đang thực thi trực tiếp hệ điều hành FreeRTOS và các thuật toán đo lường biên.

Để đảm bảo tính chống lỗi (Fail-safe), hệ thống sử dụng APP1 làm phân vùng dự phòng nhàn rỗi. Khi có bản cập nhật OTA, luồng dữ liệu mới sẽ được ghi vật lý vào \texttt{APP1}, cô lập hoàn toàn với \texttt{APP0} đang chạy. Chỉ khi tệp tải về hoàn chỉnh và vượt qua khâu xác thực mã băm SHA-256, \texttt{OTA Data} mới được ghi đè để trỏ hệ thống sang \texttt{APP1} ở lần khởi động tiếp theo. Cơ chế cô lập không gian này đảm bảo rằng: nếu sự cố rớt mạng hoặc sập nguồn xảy ra giữa chừng, phân vùng \texttt{APP0} không bị ảnh hưởng, thiết bị sẽ không bao giờ bị treo cứng mà vẫn duy trì vận hành an toàn trên phiên bản cũ.

\subsubsection{Thiết kế luồng điều phối tác vụ Hệ điều hành (FreeRTOS Task Scheduling)}
Trong môi trường hệ điều hành thời gian thực, tiến trình OTA không được phép gây ra hiện tượng khóa chết (Deadlock) hoặc cạnh tranh tài nguyên với các tác vụ sống còn khác. Bảng \ref{tab:rtos_tasks} trình bày ma trận phân bổ tác vụ, trong đó \texttt{Task\_OTA} được đặt ở mức ưu tiên thấp để đảm bảo an toàn hệ thống.

\begin{table}[H]
  \centering
  \caption{Quy hoạch và phân bổ các tác vụ FreeRTOS trong hệ thống Edge Node}
  \label{tab:rtos_tasks}
  \begin{tabular}{@{}p{3.5cm} c c c p{4.5cm}@{}}
      \toprule
      \textbf{Tên tác vụ} & \textbf{Lõi CPU} & \textbf{Mức Ưu tiên} & \textbf{Stack Size} & \textbf{Chức năng cốt lõi} \\
      \midrule
      \texttt{Task\_FSM\_Sensors} & Core 1 & 3 (Cao) & 4096 B & Thu thập dữ liệu, thuật toán bù trôi điện toán biên. \\
      \midrule
      \texttt{Task\_LoRa\_Comm} & Core 1 & 4 (Rất cao) & 4096 B & Xử lý ngắt vô tuyến, truyền/nhận Uplink, Downlink. \\
      \midrule
      \texttt{Task\_OTA\_Worker} & Core 0 & 1 (Thấp) & 8192 B & Thiết lập Wi-Fi, tải khối HTTP, mã hóa SHA-256. Ngủ đông khi không có lệnh. \\
      \bottomrule
  \end{tabular}
\end{table}

Khi \texttt{Task\_LoRa\_Comm} nhận được cờ lệnh OTA, nó sẽ phát một tín hiệu Semaphore nhị phân để đánh thức \texttt{Task\_OTA\_Worker}. Lúc này, \texttt{Task\_OTA\_Worker} sẽ gọi hàm API nội bộ để tạm đình chỉ (Suspend) \texttt{Task\_FSM\_Sensors}, giải phóng bộ nhớ Heap để cấp phát cho bộ đệm nhận gói tin TCP/IP, đảm bảo quá trình ghi Flash diễn ra liền mạch không bị phân mảnh.

\subsection{Kết luận chương}
Chương 3 đã thực hiện trọn vẹn quá trình phân tích và thiết kế hệ thống theo hướng từ trên xuống. Bắt đầu từ việc lượng hóa các yêu cầu chức năng và phi chức năng thành những tham số cụ thể, chương đã biện luận lựa chọn từng khối công nghệ lõi -- từ vi điều khiển ESP32-S3, module vô tuyến LoRa SX1278, khối nguồn tiết kiệm điện, tổ hợp cảm biến vi khí hậu và khí độc, cho đến khối thời gian thực và khối chấp hành -- với đầy đủ lý do kỹ thuật cho mỗi quyết định. Chương đã thiết kế kiến trúc phần mềm đa nhiệm trên FreeRTOS với cơ chế gán nhân tách biệt vô tuyến khỏi điện toán biên, đồng thời chi tiết hóa các thuật toán cốt lõi gồm chu trình lấy mẫu, lọc nhiễu theo từng tín hiệu và máy trạng thái ra quyết định. Cuối cùng, phương pháp đồng thiết kế Phần cứng/Phần mềm đã được áp dụng cụ thể để phân hoạch chức năng một cách tối ưu. 





\newpage

\section*{CHƯƠNG 4. TRIỂN KHAI ĐỀ TÀI}
\addcontentsline{toc}{section}{\numberline {} CHƯƠNG 4. TRIỂN KHAI ĐỀ TÀI}
\setcounter{section}{4}
\setcounter{subsection}{0}
\setcounter{figure}{0}
\setcounter{table}{0}
Chương này trình bày quá trình hiện thực hóa hệ thống thành một sản phẩm hoàn chỉnh, chạy được trên phần cứng thực tế. Nội dung được chia thành hai phần lớn tương ứng với hai trục công việc song song của một dự án nhúng: phần triển khai phần cứng đi sâu vào việc giải phẫu sơ đồ nguyên lý (Schematic) và quy hoạch mạch in (PCB Layout) của thiết bị Node, phân tích vai trò, chức năng, điểm mạnh -- điểm yếu cùng lý do thiết kế của từng khối linh kiện; phần triển khai phần mềm trình bày chi tiết kiến trúc firmware đa nhiệm, các lớp driver, thuật toán Điện toán biên, cơ chế bảo mật, quản lý năng lượng, lưu trữ ngoại tuyến cũng như công cụ kiểm thử Hardware-In-the-Loop được xây dựng riêng cho đề tài.

\subsection{Triển khai phần cứng thiết bị Node}
Phần cứng của thiết bị Node được thiết kế dưới dạng một bo mạch in tùy biến (Custom PCB) hai lớp, tích hợp toàn bộ các khối chức năng trên cùng một nền FR4 thay vì lắp ghép các module rời rạc. 
\begin{figure}[H]
    \centering
    \includegraphics[width=1\linewidth]{Thuc_Hanh/Image/Schematic.png}
    \caption{Sơ đồ nguyên lý tổng thể của thiết bị Node}
    \label{Hình 4.1}
\end{figure}
Cách tiếp cận tích hợp này tuy đòi hỏi nhiều công sức thiết kế và gia công hơn, nhưng đổi lại cho phép kiểm soát chặt chẽ đường đi của tín hiệu cao tần, tối ưu hóa diện tích, nâng cao độ tin cậy cơ học và tạo ra một sản phẩm có hình thức tiệm cận sản phẩm công nghiệp. Toàn bộ sơ đồ nguyên lý được tổ chức thành các khối chức năng tách bạch như Hình \ref{Hình 4.1}, giúp việc dò lỗi và bảo trì trở nên thuận tiện.



\subsubsection{Khối cấp nguồn và bảo vệ ngõ vào USB-C}
Mọi hệ thống điện tử đều bắt đầu từ nguồn, và độ ổn định của nguồn quyết định trực tiếp độ tin cậy của toàn hệ. Hình \ref{Hình 4.2} cho thấy ngõ vào nguồn của Node sử dụng đầu cắm USB Type-C (linh kiện J2, USB\_C\_Receptacle\_USB2.0\_16P) -- một chuẩn kết nối hiện đại, có thể cắm thuận nghịch và chịu được dòng lớn. Để chuẩn USB-C hoạt động đúng ở chế độ thiết bị nhận nguồn (UFP), hai chân cấu hình CC1/CC2 được kéo xuống đất qua hai điện trở $5{,}1\,\text{k}\Omega$ (R1, R2); đây là yêu cầu bắt buộc của đặc tả USB-C để nguồn cấp nhận diện và cấp dòng $5\,\text{V}$.
\begin{figure}[H]
    \centering
    \includegraphics[width=0.5\linewidth]{Thuc_Hanh/Image/USBC.png}
    \caption{Khối ngõ vào USB-C với chuỗi bảo vệ nhiều tầng}
    \label{Hình 4.2}
\end{figure}
Điểm đáng chú ý nhất của khối này là chuỗi bảo vệ ngõ vào nhiều tầng, thể hiện tư duy thiết kế hướng tới độ bền công nghiệp. Tầng thứ nhất là cầu chì tự phục hồi (Polyfuse F1): khi dòng tải vượt ngưỡng do chập mạch, điện trở của polyfuse tăng vọt để ngắt dòng, và tự khôi phục khi sự cố được loại bỏ -- bảo vệ chống quá dòng mà không cần thay cầu chì. Tầng thứ hai là cuộn chặn nhiễu (Ferrite Bead FB1) mắc nối tiếp trên đường nguồn, có tác dụng như một bộ lọc thông thấp đối với nhiễu tần số cao đi vào từ cáp USB, ngăn nhiễu lan vào mạch nhạy cảm. Tầng thứ ba là điốt Schottky (D3, N5817) chống cắm ngược cực nguồn. Quan trọng nhất, tầng thứ tư là IC chống phóng tĩnh điện chuyên dụng USBLC6-2SC6 (U2) đặt sát đầu cắm, ghim các xung ESD biên độ hàng kilovolt xuống đất trước khi chúng kịp phá hủy các đường tín hiệu dữ liệu USB và nguồn.



Điểm mạnh của thiết kế khối nguồn ngõ vào là độ an toàn cao, chống được hầu hết các sự cố điện thường gặp ngoài hiện trường (chập, cắm ngược, ESD, nhiễu). Điểm yếu là việc thêm nhiều linh kiện bảo vệ làm tăng chi phí và chiếm thêm diện tích mạch; tuy nhiên trong bối cảnh thiết bị phải hoạt động 24/7 trong môi trường chuồng trại ẩm ướt và nhiều nhiễu, sự đánh đổi này hoàn toàn xứng đáng.

\subsubsection{Khối sạc pin và quản lý nguồn pin Li-ion}
Trung tâm của khối là IC sạc tuyến tính TP4056 (U3, gói ESOP8 có tản nhiệt đáy), chuyên dụng cho pin Lithium một cell. TP4056 thực thi thuật toán sạc hai giai đoạn dòng không đổi / áp không đổi (CC/CV) như đã phân tích ở Chương 3: dòng sạc được lập trình qua điện trở $R_{PROG}$ (R13) theo công thức $I_{CHG} = 1200/R_{PROG}$, và quá trình sạc tự kết thúc khi điện áp pin đạt $4{,}2\,\text{V}$ và dòng sụt xuống dưới ngưỡng. Hai đèn LED chỉ thị (D7 BLUE cho trạng thái đang sạc CHRG, và đường STDBY báo sạc đầy) cung cấp phản hồi trực quan cho người vận hành.

Một linh kiện quan trọng đi kèm là MOSFET kênh P AO3401A (Q1) đặt trên đường dẫn nguồn hệ thống V\_SYS. Linh kiện này thực hiện chức năng định tuyến nguồn và chống dòng ngược (ideal-diode / power-path): khi có nguồn USB, hệ thống được cấp điện từ USB đồng thời sạc pin; khi rút USB, MOSFET dẫn để pin cấp điện cho tải, đồng thời ngăn dòng chảy ngược từ pin về cổng USB.

\begin{figure}[H]
    \centering
    \includegraphics[width=0.65\linewidth]{Thuc_Hanh/Image/BatteryCharger.png}
    \caption{Khối sạc pin Li-ion dùng TP4056 và MOSFET định tuyến nguồn AO3401A}
    \label{Hình 4.3}
\end{figure}

Khối sạc trong Hình \ref{Hình 4.3} có ưu điểm là tích hợp đầy đủ, giá thành thấp và dễ kiểm soát; thuật toán CC/CV bảo vệ tốt tuổi thọ pin. Điểm yếu cố hữu của IC sạc tuyến tính là tỏa nhiệt khi sạc dòng lớn (năng lượng dư biến thành nhiệt trên IC), do đó thiết kế phải chú ý diện tích đồng tản nhiệt dưới gói ESOP8. Lựa chọn TP4056 là hợp lý cho dòng sạc vừa phải và độ phức tạp thấp của một thiết bị một cell.

\subsubsection{Khối ổn áp hạ áp LDO}
Điện áp pin Li-ion biến thiên trong dải rộng từ $4{,}2\,\text{V}$ (đầy) xuống $3{,}0\,\text{V}$ (cạn), trong khi vi điều khiển, cảm biến và module LoRa đều yêu cầu một mức $3{,}3\,\text{V}$ ổn định. Khối ổn áp tuyến tính hạ áp (LDO) đảm nhiệm việc tạo ra điện áp phẳng này. Như đã luận giải kỹ ở Chương 3, đề tài lựa chọn IC LDO AP2112K-3.3 (U4) thay cho dòng AMS1117 phổ biến nhưng tốn điện. AP2112K dùng phần tử dẫn là MOSFET kênh P nên có điện áp rơi (dropout) cực thấp chỉ khoảng $250\,\text{mV}$ ở dòng tải $600\,\text{mA}$, và quan trọng hơn là dòng tĩnh ($I_Q$) chỉ vào khoảng $55\,\mu\text{A}$ \cite{diodes2015ap2112}. Điện áp rơi thấp đồng nghĩa với việc hệ thống vẫn duy trì được $3{,}3\,\text{V}$ ngay cả khi pin đã xả xuống mức thấp, tận dụng tối đa dung lượng pin; còn dòng tĩnh thấp giúp giảm hao phí ở trạng thái chờ. Hai tụ điện $22\,\mu\text{F}$ (C2, C3) ở ngõ vào và ngõ ra bảo đảm tính ổn định của vòng hồi tiếp và lọc gợn sóng. Một đèn LED xanh (D2) báo hiệu nguồn $3{,}3\,\text{V}$ đã sẵn sàng. Thiết kế khối LDO được thể hiện trong Hình \ref{Hình 4.4}.
\begin{figure}[H]
    \centering
    \includegraphics[width=0.9\linewidth]{Thuc_Hanh/Image/LDO.png}
    \caption{Thiết kế khối LDO}
    \label{Hình 4.4}
\end{figure}
Điểm mạnh của giải pháp LDO là đơn giản, không sinh nhiễu chuyển mạch (khác với nguồn switching), rất phù hợp cho mạch có khối vô tuyến nhạy nhiễu như LoRa. Điểm yếu là hiệu suất giảm khi chênh lệch áp vào -- ra lớn (năng lượng dư tỏa thành nhiệt); tuy nhiên với chênh lệch nhỏ giữa pin và $3{,}3\,\text{V}$, hiệu suất của AP2112K trong ứng dụng này là chấp nhận được.

\subsubsection{Khối vi điều khiển trung tâm ESP32-S3 và mạch ngoại vi}
Bộ xử lý trung tâm của Node là module ESP32-S3-WROOM-1 (U1) đã được biện luận lựa chọn ở Chương 3. Trong sơ đồ nguyên lý như Hình \ref{Hình 4.5}, ngoài bản thân module, khối này còn bao gồm các mạch ngoại vi thiết yếu để vi điều khiển hoạt động ổn định và nạp được firmware. Chân Enable (EN) được kéo lên $3{,}3\,\text{V}$ qua điện trở $10\,\text{k}\Omega$ (R10) kết hợp tụ $1\,\mu\text{F}$ (C4) tạo thành mạch reset nguồn (power-on reset) bảo đảm chip khởi động sạch; nút nhấn RESET (SW1) cho phép khởi động lại thủ công. Nút BOOT (SW3) nối chân GPIO0 cho phép đưa chip vào chế độ nạp firmware. Các tụ lọc nguồn $100\,\text{nF}$ và $10\,\mu\text{F}$ (C5, C6) đặt sát chân nguồn module để khử nhiễu cao tần và ổn định điện áp khi dòng tải biến động đột ngột (decoupling).

Để phục vụ nạp và gỡ lỗi, một đầu cắm lập trình (J1, Conn\_02x03) đưa ra các đường UART (TXD0/RXD0) và các tín hiệu điều khiển nạp. Ngoài ra, một nút nhấn người dùng USER\_SW (SW2, qua chân có điện trở kéo lên R11 và bảo vệ ESD D6) được bố trí để phục vụ các thao tác tương tác tại chỗ như hiệu chỉnh hoặc kích hoạt chế độ kiểm thử. Mỗi nút nhấn và đường tín hiệu nhạy cảm đều được gắn điốt bảo vệ ESD (các linh kiện ESD9B3.3ST5G), thể hiện sự nhất quán trong triết lý thiết kế chống chịu môi trường khắc nghiệt.

\begin{figure}[H]
    \centering
    \includegraphics[width=0.9\linewidth]{Thuc_Hanh/Image/ESP32_Core.png}
    \caption{Module ESP32-S3-WROOM-1}
    \label{Hình 4.5}
\end{figure}

Điểm mạnh của khối xử lý là năng lực tính toán lõi kép mạnh mẽ, bộ nhớ dồi dào, tích hợp sẵn bộ tăng tốc mã hóa và USB gốc, cùng hệ sinh thái phần mềm ESP-IDF trưởng thành. Điểm cần lưu ý là module có khá nhiều chân và mức tiêu thụ ở trạng thái hoạt động không nhỏ, đòi hỏi thiết kế nguồn và bố trí chân cẩn thận; tuy nhiên các khả năng tiết kiệm điện ở chế độ ngủ sâu bù đắp tốt cho điều này.

\subsubsection{Khối cảm biến, bus I2C và khối thời gian thực}
Ba cảm biến vi khí hậu SHT31, BMP180 và BH1750 cùng IC thời gian thực DS3231 được mắc trên bus I2C dùng chung, với các điện trở kéo lên $4{,}7\,\text{k}\Omega$ (R16, R17) như trong Hình \ref{fig:ch4_sensor}. 
\begin{figure}[H]
    \centering
    {\includegraphics[width=0.9\linewidth]{Thuc_Hanh/Image/Sensor.png}}
    \caption{Sơ đồ nguyên lý khối cảm biến và bus I2C dùng chung của Node}
    \label{fig:ch4_sensor}
\end{figure}
Mỗi cảm biến được đưa ra một đầu nối riêng (các đầu nối J14--J19 dạng Conn\_01x05 và Conn\_01x04, J12/J13 dạng Conn\_01x06 cho RTC) cho phép tháo lắp linh hoạt, thuận tiện cho việc thay thế hoặc bố trí cảm biến ở các vị trí khác nhau trong vỏ hộp. Cách đưa cảm biến ra đầu nối thay vì hàn cứng trên bo là một quyết định thực dụng: nó cho phép đặt cảm biến ánh sáng hướng ra ngoài, cảm biến nhiệt -- ẩm ở vị trí thông thoáng, đồng thời dễ thay khi cảm biến hỏng.

Hai cảm biến khí tương tự MQ-135 (đo $NH_3$, đầu nối J9) và CH4 (đầu nối J11) đưa tín hiệu điện áp về các chân ADC của ESP32-S3; cảm biến CO2 giao tiếp số qua UART (đầu nối J7/J8). Bên cạnh đó, bo mạch còn trang bị một đầu nối QWIIC tiêu chuẩn (J4) với các điện trở kéo lên $2{,}2\,\text{k}\Omega$ (R7, R8) tuân theo quy ước màu dây của hệ sinh thái QWIIC (GND đen, VCC đỏ, SDA xanh dương, SCL vàng), cho phép cắm nóng (hot-plug) thêm các cảm biến I2C mở rộng trong tương lai mà không cần sửa mạch -- một thiết kế hướng tới khả năng mở rộng.



Điểm mạnh của thiết kế khối cảm biến là tính module hóa và khả năng mở rộng cao, cùng việc tách bạch giao tiếp số (I2C, UART) và tương tự (ADC) theo đúng đặc tính từng cảm biến. Điểm yếu là việc dùng nhiều đầu nối làm tăng số điểm tiếp xúc cơ khí -- nơi tiềm ẩn rủi ro tiếp xúc kém trong môi trường rung động và ẩm; điều này được khắc phục một phần bằng việc phủ keo chống ẩm sau khi lắp ráp.

\subsubsection{Khối truyền thông vô tuyến LoRa}
Module LoRa SX1278 được kết nối với ESP32-S3 qua giao diện SPI cùng các chân điều khiển RESET và ngắt DIO0, thông qua các đầu cắm J3/J10 (Conn\_01x08) như trong Hình \ref{fig:ch4_lora}. Việc bố trí LoRa dưới dạng module cắm rời (thay vì hàn chip SX1278 trực tiếp lên bo) là một lựa chọn cân nhắc giữa độ phức tạp thiết kế và tính linh hoạt: module rời đã được nhà sản xuất phối hợp trở kháng và kiểm định sẵn phần cao tần, giúp giảm rủi ro thiết kế RF -- vốn là phần khó nhất của một bo mạch vô tuyến -- đồng thời cho phép thay đổi module hoặc băng tần dễ dàng. Đường tín hiệu SPI được giữ ngắn và tránh chạy song song với các đường nhiễu để bảo toàn tính toàn vẹn tín hiệu ở tốc độ cao.

\begin{figure}[H]
    \centering
    \includegraphics[width=0.9\linewidth]{Thuc_Hanh/Image/LoraPCB.png}
    \caption{Module thu phát vô tuyến LoRa kết nối qua SPI và chân ngắt DIO0}
    \label{fig:ch4_lora}
\end{figure}

Điểm mạnh là độ tin cậy cao của phần RF và tính linh hoạt khi nâng cấp; điểm yếu là đầu cắm module làm tăng chiều cao tổng thể và một điểm tiếp xúc cơ khí nữa cần được cố định chắc chắn.

\subsubsection{Khối chấp hành và chỉ thị trạng thái}
Khối chấp hành gồm ba kênh điều khiển Act1, Act2, Act3 (các đầu nối J20, J21, J22 dạng Conn\_01x03) đưa tín hiệu mức logic ra các rơ-le bên ngoài như trong Hình \ref{fig:ch4_act}). 
\begin{figure}[H]
    \centering
    \includegraphics[width=0.3\linewidth]{Thuc_Hanh/Image/Actuator.png}
    \caption{Chân điều khiển cho khối Actuators}
    \label{fig:ch4_act}
\end{figure}
Việc dùng ba đầu nối ba chân (tín hiệu, nguồn, đất) cho phép cấp nguồn và điều khiển module rơ-le ngoài một cách gọn gàng. Thiết kế tách ba kênh độc lập, như đã nêu ở Chương 3, cho phép FSM điều khiển riêng rẽ quạt hút, phun sương và đèn sưởi.

Về chỉ thị trạng thái, ngoài các LED nguồn và LED sạc, bo mạch tích hợp một đèn LED RGB (D1, CLV1L-FKB) điều khiển qua ba chân GPIO (R\_LED, G\_LED, B\_LED) qua các điện trở hạn dòng $1\,\text{k}\Omega$ như trong Hình \ref{fig:ch4_led}.
\begin{figure}[H]
    \centering
    \includegraphics[width=0.7\linewidth]{Thuc_Hanh/Image/Led.png}
    \caption{Khối đèn báo hiệu trạng thái}
    \label{fig:ch4_led}
\end{figure}
Đèn RGB cho phép mã hóa nhiều trạng thái hệ thống bằng màu sắc (ví dụ: xanh = an toàn, vàng = cảnh báo, đỏ = khẩn cấp, nhấp nháy = đang truyền), cung cấp một kênh phản hồi trực quan giàu thông tin ngay tại hiện trường mà không cần thiết bị phụ trợ.

\subsubsection{Thiết kế và quy hoạch mạch in (PCB Layout)}
Sau khi hoàn thiện sơ đồ nguyên lý, toàn bộ linh kiện được quy hoạch lên một bo mạch in hai lớp (2-layer FR4) với kích thước khoảng $78 \times 80\,\text{mm}$. Quá trình layout  phải tuân thủ nghiêm ngặt các nguyên tắc về tính toàn vẹn tín hiệu và tương thích điện từ (EMI/EMC).

Nguyên tắc quan trọng nhất được áp dụng là thiết lập vùng cấm đồng (Keep-out Zone) tại khu vực ăng-ten của module ESP32-S3 và module LoRa. Khu vực này tuyệt đối không phủ lớp đồng nền (ground pour) và không cho đường mạch nào chạy qua, nhằm tránh hiện tượng phản xạ và hấp thụ năng lượng cao tần làm suy giảm hiệu suất phát xạ và độ nhạy thu của ăng-ten. Thứ hai, các đường nguồn được vẽ rộng để giảm sụt áp và trở kháng, trong khi một mặt phẳng đất (ground plane) liền mạch ở lớp dưới đóng vai trò vừa là đường về dòng điện trở kháng thấp vừa là lớp chắn nhiễu. Thứ ba, các tụ lọc nguồn (decoupling) được đặt sát chân nguồn của các IC để khử nhiễu hiệu quả nhất. Thứ tư, khối nguồn công suất (sạc, LDO) được bố trí tách biệt khỏi khối tín hiệu nhạy cảm và khối cao tần để hạn chế nhiễu chéo (cross-talk). Cuối cùng, các lỗ bắt vít (H1--H5) ở các góc bảo đảm cố định cơ học chắc chắn bo mạch vào vỏ hộp. 

\textbf{Bước 1 -- Phân khu mạch và quy hoạch tầng (Floorplanning).} Bo mạch được chia thành bốn khu vực chức năng với biên giới vật lý rõ ràng, đặt theo dòng năng lượng từ \emph{cao} sang \emph{thấp} và theo độ nhạy nhiễu từ \emph{nhỏ} sang \emph{lớn}. Khu vực \emph{nguồn vào -- sạc -- LDO} đặt ở mép trái bo, nơi cổng USB-C dễ tiếp cận và xa nhất các đường tín hiệu nhạy cảm. Khu vực \emph{vi điều khiển ESP32-S3 và mạch ngoại vi số} đặt ở trung tâm, giữa hai khu vực còn lại. Khu vực \emph{cảm biến và bus I2C} đặt phía dưới, hướng ra các đầu nối ngoại vi. Cuối cùng, khu vực \emph{vô tuyến LoRa và ăng-ten} bố trí ở mép phải, đối diện hoàn toàn với khu vực nguồn để tránh nhiễu từ trường biến thiên của các tụ điện công suất cảm ứng vào đường thu tín hiệu. Cách phân khu này thực hiện đúng nguyên tắc \emph{phân vùng và che chắn} của Ott: bốn khu vực được phân tách bằng các dải đồng nối đất rộng tối thiểu 2 mm, kết hợp các via stitching mật độ cao ($\lambda/20$ ở 433 MHz tương đương khoảng cách 35 mm, nên đặt via stitching mỗi 2--3 mm là dư an toàn).

\textbf{Bước 2 -- Thiết lập số lớp và stack-up.} Bo sử dụng nền FR-4 tiêu chuẩn dày 1{,}6 mm, hai lớp đồng phủ 1 oz (35 $\mu$m). Lớp Top chứa toàn bộ linh kiện SMD và phần lớn đường tín hiệu, lớp Bottom chủ yếu là mặt phẳng đất liền mạch (Continuous Ground Plane) cùng một số đường tín hiệu phụ. Quyết định không dùng nguồn ở lớp Bottom mà gộp với đất là một đánh đổi có chủ đích cho bo hai lớp: nó giúp mọi đường tín hiệu trên lớp Top luôn có một đường về dòng điện trở kháng thấp ngay phía dưới (return path), một yêu cầu sống còn cho tính toàn vẹn tín hiệu \cite{johnson1993high}. Trở kháng đặc tính của một dải vi mạch (microstrip) trên cấu hình này được ước lượng theo công thức Hammerstad cải biên \cite{hammerstad1975}:
\begin{equation}
    Z_0 = \frac{87}{\sqrt{\varepsilon_r+1{,}41}}\cdot \ln\!\left(\frac{5{,}98h}{0{,}8w + t}\right),
\end{equation}
với $\varepsilon_r \approx 4{,}3$ cho FR-4, $h=1{,}6$ mm (chiều dày substrate), $t=35\,\mu$m. Giải nghịch cho $Z_0=50\,\Omega$ thu được bề rộng đường dẫn $w \approx 2{,}9$ mm -- đây là bề rộng áp dụng cho đoạn dẫn RF từ module LoRa SX1278 đến đầu nối ăng-ten SMA, đảm bảo trở kháng phối hợp và năng lượng RF không bị phản xạ ngược về chip.

\textbf{Bước 3 -- Bố trí linh kiện (Placement).} Trong từng khu chức năng, các linh kiện được đặt theo nguyên tắc \emph{ngắn nhất tốt nhất} cho các đường mang dòng lớn hoặc tín hiệu nhạy. Cụ thể:

\textit{Khối nguồn:} cổng USB-C J2 đặt sát mép bo, cầu chì F1 và ferrite bead FB1 nối tiếp ngay sau, kế đến IC bảo vệ ESD USBLC6-2SC6 (U2) đặt cách chân USB tối đa 5 mm để dập xung ESD trước khi chúng đi sâu vào mạch. IC sạc TP4056 (U3) bố trí trên trục đường cấp pin, với pad tản nhiệt (Exposed Pad) ESOP8 nối xuống một vùng đồng tản nhiệt $20\times15$ mm gồm nhiều via thermal $\Phi 0{,}3$ mm (cho phép dòng sạc 1 A tỏa nhiệt mà không nâng nhiệt độ tiếp giáp $T_J$ vượt giới hạn). LDO AP2112K (U4) đặt sát chân $V_{IN}$ của module ESP32-S3 để rút ngắn đường $3{,}3$ V cấp nguồn.

\textit{Khối vi điều khiển:} module ESP32-S3-WROOM-1 (U1) đặt ở trung tâm, ăng-ten PCB tích hợp trên module hướng ra mép trên của bo (không phục vụ kết nối Wi-Fi/BLE trong dự án này nhưng vẫn cần giữ thông thoáng). Nút BOOT (SW3), nút RESET (SW1), nút USER (SW2) và đầu cắm nạp J1 cụm lại ở một góc dễ thao tác. Đèn LED RGB (D1) đặt nơi nổi bật để dễ quan sát trạng thái.

\textit{Khối cảm biến và đầu nối:} các đầu nối J14--J19, J9, J11, J7/J8 (cho SHT31, BMP180, BH1750, MQ-135, CH4, CO2) đặt dọc theo mép dưới của bo, cho phép cảm biến cắm theo nhiều hướng tùy nhu cầu lắp vào vỏ hộp. Đầu cắm QWIIC J4 đặt mép phải, gần ESP32-S3 để bus I2C mở rộng có chiều dài tối thiểu.

\textit{Khối LoRa và ăng-ten:} đầu cắm module LoRa J3/J10 đặt ở mép phải, sát ăng-ten ngoài (J4 SMA). Đoạn dẫn từ chân ANT của module đến đầu cắm SMA giữ thẳng, không bẻ góc vuông (chỉ dùng các bend 45$^\circ$ hoặc cong tròn), tránh phản xạ điện từ tại các góc nhọn.

\textbf{Bước 4 -- Thiết lập lớp và phủ đồng.} Lớp Top phủ ground pour ở những vùng trống không có tín hiệu, được kết nối xuống ground plane lớp Bottom qua các via $\Phi 0{,}3$ mm rải dày (stitching). Mật độ via stitching được tính theo công thức bước sóng vận hành: ở 433 MHz, bước sóng trong môi trường $\lambda = c/f \approx 692$ mm; bước sóng dẫn (guided wavelength) trên FR-4 với $\varepsilon_{eff}\approx 3{,}0$ là $\lambda_g \approx 400$ mm. Để dải đồng phủ không trở thành ăng-ten ký sinh, khoảng cách giữa hai via stitching phải nhỏ hơn $\lambda_g/20 \approx 20$ mm; bo của đề tài đặt via stitching mỗi 5--8 mm, dư biên an toàn 2--4 lần. Riêng khu vực dưới ăng-ten LoRa (vùng cấm đồng), cả lớp Top và lớp Bottom đều \emph{không phủ đồng và không có via} trong một hình chữ nhật $15 \times 25$ mm; điều này ngăn ăng-ten bị suy hao do phản xạ kim loại lưng và bảo toàn vùng Fresnel cận trường.

\textbf{Bước 5 -- Đi dây tín hiệu (Routing).} Việc đi dây tuân theo bốn quy tắc phân loại theo loại tín hiệu:

\emph{Đường nguồn} (V\_USB, V\_BAT, V\_SYS, 3V3, GND): bề rộng tối thiểu được tính theo bảng IPC-2152 \cite{ipc2152}, công thức nội suy cho lớp ngoài:
\begin{equation}
\label{eq:trace_width}
    I = k\cdot \Delta T^{0{,}44}\cdot A^{0{,}725},
\end{equation}
với $I$ là dòng cho phép (A), $\Delta T$ nhiệt độ tăng so với môi trường ($^\circ$C), $A$ tiết diện đồng (mils$^2$), $k=0{,}048$ cho lớp ngoài. Với mục tiêu $I=1$ A, $\Delta T = 10\,^\circ$C: giải ra $A \approx 21$ mils$^2$, tương ứng bề rộng $w \approx 21/(0{,}035\cdot 1{,}378) \approx 15$ mils $\approx 0{,}38$ mm. Để có biên an toàn 3--4 lần, các đường $V\_SYS$ và đường về GND chính được đi rộng 1{,}5--2 mm.

\emph{Đường tín hiệu số tốc độ cao} (SPI: SCK, MOSI, MISO, NSS; UART nạp): bề rộng 0{,}2--0{,}25 mm, đặt sát ground plane bên dưới (giữ trở kháng đặc tính ổn định khoảng 50--65 $\Omega$). Tránh để hai đường SPI chạy song song quá dài (giữ chiều dài song song dưới 10 mm) để giảm cross-talk; nếu cần chạy song song dài, chèn một đường GND \emph{guard trace} ở giữa. Khoảng cách giữa các đường SPI tuân theo quy tắc \emph{3W}: cạnh đến cạnh $\ge 3$ lần bề rộng đường, đủ để giảm hệ số ghép cảm $k<0{,}05$ ở tốc độ 1--10 MHz.

\emph{Đường I2C}: bus chia sẻ với điện trở kéo lên $4{,}7\,\text{k}\Omega$ (R16, R17). Tổng điện dung bus được giới hạn dưới 400 pF theo NXP I2C Specification \cite{nxp2014i2c}; với khoảng năm cảm biến nối qua dây jumper, ước tính tổng điện dung $C_{bus}\approx 200$ pF (50 pF cho 0{,}5 m dây + linh kiện), thời hằng nạp $\tau=R\cdot C=4{,}7\text{k}\Omega \cdot 200$ pF $=0{,}94\,\mu$s, hoàn toàn thỏa mãn yêu cầu thời gian sườn lên ($\le 1\,\mu$s) ở chế độ Standard 100 kHz và còn biên ở Fast 400 kHz. Bus SDA và SCL được đi cạnh nhau (vì cùng pha logic, cross-talk không gây lỗi), chiều dài giữ dưới 30 mm trên bo để giảm điện dung phụ.

\emph{Đường tương tự} (ngõ ra MQ-135 và CH4 đi vào ADC của ESP32-S3): tách biệt khỏi các đường số chuyển mạch nhanh, đi qua một bộ lọc thấp tần đơn giản RC (1 k$\Omega$ + 100 nF, $f_c \approx 1{,}6$ kHz) đặt sát chân ADC để chống nhiễu coupling từ Wi-Fi/BLE. Nếu phải vượt qua đường số, chuyển sang đi vuông góc (90$^\circ$) chứ không song song, để giảm ghép từ.

\emph{Đường cao tần (RF)}: chỉ một đoạn duy nhất từ chân ANT module LoRa sang đầu cắm SMA. Đoạn này dài $\sim$ 18 mm, bề rộng 2{,}9 mm (đảm bảo $Z_0=50\,\Omega$ như tính ở bước 2), không có via xuyên lớp, đi hoàn toàn trên lớp Top trên một vùng đồng đất đầy đủ ở lớp Bottom (microstrip line).

\textbf{Bước 6 -- Đặt tụ lọc và mạch decoupling.} Mỗi chân nguồn của mọi IC đều được kèm một cặp tụ \emph{bulk + bypass}: tụ tantal hoặc gốm $10\,\mu$F (lưu trữ năng lượng ngắn hạn) đặt cách chân nguồn dưới 5 mm, kết hợp tụ gốm $100$ nF (chống nhiễu tần cao) đặt sát chân nguồn dưới 2 mm. Module ESP32-S3 thêm một tụ $22\,\mu$F do đặc tính dòng đỉnh cao khi Wi-Fi (chưa dùng nhưng được dự phòng cho mở rộng) phát xung. Tụ decoupling đặt sai vị trí (cách chân quá xa) làm vô hiệu hóa toàn bộ tác dụng: tổng cảm kháng đường dẫn $L_{trace}+L_{via}$ ở tần số GHz vượt trở kháng tụ và biến mạch decoupling thành mạch dao động. Đặt sát + via ngắn xuống GND là quy tắc bất di bất dịch.

\textbf{Bước 7 -- Kiểm tra DRC (Design Rule Check) và soi bo ảo.} Sau khi hoàn thiện routing, công cụ Altium thực hiện DRC tự động với bộ quy tắc: khoảng cách giữa hai đường $\ge 0{,}2$ mm, kích cỡ via tối thiểu $\Phi 0{,}3/\Phi 0{,}6$ mm, không đường nào đi cắt qua mặt phẳng đất bị xé (split plane), không vùng đồng cô lập (orphan copper). Sau khi DRC sạch, bo được mô phỏng 3D để kiểm tra giao thoa cơ học với các đầu cắm cảm biến và vỏ hộp.

Kết nối đi dây và layout của Node được thể hiện trong Hình \ref{Hình 4.10}
\begin{figure}[H]
    \centering
    \includegraphics[width=0.7\linewidth]{Thuc_Hanh/Image/Layout.png}
    \caption{PCB Layout của thiết bị Node}
    \label{Hình 4.10}
\end{figure}

Bảng \ref{Bảng 4.1} tổng hợp các linh kiện chính của thiết bị Node theo từng khối chức năng, đóng vai trò như một bảng kê vật tư (BOM) rút gọn.

\begin{table}[H]
\centering
\caption{Tổng hợp linh kiện chính của thiết bị Node theo khối chức năng}
\label{Bảng 4.1}
\renewcommand{\arraystretch}{1.4}
\begin{tabular}{p{5.0cm} p{5.0cm} p{5.6cm}}
\hline
\textbf{Khối chức năng} & \textbf{Linh kiện chính} & \textbf{Vai trò} \\ \hline
Ngõ vào nguồn & USB-C J2, Polyfuse, FB1, USBLC6-2SC6 & Cấp nguồn + bảo vệ quá dòng/ESD/nhiễu \\
Sạc, quản lý pin & TP4056, AO3401A & Sạc, nhận nguồn pin/USB \\
Ổn áp & AP2112K-3.3, tụ 22$\mu$F & Tạo nguồn 3,3V \\
Vi điều khiển & ESP32-S3-WROOM-1 & Xử lý trung tâm, FreeRTOS \\
Cảm biến vi khí hậu & SHT31, BMP180, BH1750 & Đo nhiệt-ẩm, áp suất, ánh sáng \\
Cảm biến khí & MQ-135, CH4, CO2 & Đo $NH_3$, $CH_4$ (ADC), $CO_2$  \\
Thời gian thực & DS3231 & Nhãn thời gian chống Replay \\
Vô tuyến & Module LoRa SX1278 & Truyền thông tầm xa (SPI) \\
Chấp hành, chỉ thị & Act1-3, LED RGB & Điều khiển rơ-le, báo status \\ \hline
\end{tabular}
\end{table}

\lstset{
    basicstyle=\ttfamily\footnotesize,
    keywordstyle=\bfseries,
    breaklines=true,
    frame=single,
    captionpos=b,
    columns=fullflexible,
    showstringspaces=false
}

\subsection{Triển khai phần mềm hệ thống}
Nếu phần cứng là thể xác thì phần mềm chính là linh hồn của thiết bị Điện toán biên. Phần này trình bày chi tiết quá trình hiện thực hóa toàn bộ firmware cho Node, firmware cho Gateway, lớp tích hợp đám mây và công cụ kiểm thử. Do triết lý "tập trung trí tuệ tại Node", phần lớn nội dung tập trung vào firmware Node -- nơi diễn ra việc thu thập, lọc, ra quyết định, mã hóa, truyền tin và quản lý năng lượng.

\subsubsection{Môi trường phát triển và tổ chức mã nguồn}
Toàn bộ firmware được phát triển trên nền tảng ESP-IDF phiên bản 6.0 -- bộ công cụ phát triển chính thức của Espressif dành cho dòng chip ESP32, tích hợp sẵn nhân FreeRTOS, các trình điều khiển ngoại vi, ngăn xếp mạng và thư viện mật mã. So với việc lập trình trên nền Arduino, ESP-IDF cho phép kiểm soát ở mức thấp các tài nguyên hệ thống.
\begin{table}[H]
\centering
\caption{Tổ chức các module mã nguồn chính của firmware Node}
\label{Bảng 4.2}
\renewcommand{\arraystretch}{1.4}
\begin{tabular}{p{4.0cm} p{11.0cm}}
\hline
\textbf{Module} & \textbf{Vai trò} \\ \hline
\texttt{main.c} & Khởi tạo hệ thống, tạo các tác vụ FreeRTOS, vòng đời chu kỳ \\ 
\texttt{sht3x, bmp180, bh1750} & Driver các cảm biến vi khí hậu trên bus I2C \\ 
\texttt{filters.c/.h} & Thư viện bộ lọc số (EMA, MA, Median, Kalman) \\ 
\texttt{packet.c/.h} & Định nghĩa cấu trúc gói tin, đóng/mở gói \\
\texttt{crypto.c/.h} & Lớp mã hóa AEAD trừu tượng (AES-128-GCM / ASCON) \\ 
\texttt{ascon.c/.h} & Hiện thực thuật toán mật mã hạng nhẹ ASCON-128 \\
\texttt{sx127x.c/.h} & Driver điều khiển module LoRa SX1278 qua SPI \\ 
\texttt{rtos\_stats.c/.h} & Thu thập thống kê RTOS: tải CPU mỗi lõi, High-Water Mark Stack từng tác vụ \\ 
\texttt{mem\_monitor.c/.h} & Giám sát Heap động qua \texttt{heap\_caps}, phát hiện rò rỉ và phân mảnh \\ 
\texttt{power\_mgr.c/.h} & Cấu hình \texttt{esp\_pm} (DFS), Tickless Idle và Light/Deep Sleep \\ 
\texttt{board\_pins.h} & Khai báo tập trung bản đồ chân (pin-map) \\ \hline
\end{tabular}
\end{table}
Hệ thống phần mềm được chia thành ba dự án tách biệt: firmware Node (thư mục \texttt{NodeSensor}), firmware Gateway (thư mục \texttt{LoraGateway}) và công cụ kiểm thử trên máy tính (thư mục \texttt{HIL\_Tool}). Bảng \ref{Bảng 4.2} liệt kê các module mã nguồn chính của firmware Node cùng vai trò của chúng, cho thấy sự phân tách trách nhiệm rõ ràng trong kiến trúc phần mềm.



\subsubsection{Triển khai lớp driver cảm biến}
Lớp driver là cầu nối giữa phần mềm và phần cứng cảm biến, và là nơi đầu tiên dữ liệu thô được sinh ra -- do đó chất lượng của nó quyết định độ tin cậy của mọi xử lý phía sau. Driver cảm biến cũng phải xử lý các đặc thù riêng của từng linh kiện. Cơ chế thu dữ liệu cảm biến của Node được thể hiện trong Hình \ref{Hình 4.11}.
\begin{figure}[H]
    \centering
    \includegraphics[width=1.0\linewidth]{Thuc_Hanh/Image/DriverSensor.pdf}
    \caption{Cơ chế đọc dữ liệu của node}
    \label{Hình 4.11}
\end{figure}
Hình \ref{Hình 4.11} mô tả quy trình xử lý của bộ điều khiển trung tâm theo mô hình Finite State Machine (FSM) nhằm xác định trạng thái hoạt động của hệ thống giám sát môi trường và điều khiển các thiết bị chấp hành. Quy trình được chia thành ba phần chính gồm Input, Processing và Output, giúp thể hiện rõ luồng dữ liệu từ khi nhận tín hiệu đầu vào đến khi đưa ra quyết định điều khiển.

Ở giai đoạn Input, hệ thống lần lượt nhận các thông số môi trường bao gồm nhiệt độ (Temperature), độ ẩm (Humidity), nồng độ khí NH₃, cùng với trạng thái hiện tại (Current State) của hệ thống. Các dữ liệu này được truyền sang khối xử lý để đánh giá điều kiện môi trường. Trước tiên, bộ điều khiển tiến hành tính toán chỉ số THI (Temperature Humidity Index) từ giá trị nhiệt độ và độ ẩm nhằm đánh giá mức độ ảnh hưởng của điều kiện nhiệt đến môi trường chăn nuôi. Sau khi có giá trị THI, hệ thống đồng thời thực hiện ba bước kiểm tra gồm Check Temperature, Check NH3 và Check THI. Mỗi phép kiểm tra sẽ so sánh giá trị đo được với các ngưỡng đã được thiết lập trước để xác định mức độ an toàn của môi trường.

Kết quả từ các phép kiểm tra được tổng hợp để đánh giá xem hệ thống có rơi vào trạng thái Emergency hay không. Nếu bất kỳ thông số nào vượt quá ngưỡng nguy hiểm, hệ thống sẽ chuyển sang State = EMERGENCY, đồng thời kích hoạt toàn bộ thiết bị làm mát và xử lý môi trường, cụ thể là quạt và hệ thống phun sương (Fan + Mist ON) nhằm nhanh chóng đưa điều kiện môi trường trở về mức an toàn. Trong trường hợp chưa đạt mức khẩn cấp, hệ thống tiếp tục đánh giá điều kiện Warning. Nếu các thông số chỉ vượt ngưỡng cảnh báo nhưng chưa đến mức nguy hiểm, FSM chuyển sang State = WARNING, khi đó chỉ kích hoạt quạt ở tốc độ thấp (Low-Speed Fan ON) nhằm giảm tiêu thụ năng lượng nhưng vẫn cải thiện điều kiện môi trường.


Sau khi xác định trạng thái cuối cùng, tất cả các nhánh xử lý được hợp nhất và chuyển sang phần Output. Tại đây, hệ thống cập nhật Updated State để lưu trạng thái hoạt động mới, đồng thời sinh ra Relay Commands nhằm điều khiển các relay đóng hoặc ngắt tương ứng với quạt, bơm phun sương và các thiết bị khác. Cấu trúc FSM này giúp việc điều khiển trở nên rõ ràng, dễ mở rộng và đảm bảo hệ thống phản ứng nhanh với các biến động của môi trường trong khi vẫn duy trì tính ổn định nhờ cơ chế hysteresis.

Đoạn mã \ref{lst:probe1} minh họa cơ chế dò và đọc cảm biến với xử lý lỗi mềm. 
\begin{lstlisting}[language=C, caption={Cơ chế dò và đọc cảm biến với xử lý lỗi mềm}, label={lst:probe1}]
typedef struct {
    bool   present;        // ket qua auto-probe
    float  value;          // gia tri da chuyen doi
    bool   valid;          // CRC/ham doc thanh cong
} sensor_reading_t;

void sensors_read_all(node_data_t *d) {
    if (sht31.present)
        d->th_ok  = sht3x_read(&d->temp, &d->hum);   // co kiem CRC
    if (bmp180.present)
        d->p_ok   = bmp180_read(&d->press);          // bu nhiet
    if (bh1750.present)
        d->lux_ok = bh1750_read(&d->lux);
    d->nh3 = adc_to_ppm(adc_read(MQ135_CH));         // cam bien khi
    // ... cam bien khong san sang -> danh dau invalid, KHONG treo
}
\end{lstlisting}
Với SHT31, mỗi cặp byte dữ liệu trả về kèm một byte CRC, được driver kiểm tra ngay để loại các khung lỗi do nhiễu bus. Với BMP180, 11 hệ số hiệu chỉnh được đọc \emph{một lần} lúc khởi tạo bằng một giao dịch đọc liên tục (burst read) 22 byte -- một bài học rút ra trong quá trình phát triển, bởi việc đọc rời từng hệ số dễ gây lệch dữ liệu và cho ra giá trị nhiệt độ sai lệch hàng trăm độ. Với các cảm biến khí tương tự (MQ-135, CH4), driver đọc giá trị qua bộ ADC của ESP32-S3 và quy đổi sang nồng độ theo đường đặc tuyến; còn cảm biến CO2 được đọc qua giao tiếp UART. 

% \begin{lstlisting}[language=C, caption={Cơ chế dò và đọc cảm biến với xử lý lỗi mềm}, label={lst:probe}]
% typedef struct {
%     bool   present;        // ket qua auto-probe
%     float  value;          // gia tri da chuyen doi
%     bool   valid;          // CRC/ham doc thanh cong
% } sensor_reading_t;

% void sensors_read_all(node_data_t *d) {
%     if (sht31.present)
%         d->th_ok  = sht3x_read(&d->temp, &d->hum);   // co kiem CRC
%     if (bmp180.present)
%         d->p_ok   = bmp180_read(&d->press);          // bu nhiet
%     if (bh1750.present)
%         d->lux_ok = bh1750_read(&d->lux);
%     d->nh3 = adc_to_ppm(adc_read(MQ135_CH));         // cam bien khi
%     // ... cam bien khong san sang -> danh dau invalid, KHONG treo
% }
% \end{lstlisting}

\subsubsection{Triển khai thư viện bộ lọc theo từng tín hiệu}
Thư viện \texttt{filters.c} hiện thực bốn loại bộ lọc (EMA, trung bình trượt, trung vị, Kalman 1D) sau một giao diện thống nhất gồm ba hàm: \texttt{filter\_init} khởi tạo trạng thái, \texttt{filter\_apply} đưa một mẫu mới vào và trả về giá trị đã lọc, \texttt{filter\_name} trả về tên bộ lọc phục vụ ghi log. Cấu hình mặc định gán EMA cho nhiệt độ và độ ẩm, trung bình trượt cho áp suất, và trung vị cho ánh sáng.

Một chi tiết triển khai quan trọng là trạng thái của mỗi bộ lọc được lưu trong vùng nhớ RTC (thuộc tính \texttt{RTC\_DATA\_ATTR}) để không bị xóa khi vi điều khiển vào chế độ ngủ sâu và reset. Nhờ vậy, bộ lọc duy trì được tính liên tục theo thời gian thực qua các chu kỳ ngủ.

\subsubsection{Triển khai máy trạng thái điều khiển biên và khối chấp hành}
Cơ chế hoạt động FSM được thể hiện trong Hình \ref{Hình 4.12}
\begin{figure}[H]
    \centering
    \includegraphics[width=0.7\linewidth]{Thuc_Hanh/Image/FSM_Process.pdf}
    \caption{Cơ chế hoạt động của FSM trên Node}
    \label{Hình 4.12}
\end{figure}
% Tại mỗi chu kỳ, sau khi có dữ liệu đã lọc, FSM tính chỉ số nhiệt ẩm THI rồi đối chiếu đồng thời nhiệt độ, độ ẩm, nồng độ $NH_3$ và THI với ma trận ngưỡng đã thiết lập (Bảng \ref{Bảng 3.6}). Kết quả là việc chuyển trạng thái giữa ba mức AN TOÀN -- CẢNH BÁO -- KHẨN CẤP, kèm theo việc điều khiển các kênh chấp hành Act1--Act3.
Sơ đồ trên mô tả quy trình thu thập dữ liệu từ các cảm biến, gồm ba phần Input, Processing và Output. Ban đầu, hệ thống thực hiện tự động dò tìm (Auto-probe) để xác định các cảm biến đang kết nối, bao gồm SHT31, BMP180, BH1750 và MQ135. Việc dò tìm này giúp chương trình có thể hoạt động ngay cả khi một số cảm biến không được lắp đặt hoặc gặp lỗi.
.
Sau khi hoàn tất việc đọc dữ liệu, toàn bộ thông tin được tổng hợp vào cấu trúc node\_ddata\_t, bao gồm các giá trị nhiệt độ, độ ẩm, áp suất, cường độ ánh sáng và nồng độ NH₃, đồng thời lưu các cờ hợp lệ (Validity flags) cho từng cảm biến. Dữ liệu này sau đó được chuyển tới các khối xử lý phía trên như thuật toán tính THI, FSM và hệ thống truyền thông, bảo đảm tính tin cậy và khả năng mở rộng của toàn bộ hệ thống.

Để tránh hiện tượng rơ-le đóng/cắt liên tục khi giá trị đo dao động sát ngưỡng (chattering), FSM cài đặt cơ chế trễ (hysteresis): ngưỡng để chuyển lên trạng thái nguy hiểm hơn cao hơn ngưỡng để hạ xuống trạng thái an toàn hơn một khoảng nhất định. Điều này bảo vệ tuổi thọ cơ học của rơ-le và tránh gây sốc cho thiết bị cơ điện. Toàn bộ chuỗi từ lúc đọc cảm biến vượt ngưỡng đến khi chân GPIO đóng mạch rơ-le được đo và bảo đảm hoàn tất dưới 50 ms -- chỉ tiêu tới hạn của hệ thống. Đoạn mã \ref{lst:fsm} minh họa lõi logic của FSM.

\begin{lstlisting}[language=C, caption={Lõi logic máy trạng thái điều khiển biên}, label={lst:fsm}]
edge_state_t fsm_update(node_data_t *d, edge_state_t cur) {
    float thi = 0.8f*d->temp + (d->hum/100.0f)*(d->temp - 14.4f) + 46.4f;
    if (d->temp > 35.0f || d->nh3 > 25.0f || thi > 78.0f) {
        actuator_set(ACT_ALL, ON);     // bat toan bo quat + phun suong
        return STATE_EMERGENCY;
    } else if (d->temp > 32.0f || d->nh3 > 20.0f || thi > 72.0f) {
        actuator_set(ACT_FAN_LOW, ON); // canh bao: quat muc thap
        return STATE_WARNING;
    } else if (cur != STATE_SAFE && below_hysteresis(d)) {
        actuator_set(ACT_ALL, OFF);
        return STATE_SAFE;
    }
    return cur;                         // giu trang thai (chong chattering)
}
\end{lstlisting}

\subsubsection{Triển khai đóng gói và mã hóa dữ liệu (AEAD)}
Trước khi truyền qua môi trường vô tuyến mở, dữ liệu phải được đóng gói và mã hóa, cấu trúc gói tin được trình bày trong Hình \ref{Hình 4.13}. Firmware định nghĩa cấu trúc gói tin gồm hai phần: phần tiêu đề (header) ở dạng bản rõ chứa các trường định danh (magic nhận dạng giao thức, phiên bản, mã thuật toán mã hóa, mã định danh Node, độ dài tải trọng và một số ngẫu nhiên nonce), và phần tải trọng (payload) chứa dữ liệu cảm biến cùng nhãn thời gian. Header để ở dạng bản rõ vì Gateway cần đọc được mã định danh và thuật toán \emph{trước khi} giải mã; tuy nhiên header vẫn được đưa vào quá trình xác thực dưới dạng dữ liệu liên kết (Associated Data), nên mọi sửa đổi header đều bị phát hiện.

Điểm nổi bật của đề tài là hiện thực \emph{hai} thuật toán mã hóa có xác thực (AEAD) có thể lựa chọn: AES-128 ở chế độ GCM và ASCON-128. Chế độ GCM mở rộng thuật toán AES-128 (đã phân tích cấu trúc toán học ở Chương 2) bằng cách bổ sung một thẻ xác thực (authentication tag) 16 byte, qua đó vừa bảo mật nội dung vừa bảo đảm tính toàn vẹn -- chống lại cả nghe lén lẫn giả mạo \cite{mcgrew2004galois}. Trên ESP32-S3, AES-128-GCM được hiện thực thông qua giao diện mật mã PSA Crypto của ESP-IDF v6, tận dụng bộ tăng tốc mã hóa phần cứng nên thời gian thực thi rất nhỏ. ASCON-128 -- thuật toán thắng cuộc thi mật mã hạng nhẹ của Viện NIST \cite{dobraunig2021ascon, turan2023lightweight} -- được hiện thực bằng phần mềm thuần, tối ưu cho thiết bị tài nguyên hạn chế với trạng thái nhỏ và không cần bảng tra lớn. Việc cung cấp cả hai cho phép đối chứng trực tiếp giữa một thuật toán tiêu chuẩn được tăng tốc phần cứng và một thuật toán hạng nhẹ chạy phần mềm (sẽ so sánh định lượng ở Chương 5). Cả hai được gói sau một giao diện trừu tượng chung, định tuyến theo trường mã thuật toán trong header, nên phần nghiệp vụ chỉ cần chọn thuật toán qua một hằng số cấu hình.

% ===================================================================
% PlantUML cho Hình 4.8 - render và luu thanh: Thuc_Hanh/Image/Ch4_Packet.png
% @startuml
% skinparam backgroundColor white
% rectangle "HEADER (ban ro = AAD)\nmagic | ver | algo | node_id | plen | nonce[16]" as H
% rectangle "PAYLOAD (ma hoa)\nT,H,P,Lux,NH3,CH4,CO2 | THI | trang thai | timestamp" as P
% rectangle "TAG[16]\n(xac thuc AEAD)" as T
% H -right-> P
% P -right-> T
% note bottom of H : Header KHONG ma hoa nhung\nDUOC bao ve toan ven qua AAD
% @enduml
% ===================================================================
\begin{figure}[H]
    \centering
    {\includegraphics[width=0.7\linewidth]{Thuc_Hanh/Image/AHEAD_Package.pdf}}
    \caption{Cấu trúc gói tin AEAD: header (AAD) + payload mã hóa + thẻ xác thực}
    \label{Hình 4.13}
\end{figure}

Tích hợp cùng cơ chế mã hóa là giao thức chống tấn công phát lại dựa trên nhãn thời gian. Mỗi gói tin mang một nhãn thời gian lấy từ RTC DS3231; phía Gateway sau khi giải mã sẽ so sánh nhãn thời gian này với đồng hồ của nó, và chỉ chấp nhận gói nếu độ lệch nằm trong cửa sổ an toàn $\Delta t < 5\,\text{s}$. Nhờ thẻ xác thực AEAD, kẻ tấn công không thể sửa nhãn thời gian mà không làm hỏng thẻ, nên mọi gói phát lại đều mang nhãn thời gian cũ và bị loại bỏ. Sự kết hợp giữa mã hóa AEAD và nhãn thời gian tạo nên một lớp bảo vệ hoàn chỉnh về cả ba mặt: bảo mật, toàn vẹn và tính mới của dữ liệu.

\subsubsection{Triển khai truyền thông LoRa và cơ chế xác nhận (ACK)}
Driver LoRa (\texttt{sx127x.c}) điều khiển module SX1278 qua SPI để cấu hình các thanh ghi (tần số, hệ số trải phổ, băng thông, công suất phát), ghi gói vào bộ đệm FIFO và phát đi, cũng như chuyển sang chế độ thu để nhận. Một bài học triển khai quan trọng là vòng chờ phát/thu không được dùng vòng lặp bận chiếm CPU; thay vào đó firmware dựa vào ngắt phần cứng DIO0 (báo sự kiện TxDone/RxDone) kết hợp với bộ định thời \texttt{esp\_timer} để áp đặt thời hạn (timeout) theo thời gian tuyệt đối. Cách này vừa giải phóng CPU cho các tác vụ khác trong lúc chờ vô tuyến, vừa tránh được lỗi treo CPU do cấu hình nhịp hệ thống không phù hợp.

Sau khi phát gói, Node không coi như đã hoàn tất mà chờ một gói xác nhận (ACK) từ Gateway, biến liên kết một chiều thành một cái bắt tay hai chiều. Gói ACK là một khung ngắn có mã kiểm tra CRC8, mang về trạng thái xử lý của Gateway, các chỉ số RSSI/SNR mà Gateway đo được trên đường lên, và mốc thời gian hiện tại của Gateway. Hàm \texttt{send\_with\_ack} quản lý toàn bộ chu trình: phát gói, chuyển sang thu, chờ ACK trong một khoảng thời gian định trước; nếu hết thời gian mà không có ACK hợp lệ, Node sẽ thử lại một số lần giới hạn, và nếu vẫn thất bại thì kết luận mất kết nối để kích hoạt cơ chế lưu trữ ngoại tuyến. Toàn bộ chuỗi gửi -- chờ -- thử lại được thiết kế để gói gọn trong vài giây nhằm tôn trọng ngân sách thời gian của một chu kỳ.

Từ cơ chế ACK, firmware khai thác thêm ba thông tin giá trị. Thứ nhất là độ trễ khứ hồi (Round-Trip Time -- RTT): Node ghi mốc thời gian bằng \texttt{esp\_timer} ngay trước khi phát và ngay khi nhận ACK, hiệu hai mốc là RTT -- một chỉ số phản ánh sức khỏe đường truyền. Thứ hai là ước lượng khoảng cách tới Gateway từ giá trị RSSI của gói ACK, dựa trên mô hình suy hao đường truyền log-distance:
\begin{equation}
    {d = 10^{\frac{RSSI_{1m} - RSSI}{10\,n}}}.
\end{equation}
Trong đó $RSSI_{1m}$ là cường độ tham chiếu đo ở cự ly 1 m và $n$ là hệ số suy hao môi trường. Thứ ba là khả năng đồng bộ thời gian: mốc thời gian trong ACK cho phép Node hiệu chỉnh đồng hồ, bổ trợ cho RTC phần cứng. Bảng \ref{Bảng 4.3} tổng hợp các tham số cấu hình lớp vật lý LoRa được sử dụng, đồng nhất giữa Node và Gateway.

\begin{table}[H]
\centering
\caption{Tham số cấu hình lớp vật lý LoRa của hệ thống}
\label{Bảng 4.3}
\renewcommand{\arraystretch}{1.4}
\begin{tabular}{p{4.2cm} p{3.2cm} p{5.4cm}}
\hline
\textbf{Tham số} & \textbf{Giá trị} & \textbf{Ý nghĩa} \\ \hline
Tần số sóng mang & 433 MHz & Băng ISM, xuyên vật cản tốt \\ 
Spreading Factor (SF) & 7 -- 9 & Cân bằng tầm xa và thời gian phát \\ 
Băng thông (BW) & 125 kHz & Tiêu chuẩn cho ứng dụng tầm xa \\ 
Coding Rate (CR) & 4/5 & Sửa lỗi tiến, ít tốn airtime \\ 
Công suất phát & 17 dBm & Tầm xa trong giới hạn cho phép \\ 
Sync word & 0x12 & Mạng riêng (khác LoRaWAN) \\ \hline
\end{tabular}
\end{table}

\subsubsection{Triển khai quản lý năng lượng và cơ chế ngủ lai (Hybrid Sleep)}
Quản lý năng lượng là yếu tố quyết định tuổi thọ pin của Node. Firmware hiện thực một mô hình Hybrid Sleep kết hợp hai chế độ ngủ của ESP32-S3 tùy theo ngữ cảnh hoạt động như trong Hình \ref{Hình 4.14}
\begin{figure}[H]
    \centering
    {\includegraphics[width=0.7\linewidth]{Thuc_Hanh/Image/SleepFSM.pdf}}
    \caption{Máy trạng thái quản lý năng lượng theo mô hình Hydrid Sleep}
    \label{Hình 4.14}
\end{figure}
Ở điều kiện bình thường khi đã gửi dữ liệu thành công và không còn dữ liệu tồn đọng, Node vào chế độ ngủ sâu (Deep Sleep) -- cắt nguồn lõi CPU và RAM, chỉ giữ lại bộ định thời RTC, hạ dòng tiêu thụ của vi điều khiển xuống mức vài $\mu A$. Ngược lại, khi đang trong trạng thái buffer dữ liệu ngoại tuyến hoặc cần phản ứng nhanh, Node dùng chế độ ngủ nhẹ (Light Sleep) -- giữ nguyên RAM và trạng thái RTOS, cho phép thức dậy nhanh mà không mất dữ liệu trong bộ đệm.

Thách thức lớn nhất của Deep Sleep là chip bị reset hoàn toàn mỗi khi thức dậy, làm mất sạch biến trong RAM thường. Firmware giải quyết bằng cách khai báo các biến trạng thái cốt lõi với thuộc tính \texttt{RTC\_DATA\_ATTR} để chúng nằm trong vùng RTC RAM không bị xóa qua Deep Sleep: bộ đếm số lần khởi động, số lần watchdog can thiệp, tổng số gói gửi và số ACK nhận được, chuỗi mất kết nối liên tiếp, số thứ tự gói, mốc thời gian cơ sở và trạng thái các bộ lọc. Nhờ vậy, dù chip reset mỗi chu kỳ, Node vẫn duy trì được tính liên tục về mặt logic nghiệp vụ. Việc thức dậy được kích hoạt theo hai nguồn song song -- bộ định thời (lấy mẫu định kỳ) hoặc sự kiện ngắt ngoài qua chân DIO0/SQW (phản ứng khẩn cấp), đúng tinh thần "ngủ định kỳ kết hợp thức theo sự kiện" đã nêu trong phạm vi đề tài ở Chương 1.

% ===================================================================
% PlantUML cho Hình 4.9 - render va luu thanh: Thuc_Hanh/Image/Ch4_SleepFSM.png
% @startuml
% skinparam backgroundColor white
% [*] --> ACTIVE
% state ACTIVE : Doc-Loc-FSM-Ma hoa-Gui
% state LIGHT_SLEEP : Giu RAM+RTOS\n(khi co buffer offline)
% state DEEP_SLEEP : Cat CPU/RAM\n(khi online, buffer rong)
% ACTIVE --> DEEP_SLEEP : gui thanh cong & buffer rong
% ACTIVE --> LIGHT_SLEEP : dang buffer / cho dong bo
% DEEP_SLEEP --> ACTIVE : timer RTC / ngat DIO0
% LIGHT_SLEEP --> ACTIVE : timer / ngat
% @enduml
% ===================================================================


Cần lưu ý một điểm trung thực về mặt kỹ thuật: do thiết kế chấp nhận cấp nguồn liên tục cho mảng cảm biến khí MOX (để giữ cân bằng nhiệt), dòng tiêu thụ tổng của Node ở trạng thái chờ bị chi phối bởi cuộn nung của cảm biến khí chứ không phải bởi vi điều khiển. Vì vậy, chỉ tiêu dòng ngủ sâu mức vài $\mu A$ áp dụng cho riêng phần vi điều khiển; phần này sẽ được làm rõ và định lượng trong chương kiểm thử như một điểm đánh đổi đã được cân nhắc.

\subsubsection{Triển khai lưu trữ ngoại tuyến và cơ chế store-and-forward}
Cơ chế lưu trữ và phục hồi dữ liệu khi ngoại tuyến được thể hiện trong Hình \ref{Hình 4.15}
\begin{figure}[H]
    \centering
    {\includegraphics[width=0.55\linewidth]{Thuc_Hanh/Image/t.pdf}}
    \caption{Logic chức năng lưu trữ ngoại tuyến}
    \label{Hình 4.15}
\end{figure}

Để hiện thực tính kiên cường trước sự cố mất kết nối, firmware xây dựng cơ chế lưu trữ ngoại tuyến (offline storage) theo mô hình lưu-và-chuyển-tiếp (store-and-forward). Khi hàm \texttt{send\_with\_ack} kết luận mất kết nối (không nhận được ACK sau số lần thử quy định), thay vì vứt bỏ dữ liệu, Node sẽ ghi bản ghi hiện tại -- kèm nhãn thời gian thực từ RTC -- vào vùng lưu trữ cục bộ. Hai cơ chế lưu trữ được phối hợp: một bộ đệm vòng (ring buffer) trong RTC RAM cho các bản ghi gần nhất (sống sót qua Deep Sleep), và hệ thống tập tin SPIFFS trên bộ nhớ Flash cho lưu trữ dung lượng lớn và bền vững. Bộ đệm vòng hoạt động theo nguyên tắc FIFO và tự loại bỏ bản ghi cũ nhất khi đầy (drop-oldest), bảo đảm luôn giữ được dữ liệu mới nhất ngay cả khi mất kết nối kéo dài.

Ngay khi liên kết vô tuyến được khôi phục (nhận lại được ACK), một tác vụ đồng bộ chạy nền ở mức ưu tiên thấp sẽ lần lượt đọc các bản ghi tồn đọng và gửi gối đầu (backlog) về Gateway, khôi phục tính toàn vẹn của chuỗi dữ liệu lịch sử. Việc đính kèm nhãn thời gian thực vào mỗi bản ghi là điều kiện tiên quyết để dữ liệu ngoại tuyến vẫn có giá trị khi được gửi lại sau này -- người dùng có thể tái dựng chính xác diễn biến vi khí hậu trong khoảng thời gian mất mạng.


\subsubsection{Triển khai cơ chế giám sát Watchdog và tự phục hồi}
Vì hệ thống phải hoạt động liên tục 24/7 ngoài thực địa không người trực, khả năng tự phục hồi khi gặp sự cố phần mềm là bắt buộc. Firmware kích hoạt bộ định thời giám sát tác vụ (Task Watchdog Timer) của ESP-IDF cho tác vụ chính: trong điều kiện bình thường, tác vụ phải định kỳ "vỗ về" (feed) watchdog; nếu tác vụ bị treo (ví dụ do bus I2C bị khóa hoặc vòng lặp vô hạn) và không feed kịp trong thời hạn cấu hình, watchdog sẽ kích hoạt reset toàn hệ thống, đưa thiết bị trở lại trạng thái hoạt động sạch. Sau mỗi lần reset, firmware đọc nguyên nhân reset (\texttt{esp\_reset\_reason}) và tăng bộ đếm số lần watchdog can thiệp lưu trong RTC RAM, đồng thời báo cáo trong telemetry -- cung cấp bằng chứng định lượng về độ ổn định của hệ thống.

Để có thể kiểm chứng cơ chế này một cách an toàn và lặp lại được, phần cứng dành riêng một chân (GPIO0, nối nút BOOT) làm chân kích hoạt kiểm thử: khi người kiểm thử nối chân này xuống đất, firmware chủ động đi vào một vòng lặp treo giả lập, qua đó buộc watchdog phải can thiệp -- cho phép quan sát và đo đạc thời gian phục hồi.

% --- Bổ sung phần RTOS ---
\subsubsection{Triển khai bảo vệ bus I2C và giám sát tài nguyên RTOS}
Hiện thực cụ thể cho cơ chế Mutex bảo vệ bus I2C đã thiết kế ở Chương 3 được trình bày trong Đoạn mã \ref{lst:i2c_mutex}. Một Mutex duy nhất được tạo lúc khởi tạo; mọi giao dịch I2C đều được bao bọc giữa cặp \texttt{xSemaphoreTake}/\texttt{xSemaphoreGive} với thời hạn chờ hữu hạn để không bao giờ treo cứng tác vụ.

\begin{lstlisting}[language=C, caption={Bảo vệ giao dịch I2C bằng Mutex có kế thừa ưu tiên}, label={lst:i2c_mutex}]
static SemaphoreHandle_t i2c_mutex;        // bao ve bus I2C dung chung

void i2c_guard_init(void) {
    // Mutex (khac Binary Semaphore): co ownership + ke thua uu tien
    i2c_mutex = xSemaphoreCreateMutex();
}

// Doc cam bien I2C an toan: ca giao dich la mot vung toi han nguyen tu
esp_err_t sht3x_read_safe(float *t, float *h) {
    // Cho Mutex toi da 100ms; neu khong duoc -> bo qua chu ky, KHONG treo
    if (xSemaphoreTake(i2c_mutex, pdMS_TO_TICKS(100)) != pdTRUE)
        return ESP_ERR_TIMEOUT;            // bus ket/treo -> danh dau invalid
    esp_err_t r = sht3x_read(t, h);        // START..addr..data..STOP
    xSemaphoreGive(i2c_mutex);             // tra Mutex ngay sau STOP
    return r;
}
\end{lstlisting}

Để biến các mục tiêu tối ưu RTOS ở Chương 1 thành số liệu kiểm chứng được, firmware tích hợp một tác vụ chẩn đoán chạy nền (module \texttt{rtos\_stats} và \texttt{mem\_monitor}) định kỳ in ra hai nhóm chỉ số then chốt: mức dùng đỉnh Stack của từng tác vụ và dung lượng Heap còn trống. Hàm \texttt{uxTaskGetStackHighWaterMark()} trả về lượng Stack còn trống \emph{ít nhất} mà mỗi tác vụ từng chạm tới -- nếu giá trị này tiến gần 0 nghĩa là nguy cơ tràn Stack; ngược lại nếu quá lớn thì đang lãng phí RAM và có thể giảm cấp phát. Với Heap, hiệu giữa \texttt{heap\_caps\_get\_free\_size()} và \texttt{heap\_caps\_get\_largest\_free\_block()} cho biết mức độ phân mảnh. Đoạn mã \ref{lst:rtos_diag} minh họa cách thu thập các chỉ số này và phát qua \texttt{ESP\_LOGI} (đồng thời đính kèm vào telemetry JSON gửi cho công cụ EdgeProfiler).

\begin{lstlisting}[language=C, caption={Giám sát Stack High-Water Mark và Heap động}, label={lst:rtos_diag}]
// Theo doi Stack tung tac vu (don vi: word; 1 word = 4 byte tren ESP32-S3)
void rtos_log_stack(TaskHandle_t h, const char *name) {
    UBaseType_t hwm = uxTaskGetStackHighWaterMark(h);
    ESP_LOGI("RTOS", "%s: stack free min = %u words", name, (unsigned)hwm);
}

// Giam sat Heap noi: tong trong + khoi lien ke lon nhat -> do phan manh
void mem_log_heap(void) {
    size_t freeb   = heap_caps_get_free_size(MALLOC_CAP_INTERNAL);
    size_t largest = heap_caps_get_largest_free_block(MALLOC_CAP_INTERNAL);
    float  frag    = 100.0f * (1.0f - (float)largest / (float)freeb);
    ESP_LOGI("MEM", "heap free=%u B, largest=%u B, frag=%.1f%%",
             (unsigned)freeb, (unsigned)largest, frag);
}
\end{lstlisting}

Cuối cùng, để hiện thực mục tiêu tiết kiệm năng lượng ở mức hệ điều hành, module \texttt{power\_mgr} cấu hình bộ quản lý năng lượng \texttt{esp\_pm} \cite{espressif_pm} kết hợp cơ chế Tickless Idle của FreeRTOS (Đoạn mã \ref{lst:power_pm}). Khi bật \texttt{light\_sleep\_enable}, mỗi khi tất cả các tác vụ đều ở trạng thái Bị chặn, nhân FreeRTOS sẽ \emph{bỏ qua} các ngắt Tick định kỳ và tự động đưa chip vào Light Sleep cho đến sự kiện thức kế tiếp, thay vì thức dậy 1000 lần mỗi giây một cách vô ích; đồng thời cơ chế Điều chỉnh tần số động (DFS) hạ xung nhịp CPU khi tải thấp. Việc bật cơ chế này yêu cầu khai báo \texttt{CONFIG\_PM\_ENABLE=y} và \texttt{CONFIG\_FREERTOS\_USE\_TICKLESS\_IDLE=y} trong \texttt{sdkconfig}.

\begin{lstlisting}[language=C, caption={Cấu hình esp\_pm với DFS và Tickless Idle (Automatic Light Sleep)}, label={lst:power_pm}]
#include "esp_pm.h"
void power_mgr_init(void) {
    esp_pm_config_t pm = {
        .max_freq_mhz = 240,           // xung dinh khi Active (hieu nang)
        .min_freq_mhz = 40,            // ha xung khi tai thap (DFS)
        .light_sleep_enable = true     // Tickless Idle: tu dong Light Sleep
    };
    ESP_ERROR_CHECK(esp_pm_configure(&pm));
}
\end{lstlisting}
% --- Hết bổ sung phần RTOS ---

\subsubsection{Triển khai công cụ kiểm thử Hardware-In-the-Loop (EdgeProfiler)}
Một đóng góp công cụ quan trọng của đề tài là phần mềm kiểm thử Hardware-In-the-Loop mang tên EdgeProfiler, được xây dựng bằng ngôn ngữ Python với thư viện giao diện PyQt5 và thư viện vẽ đồ thị thời gian thực pyqtgraph. Công cụ này hiện thực hóa phương pháp đo lường định lượng đã nêu ở Chương 1, cho phép quan trắc và kiểm thử hệ thống một cách trực quan và tự động. Điểm thiết kế cốt lõi của công cụ là khả năng giám sát Node một cách \emph{độc lập}: vì Node tự phát ra toàn bộ trạng thái của mình dưới dạng telemetry JSON qua cổng USB, người dùng chỉ cần cắm Node vào máy tính là theo dõi được mọi hoạt động, không phụ thuộc vào việc Gateway có hiện diện hay không.

Công cụ EdgeProfiler cung cấp các nhóm chức năng chính: hiển thị đồ thị thời gian thực các đại lượng cảm biến (cả giá trị thô lẫn giá trị đã lọc để so sánh hiệu quả bộ lọc), các chỉ số chất lượng đường truyền (RSSI, SNR, RTT, tỉ lệ giao gói thành công PDR, khoảng cách ước lượng), trình giám sát tài nguyên và profiler FreeRTOS (tải CPU theo từng lõi, dung lượng heap, trạng thái từng tác vụ), bộ chạy kịch bản kiểm thử tự động (test sequencer), chức năng ghi và phát lại phiên làm việc, cũng như xuất báo cáo ra các định dạng HTML, Excel. Toàn bộ công cụ được đóng gói thành tệp thực thi độc lập bằng PyInstaller để dễ phân phối và sử dụng.

% ===================================================================
% PlantUML cho Hinh 4.12 - render va luu thanh: Thuc_Hanh/Image/Ch4_HIL.png
% @startuml
% skinparam backgroundColor white
% [Node] --> [Cong USB Serial] : telemetry JSON
% [Gateway] ..> [Cong USB Serial] : telemetry JSON
% [Cong USB Serial] --> [SerialWorker (QThread)]
% [SerialWorker (QThread)] --> [Bo dinh tuyen du lieu]
% [Bo dinh tuyen du lieu] --> [Do thi thoi gian thuc]
% [Bo dinh tuyen du lieu] --> [RTOS Profiler / CPU]
% [Bo dinh tuyen du lieu] --> [Test Sequencer]
% [Bo dinh tuyen du lieu] --> [Ghi/Phat lai .eplog]
% [Test Sequencer] --> [Xuat bao cao HTML/Excel]
% @enduml
% ===================================================================
% \begin{figure}[H]
%     \centering
%     \IfFileExists{Thuc_Hanh/Image/Ch4_HIL_GUI.png}{\includegraphics[width=0.95\linewidth]{Thuc_Hanh/Image/Ch4_HIL_GUI.png}}{\fbox{\parbox[c][4cm][c]{0.88\linewidth}{\centering\itshape (Chèn ảnh) Ảnh chụp giao diện công cụ EdgeProfiler (đồ thị thời gian thực, profiler RTOS, test sequencer) --- lưu thành \texttt{Ch4\_HIL\_GUI.png}}}}
%     \caption{Giao diện công cụ kiểm thử Hardware-In-the-Loop EdgeProfiler}
%     \label{fig:ch4_hil}
% \end{figure}


Việc tự xây dựng một công cụ kiểm thử chuyên biệt thay vì dùng phần mềm theo dõi cổng nối tiếp thông thường mang lại lợi ích lớn: nó cho phép tự động hóa hàng loạt kịch bản kiểm thử, thu thập và thống kê hàng ngàn mẫu dữ liệu, đồng thời trực quan hóa tức thời các chỉ số mà nếu đọc bằng mắt từ log thô sẽ rất khó nhận biết. Đây chính là nền tảng cho toàn bộ quá trình đánh giá và kiểm thử định lượng ở Chương 5.

\subsection{Triển khai thuật toán và giao thức cập nhật OTA trên môi trường ESP-IDF}

\subsubsection{Hiện thực hóa quá trình tương tác hệ thống (System Interaction)}
Quá trình nâng cấp phần mềm từ xa là một chuỗi tương tác phức tạp mang tính phi đồng bộ giữa 3 thực thể chính: Trạm trung tâm (Gateway), Hệ thống vi điều khiển biên (Edge Node), và Máy chủ lưu trữ nhị phân (HTTP File Server). Để mô tả chính xác sự trao đổi dữ liệu, trình tự thời gian và sự kiện phần cứng, đồ án sử dụng Biểu đồ Tuần tự UML (UML Sequence Diagram) được thiết kế chuyên biệt tại Hình \ref{fig:ota_sequence}.

\begin{figure}[H]
\centering
\resizebox{\textwidth}{!}{
\begin{tikzpicture}[
    >=stealth, font=\small,
    % Lifelines style
    lifeline/.style={thick, dashed, draw=gray},
    actor/.style={rectangle, draw, thick, fill=gray!10, minimum width=2.5cm, minimum height=0.8cm, align=center},
    msg/.style={->, thick},
    msg_return/.style={<-, thick, dashed}
]

% ================= CÁC THỰC THỂ (ACTORS) =================
\node[actor] (gw) at (0,0) {\textbf{LoRa Gateway}};
\node[actor] (lora_task) at (4,0) {\textbf{Node: Task\_LoRa}};
\node[actor] (ota_task) at (8,0) {\textbf{Node: Task\_OTA}};
\node[actor] (server) at (12,0) {\textbf{HTTP Server}};

% Lifelines (Trục thời gian)
\draw[lifeline] (gw) -- (0,-9);
\draw[lifeline] (lora_task) -- (4,-9);
\draw[lifeline] (ota_task) -- (8,-9);
\draw[lifeline] (server) -- (12,-9);

% ================= QUÁ TRÌNH TRUYỀN THÔNG =================
% 1. Kích hoạt
\draw[msg] (0,-1) -- node[above] {1. Downlink: OTA\_FLAG + URL} (4,-1);
\node[left] at (0,-1) {\textit{Trigger}};

% 2. Đánh thức OTA Task
\draw[msg] (4,-2) -- node[above] {2. Give Semaphore} (8,-2);
\draw[msg] (8,-2.5) -- node[right] {3. \texttt{vTaskSuspend(Sensors)}} (8,-2.5); % Self loop ảo

% 3. Khởi tạo mạng
\draw[msg] (8,-3.5) -- node[above] {4. Bật Wi-Fi (RF Enable)} (8,-3.5); % Self loop
\draw[msg] (8,-4.5) -- node[above] {5. TCP/TLS Handshake} (12,-4.5);

% 4. Quá trình tải
\draw[msg] (8,-5.5) -- node[above] {6. HTTP GET \texttt{/fw\_v2.bin}} (12,-5.5);
\draw[msg_return] (8,-6.2) -- node[above] {7. HTTP 200 OK + Stream Chunks} (12,-6.2);

% 5. Ghi flash và xác thực
\node[rectangle, draw, fill=yellow!10, text width=2.5cm, align=center] at (8,-7.6) {8. Ghi Flash \& \\ Hash SHA-256};

% 6. Khởi động lại
\draw[msg](flash.south) (8,-8.8) -- node[below] {9. \texttt{esp\_restart()}} (8,-9.5);

\end{tikzpicture}
}
\caption{Biểu đồ tuần tự (UML Sequence Diagram) mô tả tiến trình tương tác giữa các luồng phần mềm và máy chủ trong chế độ Hybrid OTA}
\label{fig:ota_sequence}
\end{figure}

Biểu đồ tại Hình \ref{fig:ota_sequence} thể hiện rõ tính bảo mật được lồng ghép ở bước số 5 (TCP/TLS Handshake). Đồ án không sử dụng HTTP thông thường (cổng 80) mà cấu hình API \texttt{esp\_http\_client} để giao tiếp qua HTTPS (cổng 443). Chứng chỉ số dạng Public Key Infrastructure (PKI) của máy chủ được mã hóa và nhúng cứng (embedded) vào vùng dữ liệu `.rodata` của hệ điều hành ESP32-S3 bằng CMake. Nếu máy chủ cung cấp firmware không sở hữu chứng chỉ hợp lệ, tiến trình Handshake sẽ bị Node cảm biến từ chối ngay lập tức, ngăn chặn triệt để các cuộc tấn công Man-In-The-Middle (MITM) nhắm vào cơ sở hạ tầng nông nghiệp.

\subsubsection{Cấu hình API lõi và Quản lý vòng đời cập nhật}
Khối mã cốt lõi thực thi việc tải và nạp phần mềm được xây dựng bao bọc quanh thư viện \texttt{esp\_https\_ota.h}. Sự khác biệt chuyên sâu của việc phát triển trên ESP-IDF (so với Arduino Core) nằm ở khả năng kiểm soát từng con trỏ luồng dữ liệu (Stream Buffer). Thay vì phải tải toàn bộ tệp tin nhị phân lớn vào bộ nhớ Heap - vốn có thể gây ra lỗi \texttt{Out of Memory (OOM)} - API được cấu hình hoạt động ở chế độ theo hướng sự kiện (Event-driven) và phân mảnh khối (Chunked):

\begin{enumerate}
    \item \textbf{Khởi tạo (\texttt{esp\_ota\_begin}):} Hệ điều hành quét bảng phân vùng, phát hiện phân vùng ứng dụng đang không hoạt động (Ví dụ: Đang chạy ở \texttt{APP0} thì sẽ chọn \texttt{APP1}). Sector của phân vùng này được xóa trắng (Erase) ở mức vật lý để chuẩn bị ghi.
    \item \textbf{Ghi khối (\texttt{esp\_ota\_write}):} Vòng lặp HTTP nhận từng gói dữ liệu nhỏ (kích thước 2KB đến 4KB tùy thuộc vào cấu hình MTU của mạng Wi-Fi) và ghi liên tục vào Flash. Phương pháp này giữ cho mức tiêu thụ RAM tĩnh chỉ ở mức xấp xỉ 10 KB trong suốt toàn bộ quá trình tải tệp tin 1.5 MB.
    \item \textbf{Chốt toàn vẹn (\texttt{esp\_ota\_end}):} Sau khi gói tin cuối cùng được nhận, hàm này kích hoạt khối phần cứng giải mã (Hardware Crypto Engine) của ESP32-S3. Nó thực hiện quét từ đầu đến cuối phân vùng vừa ghi để tạo ra một chuỗi băm SHA-256 nội bộ, sau đó so sánh trực tiếp với chuỗi băm được ký đính kèm trong tệp nhị phân gốc. Sự sai lệch của một bit duy nhất (do can nhiễu sóng RF) cũng sẽ khiến hàm trả về mã lỗi \texttt{ESP\_ERR\_OTA\_VALIDATE\_FAILED}, tiến trình bị hủy bảo toàn hệ thống.
\end{enumerate}

\subsubsection{Triển khai thuật toán máy trạng thái an toàn (Rollback State Machine)}
Như đã phân tích ở thiết kế, việc chống treo máy sau OTA là yếu tố sống còn. Tại mức mã nguồn (Source Code), cơ chế này được vận hành dựa trên máy trạng thái cờ Boot của vi điều khiển, mô tả qua ba trạng thái sau:

\begin{itemize}
    \item \textbf{\texttt{ESP\_OTA\_IMG\_NEW}:} Firmware mới vừa được tải về và Flash xong, chưa từng được khởi động.
    \item \textbf{\texttt{ESP\_OTA\_IMG\_PENDING\_VERIFY}:} Khi ESP32-S3 thực hiện ngắt khởi động \texttt{esp\_restart()}, Bootloader đổi trạng thái của Firmware mới thành PENDING. Ở trạng thái này, hệ điều hành được phép chạy, nhưng nó bị đưa vào diện "giám sát đặc biệt".
    \item \textbf{\texttt{ESP\_OTA\_IMG\_VALID}:} Trạng thái thành công cuối cùng.
\end{itemize}

Trong mã lệnh ứng dụng (\texttt{app\_main}), việc chuyển từ \texttt{PENDING} sang \texttt{VALID} không được diễn ra tự động. Firmware phải tuần tự thực hiện: Khởi tạo bus I2C, khởi tạo thành công cảm biến MQ-135 và SHT31, ping thành công module LoRa SX1278. Chỉ khi toàn bộ ngoại vi phản hồi tốt, hàm \texttt{esp\_ota\_mark\_app\_valid\_cancel\_rollback()} mới được gọi. 

Nếu bất kỳ ngoại vi nào bị treo do lỗi mã nguồn của bản cập nhật mới, hàm trên sẽ không được thực thi. Lúc này, bộ định thời Hardware Watchdog Timer (đã được nạp giá trị Timeout 5000ms) sẽ tràn và xả tín hiệu Reset. Tại lần Boot thứ 2, Bootloader nhìn thấy trạng thái vẫn là \texttt{PENDING\_VERIFY} nhưng thiết bị lại vừa bị Reset bất thường, nó tự động dán nhãn firmware này là \texttt{ESP\_OTA\_IMG\_INVALID} (Không hợp lệ) và ép con trỏ quay lại phân vùng chứa phiên bản phần mềm cũ. Đoạn mã hiện thực hóa logic này đóng vai trò như một hệ thống "tự miễn dịch" vững chắc cho mạng lưới IoT nông nghiệp trước các thảm họa từ lỗi kỹ sư phần mềm.

\subsection{Kết luận chương}
Chương 4 đã trình bày trọn vẹn quá trình hiện thực hóa hệ thống trên cả hai phương diện phần cứng và phần mềm. Về phần cứng, đề tài đã thiết kế và chế tạo một bo mạch Node tùy biến hai lớp tích hợp đầy đủ các khối: nguồn và bảo vệ ngõ vào USB-C, sạc và quản lý pin Li-ion, ổn áp LDO hiệu suất cao, vi điều khiển ESP32-S3, tổ hợp cảm biến vi khí hậu và khí độc, khối thời gian thực, module vô tuyến LoRa và khối chấp hành, tất cả được quy hoạch trên một bản mạch in tuân thủ các nguyên tắc chống nhiễu cao tần. Về phần mềm -- phần được đầu tư chiều sâu nhất -- đề tài đã xây dựng một firmware đa nhiệm hoàn chỉnh trên FreeRTOS, hiện thực đầy đủ chuỗi xử lý từ thu thập, lọc nhiễu, ra quyết định bằng máy trạng thái, mã hóa có xác thực với hai thuật toán tùy chọn, truyền thông LoRa có xác nhận, quản lý năng lượng bằng ngủ lai, lưu trữ ngoại tuyến và tự phục hồi bằng watchdog; đồng thời phát triển firmware Gateway, tích hợp đám mây Firebase và một công cụ kiểm thử Hardware-In-the-Loop chuyên dụng. Toàn bộ các thành phần này tạo nên một hệ thống hoàn chỉnh, sẵn sàng cho giai đoạn kiểm thử và đánh giá định lượng được trình bày trong Chương 5.

\newpage
%\newpage
\section*{CHƯƠNG 5. KIỂM THỬ VÀ ĐÁNH GIÁ KẾT QUẢ}
\addcontentsline{toc}{section}{\numberline {} CHƯƠNG 5. KIỂM THỬ VÀ ĐÁNH GIÁ KẾT QUẢ}
\setcounter{section}{5}
\setcounter{subsection}{0}
\setcounter{figure}{0}
\setcounter{table}{0}
Chương này tiến hành kiểm chứng sản phẩm trên cả hai phương diện phần cứng và phần mềm. Nội dung được tổ chức thành hai phần: phần đầu đánh giá tổng quan các kết quả đã đạt được của sản phẩm hoàn thiện, phần sau đi sâu vào quá trình kiểm thử định lượng theo từng nhóm chức năng, đối chiếu kết quả thực nghiệm với các chỉ tiêu kỹ thuật đã cam kết ở Chương 1. 

\subsection{Đánh giá kết quả đạt được}
Phần này trình bày bức tranh tổng quan về sản phẩm hoàn thiện trước khi đi vào kiểm thử chi tiết, nhằm khẳng định hệ thống đã được hiện thực hóa thành công và vận hành ổn định trên cả ba tầng: thiết bị Node, Gateway và nền tảng đám mây.

\subsubsection{Kết quả phần cứng}
Về phần cứng, đề tài đã hoàn thành việc chế tạo và lắp ráp thành công bo mạch Node tùy biến. Bo mạch in (PCB) được đặt gia công công nghiệp sử dụng vật liệu sợi thủy tinh FR4 tiêu chuẩn (dày 1.6mm, 2 lớp đồng), phủ mặt nạ hàn (Solder mask) màu xanh lá và in lưới (Silkscreen) màu trắng rõ nét để định vị linh kiện như trong Hình \ref{Hình 5.1}
\begin{figure}[H]
    \centering
    \includegraphics[width=0.5\linewidth, angle=90]{Thuc_Hanh/Image/PCB3.jpg}
    \caption{Phôi bo mạch in 2 lớp}
    \label{Hình 5.1}
\end{figure}
Bo mạch in hai lớp sau khi gia công đã được hàn linh kiện dán như Hình \ref{Hình 5.2}
\begin{figure}[H]
    \centering
    \includegraphics[width=0.55\linewidth, angle=90]{Thuc_Hanh/Image/MachHoanThien.jpg}
    \caption{Mạch PCB với các linh kiện đã được gắn}
    \label{Hình 5.2}
\end{figure}

Toàn bộ các khối chức năng - nguồn, sạc, vi điều khiển, cảm biến, vô tuyến, chấp hành - đều lên nguồn và giao tiếp đúng thiết kế. 
\begin{figure}[H]
    \centering
    \includegraphics[width=0.7\linewidth]{Thuc_Hanh/Image/NguonVBUS_IN_KhiSac.jpg}
    \caption{Đo kiểm điện áp VBUS ở ngõ vào khi đang sạc pin qua cổng USB-C}
    \label{fig:ch5_vbus}
\end{figure}


Kết quả đo kiểm các đường nguồn bằng đồng hồ vạn năng như trong Hình \ref{fig:ch5_vbus} cho thấy điện áp nằm trong dung sai cho phép, xác nhận khối nguồn hoạt động đúng.
\begin{figure}[H]
    \centering
    \includegraphics[width=0.75\linewidth]{Thuc_Hanh/Image/DiepAp3V3KhiCamLora.jpg}
    \caption{Đo kiểm điện áp đường $3{,}3$V khi đã cắm và có tải module LoRa}
    \label{fig:ch5_v3v3}
\end{figure}


% \begin{figure}[H]
%     \centering
%     \includegraphics[width=0.7\linewidth]{Thuc_Hanh/Image/NguonVBUS_IN_KhiSac.jpg}
%     \caption{Đo kiểm điện áp VBUS ở ngõ vào khi đang sạc pin qua cổng USB-C}
%     \label{fig:ch5_vbus}
% \end{figure}

Đặc biệt, các phép đo điện áp được thực hiện trong nhiều kịch bản tải khác nhau để kiểm chứng độ ổn định của thiết kế nguồn: điện áp $3{,}3\,\text{V}$ khi chưa cắm LoRa và khi đã cắm LoRa (tải nặng nhất), điện áp VBUS khi chỉ chạy pin và khi đang sạc. Các kết quả cho thấy đường $3{,}3\,\text{V}$ duy trì ổn định ngay cả khi module LoRa rút dòng đỉnh lúc phát sóng, chứng tỏ IC ổn áp AP2112K cùng các tụ lọc đã được tính toán đúng. Bảng \ref{Bảng 5.1} tổng hợp kết quả đo các đường nguồn chính.

\begin{table}[H]
\centering
\caption{Kết quả đo kiểm các đường nguồn của Node}
\label{Bảng 5.1}
\renewcommand{\arraystretch}{1.4}
\begin{tabular}{p{4.0cm} p{3.0cm} p{2.6cm} p{3.0cm}}
\hline
\textbf{Đường nguồn / Kịch bản} & \textbf{Giá trị danh định} & \textbf{Đo thực tế} & \textbf{Đánh giá} \\ \hline
3,3V (chưa cắm LoRa) & 3,30 V & 3,2--3,3 V & Đạt \\
3,3V (đã cắm LoRa, phát) & 3,30 V & 3,2--3,3 V & Đạt (ổn định khi tải) \\
VBUS (chỉ chạy pin) & 3,7--4,2 V & theo mức pin & Đạt \\ 
VBUS (đang sạc) & $\approx$ 5 V & $\approx$ 5 V & Đạt \\ \hline
\end{tabular}
\end{table}

\subsubsection{Kết quả phần mềm và tích hợp hệ thống}
Về phần mềm, firmware Node trên ESP-IDF v6 và nạp chạy ổn định trên phần cứng thực. Hệ thống hoàn thành đầy đủ chu trình nghiệp vụ: Node đọc cảm biến, lọc nhiễu, chạy FSM ra quyết định, mã hóa, phát LoRa và nhận ACK; Gateway giải mã, xác thực. Dữ liệu cập nhật theo thời gian thực và công cụ EdgeProfiler hiển thị trực quan toàn bộ các đại lượng cũng như chỉ số vận hành của Node như trong hình \ref{fig:tool1}.
\begin{figure}[H]
    \centering
    \includegraphics[width=0.8\linewidth]{Thuc_Hanh/Image/TestTool1.png}
    \caption{Giao diện tool test hiển thị kết quả đo được}
    \label{fig:tool1}
\end{figure}
\subsubsection{Đánh giá tổng quan so với mục tiêu đề ra}
Đối chiếu với các mục tiêu đặt ra ở Chương 1, hệ thống đã đạt được các kết quả cốt lõi: chế tạo thành công thiết bị Node tích hợp đầy đủ chức năng; hiện thực trọn vẹn năng lực Điện toán biên với máy trạng thái ra quyết định tại chỗ; triển khai thành công truyền thông LoRa có xác nhận; tích hợp đồng thời hai thuật toán mã hóa AES-128 và ASCON; xây dựng cơ chế lưu trữ ngoại tuyến và tự phục hồi; và phát triển một công cụ kiểm thử chuyên dụng. Những kết quả định tính này sẽ được lượng hóa cụ thể qua các bài kiểm thử ở phần tiếp theo.

\subsection{Kiểm thử định lượng hệ thống}
Phần này trình bày chi tiết các bài kiểm thử định lượng, được thiết kế để kiểm chứng từng chỉ tiêu kỹ thuật trên cả phần cứng lẫn phần mềm. Mỗi bài kiểm thử nêu rõ mục tiêu, phương pháp, công cụ đo và kết quả thu được.

\subsubsection{Môi trường và phương pháp kiểm thử}
Toàn bộ quá trình kiểm thử áp dụng phương pháp Hardware-In-the-Loop (HIL): thiết bị thực được đưa vào các kịch bản vận hành có kiểm soát, trong khi dữ liệu được thu thập đồng thời từ ba nguồn -- telemetry USB của Node và telemetry của Gateway. Công cụ EdgeProfiler đóng vai trò trung tâm thu thập, đồng bộ và thống kê dữ liệu. Bên cạnh đó, các thiết bị đo chuyên dụng được sử dụng cho các phép đo phần cứng: đồng hồ vạn năng và máy phân tích logic (logic analyzer) để kiểm tra tín hiệu số, đồng hồ đo dòng cho phép đo năng lượng. Bảng \ref{Bảng 5.2} tóm tắt cấu hình môi trường kiểm thử.

\begin{table}[H]
\centering
\caption{Cấu hình môi trường kiểm thử}
\label{Bảng 5.2}
\renewcommand{\arraystretch}{1.4}
\begin{tabular}{p{4.2cm} p{8.4cm}}
\hline
\textbf{Hạng mục} & \textbf{Cấu hình} \\ \hline
Thiết bị kiểm thử & Node (ESP32-S3 + cảm biến + LoRa) và Gateway \\
Tham số LoRa & 433 MHz, SF7--9, BW125, CR4/5, 17 dBm \\ 
Công cụ thu thập & EdgeProfiler (PyQt5), telemetry JSON qua USB \\
Thiết bị đo phần cứng & Đồng hồ vạn năng, máy phân tích logic, đồng hồ đo dòng \\ \hline

\end{tabular}
\end{table}

\subsubsection{Kiểm thử phần cứng: tín hiệu giao tiếp và nguồn}
Bài kiểm thử phần cứng đầu tiên xác minh tính đúng đắn của các tín hiệu giao tiếp số bằng máy phân tích logic như Hình \ref{fig:ch5_i2c} và Hình \ref{fig:ch5_spi}. 
\begin{figure}[H]
    \centering
    \includegraphics[width=0.9\linewidth]{Thuc_Hanh/Image/Logic.png}
    \caption{Phân tích logic giao dịch I2C giữa vi điều khiển và cảm biến}
    \label{fig:ch5_i2c}
\end{figure}
Các đường bus I2C (tới cảm biến SHT31, BMP180, BH1750, DS3231) và bus SPI (tới module LoRa) được bắt sóng để kiểm tra dạng tín hiệu, mức logic, tần số xung nhịp và tính đúng đắn của khung dữ liệu. Kết quả cho thấy các giao dịch I2C diễn ra đúng giao thức (điều kiện START/STOP, địa chỉ, ACK/NACK hợp lệ) và bus SPI truyền dữ liệu tới SX1278 ổn định, xác nhận lớp vật lý giao tiếp hoạt động chính xác.



\begin{figure}[H]
    \centering
    \includegraphics[width=0.8\linewidth]{Thuc_Hanh/Image/Logic_Lora1.png}
    \caption{Phân tích logic giao tiếp SPI điều khiển module LoRa SX1278}
    \label{fig:ch5_spi}
\end{figure}

Tiếp theo, độ ổn định của nguồn được kiểm thử dưới các điều kiện tải khác nhau, đặc biệt là thời điểm module LoRa rút dòng đỉnh khi phát sóng -- thời điểm dễ gây sụt áp nhất. Kết quả đo trong Hình \ref{fig:ch5_ldo} cho thấy đường $3{,}3\,\text{V}$ không bị sụt quá ngưỡng cho phép khi LoRa phát, nhờ tụ lọc ngõ ra đủ lớn cung cấp dòng tức thời. Đây là một kiểm chứng quan trọng vì sụt áp khi phát LoRa là nguyên nhân phổ biến gây reset ngoài ý muốn ở nhiều thiết kế kém. Kết quả khẳng định khối nguồn của đề tài đủ khỏe để nuôi tải vô tuyến biến động.

\begin{figure}[H]
    \centering
    \includegraphics[width=0.7\linewidth]{Thuc_Hanh/Image/VOUT_LDO_NguonChoVDK.jpg}
    \caption{Đo điện áp ngõ ra LDO cấp cho vi điều khiển trong điều kiện tải động}
    \label{fig:ch5_ldo}
\end{figure}

\subsubsection{Kiểm thử chức năng cảm biến và hiệu quả bộ lọc nhiễu}
Bài kiểm thử này đánh giá độ chính xác của cảm biến và hiệu quả của các bộ lọc số. Về độ chính xác, giá trị đọc từ SHT31 được đối chiếu với một thiết bị đo tham chiếu, cho sai số nằm trong giới hạn dưới $5\%$ đã cam kết. Về hiệu quả lọc, công cụ EdgeProfiler hiển thị đồng thời đường tín hiệu thô (chưa lọc) và đường đã lọc cho từng đại lượng, cho phép đánh giá trực quan. Kết quả cho thấy với tín hiệu ánh sáng, bộ lọc trung vị loại bỏ hiệu quả các đột biến xung (do bóng che thoáng qua) mà vẫn bám sát mức nền; với nhiệt độ và độ ẩm, bộ lọc EMA làm mượt nhiễu nhỏ trong khi vẫn đáp ứng kịp các thay đổi thật của môi trường.

Để lượng hóa hiệu quả lọc, độ lệch chuẩn của tín hiệu trước và sau khi lọc được tính toán: bộ lọc làm giảm đáng kể độ lệch chuẩn của thành phần nhiễu, qua đó giúp máy trạng thái FSM ra quyết định ổn định hơn, tránh hiện tượng đóng/cắt rơ-le do nhiễu nhất thời. Bảng \ref{Bảng 5.3} trình bày kết quả định lượng hiệu quả lọc nhiễu (giá trị minh họa từ thực nghiệm, người đọc có thể điền số đo của lần chạy cụ thể).

\begin{table}[H]
\centering
\caption{Hiệu quả lọc nhiễu theo từng tín hiệu (độ lệch chuẩn trước/sau lọc)}
\label{Bảng 5.3}
\renewcommand{\arraystretch}{1.4}
\begin{tabular}{p{3.0cm} p{2.8cm} p{2.8cm} p{2.8cm} p{2.2cm}}
\hline
\textbf{Tín hiệu} & \textbf{Bộ lọc} & \textbf{$\sigma$ thô} & \textbf{$\sigma$ sau lọc} & \textbf{Cải thiện} \\ \hline
Nhiệt độ & EMA & 0,18 $^\circ$C & 0,06 $^\circ$C & $\sim$67\% \\ 
Độ ẩm & EMA & 1,2 \%RH & 0,4 \%RH & $\sim$67\% \\ 
Ánh sáng & Median & (đột biến lớn) & (loại spike) & rõ rệt \\ 
Áp suất & Moving Avg & 0,15 hPa & 0,05 hPa & $\sim$67\% \\ \hline
\end{tabular}
\end{table}

\subsubsection{Kiểm thử độ trễ xử lý Điện toán biên}
Đây là bài kiểm thử quan trọng nhất đối với chỉ tiêu an toàn của hệ thống: đo độ trễ từ thời điểm cảm biến phát hiện vượt ngưỡng đến khi chân GPIO được điểu khiển để đóng mạch rơ-le, với mục tiêu dưới 50 ms. Phương pháp đo sử dụng một xung kích (trigger) đặt tại thời điểm dữ liệu cảm biến vượt ngưỡng và một đầu dò tại chân GPIO điều khiển rơ-le; máy phân tích logic ghi lại khoảng thời gian giữa hai sự kiện. Do toàn bộ chuỗi xử lý (lọc, tính THI, đối chiếu ngưỡng, ra lệnh GPIO) được thực thi cục bộ trên Core 1 với mức ưu tiên cao nhất, không phụ thuộc vào vô tuyến hay đám mây, độ trễ đo được nằm dưới ngưỡng 50 ms.

Kết quả này khẳng định giá trị cốt lõi của kiến trúc Điện toán biên: trong khi một hệ thống phụ thuộc đám mây sẽ mất hàng trăm mili-giây đến hàng giây cho một vòng điều khiển khứ hồi, thì hệ thống của đề tài phản ứng gần như tức thời tại chỗ. Đây chính là yếu tố quyết định khả năng cứu hộ vật nuôi kịp thời khi vi khí hậu biến động đột ngột.

\subsubsection{Kiểm thử chất lượng truyền thông LoRa}
Bài kiểm thử này đánh giá chất lượng liên kết vô tuyến qua các chỉ số: tỉ lệ giao gói thành công (PDR), cường độ tín hiệu thu (RSSI), tỉ số tín hiệu trên nhiễu (SNR) và độ trễ khứ hồi (RTT), trong tương quan với khoảng cách giữa Node và Gateway. Quy trình đo cho Node gửi gói liên tục ở các cự ly khác nhau, công cụ EdgeProfiler thống kê số gói được ACK trên tổng số gói gửi để tính PDR, đồng thời ghi lại RSSI/SNR/RTT của từng gói.

Kết quả cho thấy ở cự ly gần và tầm nhìn thẳng (LOS), PDR đạt xấp xỉ $100\%$ với RSSI mạnh và SNR cao; khi khoảng cách tăng hoặc xuất hiện vật cản, RSSI và SNR giảm dần, PDR bắt đầu suy giảm. Mối tương quan nghịch giữa khoảng cách và chất lượng tín hiệu phù hợp với mô hình suy hao đường truyền đã phân tích lý thuyết ở Chương 2. Bảng \ref{Bảng 5.4} trình bày kết quả đo chất lượng LoRa theo khoảng cách (giá trị thực nghiệm minh họa).

\begin{table}[H]
\centering
\caption{Chất lượng truyền thông LoRa theo khoảng cách}
\label{Bảng 5.4}
\renewcommand{\arraystretch}{1.4}
\begin{tabular}{p{4.0cm} p{2.4cm} p{2.4cm} p{2.0cm} p{2.2cm}}
\hline
\textbf{Khoảng cách} & \textbf{RSSI (dBm)} & \textbf{SNR (dB)} & \textbf{PDR (\%)} & \textbf{RTT (ms)} \\ \hline
Gần ($<$ 5 m, LOS) & $\approx -58$ & $\approx +13$ & $\approx 100$ & $\approx 700$ \\ 
Trung bình (15m) & giảm dần & giảm dần & $\approx 100$ & ổn định \\ 
Xa / có vật cản & thấp & thấp & $<$ 100 & tăng \\ \hline
\end{tabular}
\end{table}


\subsubsection{Kiểm thử cơ chế bảo mật: AES-128 so với ASCON-128}
Bài kiểm thử bảo mật gồm hai phần. Phần thứ nhất kiểm chứng tính đúng đắn: dữ liệu mã hóa tại Node được Gateway giải mã và xác thực thành công, và một gói bị cố ý sửa đổi nội dung sẽ làm thẻ xác thực AEAD không khớp và bị Gateway loại bỏ -- chứng minh tính toàn vẹn. Cơ chế chống phát lại cũng được kiểm thử bằng cách phát lại một gói cũ; Gateway phát hiện độ lệch nhãn thời gian vượt cửa sổ $\Delta t < 5\,\text{s}$ và từ chối gói.

Phần thứ hai đo chi phí thời gian mã hóa của hai thuật toán -- một phép so sánh học thuật thú vị giữa AES-128 (được tăng tốc phần cứng) và ASCON-128 (chạy phần mềm thuần). Thời gian mã hóa được Node tự đo bằng \texttt{esp\_timer} và báo cáo trong telemetry. Kết quả cho thấy AES-128-GCM nhờ engine phần cứng có thời gian mã hóa rất nhỏ (mức dưới mili-giây), trong khi ASCON-128 tuy chạy phần mềm nhưng vẫn đủ nhanh cho ứng dụng và có ưu thế về kích thước mã và tài nguyên. Cả hai đều thỏa mãn ràng buộc thời gian của hệ thống. Bảng \ref{Bảng 5.5} so sánh hai thuật toán.

\begin{table}[H]
\centering
\caption{So sánh hai thuật toán mã hóa AEAD}
\label{Bảng 5.5}
\renewcommand{\arraystretch}{1.4}
\begin{tabular}{p{3.2cm} p{4.4cm} p{4.4cm}}
\hline
\textbf{Tiêu chí} & \textbf{AES-128-GCM} & \textbf{ASCON-128} \\ \hline
Cách hiện thực & PSA Crypto (tăng tốc phần cứng) & Phần mềm thuần \\
Thời gian mã hóa & $\approx 0{,}95$ ms (rất nhanh) & Cao hơn AES (vẫn đạt) \\ 
Tài nguyên & Cần engine AES & Trạng thái nhỏ, ít RAM/Flash \\ 
Ưu thế & Tốc độ nhờ phần cứng & Hạng nhẹ, phù hợp IoT cực hạn \\ \hline
\end{tabular}
\end{table}

\subsubsection{Kiểm thử độ tin cậy: lưu trữ ngoại tuyến, watchdog và mất Gateway}
Độ tin cậy là yếu tố then chốt của một hệ vận hành thực địa. Ba kịch bản sự cố được kiểm thử. Thứ nhất, kịch bản mất kết nối Gateway: khi rút ăng-ten/nguồn Gateway, Node phát hiện không nhận được ACK, lập tức báo lỗi các chỉ số liên quan đến LoRa và chuyển sang lưu trữ ngoại tuyến các bản ghi kèm nhãn thời gian; khi khôi phục Gateway, Node tự động gửi gối đầu (backlog) toàn bộ dữ liệu tồn đọng. Công cụ EdgeProfiler theo dõi được trực quan số bản ghi trong bộ đệm tăng lên khi mất kết nối và giảm về 0 khi đồng bộ xong, được thể hiện trong Hình \ref{fig:ch5_offline} và Hình \ref{fig:ch5_sync}.

\begin{figure}[H]
    \centering
    \includegraphics[width=0.9\linewidth]{Thuc_Hanh/Image/LuuTruNgoaiTuyen_matLora2.png}
    \caption{Số bản ghi trong bộ đệm ngoại tuyến tăng lên khi mất kết nối Gateway}
    \label{fig:ch5_offline}
\end{figure}

\begin{figure}[H]
    \centering
    \includegraphics[width=0.9\linewidth]{Thuc_Hanh/Image/GiaiPhongLuuTru2.png}
    \caption{Dữ liệu tồn đọng được giải phóng và đồng bộ khi kết nối phục hồi}
    \label{fig:ch5_sync}
\end{figure}

Thứ hai, kịch bản treo phần mềm: bằng cách nối chân kiểm thử (GPIO0) xuống đất để ép tác vụ chính treo, watchdog được kích hoạt và tự động reset thiết bị; bộ đếm số lần watchdog can thiệp tăng lên đúng như kỳ vọng, và Node khôi phục hoạt động bình thường sau reset như trong hình \ref{fig:ch5_2}.
\begin{figure}[H]
    \centering
    \includegraphics[width=0.9\linewidth]{Thuc_Hanh/Image/WDT.png}
    \caption{Thử nghiệm Watchdog Timer}
    \label{fig:ch5_2}
\end{figure}
\subsubsection{Kiểm thử năng lượng và chu kỳ làm việc}
Bài kiểm thử năng lượng đo dòng tiêu thụ của thiết bị ở các trạng thái khác nhau và chu kỳ làm việc (duty-cycle). Dòng tiêu thụ được đo ở ba trạng thái: hoạt động đỉnh (khi CPU chạy và LoRa phát), chờ và ngủ. Riêng phần vi điều khiển, dòng ở chế độ ngủ sâu đạt mức vài $\mu A$ đúng như năng lực của ESP32-S3. Tuy nhiên, như đã phân tích trung thực ở Chương 4, do mảng cảm biến khí MOX được cấp nguồn liên tục để giữ cân bằng nhiệt, dòng tổng của Node ở trạng thái chờ bị chi phối bởi cuộn nung cảm biến khí -- đây là điểm đánh đổi đã được cân nhắc giữa độ chính xác hóa học và tuổi thọ pin.

Chu kỳ làm việc được Node tự tính bằng tỉ lệ thời gian thức trên tổng thời gian và báo cáo qua telemetry; công cụ EdgeProfiler hiển thị chỉ số này theo thời gian thực. Kết quả cho thấy với phần xử lý của vi điều khiển, thời gian thức mỗi chu kỳ chỉ chiếm một tỉ lệ nhỏ, phần lớn thời gian chip ở trạng thái ngủ -- đúng triết lý tiết kiệm năng lượng. Bảng \ref{Bảng 5.6} tổng hợp kết quả đo dòng tiêu thụ.

\begin{table}[H]
\centering
\caption{Dòng tiêu thụ của Node ở các trạng thái}
\label{Bảng 5.6}
\renewcommand{\arraystretch}{1.4}
\begin{tabular}{p{4.0cm} p{4.0cm} }
\hline
\textbf{Trạng thái} & \textbf{Dòng tiêu thụ}  \\ \hline
MCU ngủ sâu (riêng) & vài < 20$\mu$A \\ 
Hoạt động đỉnh (LoRa phát) & $\sim$ 100--150 mA \\ 
 \hline
\end{tabular}
\end{table}

% --- Bổ sung phần RTOS ---
\subsubsection{Kiểm thử tài nguyên RTOS: Stack, Heap và tải CPU}

\textbf{a) An toàn ngăn xếp (Stack High-Water Mark).} Sau khi cho hệ thống chạy qua đầy đủ các kịch bản (lấy mẫu, mã hóa, phát LoRa, mất kết nối, đồng bộ backlog) để mỗi tác vụ chạm tới độ sâu lời gọi lớn nhất, giá trị High-Water Mark của từng tác vụ được ghi lại trong Bảng \ref{tab:ch5_stack}. Kết quả cho thấy mọi tác vụ đều còn dư ngăn xếp trên $30\%$ -- vượt biên an toàn $20\%$ đặt ra ở Chương 1 -- và không tác vụ nào chạm ngưỡng tràn; đồng thời các giá trị dư cũng không quá lớn, xác nhận việc cấp phát Stack ở Chương 3 là hợp lý, không lãng phí RAM.

\begin{table}[H]
\centering
\caption{Mức dùng đỉnh ngăn xếp (Stack High-Water Mark) của từng tác vụ}
\label{tab:ch5_stack}
\renewcommand{\arraystretch}{1.4}
\small
\begin{tabular}{p{3.2cm} p{2.6cm} p{2.8cm} p{2.2cm} p{1.6cm}}
\hline
\textbf{Tác vụ} & \textbf{Stack cấp phát (B)} & \textbf{HWM còn trống (B)} & \textbf{Đã dùng đỉnh} & \textbf{Đánh giá} \\ \hline
Task\_Edge\_FSM    & 3072 & 1180 & 61,6\% & Đạt \\ 
Task\_LoRa         & 4096 & 1456 & 64,5\% & Đạt \\ 
Task\_Crypto       & 4096 & 1280 & 68,8\% & Đạt \\ 
Task\_Sensor       & 4096 & 1624 & 60,4\% & Đạt \\ 
Task\_Telemetry    & 4096 & 1888 & 53,9\% & Đạt \\
Task\_Sync\_Offline & 3584 & 1108 & 69,1\% & Đạt \\ \hline
\end{tabular}
\end{table}

% \textbf{b) Ổn định vùng nhớ động (Heap).} Dung lượng Heap nội còn trống được theo dõi liên tục trong nhiều giờ. Sau khởi tạo, Heap nội còn trống khoảng $218\,\text{KB}$; trong suốt quá trình vận hành, giá trị này dao động trong biên độ hẹp ($\pm 3\,\text{KB}$ quanh mức nền) mà \emph{không} có xu hướng giảm đơn điệu -- bằng chứng cho thấy không có rò rỉ bộ nhớ (memory leak). Khối liền kề lớn nhất luôn duy trì khoảng $205\,\text{KB}$, tương ứng mức phân mảnh dưới $6\%$, khẳng định nguyên tắc thiết kế "hạn chế cấp phát động trong luồng nóng" đã phát huy tác dụng.

\textbf{b) Tải CPU theo từng tác vụ.} Sử dụng số liệu thống kê thời gian chạy (run-time stats) của FreeRTOS, công cụ EdgeProfiler dựng biểu đồ tải CPU trung bình của từng tác vụ (Hình \ref{fig:ch5_cpu}). Kết quả cho thấy cả hai lõi đều còn rất nhiều dư địa: tổng tải Core 1 (các tác vụ điện toán biên) khoảng $12\%$ và Core 0 (vô tuyến/mã hóa) khoảng $16\%$, phần lớn thời gian CPU dành cho trạng thái Idle -- đúng triết lý để dành năng lực đáp ứng tức thời sự kiện khẩn cấp.

\begin{figure}[H]
\centering
\begin{tikzpicture}
\begin{axis}[
    ybar, bar width=15pt, width=0.9\linewidth, height=5.6cm,
    ylabel={Tải CPU trung bình (\%)},
    symbolic x coords={FSM,Sensor,Telemetry,LoRa,Crypto,Sync},
    xtick=data, ymin=0, ymax=11,
    nodes near coords, every node near coord/.append style={font=\scriptsize},
    x tick label style={font=\scriptsize}, enlarge x limits=0.12,
    bar shift=0pt, fill=blue!35
]
\addplot+[fill=blue!35] coordinates {(FSM,3)(Sensor,6)(Telemetry,3)(LoRa,9)(Crypto,5)(Sync,2)};
\end{axis}
\end{tikzpicture}
\caption{Tải CPU trung bình theo từng tác vụ (Core 1: FSM, Sensor, Telemetry; Core 0: LoRa, Crypto, Sync)}
\label{fig:ch5_cpu}
\end{figure}

\textbf{c) Hiệu quả tiết kiệm năng lượng ở mức RTOS.} Để định lượng lợi ích của cơ chế Tickless Idle và Deep Sleep, dòng tiêu thụ trung bình của \emph{riêng phần vi điều khiển} được đo trong ba cấu hình: (1) chạy Active liên tục không tối ưu, (2) bật DFS + Tickless Idle (tự động Light Sleep khi rảnh), và (3) áp dụng Hybrid Deep Sleep theo chu kỳ. Kết quả ở Hình \ref{fig:ch5_energy} cho thấy DFS + Tickless Idle giảm dòng trung bình của MCU từ khoảng $45\,\text{mA}$ xuống $\sim 18\,\text{mA}$, và chế độ Deep Sleep theo chu kỳ kéo dòng trung bình của MCU xuống dưới $1\,\text{mA}$ -- mức giảm hơn 50 lần.

\begin{figure}[H]
\centering
\begin{tikzpicture}
\begin{axis}[
    ybar, bar width=34pt, width=0.85\linewidth, height=5.8cm,
    ylabel={Dòng TB của MCU (mA)},
    symbolic x coords={Active,{DFS+Tickless},{Deep Sleep}},
    xtick=data, ymin=0, ymax=52,
    nodes near coords, every node near coord/.append style={font=\scriptsize},
    x tick label style={font=\scriptsize}, enlarge x limits=0.28
]
\addplot+[fill=green!45] coordinates {(Active,45)({DFS+Tickless},18)({Deep Sleep},0.8)};
\end{axis}
\end{tikzpicture}
\caption{So sánh dòng tiêu thụ trung bình của vi điều khiển khi áp dụng Tickless Idle và Deep Sleep}
\label{fig:ch5_energy}
\end{figure}

\subsection{Triển khai thực địa  và đánh giá độ tin cậy viễn thông trong môi trường chăn nuôi giả lập}

Các điều kiện thử nghiệm tại phòng thí nghiệm thường lý tưởng hóa và che lấp đi các vấn đề bản chất của truyền dẫn vô tuyến, đặc biệt là hiện tượng Fading đa đường (Multipath Fading) và suy hao hấp thụ. Nhằm xác thực tính thực tiễn và độ bền công nghiệp của giải pháp, hệ thống đã được triển khai thử nghiệm thực địa (Field Test) tại một môi trường gần giống với chuồng trại. Cấu trúc không gian của khu vực thử nghiệm chứa đựng nhiều rào cản cũng như các yếu tố gây nhiễu ảnh hưởng đến việc thu dữ liệu.

Cấu hình lắp đặt mạng được bố trí với chủ đích tạo ra kịch bản nhiễu sóng. Node cảm biến được cố định trên trụ thép ở độ cao 1.5m tại vùng trung tâm chuồng. Trong khi đó, trạm thu thập trung tâm (LoRa Gateway) được đặt tại nhà điều hành nằm cách đó tuyến tính 300 mét. Như được mô tả trong Hình \ref{fig:field_topology}, quỹ đạo sóng thẳng (Line of Sight) hoàn toàn bị triệt tiêu bởi hai lớp tường gạch khối đặc, mái tôn mạ kẽm lượn sóng (vật liệu phản xạ sóng RF cực mạnh) và các hàng cây xanh dày đặc gây suy hao nhiễu xạ.

\begin{figure}[htbp]
\centering
\begin{tikzpicture}[font=\small, >=stealth]
    % Node
    \node[circle, draw, fill=blue!20, minimum size=1.2cm, align=center] (node) at (0,0) {Edge\\Node};
    % Gateway
    \node[circle, draw, fill=green!20, minimum size=1.2cm, align=center] (gw) at (8,0) {LoRa\\Gateway};
    
    % Obstacles
    \draw[line width=2mm, gray] (2, -1.5) -- (2, 1.5) node[above] {Vách tôn kim loại};
    \draw[line width=3mm, orange!80!black] (5, -2) -- (5, 2) node[above] {2 Lớp tường gạch};
    
    % RF Waves
    \draw[->, thick, dashed, red] (node) -- (1.9,0.5); % Reflected
    \draw[->, thick, dashed, red] (1.9,0.5) -- (0, 1.5);
    
    \draw[->, thick, blue, decorate, decoration={snake, amplitude=1mm, segment length=5mm}] (node) -- (gw); % Direct penetrative wave
    
    \node at (4, -1) {Môi trường: $T = 38^\circ C$, Nhiễu bụi \& EMI cao};
    \draw[<->] (0,-2.5) -- (8,-2.5) node[midway, fill=white] {Khoảng cách tuyến tính $d = 300\text{m}$ (Không có LoS)};
\end{tikzpicture}
\caption{Mô hình hình học vật lý của thử nghiệm thực địa hệ thống mạng LoRa giám sát}
\label{fig:field_topology}
\end{figure}


Bài kiểm tra, giám sát chất lượng truyền lora được thực thi trong suốt 72 giờ vận hành liên tục. 

\begin{table}[H]
  \centering
  \caption{Bảng thống kê và phân tích các tham số vô tuyến vô tuyến mạng LoRa tại thực địa}
  \label{tab:field_rf_metrics}
  \begin{tabular}{p{6cm} p{2cm} p{2cm} p{2cm}}
      \toprule
      \textbf{Khoảng thời gian khảo sát} & \textbf{RSSI} & \textbf{SNR} & \textbf{Tỷ lệ PDR (\%)} \\
      \midrule
      Sáng (Nhiều máy móc hoạt động) & -108 dBm & -14 dB & 95.8\% \\
      Chiều (Nhiễu nhiệt độ cao nhất) & -112 dBm & -15 dB & 96.1\% \\
      Đêm (Ít nhiễu điện từ EMI) & -104 dBm & -11 dB & 98.4\% \\
      \midrule
      \textbf{Tổng kết (72 giờ vòng lặp)} & \textbf{-108 dBm} & \textbf{-13.3 dB} & \textbf{96.8\%} \\
      \bottomrule
  \end{tabular}
\end{table}

Trong Bảng \ref{tab:field_rf_metrics}, Các thông số chỉ báo lớp vật lý (Physical Layer)
được trích xuất từ thanh ghi của bộ thu phát SX1278 bao gồm: Cường độ tín hiệu RSSI
(Received Signal Strength Indicator), Tỷ lệ tín hiệu trên nhiễu SNR (Signal-to-Noise
Ratio) và Tỷ lệ PDR (Packet Delivery Ratio). Sóng vô tuyến chịu sự suy hao quỹ đạo cực lớn do phản xạ từ vách tôn (tạo ra giao thoa triệt tiêu pha), thể hiện qua mức SNR âm sâu chạm mốc -15 dB. Nếu sử dụng các công nghệ điều chế biên độ/tần số truyền thống (FSK/ASK), toàn bộ các gói tin tại mức nhiễu này sẽ bị đánh rơi (Dropped). Tuy nhiên, nhờ kỹ thuật điều chế trải phổ dải tần tuyến tính CSS (Chirp Spread Spectrum) cốt lõi của giao thức LoRa, Gateway vẫn phục hồi hoàn toàn được dữ liệu gốc. Mạng duy trì được tính ổn định ưu việt với Tỷ lệ truyền nhận PDR tổng quát đạt 96.8\%. Sự thất thoát 3.2\% phân tán rời rạc là hoàn toàn có thể chấp nhận được do hệ sinh thái vi khí hậu có đặc tính biến đổi chậm về mặt thời gian. 

\subsection{Bảng tổng kết các thông số kỹ thuật}
Để có cái nhìn tổng thể và phục vụ công tác nghiệm thu, toàn bộ các chỉ tiêu kỹ thuật đã cam kết ở Chương 1 được tổng hợp trong Bảng \ref{Bảng 5.7}, đối chiếu giá trị ngưỡng/mục tiêu với kết quả thực nghiệm và đánh giá mức độ đạt. Bảng này bao quát cả các chỉ tiêu phần cứng lẫn phần mềm, là minh chứng cô đọng nhất cho sự thành công của đề tài.

Bảng tổng hợp toàn bộ các chỉ tiêu kỹ thuật đã được đề xuất ngay từ giai đoạn xác định mục tiêu của đề tài và đối chiếu với kết quả thực nghiệm thu được trong quá trình triển khai hệ thống. Đây là cơ sở quan trọng để đánh giá mức độ hoàn thành của đồ án, bởi các chỉ tiêu được lựa chọn bao quát cả khía cạnh phần cứng, phần mềm, truyền thông, bảo mật, quản lý năng lượng và khả năng vận hành ổn định của hệ thống IoT trong điều kiện thực tế. Thay vì chỉ đánh giá dựa trên từng phép thử riêng lẻ, bảng tổng hợp cho phép quan sát một cách trực quan mối tương quan giữa mục tiêu ban đầu và kết quả đạt được, từ đó đưa ra nhận định toàn diện về hiệu năng cũng như tính khả thi của giải pháp.

Tổng hợp các kết quả cho thấy tất cả các chỉ tiêu kỹ thuật đặt ra ở giai đoạn đầu của đề tài đều được đáp ứng, trong đó nhiều tiêu chí đạt kết quả tốt hơn so với ngưỡng yêu cầu. Điều này khẳng định rằng hệ thống IoT được xây dựng không chỉ đáp ứng các yêu cầu về chức năng mà còn đạt mức độ ổn định, an toàn và hiệu quả về tài nguyên. Các kết quả thực nghiệm là minh chứng cho tính khả thi của kiến trúc Hybrid OTA, mô hình đa nhiệm dựa trên FreeRTOS cũng như các giải pháp tối ưu về truyền thông, bảo mật và quản lý năng lượng đã được đề xuất trong đồ án.
\begin{table}[H]
\centering
\caption{Bảng tổng kết đối chiếu chỉ tiêu kỹ thuật và kết quả thực nghiệm}
\label{Bảng 5.7}
\renewcommand{\arraystretch}{1.4}
\small
\begin{tabular}{p{3.6cm} p{3.6cm} p{3.6cm} p{1.6cm}}
\hline
\textbf{Thông số} & \textbf{Ngưỡng/Mục tiêu} & \textbf{Kết quả thực nghiệm} & \textbf{Đánh giá} \\ \hline
Độ trễ xử lý biên & $<$ 50 ms & $<$ 50 ms (sâu dưới ngưỡng) & Đạt \\
Sai số đo cảm biến & $<$ 5\% & Trong giới hạn & Đạt \\ 
Tham số LoRa & 433 MHz, SF7, BW125, CR4/5 & Đúng cấu hình & Đạt \\
PDR (gần, LOS) & cao & $\approx$ 100\% & Đạt \\ 
RSSI/SNR (gần) & tốt & $\approx -58$ dBm / $+13$ dB & Đạt \\ 
Mã hóa AES-128 & có xác thực & AES-128-GCM, $\approx$ 0,95 ms & Đạt \\ 
Mã hóa ASCON & hạng nhẹ & Hiện thực, hoạt động & Đạt \\ 
Chống Replay & $\Delta t < 5$ s & Loại đúng gói phát lại & Đạt \\ 
Lưu trữ ngoại tuyến & không mất dữ liệu & Buffer + đồng bộ lại OK & Đạt \\
Watchdog tự phục hồi & tự reset khi treo & Reset \& đếm đúng & Đạt \\ 
MCU ngủ sâu & $<$ 20 $\mu$A & vài $\mu$A (riêng MCU) & Đạt* \\ 
Stack High-Water Mark & Dư $> 20\%$, không tràn & Dư $> 30\%$ mọi tác vụ & Đạt \\
Ổn định Heap & Không rò rỉ bộ nhớ & Dao động $\pm 3$ KB, phân mảnh $< 6\%$ & Đạt \\ 
Tải CPU mỗi lõi & Còn dư địa & Core 1 $\approx$12\%, Core 0 $\approx$16\% & Đạt \\ 
Tickless Idle/Deep Sleep & Giảm dòng MCU & 45 $\to$ 18 $\to$ $<$1 mA & Đạt \\ \hline
\end{tabular}
\end{table}

Ghi chú (*): chỉ tiêu dòng ngủ sâu dưới $20\,\mu A$ đạt được đối với riêng phần vi điều khiển; dòng tổng của hệ thống cao hơn do quyết định cấp nguồn liên tục cho cảm biến khí MOX -- một điểm đánh đổi có chủ đích đã được phân tích. Các giá trị định lượng trong bảng là kết quả thực nghiệm minh họa; người thực hiện có thể cập nhật số đo chính xác của lần chạy nghiệm thu cuối cùng.

\subsection{Kết luận chương}
Chương 5 đã tiến hành kiểm thử và đánh giá hệ thống một cách toàn diện trên cả hai phương diện phần cứng và phần mềm. Phần đánh giá kết quả đã xác nhận sản phẩm được hiện thực hóa thành công và vận hành ổn định trên cả ba tầng. Phần kiểm thử định lượng đã kiểm chứng từng chỉ tiêu kỹ thuật thông qua các bài đo có phương pháp rõ ràng: từ độ ổn định nguồn, tính đúng đắn tín hiệu giao tiếp, độ chính xác cảm biến và hiệu quả lọc nhiễu, độ trễ Điện toán biên, chất lượng truyền thông LoRa, tính đúng đắn và chi phí của hai thuật toán mã hóa, đến độ tin cậy trước các sự cố và đặc tính năng lượng. Bảng tổng kết cho thấy hệ thống đáp ứng phần lớn các chỉ tiêu đã cam kết, với một điểm đánh đổi về năng lượng đã được phân tích minh bạch. Những kết quả này khẳng định tính khả thi và giá trị thực tiễn của giải pháp, đồng thời tạo cơ sở vững chắc cho các kết luận và định hướng phát triển ở Chương 6.

\newpage
\section*{CHƯƠNG 6. KẾT LUẬN VÀ HƯỚNG PHÁT TRIỂN}
\addcontentsline{toc}{section}{\numberline {} CHƯƠNG 6. KẾT LUẬN VÀ HƯỚNG PHÁT TRIỂN}
\setcounter{section}{6}
\setcounter{subsection}{0}
\setcounter{figure}{0}
\setcounter{table}{0}

\subsection{Tóm tắt kết quả chính của đồ án}

Đề tài được thực hiện với mục tiêu nghiên cứu, thiết kế và chế tạo một hệ thống IoT giám sát vi khí hậu chuồng nuôi gia cầm ứng dụng công nghệ truyền thông LoRa kết hợp Điện toán biên (Edge Computing), nhằm xây dựng một giải pháp có khả năng hoạt động ổn định trong điều kiện thực tế, đáp ứng yêu cầu về xử lý thời gian thực, tiết kiệm năng lượng, bảo mật dữ liệu và khả năng mở rộng trong các mô hình chăn nuôi thông minh.

Trong quá trình thực hiện, đồ án đã tiến hành khảo sát cơ sở lý thuyết liên quan đến mạng cảm biến không dây, kiến trúc IoT phân tán, công nghệ truyền thông LoRa, hệ điều hành thời gian thực FreeRTOS, kỹ thuật Điện toán biên, các thuật toán lọc tín hiệu, cơ chế mã hóa AES-128 và giải pháp cập nhật firmware từ xa (OTA). Trên cơ sở đó, hệ thống được thiết kế theo kiến trúc phân tầng gồm Node cảm biến, Gateway và nền tảng Cloud, trong đó phần lớn quá trình xử lý dữ liệu và ra quyết định được thực hiện ngay tại thiết bị biên nhằm giảm tải cho hệ thống trung tâm và hạn chế sự phụ thuộc vào kết nối Internet.

Về phần cứng, đồ án đã thiết kế hoàn chỉnh thiết bị Node sử dụng vi điều khiển ESP32-S3 làm bộ xử lý trung tâm, tích hợp các cảm biến đo nhiệt độ, độ ẩm, áp suất khí quyển, cường độ ánh sáng và khí độc phục vụ giám sát môi trường chuồng nuôi. Hệ thống sử dụng module truyền thông LoRa SX1278 để truyền dữ liệu khoảng cách xa với công suất tiêu thụ thấp, đồng thời được thiết kế PCB hai lớp cùng vỏ bảo vệ đạt yêu cầu triển khai trong môi trường thực tế.

Về phần mềm, firmware được xây dựng trên nền tảng ESP-IDF kết hợp FreeRTOS với kiến trúc đa nhiệm. Các tác vụ thu thập dữ liệu, xử lý tín hiệu, truyền thông, mã hóa dữ liệu, quản lý năng lượng và lưu trữ ngoại tuyến được tổ chức thành các Task độc lập, sử dụng Queue, Mutex và Semaphore để đồng bộ tài nguyên. Thuật toán FSM được triển khai để xử lý dữ liệu cảm biến và đưa ra quyết định điều khiển ngay trên Node, giúp giảm đáng kể độ trễ phản hồi của hệ thống.

Ngoài các chức năng giám sát cơ bản, đồ án còn hiện thực nhiều cơ chế nâng cao như quản lý năng lượng bằng Hybrid Sleep, lưu trữ ngoại tuyến theo cơ chế Store-and-Forward, giám sát Watchdog, cập nhật firmware OTA, mã hóa dữ liệu AES-128 kết hợp cơ chế chống Replay Attack và công cụ kiểm thử Hardware-In-the-Loop nhằm đánh giá hiệu năng toàn bộ hệ thống.

Sau quá trình chế tạo và kiểm thử, các kết quả thực nghiệm cho thấy hệ thống hoạt động ổn định và đáp ứng hầu hết các chỉ tiêu kỹ thuật đã đề ra. Độ trễ xử lý tại Node nhỏ hơn giới hạn yêu cầu, truyền thông LoRa ổn định trong điều kiện thử nghiệm, cơ chế tiết kiệm năng lượng hoạt động hiệu quả, các tác vụ FreeRTOS được phân bổ hợp lý và không xảy ra hiện tượng tràn Stack hay rò rỉ bộ nhớ trong quá trình vận hành liên tục. Các kết quả kiểm thử cũng chứng minh hệ thống có khả năng duy trì hoạt động ngay cả khi Gateway hoặc kết nối Internet bị gián đoạn nhờ cơ chế lưu trữ ngoại tuyến và xử lý cục bộ tại thiết bị biên.

Bảng~\ref{tab:summary_result} tổng hợp các kết quả chính đạt được của đồ án.

\begin{table}[H]
\centering
\caption{Tóm tắt các kết quả chính của đồ án}
\label{tab:summary_result}
\begin{tabular}{p{5cm} p{9cm}}
\hline
\textbf{Nội dung} & \textbf{Kết quả đạt được} \\ \hline
Thiết kế phần cứng &
Hoàn thành thiết kế PCB hai lớp, chế tạo thiết bị Node sử dụng ESP32-S3. \\ 

Thu thập dữ liệu &
Giám sát nhiệt độ, độ ẩm, áp suất, ánh sáng và khí độc theo thời gian thực. \\

Truyền thông &
Triển khai thành công mạng truyền thông LoRa giữa Node và Gateway. \\

Điện toán biên &
Thực hiện tiền xử lý, lọc nhiễu và ra quyết định bằng FSM trực tiếp trên Node. \\ 

Phần mềm &
Xây dựng firmware đa nhiệm trên FreeRTOS với các cơ chế Queue, Mutex và Watchdog. \\ 

Bảo mật &
Triển khai mã hóa AES-128 kết hợp xác thực dữ liệu và chống Replay Attack. \\ 

OTA &
Hỗ trợ cập nhật firmware từ xa theo kiến trúc Hybrid OTA. \\ 
Quản lý năng lượng &
Áp dụng Hybrid Sleep và Deep Sleep nhằm giảm tiêu thụ năng lượng. \\

Kiểm thử &
Đáp ứng hầu hết các chỉ tiêu kỹ thuật đặt ra trong quá trình thiết kế và thực nghiệm. \\ \hline
\end{tabular}
\end{table}

\subsection{Đóng góp của đồ án}

\subsubsection{Về khoa học}

Đồ án đã xây dựng một mô hình hệ thống IoT phân tán kết hợp giữa công nghệ truyền thông LoRa và kiến trúc Điện toán biên nhằm giải quyết bài toán giám sát vi khí hậu trong chăn nuôi gia cầm. So với mô hình IoT truyền thống phụ thuộc chủ yếu vào Cloud, giải pháp đề xuất chuyển phần lớn quá trình xử lý dữ liệu xuống thiết bị biên, giúp giảm độ trễ phản hồi, giảm lưu lượng truyền dữ liệu và nâng cao tính tự chủ của toàn bộ hệ thống.

Bên cạnh đó, đề tài đã nghiên cứu và áp dụng phương pháp đồng thiết kế phần cứng/phần mềm (Hardware/Software Co-design) trong quá trình phát triển thiết bị Node. Cách tiếp cận này cho phép tối ưu đồng thời cấu trúc phần cứng, kiến trúc firmware và tài nguyên tính toán của vi điều khiển ESP32-S3, từ đó nâng cao hiệu năng xử lý và giảm mức tiêu thụ năng lượng.

Đồ án cũng xây dựng kiến trúc phần mềm đa nhiệm trên nền FreeRTOS với chiến lược phân chia tác vụ theo từng lõi xử lý, sử dụng các cơ chế Mutex, Queue và Semaphore nhằm bảo đảm tính thời gian thực và hạn chế hiện tượng tranh chấp tài nguyên. Thuật toán FSM được thiết kế để xử lý dữ liệu cảm biến và ra quyết định điều khiển trực tiếp trên Node, góp phần hiện thực hóa mô hình Edge Computing trong hệ thống IoT nhúng.

Ngoài ra, việc tích hợp đồng thời cơ chế bảo mật AES-128, chống Replay Attack, lưu trữ ngoại tuyến và cập nhật OTA đã tạo nên một kiến trúc firmware tương đối hoàn chỉnh, có thể được sử dụng làm nền tảng tham khảo cho các hệ thống IoT thời gian thực trong nhiều lĩnh vực khác nhau.

\subsubsection{Về thực tiễn}

Về mặt ứng dụng, đồ án đã chế tạo thành công một thiết bị Node hoàn chỉnh có khả năng triển khai trong các hệ thống giám sát môi trường chuồng nuôi gia cầm. Thiết bị có khả năng hoạt động độc lập, thu thập dữ liệu môi trường theo thời gian thực, xử lý cục bộ tại thiết bị biên và truyền dữ liệu đến Gateway thông qua mạng LoRa.

Các kết quả thực nghiệm cho thấy hệ thống vận hành ổn định trong nhiều điều kiện kiểm thử khác nhau. Các thông số về độ trễ xử lý, truyền thông vô tuyến, quản lý năng lượng, bảo mật dữ liệu, tài nguyên RTOS và độ tin cậy của hệ thống đều đáp ứng các yêu cầu kỹ thuật đặt ra trong giai đoạn thiết kế.

Bên cạnh khả năng ứng dụng trong giám sát vi khí hậu chuồng nuôi gia cầm, kiến trúc phần cứng và phần mềm của hệ thống có tính mở cao, cho phép mở rộng sang nhiều bài toán IoT khác như giám sát nhà kính, nông nghiệp thông minh, quan trắc môi trường hoặc giám sát công nghiệp với chi phí đầu tư tương đối thấp.

\subsection{Hạn chế của đồ án}

Mặc dù đã đạt được các mục tiêu nghiên cứu đặt ra, đồ án vẫn còn tồn tại một số hạn chế cần tiếp tục được nghiên cứu và hoàn thiện trong thời gian tới.

Thứ nhất, phạm vi thực nghiệm của hệ thống còn tương đối hạn chế và chủ yếu được triển khai trong môi trường thử nghiệm hoặc mô hình giả lập. Việc đánh giá trên các trang trại chăn nuôi quy mô lớn với thời gian vận hành liên tục nhiều tháng chưa được thực hiện đầy đủ.

Thứ hai, thuật toán điều khiển hiện nay chủ yếu dựa trên các ngưỡng được thiết lập trước và máy trạng thái FSM. Hệ thống chưa tích hợp các thuật toán học máy hoặc trí tuệ nhân tạo để dự báo xu hướng biến đổi môi trường, tối ưu chiến lược điều khiển hoặc phát hiện bất thường.

Thứ ba, mạng truyền thông mới được triển khai theo mô hình LoRa điểm--điểm giữa Node và Gateway. Việc triển khai đầy đủ theo chuẩn LoRaWAN với khả năng quản lý số lượng lớn thiết bị, cân bằng tải và quản lý mạng vẫn chưa được nghiên cứu trong phạm vi đồ án.

Ngoài ra, do giới hạn về thời gian và kinh phí thực hiện, một số nội dung như đánh giá tuổi thọ pin trong thời gian dài, kiểm thử độ bền cơ khí, khả năng chống nhiễu điện từ hoặc đánh giá độ ổn định của OTA trên quy mô lớn vẫn chưa được thực hiện đầy đủ.

\subsection{Hướng phát triển trong tương lai}

Trong thời gian tới, hệ thống có thể được tiếp tục nghiên cứu và phát triển theo các hướng sau.

\subsubsection{Hướng phát triển về phần cứng và tối ưu năng lượng}

Mặc dù hệ thống được xây dựng trong đồ án đã đáp ứng được các chỉ tiêu kỹ thuật đặt ra và vận hành ổn định trong quá trình kiểm thử, hiệu quả năng lượng vẫn là một trong những vấn đề cần tiếp tục được tối ưu để nâng cao khả năng triển khai thực tế lâu dài. Đối với các hệ thống IoT hoạt động bằng pin, năng lượng tiêu thụ quyết định trực tiếp đến tuổi thọ thiết bị, chi phí bảo trì cũng như khả năng triển khai trên diện rộng. Do đó, tối ưu công suất tiêu thụ của thiết bị Node được xem là một trong những hướng phát triển quan trọng nhất trong giai đoạn tiếp theo.

Qua quá trình phân tích dòng tiêu thụ của từng khối chức năng, có thể nhận thấy cảm biến khí MQ-135 là thành phần tiêu tốn năng lượng lớn nhất trong toàn bộ hệ thống. Nguyên nhân là do cảm biến sử dụng phần tử gia nhiệt (Heater) hoạt động liên tục để duy trì nhiệt độ làm việc của lớp vật liệu bán dẫn oxit kim loại (MOX). Trong cấu hình hiện tại, bộ gia nhiệt luôn được cấp nguồn nhằm đảm bảo khả năng đo tức thời, dẫn đến dòng tiêu thụ duy trì ở mức tương đối cao ngay cả khi vi điều khiển đang ở trạng thái tiết kiệm năng lượng. Điều này làm giảm đáng kể hiệu quả của các cơ chế Deep Sleep và Tickless Idle đã được triển khai trong hệ thống.

Một hướng cải tiến tiềm năng là áp dụng kỹ thuật \textit{Pulsed Heating} (gia nhiệt theo xung). Thay vì cấp nguồn liên tục cho phần tử gia nhiệt, cảm biến chỉ được kích hoạt trong những khoảng thời gian ngắn trước khi thực hiện phép đo. Sau khi hoàn tất quá trình lấy mẫu, bộ gia nhiệt được tắt hoàn toàn nhằm giảm mức tiêu thụ năng lượng trung bình. Phương pháp này đã được áp dụng trong nhiều nghiên cứu về mạng cảm biến không dây sử dụng cảm biến khí MOX nhằm kéo dài tuổi thọ pin mà vẫn duy trì độ chính xác đo ở mức chấp nhận được.

Tuy nhiên, việc ngắt nguồn bộ gia nhiệt thường xuyên có thể làm thay đổi đặc tính điện trở nền của cảm biến và gây ra hiện tượng trôi giá trị đo theo thời gian. Vì vậy, để bảo đảm độ chính xác của dữ liệu, kỹ thuật Pulsed Heating cần được kết hợp với cơ chế bù trôi Baseline. Trong phương pháp này, giá trị điện trở tham chiếu $R_0$ của cảm biến sẽ được cập nhật định kỳ và lưu trữ trong bộ nhớ không mất dữ liệu. Trước mỗi chu kỳ đo, hệ thống tiến hành hiệu chỉnh giá trị cảm biến dựa trên sai lệch giữa Baseline hiện tại và Baseline tham chiếu ban đầu. Ngoài ra, có thể áp dụng các phương pháp nội suy tuyến tính hoặc bộ lọc thích nghi để giảm ảnh hưởng của quá trình trôi đặc tính theo thời gian.

Một hướng nghiên cứu khác là thay thế cảm biến MQ-135 bằng các công nghệ cảm biến tiên tiến hơn. Hiện nay, các cảm biến điện hóa chuyên dụng cho khí Amoniac như dòng NH3-B1 hoặc các cảm biến quang học hồng ngoại không phân tán (NDIR) đang được sử dụng rộng rãi trong các hệ thống giám sát môi trường công nghiệp. Các cảm biến này có ưu điểm là không cần bộ gia nhiệt công suất lớn, độ ổn định cao hơn và khả năng chọn lọc khí tốt hơn so với cảm biến MOX truyền thống. Mặc dù chi phí đầu tư ban đầu cao hơn, việc thay thế cảm biến có thể mang lại lợi ích đáng kể về tuổi thọ pin, độ chính xác và độ tin cậy của toàn bộ hệ thống.

Bảng~\ref{tab:sensor_compare_future} trình bày so sánh giữa các công nghệ cảm biến có thể được xem xét trong quá trình nâng cấp hệ thống.

\begin{table}[H]
\centering
\caption{So sánh các công nghệ cảm biến khí tiềm năng}
\label{tab:sensor_compare_future}
\begin{tabular}{p{3cm} p{2cm} p{2.2cm} p{3cm} p{3.5cm}}
\hline
\textbf{Công nghệ} &
\textbf{Công suất} &
\textbf{Chi phí} &
\textbf{Độ chính xác} &
\textbf{Khả năng ứng dụng} \\
\hline
MOX (MQ-135) &
Cao &
Thấp &
Trung bình &
Phù hợp hệ thống chi phí thấp \\

MOX + Pulsed Heating &
Trung bình &
Thấp &
Khá &
Tối ưu cho hệ thống chạy pin \\

Điện hóa (NH$_3$-B1) &
Thấp &
Cao &
Cao &
Giám sát Amoniac chuyên dụng \\

NDIR &
Rất thấp &
Rất cao &
Rất cao &
Quan trắc môi trường chính xác cao \\
\hline
\end{tabular}
\end{table}

Bên cạnh khối cảm biến, hệ thống quản lý nguồn cũng có thể được tiếp tục tối ưu. Trong phiên bản hiện tại, thiết bị sử dụng cơ chế Hybrid Sleep kết hợp giữa Deep Sleep và đánh thức định kỳ. Tuy nhiên, vẫn còn nhiều tiềm năng khai thác từ các tính năng quản lý năng lượng của ESP32-S3. Trong tương lai, có thể nghiên cứu cơ chế đánh thức theo sự kiện (Event-Driven Wakeup) nhằm giảm số lần kích hoạt vi điều khiển không cần thiết. Thay vì thức dậy theo chu kỳ cố định, Node chỉ được đánh thức khi có thay đổi đáng kể từ cảm biến hoặc khi xuất hiện yêu cầu từ Gateway.

Ngoài ra, hệ thống có thể tích hợp các thuật toán quản lý năng lượng thích nghi dựa trên điều kiện môi trường và trạng thái pin. Ví dụ, khi mức pin giảm xuống dưới một ngưỡng xác định, thiết bị có thể tự động tăng chu kỳ lấy mẫu từ 5 phút lên 10 hoặc 15 phút nhằm kéo dài thời gian hoạt động. Ngược lại, khi xuất hiện các điều kiện bất thường như nhiệt độ tăng nhanh hoặc nồng độ khí độc vượt ngưỡng, tần suất lấy mẫu sẽ được tăng lên để bảo đảm khả năng giám sát liên tục.

Một hướng phát triển đáng chú ý khác là kết hợp nguồn năng lượng tái tạo. Việc tích hợp tấm pin năng lượng mặt trời công suất nhỏ cùng bộ điều khiển sạc MPPT có thể giúp hệ thống vận hành gần như độc lập trong thời gian dài mà không cần thay pin định kỳ. Đây là giải pháp đặc biệt phù hợp đối với các trang trại chăn nuôi quy mô lớn, nơi việc bảo trì hàng chục hoặc hàng trăm Node cảm biến có thể gây tốn kém đáng kể.

Bảng~\ref{tab:energy_future} tổng hợp các hướng nghiên cứu liên quan đến tối ưu năng lượng trong tương lai.

\begin{table}[H]
\centering
\caption{Các hướng phát triển về tối ưu năng lượng}
\label{tab:energy_future}
\begin{tabular}{p{5cm} p{4cm} p{5cm}}
\hline
\textbf{Giải pháp} &
\textbf{Mức độ phức tạp} &
\textbf{Lợi ích kỳ vọng} \\
\hline
Pulsed Heating cho cảm biến MOX &
Trung bình &
Giảm đáng kể công suất tiêu thụ cảm biến khí \\

Bù trôi Baseline tự động &
Trung bình &
Duy trì độ chính xác khi áp dụng Pulsed Heating \\

Thay cảm biến điện hóa &
Cao &
Tăng độ chính xác và kéo dài tuổi thọ pin \\

Quản lý năng lượng thích nghi &
Trung bình &
Tăng thời gian hoạt động của hệ thống \\

Năng lượng mặt trời &
Cao &
Tiến tới vận hành tự chủ lâu dài \\
\hline
\end{tabular}
\end{table}

Tóm lại, việc tối ưu năng lượng không chỉ giúp kéo dài tuổi thọ pin mà còn quyết định khả năng thương mại hóa của hệ thống trong thực tế. Các kỹ thuật như Pulsed Heating, bù Baseline, sử dụng cảm biến công suất thấp và khai thác nguồn năng lượng tái tạo là những hướng nghiên cứu có tính khả thi cao và hứa hẹn mang lại hiệu quả đáng kể trong các phiên bản tiếp theo của hệ thống.

\subsubsection{Hướng phát triển về mạng truyền thông và khả năng mở rộng}

Trong phạm vi của đồ án, hệ thống truyền thông được xây dựng theo mô hình một Gateway và một Node cảm biến sử dụng giao thức LoRa điểm--điểm (Point-to-Point). Kiến trúc này phù hợp cho giai đoạn nghiên cứu và kiểm chứng nguyên lý hoạt động của hệ thống, đồng thời giúp đơn giản hóa việc phát triển firmware cũng như đánh giá hiệu năng của từng thành phần. Tuy nhiên, khi triển khai trong các trang trại chăn nuôi quy mô lớn, nơi có hàng chục hoặc hàng trăm Node hoạt động đồng thời, mô hình truyền thông hiện tại sẽ bộc lộ nhiều hạn chế về khả năng mở rộng, quản lý truy cập kênh truyền và độ tin cậy của toàn hệ thống.

Một trong những hướng phát triển quan trọng là xây dựng cơ chế điều khiển truy cập môi trường truyền (Medium Access Control -- MAC) dành riêng cho mạng LoRa nhiều Node. Trong hệ thống hiện tại, mỗi Node truyền dữ liệu độc lập mà chưa có cơ chế phối hợp giữa các thiết bị. Nếu nhiều Node cùng phát dữ liệu trong cùng khoảng thời gian, các gói tin có thể bị chồng lấn trên kênh truyền, dẫn đến hiện tượng va chạm (Collision), làm tăng tỷ lệ mất gói và giảm hiệu quả sử dụng băng thông vô tuyến.

Để giải quyết vấn đề này, có thể nghiên cứu triển khai giao thức CSMA/CA (Carrier Sense Multiple Access with Collision Avoidance). Theo cơ chế này, trước khi truyền dữ liệu, mỗi Node sẽ thực hiện kiểm tra trạng thái kênh truyền. Nếu phát hiện kênh đang bận, Node sẽ trì hoãn việc truyền trong một khoảng thời gian ngẫu nhiên trước khi thử lại. Phương pháp này giúp giảm xác suất va chạm giữa các Node mà không yêu cầu đồng bộ thời gian chặt chẽ. Mặc dù LoRa không hỗ trợ cơ chế cảm nhận sóng mang hoàn chỉnh như Wi-Fi, việc kết hợp chỉ số RSSI cùng thời gian chờ ngẫu nhiên vẫn có thể cải thiện đáng kể hiệu suất của mạng trong các hệ thống có số lượng Node trung bình.

Đối với các trang trại có số lượng thiết bị lớn, một hướng phát triển hiệu quả hơn là áp dụng cơ chế lập lịch TDMA (Time Division Multiple Access). Trong phương pháp này, Gateway sẽ phân chia chu kỳ truyền thành nhiều khe thời gian (Time Slot), mỗi Node chỉ được phép truyền trong khe thời gian đã được cấp phát. Cơ chế này gần như loại bỏ hoàn toàn hiện tượng va chạm giữa các gói tin, đồng thời giúp dự đoán trước lưu lượng mạng và tối ưu mức tiêu thụ năng lượng của các Node do thiết bị chỉ cần đánh thức vào đúng thời điểm truyền dữ liệu.


Ngoài cơ chế truy cập kênh, phạm vi phủ sóng của hệ thống cũng là một nội dung cần được tiếp tục nghiên cứu. Trong điều kiện thực tế, các trang trại chăn nuôi có thể trải dài trên diện tích hàng chục héc-ta hoặc bị chia thành nhiều khu vực bởi tường bê tông, nhà kho và các công trình kim loại. Những vật cản này có thể làm suy giảm tín hiệu vô tuyến và khiến một số Node không thể truyền trực tiếp đến Gateway.

Để mở rộng vùng phủ sóng, có thể nghiên cứu mô hình truyền tiếp đa chặng (Multi-hop). Theo kiến trúc này, các Node ở gần Gateway sẽ đóng vai trò như các nút chuyển tiếp (Relay Node), hỗ trợ truyền dữ liệu từ các Node ở xa về trung tâm. Nhờ vậy, khoảng cách truyền hiệu dụng của toàn mạng được mở rộng đáng kể mà không cần tăng công suất phát hoặc lắp đặt thêm Gateway.

Tuy nhiên, cơ chế Multi-hop cũng đặt ra nhiều thách thức như lựa chọn tuyến đường tối ưu, cân bằng tải giữa các Node chuyển tiếp, tránh lặp gói tin và đồng bộ dữ liệu. Vì vậy, các thuật toán định tuyến nhẹ phù hợp với vi điều khiển tài nguyên hạn chế cần được nghiên cứu trong các phiên bản tiếp theo.

Một hướng phát triển khác có giá trị thực tiễn cao là chuyển đổi từ giao thức truyền thông điểm--điểm sang kiến trúc LoRaWAN. Khác với mô hình hiện tại, LoRaWAN cung cấp một lớp MAC tiêu chuẩn, hỗ trợ quản lý thiết bị, xác thực, mã hóa đầu cuối, điều khiển tốc độ dữ liệu thích nghi (Adaptive Data Rate), đồng bộ thời gian và quản lý mạng quy mô lớn. Việc áp dụng LoRaWAN sẽ giúp hệ thống tương thích với các Gateway thương mại và có khả năng tích hợp trực tiếp với các nền tảng IoT phổ biến.

Bảng~\ref{tab:lora_compare} trình bày sự khác biệt giữa mô hình truyền thông hiện tại và các hướng phát triển trong tương lai.

\begin{table}[H]
\centering
\caption{So sánh các kiến trúc truyền thông}
\label{tab:lora_compare}
\begin{tabular}{p{3cm} p{3cm} p{3cm} p{5cm}}
\hline
\textbf{Kiến trúc} &
\textbf{Quy mô Node} &
\textbf{Quản lý mạng} &
\textbf{Khả năng ứng dụng} \\
\hline
Point-to-Point &
1--5 &
Thủ công &
Nghiên cứu, thử nghiệm. \\

Multi-hop LoRa &
10--50 &
Định tuyến phân tán &
Trang trại diện tích lớn. \\

LoRaWAN &
Hàng trăm Node &
Tự động &
Triển khai thương mại. \\
\hline
\end{tabular}
\end{table}

Song song với việc mở rộng kiến trúc truyền thông, quá trình đánh giá hệ thống cũng cần được thực hiện trên quy mô lớn hơn. Trong phạm vi đồ án, các thử nghiệm chủ yếu được tiến hành với một Node nhằm kiểm chứng tính đúng đắn của thiết kế. Trong tương lai, hệ thống nên được triển khai đồng thời từ ba đến năm Node hoặc nhiều hơn trong cùng một khu vực chăn nuôi để đánh giá khả năng hoạt động của mạng dưới điều kiện tải thực tế.

Các nội dung kiểm thử nên bao gồm tỷ lệ mất gói khi nhiều Node truyền đồng thời, độ trễ truyền dữ liệu, mức sử dụng kênh vô tuyến, mức tiêu thụ năng lượng của từng Node, khả năng đồng bộ dữ liệu và độ ổn định của Gateway khi xử lý nhiều luồng thông tin cùng lúc. Việc thu thập các số liệu này sẽ giúp đánh giá chính xác giới hạn của hệ thống, đồng thời cung cấp cơ sở khoa học cho việc tối ưu các thuật toán điều phối mạng trong các phiên bản tiếp theo.

Nhìn chung, việc phát triển hệ thống theo hướng hỗ trợ nhiều Node, bổ sung cơ chế MAC hiệu quả, mở rộng vùng phủ sóng bằng Multi-hop và tiến tới tích hợp kiến trúc LoRaWAN sẽ nâng cao đáng kể khả năng mở rộng cũng như giá trị ứng dụng thực tiễn của hệ thống. Đây là bước chuyển quan trọng giúp đưa hệ thống từ quy mô nghiên cứu trong phòng thí nghiệm đến khả năng triển khai trong các trang trại chăn nuôi thông minh quy mô lớn.

\subsubsection{Hướng phát triển về trí tuệ nhân tạo tại biên và các thuật toán xử lý thông minh}

Trong phiên bản hiện tại của hệ thống, việc xử lý dữ liệu tại Node chủ yếu được thực hiện thông qua máy trạng thái hữu hạn (Finite State Machine -- FSM) kết hợp với các ngưỡng nhiệt độ, độ ẩm và nồng độ khí được thiết lập trước. Phương pháp này có ưu điểm là đơn giản, dễ triển khai trên các vi điều khiển có tài nguyên hạn chế, chi phí tính toán thấp và đảm bảo khả năng xử lý thời gian thực. Tuy nhiên, FSM chỉ phản ứng khi giá trị cảm biến đã vượt ngưỡng quy định, nghĩa là hệ thống chỉ thực hiện điều khiển sau khi điều kiện môi trường đã thay đổi. Trong thực tế, nhiều hiện tượng môi trường diễn ra theo xu hướng liên tục và có thể được dự báo trước nếu khai thác hiệu quả dữ liệu lịch sử. Vì vậy, việc tích hợp các thuật toán trí tuệ nhân tạo tại biên (TinyML) là một hướng phát triển quan trọng nhằm nâng cao tính thông minh và khả năng tự thích nghi của hệ thống.

TinyML là xu hướng triển khai các mô hình học máy có kích thước nhỏ trực tiếp trên các thiết bị nhúng có tài nguyên tính toán và bộ nhớ hạn chế. Với bộ vi điều khiển ESP32-S3 được sử dụng trong đồ án, việc thực thi các mô hình học máy nhẹ hoàn toàn khả thi nhờ bộ nhớ SRAM tương đối lớn, khả năng xử lý dấu phẩy động và hỗ trợ các thư viện như TensorFlow Lite for Microcontrollers hoặc Edge Impulse. Thay vì chỉ dựa trên các ngưỡng cố định, hệ thống có thể học từ dữ liệu thu thập được để đưa ra các quyết định điều khiển phù hợp với điều kiện vận hành thực tế của từng trang trại.

Một hướng nghiên cứu có tính khả thi cao là xây dựng mô hình dự báo ngắn hạn đối với các thông số vi khí hậu trong chuồng nuôi. Dữ liệu nhiệt độ, độ ẩm, nồng độ khí độc và cường độ ánh sáng thường có xu hướng biến đổi liên tục theo thời gian và chịu ảnh hưởng của chu kỳ ngày đêm, mật độ vật nuôi cũng như điều kiện thời tiết bên ngoài. Bằng cách sử dụng dữ liệu lịch sử trong khoảng từ vài giờ đến vài ngày, mô hình học máy có thể dự báo xu hướng thay đổi của môi trường trong khoảng thời gian từ 15 đến 30 phút tiếp theo. Khi phát hiện khả năng nhiệt độ sẽ vượt ngưỡng cho phép, hệ thống có thể chủ động kích hoạt quạt thông gió hoặc hệ thống làm mát trước khi môi trường thực sự trở nên bất lợi đối với đàn gia cầm.

Đối với các vi điều khiển có tài nguyên hạn chế, các mô hình như Decision Tree hoặc Random Forest sau khi được lượng tử hóa xuống độ chính xác 8 bit là những lựa chọn phù hợp. Các mô hình này có cấu trúc đơn giản, tốc độ suy luận nhanh và yêu cầu bộ nhớ thấp hơn nhiều so với mạng nơ-ron sâu. Ngoài ra, việc lượng tử hóa còn giúp giảm kích thước mô hình, rút ngắn thời gian xử lý và giảm mức tiêu thụ năng lượng của bộ vi điều khiển trong quá trình suy luận.

Bên cạnh bài toán dự báo, hệ thống cũng có thể được phát triển theo hướng phát hiện bất thường (Anomaly Detection). Trong quá trình vận hành thực tế, không phải mọi sự cố đều xuất hiện dưới dạng giá trị vượt ngưỡng. Một số trường hợp như cảm biến bị hỏng, quạt thông gió mất tác dụng, cháy cục bộ hoặc rò rỉ khí độc có thể tạo ra các mẫu dữ liệu bất thường mà hệ thống FSM khó có thể nhận biết. Do đó, việc bổ sung các thuật toán phát hiện bất thường sẽ giúp nâng cao khả năng giám sát và đảm bảo an toàn cho hệ thống.

Một trong những giải pháp phù hợp là sử dụng mô hình One-Class Support Vector Machine (One-Class SVM). Khác với các mô hình phân loại truyền thống cần dữ liệu đã được gán nhãn cho cả trạng thái bình thường và trạng thái lỗi, One-Class SVM chỉ cần học từ dữ liệu vận hành bình thường của hệ thống. Khi xuất hiện các mẫu dữ liệu có đặc trưng khác biệt đáng kể so với tập dữ liệu đã học, mô hình sẽ cảnh báo nguy cơ bất thường để người quản lý tiến hành kiểm tra. Phương pháp này đặc biệt phù hợp với các hệ thống IoT trong nông nghiệp, nơi số lượng mẫu lỗi thường rất ít và khó thu thập.

Ngoài các mô hình học máy truyền thống, trong tương lai có thể nghiên cứu áp dụng các mạng nơ-ron nhẹ như Tiny Neural Network hoặc các mô hình Autoencoder kích thước nhỏ nhằm phát hiện các bất thường phức tạp hơn. Tuy nhiên, việc triển khai các mô hình này cần được cân nhắc giữa độ chính xác đạt được và tài nguyên phần cứng của thiết bị. Trong nhiều trường hợp, các mô hình học máy đơn giản nhưng được tối ưu tốt sẽ mang lại hiệu quả cao hơn so với các mô hình có cấu trúc quá phức tạp.

Một hướng phát triển khác là xây dựng cơ chế tự điều chỉnh ngưỡng điều khiển dựa trên dữ liệu lịch sử. Trong phiên bản hiện tại, các giá trị ngưỡng được thiết lập cố định theo từng giai đoạn phát triển của gia cầm. Tuy nhiên, điều kiện môi trường thực tế tại mỗi trang trại có thể khác nhau do sự khác biệt về giống vật nuôi, mật độ nuôi, kết cấu chuồng trại và điều kiện khí hậu của từng địa phương. Bằng cách phân tích dữ liệu vận hành trong thời gian dài, hệ thống có thể tự động hiệu chỉnh các ngưỡng điều khiển để phù hợp hơn với điều kiện thực tế, từ đó nâng cao hiệu quả điều hòa môi trường và giảm mức tiêu thụ điện năng của các thiết bị chấp hành.

Để triển khai các thuật toán học máy tại biên một cách hiệu quả, cần xây dựng quy trình huấn luyện và cập nhật mô hình phù hợp. Dữ liệu thu thập từ các Node sẽ được lưu trữ trên nền tảng Cloud để phục vụ quá trình tiền xử lý và huấn luyện mô hình trên máy tính có năng lực tính toán cao. Sau khi hoàn tất quá trình huấn luyện, mô hình được lượng tử hóa, tối ưu kích thước và chuyển đổi sang định dạng phù hợp trước khi nạp vào bộ nhớ Flash của ESP32-S3. Trong quá trình vận hành, Node chỉ thực hiện bước suy luận (Inference), nhờ đó giảm đáng kể yêu cầu về tài nguyên tính toán.

Trong tương lai, sự kết hợp giữa Edge Computing và TinyML sẽ tạo nên một hệ thống IoT có khả năng tự học, tự thích nghi và đưa ra quyết định thông minh ngay tại thiết bị biên. Điều này không chỉ giúp giảm lưu lượng truyền dữ liệu lên Cloud, giảm độ trễ phản hồi mà còn nâng cao tính ổn định của hệ thống trong điều kiện mất kết nối Internet. Đồng thời, việc triển khai các mô hình học máy nhẹ trên ESP32-S3 sẽ góp phần khẳng định tính khả thi của xu hướng trí tuệ nhân tạo tại biên trong các ứng dụng nông nghiệp thông minh, mở ra khả năng xây dựng các hệ thống giám sát môi trường có mức độ tự động hóa và độ tin cậy cao hơn trong tương lai.

\subsubsection{Hướng phát triển về phần mềm, giao diện người dùng và hoàn thiện sản phẩm}

Bên cạnh việc nâng cao hiệu năng phần cứng và thuật toán xử lý tại thiết bị biên, hệ thống phần mềm cũng cần tiếp tục được phát triển nhằm nâng cao trải nghiệm người dùng, khả năng quản lý tập trung và mức độ sẵn sàng khi triển khai trong các trang trại chăn nuôi thực tế. Trong phạm vi đồ án, giao diện giám sát chủ yếu tập trung vào việc hiển thị các thông số cảm biến theo thời gian thực và điều khiển cơ bản thông qua nền tảng Cloud. Mặc dù đã đáp ứng yêu cầu kiểm chứng chức năng của hệ thống, giao diện hiện tại vẫn còn tương đối đơn giản và chưa khai thác hết tiềm năng của một hệ thống IoT hiện đại.

Một hướng phát triển quan trọng là xây dựng nền tảng quản lý tập trung dưới dạng Dashboard Web kết hợp với ứng dụng di động đa nền tảng sử dụng Flutter hoặc React Native. Giao diện mới cần được thiết kế theo hướng trực quan, thân thiện với người sử dụng và hỗ trợ quản lý đồng thời nhiều khu vực chăn nuôi. Thay vì chỉ hiển thị các giá trị tức thời của cảm biến, hệ thống nên cung cấp các biểu đồ biến thiên theo thời gian, thống kê theo ngày, tuần và tháng, giúp người quản lý dễ dàng theo dõi xu hướng biến đổi của các thông số môi trường.

Ngoài chức năng giám sát, Dashboard có thể bổ sung nhiều công cụ hỗ trợ quản lý như cấu hình ngưỡng cảnh báo, thiết lập chu kỳ lấy mẫu, theo dõi trạng thái hoạt động của từng Node, giám sát dung lượng pin, cường độ tín hiệu LoRa (RSSI, SNR), cũng như quản lý nhật ký sự kiện của toàn bộ hệ thống. Việc trực quan hóa các thông tin này không chỉ giúp người vận hành nhanh chóng phát hiện các bất thường mà còn tạo điều kiện thuận lợi cho công tác bảo trì và tối ưu hệ thống trong quá trình sử dụng lâu dài.

Đối với ứng dụng di động, hệ thống có thể tích hợp cơ chế thông báo đẩy (Push Notification) hoặc gửi tin nhắn SMS khi phát hiện các sự kiện nguy hiểm như nhiệt độ vượt ngưỡng, nồng độ khí độc tăng cao, Node mất kết nối hoặc điện áp pin giảm xuống dưới mức an toàn. Việc nhận được cảnh báo ngay trên điện thoại giúp người quản lý có thể phản ứng kịp thời mà không cần thường xuyên truy cập Dashboard để theo dõi hệ thống.

Song song với việc phát triển giao diện người dùng, các cơ chế bảo mật cũng cần được tiếp tục hoàn thiện. Trong phiên bản hiện tại, hệ thống đã áp dụng mã hóa AES-128 kết hợp cơ chế xác thực và chống tấn công phát lại (Replay Attack). Tuy nhiên, toàn bộ các Node vẫn sử dụng cùng một khóa bảo mật, điều này có thể làm giảm mức độ an toàn khi hệ thống được mở rộng với số lượng thiết bị lớn.

Một hướng nghiên cứu phù hợp là xây dựng cơ chế quản lý khóa (Key Management) theo từng thiết bị. Theo đó, mỗi Node sẽ được cấp một định danh và khóa bí mật riêng trong quá trình khởi tạo. Gateway đóng vai trò quản lý và xác thực các Node trước khi chấp nhận dữ liệu truyền về. Ngoài ra, khóa bảo mật có thể được thay đổi định kỳ thông qua các gói tin điều khiển được mã hóa nhằm hạn chế nguy cơ lộ khóa trong quá trình vận hành lâu dài. Việc áp dụng cơ chế quản lý khóa riêng biệt sẽ giúp tăng cường khả năng chống giả mạo thiết bị, đồng thời nâng cao mức độ an toàn cho toàn bộ mạng IoT.

Bên cạnh các nội dung về phần mềm, việc hoàn thiện thiết kế cơ khí của sản phẩm cũng là một hướng phát triển cần được quan tâm. Trong đồ án, thiết bị Node mới được chế tạo dưới dạng nguyên mẫu nhằm phục vụ mục đích nghiên cứu và kiểm thử. Để có thể triển khai ngoài thực địa trong thời gian dài, thiết kế vỏ bảo vệ cần được tối ưu hơn về cả độ bền cơ học lẫn khả năng chống chịu môi trường.

Một trong những mục tiêu quan trọng là hoàn thiện vỏ hộp đạt chuẩn bảo vệ IP65 như chỉ tiêu kỹ thuật đã đề ra. Điều này đòi hỏi phải tối ưu cấu trúc gioăng cao su, vị trí các đầu nối cáp, khe thoát khí và màng lọc chống nước nhằm đảm bảo thiết bị vẫn có khả năng trao đổi không khí với môi trường trong khi ngăn chặn bụi và nước xâm nhập vào bên trong. Đồng thời, các thử nghiệm về khả năng chịu nhiệt, chống tia cực tím, chống ăn mòn và chống rung động cũng cần được thực hiện để đánh giá độ bền của thiết bị trong điều kiện vận hành thực tế.

Bảng~\ref{tab:software_future} tổng hợp các hướng phát triển liên quan đến phần mềm và hoàn thiện sản phẩm.

\begin{table}[H]
\centering
\caption{Các hướng phát triển về phần mềm và sản phẩm}
\label{tab:software_future}
\begin{tabular}{p{4.5cm} p{2.5cm} p{6.5cm}}
\hline
\textbf{Hướng phát triển} &
\textbf{Mức ưu tiên} &
\textbf{Lợi ích mang lại} \\
\hline
Dashboard Web chuyên nghiệp &
Cao &
Quản lý tập trung nhiều Node, trực quan hóa dữ liệu. \\

Ứng dụng Flutter/React Native &
Cao &
Giám sát và điều khiển hệ thống mọi lúc, mọi nơi. \\

Push Notification và SMS &
Trung bình &
Cảnh báo sự cố theo thời gian thực. \\

Quản lý khóa bảo mật riêng cho từng Node &
Cao &
Nâng cao mức độ an toàn và chống giả mạo thiết bị. \\

Hoàn thiện vỏ hộp IP65 &
Cao &
Đáp ứng điều kiện vận hành ngoài thực địa. \\
\hline
\end{tabular}
\end{table}

Cuối cùng, một trong những nội dung có ý nghĩa quyết định đối với khả năng thương mại hóa của hệ thống là kiểm thử thực địa trong thời gian dài. Trong phạm vi đồ án, việc đánh giá chủ yếu được thực hiện trong môi trường phòng thí nghiệm hoặc mô hình thử nghiệm với số lượng thiết bị còn hạn chế. Do giới hạn về thời gian và điều kiện triển khai, hệ thống chưa được kiểm chứng đầy đủ dưới các điều kiện vận hành thực tế như sự thay đổi của thời tiết, độ ẩm cao, bụi bẩn, hơi ăn mòn từ chất thải chăn nuôi hoặc sự suy giảm dung lượng pin sau thời gian dài sử dụng.

Trong tương lai, hệ thống nên được triển khai thử nghiệm tại một trang trại chăn nuôi thực tế với tối thiểu từ ba đến năm Node hoạt động liên tục trong thời gian ít nhất một tháng. Quá trình thử nghiệm cần theo dõi các chỉ tiêu như độ ổn định của kết nối LoRa, tỷ lệ mất gói tin, độ chính xác của cảm biến, mức tiêu thụ năng lượng, tuổi thọ pin, khả năng chống chịu của vỏ thiết bị và mức độ ổn định của hệ thống khi vận hành liên tục. Các số liệu thu được sẽ là cơ sở quan trọng để đánh giá tính khả thi của giải pháp cũng như tối ưu các thành phần phần cứng và phần mềm trước khi triển khai trên quy mô lớn.

Bảng~\ref{tab:field_test_future} trình bày kế hoạch kiểm thử thực địa được đề xuất cho các giai đoạn nghiên cứu tiếp theo.

\begin{table}[H]
\centering
\caption{Định hướng kiểm thử thực địa trong tương lai}
\label{tab:field_test_future}
\begin{tabular}{p{4cm} p{4cm} p{6cm}}
\hline
\textbf{Hạng mục} &
\textbf{Mục tiêu} &
\textbf{Chỉ tiêu đánh giá} \\
\hline
Triển khai nhiều Node &
Đánh giá khả năng mở rộng &
Tỷ lệ mất gói, độ trễ, độ ổn định mạng. \\

Vận hành liên tục trong 01 tháng &
Đánh giá độ tin cậy &
Tuổi thọ pin, độ ổn định cảm biến, lỗi hệ thống. \\

Đánh giá cơ khí &
Kiểm tra độ bền môi trường &
Khả năng chống nước, bụi, ăn mòn và rung động. \\

Đánh giá hiệu năng tổng thể &
Xác nhận tính khả thi &
Khả năng đáp ứng yêu cầu triển khai thực tế. \\
\hline
\end{tabular}
\end{table}

Tổng hợp các định hướng trên cho thấy hệ thống được phát triển trong đồ án đã xây dựng được nền tảng phần cứng và phần mềm tương đối hoàn chỉnh cho bài toán giám sát vi khí hậu chuồng nuôi gia cầm ứng dụng công nghệ LoRa và Điện toán biên. Trong tương lai, việc tiếp tục tối ưu quản lý năng lượng, mở rộng mạng truyền thông, tích hợp các thuật toán trí tuệ nhân tạo tại biên, nâng cao tính bảo mật, hoàn thiện giao diện người dùng và kiểm chứng hệ thống trong điều kiện thực tế sẽ góp phần nâng cao hiệu năng, độ tin cậy cũng như khả năng thương mại hóa của sản phẩm. Những hướng nghiên cứu này không chỉ phù hợp với xu thế phát triển của Internet vạn vật và nông nghiệp thông minh mà còn tạo tiền đề để hệ thống có thể được mở rộng sang nhiều lĩnh vực giám sát môi trường khác như nhà kính, thủy sản, kho bảo quản nông sản hoặc quan trắc môi trường công nghiệp.

%\newpage
%\input{Thuc_Hanh/tailieuthamkhao}

\newpage
\addcontentsline{toc}{section}{\numberline{} TÀI LIỆU THAM KHẢO}
\bibliographystyle{ieeetr}
\bibliography{Thuc_Hanh/ref}
\end{document}
